%%C070001 p

本章包含以下案例：


%%C070002 li

\begin{itemize}
\item 从二进制文件读写原始数据


%%C070003 li

\item 从二进制文件读写对象


%%C070004 li

\item 为流使用本地化设置


%%C070005 li

\item 用I/O操纵器控制流的输出


%%C070006 li

\item 使用货币I/O操纵器


%%C070007 li

\item 使用时间I/O操纵器


%%C070008 li

\item 处理文件系统路径


%%C070009 li

\item 创建、复制与删除文件和目录


%%C070010 li

\item 清除文件内容


%%C070011 li

\item 检查现有文件或目录的属性


%%C070012 li

\item 枚举目录内容


%%C070013 li

\item 查找文件

\end{itemize}


%%C070014 h2

\section{引言}


%%C070015 p

C++标准库中最重要的组成部分之一是基于流的输入输出库，它使开发者能够处理文件、内存流或其他类型的I/O设备。本章第一部分给出一些常见流操作的解决方案，例如读写数据、本地化设置，以及操控流的输入与输出。本章第二部分探讨C++17新增的\texttt{filesystem}库，它使开发者能对文件系统及其对象（例如文件和目录）执行操作。


%%C070016 h2

\section{从二进制文件读写原始数据}


%%C070017 p

程序处理的部分数据必须以多种方式持久化到磁盘文件中，既可以存入数据库，也可以写入平面文件，形式可以是文本数据，也可以是二进制数据。本案例和下一个案例关注的是把原始数据与对象持久化到二进制文件，以及从二进制文件加载。这里的原始数据指无结构数据；本案例将讨论写入和读取一个缓冲区的内容（即一段连续内存，可以是类C数组、\texttt{std::\allowbreak{}vector}或\texttt{std::\allowbreak{}array}）。


%%C070018 h3

\subsection{准备工作}


%%C070019 p

阅读本案例之前，你应当熟悉标准流输入输出库；不过下面也会给出理解本案例所必需的说明。你还应了解二进制文件与文本文件的区别。


%%C070020 p

本案例将使用\texttt{ofstream}和\texttt{ifstream}类，它们位于\texttt{<\allowbreak{}fstream>}头文件的\texttt{std}命名空间中。


%%C070021 p

在下面的示例中，我们将把以下数据写入二进制文件（随后再读回）：


%%C070022 code

\begin{codeblockbg}
\begin{Verbatim}[commandchars=\\\{\},breaklines=true,breakanywhere=true,fontsize=\codefont,xleftmargin=17bp]
\PYG{n}{std}\PYG{o}{:}\PYG{o}{:}\PYG{n}{vector}\PYG{o}{\PYGZlt{}}\PYG{k+kt}{unsigned}\PYG{+w}{ }\PYG{k+kt}{char}\PYG{o}{\PYGZgt{}}\PYG{+w}{ }\PYG{n}{output}\PYG{+w}{ }\PYG{p}{\PYGZob{}}\PYG{l+m+mi}{0}\PYG{p}{,}\PYG{+w}{ }\PYG{l+m+mi}{1}\PYG{p}{,}\PYG{+w}{ }\PYG{l+m+mi}{2}\PYG{p}{,}\PYG{+w}{ }\PYG{l+m+mi}{3}\PYG{p}{,}\PYG{+w}{ }\PYG{l+m+mi}{4}\PYG{p}{,}\PYG{+w}{ }\PYG{l+m+mi}{5}\PYG{p}{,}\PYG{+w}{ }\PYG{l+m+mi}{6}\PYG{p}{,}\PYG{+w}{ }\PYG{l+m+mi}{7}\PYG{p}{,}\PYG{+w}{ }\PYG{l+m+mi}{8}\PYG{p}{,}\PYG{+w}{ }\PYG{l+m+mi}{9}\PYG{p}{\PYGZcb{}}\PYG{p}{;}
\end{Verbatim}
\end{codeblockbg}


%%C070023 h3

\subsection{具体做法}


%%C070024 p

要把一个缓冲区的内容（在本例中是\texttt{std::\allowbreak{}vector}）写入二进制文件，应执行以下步骤：


%%C070025 p

以二进制模式打开一个用于写入的文件流，具体做法是创建


%%C070026 p

\texttt{std::\allowbreak{}ofstream}类的实例：


%%C070027 code

\begin{codeblockbg}
\begin{Verbatim}[commandchars=\\\{\},breaklines=true,breakanywhere=true,fontsize=\codefont,xleftmargin=17bp]
\PYG{n}{std}\PYG{o}{:}\PYG{o}{:}\PYG{n}{ofstream}\PYG{+w}{ }\PYG{n}{ofile}\PYG{p}{(}\PYG{l+s}{\PYGZdq{}}\PYG{l+s}{sample.bin}\PYG{l+s}{\PYGZdq{}}\PYG{p}{,}\PYG{+w}{ }\PYG{n}{std}\PYG{o}{:}\PYG{o}{:}\PYG{n}{ios}\PYG{o}{:}\PYG{o}{:}\PYG{n}{binary}\PYG{p}{)}\PYG{p}{;}
\end{Verbatim}
\end{codeblockbg}


%%C070028 p

在向文件写入数据之前，确保文件确实已经打开：


%%C070029 code

\begin{codeblockbg}
\begin{Verbatim}[commandchars=\\\{\},breaklines=true,breakanywhere=true,fontsize=\codefont,xleftmargin=17bp]
\PYG{k}{if}\PYG{p}{(}\PYG{n}{ofile}\PYG{p}{.}\PYG{n}{is\PYGZus{}open}\PYG{p}{(}\PYG{p}{)}\PYG{p}{)}
\PYG{p}{\PYGZob{}}
\PYG{+w}{  }\PYG{c+c1}{// 流文件操作}
\PYG{p}{\PYGZcb{}}
\end{Verbatim}
\end{codeblockbg}


%%C070030 p

写数据时，提供指向字符数组的指针以及要写入的


%%C070031 p

字符数：


%%C070032 code

\begin{codeblockbg}
\begin{Verbatim}[commandchars=\\\{\},breaklines=true,breakanywhere=true,fontsize=\codefont,xleftmargin=17bp]
\PYG{n}{ofile}\PYG{p}{.}\PYG{n}{write}\PYG{p}{(}\PYG{k}{reinterpret\PYGZus{}cast}\PYG{o}{\PYGZlt{}}\PYG{k+kt}{char}\PYG{o}{*}\PYG{o}{\PYGZgt{}}\PYG{p}{(}\PYG{n}{output}\PYG{p}{.}\PYG{n}{data}\PYG{p}{(}\PYG{p}{)}\PYG{p}{)}\PYG{p}{,}
\PYG{+w}{            }\PYG{n}{output}\PYG{p}{.}\PYG{n}{size}\PYG{p}{(}\PYG{p}{)}\PYG{p}{)}\PYG{p}{;}
\end{Verbatim}
\end{codeblockbg}


%%C070033 p

把流缓冲区的内容刷新到实际的磁盘文件；关闭流时这一操作会自动


%%C070034 p

完成：


%%C070035 code

\begin{codeblockbg}
\begin{Verbatim}[commandchars=\\\{\},breaklines=true,breakanywhere=true,fontsize=\codefont,xleftmargin=17bp]
\PYG{n}{ofile}\PYG{p}{.}\PYG{n}{close}\PYG{p}{(}\PYG{p}{)}\PYG{p}{;}
\end{Verbatim}
\end{codeblockbg}


%%C070036 p

要把二进制文件的全部内容读入一个缓冲区，应执行以下步骤：


%%C070037 p

以二进制模式打开一个用于读取文件的文件流，具体做法是创建


%%C070038 p

\texttt{std::\allowbreak{}ifstream}类的实例：


%%C070039 code

\begin{codeblockbg}
\begin{Verbatim}[commandchars=\\\{\},breaklines=true,breakanywhere=true,fontsize=\codefont,xleftmargin=17bp]
\PYG{n}{std}\PYG{o}{:}\PYG{o}{:}\PYG{n}{ifstream}\PYG{+w}{ }\PYG{n}{ifile}\PYG{p}{(}\PYG{l+s}{\PYGZdq{}}\PYG{l+s}{sample.bin}\PYG{l+s}{\PYGZdq{}}\PYG{p}{,}\PYG{+w}{ }\PYG{n}{std}\PYG{o}{:}\PYG{o}{:}\PYG{n}{ios}\PYG{o}{:}\PYG{o}{:}\PYG{n}{binary}\PYG{p}{)}\PYG{p}{;}
\end{Verbatim}
\end{codeblockbg}


%%C070040 p

在从文件读取数据之前，确保文件确实已经打开：


%%C070041 code

\begin{codeblockbg}
\begin{Verbatim}[commandchars=\\\{\},breaklines=true,breakanywhere=true,fontsize=\codefont,xleftmargin=17bp]
\PYG{k}{if}\PYG{p}{(}\PYG{n}{ifile}\PYG{p}{.}\PYG{n}{is\PYGZus{}open}\PYG{p}{(}\PYG{p}{)}\PYG{p}{)}
\PYG{p}{\PYGZob{}}
\PYG{+w}{  }\PYG{c+c1}{// 流文件操作}
\PYG{p}{\PYGZcb{}}
\end{Verbatim}
\end{codeblockbg}


%%C070042 p

确定文件长度：把输入位置指示器移到文件的


%%C070043 p

末尾，读出其值，再把指示器移到开头：


%%C070044 code

\begin{codeblockbg}
\begin{Verbatim}[commandchars=\\\{\},breaklines=true,breakanywhere=true,fontsize=\codefont,xleftmargin=17bp]
\PYG{n}{ifile}\PYG{p}{.}\PYG{n}{seekg}\PYG{p}{(}\PYG{l+m+mi}{0}\PYG{p}{,}\PYG{+w}{ }\PYG{n}{std}\PYG{o}{:}\PYG{o}{:}\PYG{n}{ios\PYGZus{}base}\PYG{o}{:}\PYG{o}{:}\PYG{n}{end}\PYG{p}{)}\PYG{p}{;}
\PYG{k}{auto}\PYG{+w}{ }\PYG{n}{length}\PYG{+w}{ }\PYG{o}{=}\PYG{+w}{ }\PYG{n}{ifile}\PYG{p}{.}\PYG{n}{tellg}\PYG{p}{(}\PYG{p}{)}\PYG{p}{;}
\PYG{n}{ifile}\PYG{p}{.}\PYG{n}{seekg}\PYG{p}{(}\PYG{l+m+mi}{0}\PYG{p}{,}\PYG{+w}{ }\PYG{n}{std}\PYG{o}{:}\PYG{o}{:}\PYG{n}{ios\PYGZus{}base}\PYG{o}{:}\PYG{o}{:}\PYG{n}{beg}\PYG{p}{)}\PYG{p}{;}
\end{Verbatim}
\end{codeblockbg}


%%C070045 p

为读取文件内容分配内存：


%%C070046 code

\begin{codeblockbg}
\begin{Verbatim}[commandchars=\\\{\},breaklines=true,breakanywhere=true,fontsize=\codefont,xleftmargin=17bp]
\PYG{n}{std}\PYG{o}{:}\PYG{o}{:}\PYG{n}{vector}\PYG{o}{\PYGZlt{}}\PYG{k+kt}{unsigned}\PYG{+w}{ }\PYG{k+kt}{char}\PYG{o}{\PYGZgt{}}\PYG{+w}{ }\PYG{n}{input}\PYG{p}{;}
\PYG{n}{input}\PYG{p}{.}\PYG{n}{resize}\PYG{p}{(}\PYG{k}{static\PYGZus{}cast}\PYG{o}{\PYGZlt{}}\PYG{k+kt}{size\PYGZus{}t}\PYG{o}{\PYGZgt{}}\PYG{p}{(}\PYG{n}{length}\PYG{p}{)}\PYG{p}{)}\PYG{p}{;}
\end{Verbatim}
\end{codeblockbg}


%%C070047 p

读取文件内容到已分配的缓冲区：提供指向用于接收数据的字符数组的


%%C070048 p

指针，以及要读取的字符数：


%%C070049 code

\begin{codeblockbg}
\begin{Verbatim}[commandchars=\\\{\},breaklines=true,breakanywhere=true,fontsize=\codefont,xleftmargin=17bp]
\PYG{n}{ifile}\PYG{p}{.}\PYG{n}{read}\PYG{p}{(}\PYG{k}{reinterpret\PYGZus{}cast}\PYG{o}{\PYGZlt{}}\PYG{k+kt}{char}\PYG{o}{*}\PYG{o}{\PYGZgt{}}\PYG{p}{(}\PYG{n}{input}\PYG{p}{.}\PYG{n}{data}\PYG{p}{(}\PYG{p}{)}\PYG{p}{)}\PYG{p}{,}\PYG{+w}{ }\PYG{n}{length}\PYG{p}{)}\PYG{p}{;}
\end{Verbatim}
\end{codeblockbg}


%%C070050 p

检查读取操作是否成功完成：


%%C070051 code

\begin{codeblockbg}
\begin{Verbatim}[commandchars=\\\{\},breaklines=true,breakanywhere=true,fontsize=\codefont,xleftmargin=17bp]
\PYG{k}{auto}\PYG{+w}{ }\PYG{n}{success}\PYG{+w}{ }\PYG{o}{=}\PYG{+w}{ }\PYG{o}{!}\PYG{n}{ifile}\PYG{p}{.}\PYG{n}{fail}\PYG{p}{(}\PYG{p}{)}\PYG{+w}{ }\PYG{o}{\PYGZam{}}\PYG{o}{\PYGZam{}}\PYG{+w}{ }\PYG{n}{length}\PYG{+w}{ }\PYG{o}{=}\PYG{o}{=}\PYG{+w}{ }\PYG{n}{ifile}\PYG{p}{.}\PYG{n}{gcount}\PYG{p}{(}\PYG{p}{)}\PYG{p}{;}
\end{Verbatim}
\end{codeblockbg}


%%C070052 p

最后，关闭文件流：


%%C070053 code

\begin{codeblockbg}
\begin{Verbatim}[commandchars=\\\{\},breaklines=true,breakanywhere=true,fontsize=\codefont,xleftmargin=17bp]
\PYG{n}{ifile}\PYG{p}{.}\PYG{n}{close}\PYG{p}{(}\PYG{p}{)}\PYG{p}{;}
\end{Verbatim}
\end{codeblockbg}


%%C070054 h3

\subsection{工作原理}


%%C070055 p

标准基于流的输入输出库提供了多个类，用于实现高层的输入、输出或同时输入输出的文件流、字符串流与字符数组操作，还提供控制这些流行为的操纵器，以及若干预定义的流对象（\texttt{cin}/\texttt{wcin}、\texttt{cout}/\texttt{wcout}、\texttt{cerr}/\texttt{wcerr}和\texttt{clog}/\texttt{wclog}）。


%%C070056 p

这些流以类模板实现；针对文件，该库提供了以下几个类：


%%C070057 li

\begin{itemize}
\item \texttt{basic\_\allowbreak{}filebuf}实现原始文件的输入输出操作，其语义与C的\texttt{FILE}流类似。


%%C070058 li

\item \texttt{basic\_\allowbreak{}ifstream}实现由\texttt{basic\_\allowbreak{}istream}流接口定义的高层文件流输入操作，内部使用一个\texttt{basic\_\allowbreak{}filebuf}对象。


%%C070059 li

\item \texttt{basic\_\allowbreak{}ofstream}实现由\texttt{basic\_\allowbreak{}ostream}流接口定义的高层文件流输出操作，内部使用一个\texttt{basic\_\allowbreak{}filebuf}对象。


%%C070060 li

\item \texttt{basic\_\allowbreak{}fstream}实现由\texttt{basic\_\allowbreak{}iostream}流接口定义的高层文件流输入输出操作，内部使用一个\texttt{basic\_\allowbreak{}filebuf}对象。

\end{itemize}


%%C070061 p

\texttt{<\allowbreak{}fstream>}头文件还定义了前面那些类模板的若干typedef。前面示例中使用的\texttt{ofstream}和\texttt{ifstream}对象就是这些类型同义词：


%%C070062 code

\begin{codeblockbg}
\begin{Verbatim}[commandchars=\\\{\},breaklines=true,breakanywhere=true,fontsize=\codefont,xleftmargin=17bp]
\PYG{k}{typedef}\PYG{+w}{ }\PYG{n}{basic\PYGZus{}ifstream}\PYG{o}{\PYGZlt{}}\PYG{k+kt}{char}\PYG{o}{\PYGZgt{}}\PYG{+w}{    }\PYG{n}{ifstream}\PYG{p}{;}
\PYG{k}{typedef}\PYG{+w}{ }\PYG{n}{basic\PYGZus{}ifstream}\PYG{o}{\PYGZlt{}}\PYG{k+kt}{wchar\PYGZus{}t}\PYG{o}{\PYGZgt{}}\PYG{+w}{ }\PYG{n}{wifstream}\PYG{p}{;}
\PYG{k}{typedef}\PYG{+w}{ }\PYG{n}{basic\PYGZus{}ofstream}\PYG{o}{\PYGZlt{}}\PYG{k+kt}{char}\PYG{o}{\PYGZgt{}}\PYG{+w}{    }\PYG{n}{ofstream}\PYG{p}{;}
\PYG{k}{typedef}\PYG{+w}{ }\PYG{n}{basic\PYGZus{}ofstream}\PYG{o}{\PYGZlt{}}\PYG{k+kt}{wchar\PYGZus{}t}\PYG{o}{\PYGZgt{}}\PYG{+w}{ }\PYG{n}{wofstream}\PYG{p}{;}
\PYG{k}{typedef}\PYG{+w}{ }\PYG{n}{basic\PYGZus{}fstream}\PYG{o}{\PYGZlt{}}\PYG{k+kt}{char}\PYG{o}{\PYGZgt{}}\PYG{+w}{     }\PYG{n}{fstream}\PYG{p}{;}
\PYG{k}{typedef}\PYG{+w}{ }\PYG{n}{basic\PYGZus{}fstream}\PYG{o}{\PYGZlt{}}\PYG{k+kt}{wchar\PYGZus{}t}\PYG{o}{\PYGZgt{}}\PYG{+w}{  }\PYG{n}{wfstream}\PYG{p}{;}
\end{Verbatim}
\end{codeblockbg}


%%C070063 p

上一节展示了如何向文件流写入原始数据、以及从文件流读取原始数据。其工作原理在此详细说明：


%%C070064 li

\begin{itemize}
\item 要向文件写入数据，我们实例化了\texttt{std::\allowbreak{}ofstream}类型的对象。在构造函数中传入要打开的文件名和流的打开模式，这里指定\texttt{std::\allowbreak{}ios::\allowbreak{}binary}表示二进制模式。以这种方式打开文件会丢弃文件原有的内容。如果要在已有文件末尾追加内容，还应使用标志\texttt{std::\allowbreak{}ios::\allowbreak{}app}（即\myttbr{std::\allowbreak{}ios::\allowbreak{}app | std::\allowbreak{}ios::\allowbreak{}binary}）。该构造函数会在其底层的原始文件对象（即\texttt{basic\_\allowbreak{}filebuf}对象）上调用\texttt{open()}。如果该操作失败，会设置失败位。为检查流是否已成功关联到文件设备，我们使用了\texttt{is\_\allowbreak{}open()}（它内部调用底层\texttt{basic\_\allowbreak{}filebuf}的同名方法）。向文件流写入数据通过\texttt{write()}方法完成，该方法接受指向要写入的字符串（字符序列）的指针和要写入的字符数。由于该方法按字符序列操作，如果数据是其他类型（例如本例中的\texttt{unsigned char}），就需要\texttt{reinterpret\_\allowbreak{}cast}。写操作失败时不会设置失败位，但可能抛出\myttbr{std::\allowbreak{}ios\_\allowbreak{}base::\allowbreak{}failure}异常。不过数据并不会直接写入文件设备，而是存放在\texttt{basic\_\allowbreak{}filebuf}对象中。要把它真正写入文件，需要刷新缓冲区，即调用\texttt{flush()}。如前面示例所示，关闭文件流时会自动完成这一操作。


%%C070065 li

\item 要从文件读取数据，我们实例化了\texttt{std::\allowbreak{}ifstream}类型的对象。在构造函数中传入与为写入而打开文件时相同的实参：文件名和打开模式，即\texttt{std::\allowbreak{}ios::\allowbreak{}binary}。构造函数会在底层的\myttbr{std::\allowbreak{}basic\_\allowbreak{}filebuf}对象上调用\texttt{open()}。为检查流是否已成功关联到文件设备，我们使用了\texttt{is\_\allowbreak{}open()}（它内部调用底层\texttt{basic\_\allowbreak{}filebuf}的同名方法）。在本例中，我们把文件的全部内容读入一块内存缓冲区，具体是一个\texttt{std::\allowbreak{}vector}。在读数据之前，必须先知道文件大小，才能分配足够大的缓冲区容纳这些数据。为此，我们用\texttt{seekg()}把输入位置指示器移到文件末尾，然后调用\texttt{tellg()}返回当前位置——在此它表示文件大小（以字节为单位），随后再把输入位置指示器移到文件开头，以便从头开始读取。也可以避免为把位置指示器移到末尾而调用\texttt{seekg()}：打开文件时直接把位置指示器移到末尾，这可以通过在构造函数（或\texttt{open()}方法）中使用\texttt{std::\allowbreak{}ios::\allowbreak{}ate}打开标志实现。为文件内容分配了足够内存之后，我们用\texttt{read()}方法把数据从文件复制到内存。该方法接受指向用于接收从流中读到的数据的字符序列的指针，以及要读取的字符数。由于流按字符操作，如果缓冲区存放的是其他类型的数据（例如本例中的\texttt{unsigned char}），就需要\texttt{reinterpret\_\allowbreak{}cast}表达式。若发生错误，该操作会抛出\myttbr{std::\allowbreak{}basic\_\allowbreak{}ios::\allowbreak{}failure}异常。要确定已成功从流中读到的字符数，可以使用\texttt{gcount()}方法。读取操作完成后，我们关闭文件流。

\end{itemize}


%%C070066 p

这些示例中展示的操作是向文件流写入数据和从文件流读取数据所必需的最小集合。不过要注意，你应当对操作是否成功做适当检查，并捕获可能发生的异常。


%%C070067 p

本案例到目前为止讨论的示例代码可以重组为两个通用函数，分别用于向文件写入数据和从文件读取数据：


%%C070068 code

\begin{codeblockbg}
\begin{Verbatim}[commandchars=\\\{\},breaklines=true,breakanywhere=true,fontsize=\codefont,xleftmargin=17bp]
\PYG{k+kt}{bool}\PYG{+w}{ }\PYG{n+nf}{write\PYGZus{}data}\PYG{p}{(}\PYG{k+kt}{char}\PYG{+w}{ }\PYG{k}{const}\PYG{+w}{ }\PYG{o}{*}\PYG{+w}{ }\PYG{k}{const}\PYG{+w}{ }\PYG{n}{filename}\PYG{p}{,}
\PYG{+w}{                }\PYG{k+kt}{char}\PYG{+w}{ }\PYG{k}{const}\PYG{+w}{ }\PYG{o}{*}\PYG{+w}{ }\PYG{k}{const}\PYG{+w}{ }\PYG{n}{data}\PYG{p}{,}
\PYG{+w}{                }\PYG{k+kt}{size\PYGZus{}t}\PYG{+w}{ }\PYG{k}{const}\PYG{+w}{ }\PYG{n}{size}\PYG{p}{)}
\PYG{p}{\PYGZob{}}
\PYG{+w}{  }\PYG{k}{auto}\PYG{+w}{ }\PYG{n}{success}\PYG{+w}{ }\PYG{o}{=}\PYG{+w}{ }\PYG{n+nb}{false}\PYG{p}{;}
\PYG{+w}{  }\PYG{n}{std}\PYG{o}{:}\PYG{o}{:}\PYG{n}{ofstream}\PYG{+w}{ }\PYG{n}{ofile}\PYG{p}{(}\PYG{n}{filename}\PYG{p}{,}\PYG{+w}{ }\PYG{n}{std}\PYG{o}{:}\PYG{o}{:}\PYG{n}{ios}\PYG{o}{:}\PYG{o}{:}\PYG{n}{binary}\PYG{p}{)}\PYG{p}{;}
\PYG{k}{if}\PYG{p}{(}\PYG{n}{ofile}\PYG{p}{.}\PYG{n}{is\PYGZus{}open}\PYG{p}{(}\PYG{p}{)}\PYG{p}{)}
\PYG{p}{\PYGZob{}}
\PYG{+w}{  }\PYG{k}{try}
\PYG{+w}{  }\PYG{p}{\PYGZob{}}
\PYG{+w}{    }\PYG{n}{ofile}\PYG{p}{.}\PYG{n}{write}\PYG{p}{(}\PYG{n}{data}\PYG{p}{,}\PYG{+w}{ }\PYG{n}{size}\PYG{p}{)}\PYG{p}{;}
\PYG{+w}{     }\PYG{n}{success}\PYG{+w}{ }\PYG{o}{=}\PYG{+w}{ }\PYG{n+nb}{true}\PYG{p}{;}
\PYG{+w}{  }\PYG{p}{\PYGZcb{}}
\PYG{+w}{  }\PYG{k}{catch}\PYG{p}{(}\PYG{n}{std}\PYG{o}{:}\PYG{o}{:}\PYG{n}{ios\PYGZus{}base}\PYG{o}{:}\PYG{o}{:}\PYG{n}{failure}\PYG{+w}{ }\PYG{o}{\PYGZam{}}\PYG{p}{)}
\PYG{+w}{  }\PYG{p}{\PYGZob{}}
\PYG{+w}{    }\PYG{c+c1}{// 处理错误}
\PYG{+w}{  }\PYG{p}{\PYGZcb{}}
\PYG{+w}{  }\PYG{n}{ofile}\PYG{p}{.}\PYG{n}{close}\PYG{p}{(}\PYG{p}{)}\PYG{p}{;}
\PYG{p}{\PYGZcb{}}
\PYG{+w}{  }\PYG{k}{return}\PYG{+w}{ }\PYG{n}{success}\PYG{p}{;}
\PYG{p}{\PYGZcb{}}
\PYG{k+kt}{size\PYGZus{}t}\PYG{+w}{ }\PYG{n+nf}{read\PYGZus{}data}\PYG{p}{(}\PYG{k+kt}{char}\PYG{+w}{ }\PYG{k}{const}\PYG{+w}{ }\PYG{o}{*}\PYG{+w}{ }\PYG{k}{const}\PYG{+w}{ }\PYG{n}{filename}\PYG{p}{,}
\PYG{+w}{                 }\PYG{n}{std}\PYG{o}{:}\PYG{o}{:}\PYG{n}{function}\PYG{o}{\PYGZlt{}}\PYG{k+kt}{char}\PYG{o}{*}\PYG{p}{(}\PYG{k+kt}{size\PYGZus{}t}\PYG{+w}{ }\PYG{k}{const}\PYG{p}{)}\PYG{o}{\PYGZgt{}}\PYG{+w}{ }\PYG{n}{allocator}\PYG{p}{)}
\PYG{p}{\PYGZob{}}
\PYG{+w}{  }\PYG{k+kt}{size\PYGZus{}t}\PYG{+w}{ }\PYG{n}{readbytes}\PYG{+w}{ }\PYG{o}{=}\PYG{+w}{ }\PYG{l+m+mi}{0}\PYG{p}{;}
\PYG{+w}{  }\PYG{n}{std}\PYG{o}{:}\PYG{o}{:}\PYG{n}{ifstream}\PYG{+w}{ }\PYG{n}{ifile}\PYG{p}{(}\PYG{n}{filename}\PYG{p}{,}\PYG{+w}{ }\PYG{n}{std}\PYG{o}{:}\PYG{o}{:}\PYG{n}{ios}\PYG{o}{:}\PYG{o}{:}\PYG{n}{ate}\PYG{+w}{ }\PYG{o}{|}\PYG{+w}{ }\PYG{n}{std}\PYG{o}{:}\PYG{o}{:}\PYG{n}{ios}\PYG{o}{:}\PYG{o}{:}\PYG{n}{binary}\PYG{p}{)}\PYG{p}{;}
\PYG{k}{if}\PYG{p}{(}\PYG{n}{ifile}\PYG{p}{.}\PYG{n}{is\PYGZus{}open}\PYG{p}{(}\PYG{p}{)}\PYG{p}{)}
\PYG{p}{\PYGZob{}}
\PYG{+w}{  }\PYG{k}{auto}\PYG{+w}{ }\PYG{n}{length}\PYG{+w}{ }\PYG{o}{=}\PYG{+w}{ }\PYG{k}{static\PYGZus{}cast}\PYG{o}{\PYGZlt{}}\PYG{k+kt}{size\PYGZus{}t}\PYG{o}{\PYGZgt{}}\PYG{p}{(}\PYG{n}{ifile}\PYG{p}{.}\PYG{n}{tellg}\PYG{p}{(}\PYG{p}{)}\PYG{p}{)}\PYG{p}{;}
\PYG{+w}{  }\PYG{n}{ifile}\PYG{p}{.}\PYG{n}{seekg}\PYG{p}{(}\PYG{l+m+mi}{0}\PYG{p}{,}\PYG{+w}{ }\PYG{n}{std}\PYG{o}{:}\PYG{o}{:}\PYG{n}{ios\PYGZus{}base}\PYG{o}{:}\PYG{o}{:}\PYG{n}{beg}\PYG{p}{)}\PYG{p}{;}
\PYG{k}{auto}\PYG{+w}{ }\PYG{n}{buffer}\PYG{+w}{ }\PYG{o}{=}\PYG{+w}{ }\PYG{n}{allocator}\PYG{p}{(}\PYG{n}{length}\PYG{p}{)}\PYG{p}{;}
\PYG{k}{try}
\PYG{p}{\PYGZob{}}
\PYG{+w}{  }\PYG{n}{ifile}\PYG{p}{.}\PYG{n}{read}\PYG{p}{(}\PYG{n}{buffer}\PYG{p}{,}\PYG{+w}{ }\PYG{n}{length}\PYG{p}{)}\PYG{p}{;}
\PYG{+w}{  }\PYG{n}{readbytes}\PYG{+w}{ }\PYG{o}{=}\PYG{+w}{ }\PYG{k}{static\PYGZus{}cast}\PYG{o}{\PYGZlt{}}\PYG{k+kt}{size\PYGZus{}t}\PYG{o}{\PYGZgt{}}\PYG{p}{(}\PYG{n}{ifile}\PYG{p}{.}\PYG{n}{gcount}\PYG{p}{(}\PYG{p}{)}\PYG{p}{)}\PYG{p}{;}
\PYG{p}{\PYGZcb{}}
\PYG{k}{catch}\PYG{+w}{ }\PYG{p}{(}\PYG{n}{std}\PYG{o}{:}\PYG{o}{:}\PYG{n}{ios\PYGZus{}base}\PYG{o}{:}\PYG{o}{:}\PYG{n}{failure}\PYG{+w}{ }\PYG{o}{\PYGZam{}}\PYG{p}{)}
\PYG{p}{\PYGZob{}}
\PYG{+w}{  }\PYG{c+c1}{// 处理错误}
\PYG{p}{\PYGZcb{}}
\PYG{+w}{  }\PYG{n}{ifile}\PYG{p}{.}\PYG{n}{close}\PYG{p}{(}\PYG{p}{)}\PYG{p}{;}
\PYG{p}{\PYGZcb{}}
\PYG{+w}{  }\PYG{k}{return}\PYG{+w}{ }\PYG{n}{readbytes}\PYG{p}{;}
\PYG{p}{\PYGZcb{}}
\end{Verbatim}
\end{codeblockbg}


%%C070069 p

\texttt{write\_\allowbreak{}data()}是一个函数，它接受文件名、指向字符数组的指针及其长度，并把该数组写入指定文件。\texttt{read\_\allowbreak{}data()}是一个函数，它接受文件名和一个用于分配缓冲区的函数，并把文件的全部内容读入该分配函数返回的缓冲区。下面是使用这两个函数的一个示例：


%%C070070 code

\begin{codeblockbg}
\begin{Verbatim}[commandchars=\\\{\},breaklines=true,breakanywhere=true,fontsize=\codefont,xleftmargin=17bp]
\PYG{n}{std}\PYG{o}{:}\PYG{o}{:}\PYG{n}{vector}\PYG{o}{\PYGZlt{}}\PYG{k+kt}{unsigned}\PYG{+w}{ }\PYG{k+kt}{char}\PYG{o}{\PYGZgt{}}\PYG{+w}{ }\PYG{n}{output}\PYG{+w}{ }\PYG{p}{\PYGZob{}}\PYG{l+m+mi}{0}\PYG{p}{,}\PYG{+w}{ }\PYG{l+m+mi}{1}\PYG{p}{,}\PYG{+w}{ }\PYG{l+m+mi}{2}\PYG{p}{,}\PYG{+w}{ }\PYG{l+m+mi}{3}\PYG{p}{,}\PYG{+w}{ }\PYG{l+m+mi}{4}\PYG{p}{,}\PYG{+w}{ }\PYG{l+m+mi}{5}\PYG{p}{,}\PYG{+w}{ }\PYG{l+m+mi}{6}\PYG{p}{,}\PYG{+w}{ }\PYG{l+m+mi}{7}\PYG{p}{,}\PYG{+w}{ }\PYG{l+m+mi}{8}\PYG{p}{,}\PYG{+w}{ }\PYG{l+m+mi}{9}\PYG{p}{\PYGZcb{}}\PYG{p}{;}
\PYG{n}{std}\PYG{o}{:}\PYG{o}{:}\PYG{n}{vector}\PYG{o}{\PYGZlt{}}\PYG{k+kt}{unsigned}\PYG{+w}{ }\PYG{k+kt}{char}\PYG{o}{\PYGZgt{}}\PYG{+w}{ }\PYG{n}{input}\PYG{p}{;}
\PYG{k}{if}\PYG{p}{(}\PYG{n}{write\PYGZus{}data}\PYG{p}{(}\PYG{l+s}{\PYGZdq{}}\PYG{l+s}{sample.bin}\PYG{l+s}{\PYGZdq{}}\PYG{p}{,}
\PYG{+w}{              }\PYG{k}{reinterpret\PYGZus{}cast}\PYG{o}{\PYGZlt{}}\PYG{k+kt}{char}\PYG{o}{*}\PYG{o}{\PYGZgt{}}\PYG{p}{(}\PYG{n}{output}\PYG{p}{.}\PYG{n}{data}\PYG{p}{(}\PYG{p}{)}\PYG{p}{)}\PYG{p}{,}
\PYG{+w}{              }\PYG{n}{output}\PYG{p}{.}\PYG{n}{size}\PYG{p}{(}\PYG{p}{)}\PYG{p}{)}\PYG{p}{)}
\PYG{p}{\PYGZob{}}
\PYG{+w}{  }\PYG{k}{if}\PYG{p}{(}\PYG{n}{read\PYGZus{}data}\PYG{p}{(}\PYG{l+s}{\PYGZdq{}}\PYG{l+s}{sample.bin}\PYG{l+s}{\PYGZdq{}}\PYG{p}{,}
\PYG{+w}{               }\PYG{p}{[}\PYG{o}{\PYGZam{}}\PYG{n}{input}\PYG{p}{]}\PYG{p}{(}\PYG{k+kt}{size\PYGZus{}t}\PYG{+w}{ }\PYG{k}{const}\PYG{+w}{ }\PYG{n}{length}\PYG{p}{)}\PYG{+w}{ }\PYG{p}{\PYGZob{}}
\PYG{+w}{    }\PYG{n}{input}\PYG{p}{.}\PYG{n}{resize}\PYG{p}{(}\PYG{n}{length}\PYG{p}{)}\PYG{p}{;}
\PYG{+w}{    }\PYG{k}{return}\PYG{+w}{ }\PYG{k}{reinterpret\PYGZus{}cast}\PYG{o}{\PYGZlt{}}\PYG{k+kt}{char}\PYG{o}{*}\PYG{o}{\PYGZgt{}}\PYG{p}{(}\PYG{n}{input}\PYG{p}{.}\PYG{n}{data}\PYG{p}{(}\PYG{p}{)}\PYG{p}{)}\PYG{p}{;}\PYG{p}{\PYGZcb{}}\PYG{p}{)}\PYG{+w}{ }\PYG{o}{\PYGZgt{}}\PYG{+w}{ }\PYG{l+m+mi}{0}\PYG{p}{)}
\PYG{+w}{    }\PYG{p}{\PYGZob{}}
\PYG{+w}{      }\PYG{n}{std}\PYG{o}{:}\PYG{o}{:}\PYG{n}{cout}\PYG{+w}{ }\PYG{o}{\PYGZlt{}}\PYG{o}{\PYGZlt{}}\PYG{+w}{ }\PYG{p}{(}\PYG{n}{output}\PYG{+w}{ }\PYG{o}{=}\PYG{o}{=}\PYG{+w}{ }\PYG{n}{input}\PYG{+w}{ }\PYG{o}{?}\PYG{+w}{ }\PYG{l+s}{\PYGZdq{}}\PYG{l+s}{equal}\PYG{l+s}{\PYGZdq{}}\PYG{o}{:}\PYG{+w}{ }\PYG{l+s}{\PYGZdq{}}\PYG{l+s}{not equal}\PYG{l+s}{\PYGZdq{}}\PYG{p}{)}
\PYG{+w}{                }\PYG{o}{\PYGZlt{}}\PYG{o}{\PYGZlt{}}\PYG{+w}{ }\PYG{n}{std}\PYG{o}{:}\PYG{o}{:}\PYG{n}{endl}\PYG{p}{;}
\PYG{+w}{    }\PYG{p}{\PYGZcb{}}
\PYG{+w}{  }\PYG{p}{\PYGZcb{}}
\end{Verbatim}
\end{codeblockbg}


%%C070071 p

另一种做法是改用动态分配的缓冲区，而不使用\texttt{std::\allowbreak{}vector}；在整个示例中需要做的改动很小：


%%C070072 code

\begin{codeblockbg}
\begin{Verbatim}[commandchars=\\\{\},breaklines=true,breakanywhere=true,fontsize=\codefont,xleftmargin=17bp]
\PYG{n}{std}\PYG{o}{:}\PYG{o}{:}\PYG{n}{vector}\PYG{o}{\PYGZlt{}}\PYG{k+kt}{unsigned}\PYG{+w}{ }\PYG{k+kt}{char}\PYG{o}{\PYGZgt{}}\PYG{+w}{ }\PYG{n}{output}\PYG{+w}{ }\PYG{p}{\PYGZob{}}\PYG{l+m+mi}{0}\PYG{p}{,}\PYG{+w}{ }\PYG{l+m+mi}{1}\PYG{p}{,}\PYG{+w}{ }\PYG{l+m+mi}{2}\PYG{p}{,}\PYG{+w}{ }\PYG{l+m+mi}{3}\PYG{p}{,}\PYG{+w}{ }\PYG{l+m+mi}{4}\PYG{p}{,}\PYG{+w}{ }\PYG{l+m+mi}{5}\PYG{p}{,}\PYG{+w}{ }\PYG{l+m+mi}{6}\PYG{p}{,}\PYG{+w}{ }\PYG{l+m+mi}{7}\PYG{p}{,}\PYG{+w}{ }\PYG{l+m+mi}{8}\PYG{p}{,}\PYG{+w}{ }\PYG{l+m+mi}{9}\PYG{p}{\PYGZcb{}}\PYG{p}{;}
\PYG{k+kt}{unsigned}\PYG{+w}{ }\PYG{k+kt}{char}\PYG{o}{*}\PYG{+w}{ }\PYG{n}{input}\PYG{+w}{ }\PYG{o}{=}\PYG{+w}{ }\PYG{n+nb}{nullptr}\PYG{p}{;}
\PYG{k+kt}{size\PYGZus{}t}\PYG{+w}{ }\PYG{n}{readb}\PYG{+w}{ }\PYG{o}{=}\PYG{+w}{ }\PYG{l+m+mi}{0}\PYG{p}{;}
\PYG{k}{if}\PYG{p}{(}\PYG{n}{write\PYGZus{}data}\PYG{p}{(}\PYG{l+s}{\PYGZdq{}}\PYG{l+s}{sample.bin}\PYG{l+s}{\PYGZdq{}}\PYG{p}{,}
\PYG{+w}{              }\PYG{k}{reinterpret\PYGZus{}cast}\PYG{o}{\PYGZlt{}}\PYG{k+kt}{char}\PYG{o}{*}\PYG{o}{\PYGZgt{}}\PYG{p}{(}\PYG{n}{output}\PYG{p}{.}\PYG{n}{data}\PYG{p}{(}\PYG{p}{)}\PYG{p}{)}\PYG{p}{,}
\PYG{+w}{              }\PYG{n}{output}\PYG{p}{.}\PYG{n}{size}\PYG{p}{(}\PYG{p}{)}\PYG{p}{)}\PYG{p}{)}
\PYG{p}{\PYGZob{}}
\PYG{+w}{  }\PYG{k}{if}\PYG{p}{(}\PYG{p}{(}\PYG{n}{readb}\PYG{+w}{ }\PYG{o}{=}\PYG{+w}{ }\PYG{n}{read\PYGZus{}data}\PYG{p}{(}
\PYG{+w}{     }\PYG{l+s}{\PYGZdq{}}\PYG{l+s}{sample.bin}\PYG{l+s}{\PYGZdq{}}\PYG{p}{,}
\PYG{+w}{     }\PYG{p}{[}\PYG{o}{\PYGZam{}}\PYG{n}{input}\PYG{p}{]}\PYG{p}{(}\PYG{k+kt}{size\PYGZus{}t}\PYG{+w}{ }\PYG{k}{const}\PYG{+w}{ }\PYG{n}{length}\PYG{p}{)}\PYG{+w}{ }\PYG{p}{\PYGZob{}}
\PYG{+w}{    }\PYG{n}{input}\PYG{+w}{ }\PYG{o}{=}\PYG{+w}{ }\PYG{k}{new}\PYG{+w}{ }\PYG{k+kt}{unsigned}\PYG{+w}{ }\PYG{k+kt}{char}\PYG{p}{[}\PYG{n}{length}\PYG{p}{]}\PYG{p}{;}
\PYG{+w}{    }\PYG{k}{return}\PYG{+w}{ }\PYG{k}{reinterpret\PYGZus{}cast}\PYG{o}{\PYGZlt{}}\PYG{k+kt}{char}\PYG{o}{*}\PYG{o}{\PYGZgt{}}\PYG{p}{(}\PYG{n}{input}\PYG{p}{)}\PYG{p}{;}\PYG{+w}{ }\PYG{p}{\PYGZcb{}}\PYG{p}{)}\PYG{p}{)}\PYG{+w}{ }\PYG{o}{\PYGZgt{}}\PYG{+w}{ }\PYG{l+m+mi}{0}\PYG{p}{)}
\PYG{+w}{    }\PYG{p}{\PYGZob{}}
\PYG{+w}{      }\PYG{k}{auto}\PYG{+w}{ }\PYG{n}{cmp}\PYG{+w}{ }\PYG{o}{=}\PYG{+w}{ }\PYG{n}{memcmp}\PYG{p}{(}\PYG{n}{output}\PYG{p}{.}\PYG{n}{data}\PYG{p}{(}\PYG{p}{)}\PYG{p}{,}\PYG{+w}{ }\PYG{n}{input}\PYG{p}{,}\PYG{+w}{ }\PYG{n}{output}\PYG{p}{.}\PYG{n}{size}\PYG{p}{(}\PYG{p}{)}\PYG{p}{)}\PYG{p}{;}
\PYG{+w}{      }\PYG{n}{std}\PYG{o}{:}\PYG{o}{:}\PYG{n}{cout}\PYG{+w}{ }\PYG{o}{\PYGZlt{}}\PYG{o}{\PYGZlt{}}\PYG{+w}{ }\PYG{p}{(}\PYG{n}{cmp}\PYG{+w}{ }\PYG{o}{=}\PYG{o}{=}\PYG{+w}{ }\PYG{l+m+mi}{0}\PYG{+w}{ }\PYG{o}{?}\PYG{+w}{ }\PYG{l+s}{\PYGZdq{}}\PYG{l+s}{equal}\PYG{l+s}{\PYGZdq{}}\PYG{o}{:}\PYG{+w}{ }\PYG{l+s}{\PYGZdq{}}\PYG{l+s}{not equal}\PYG{l+s}{\PYGZdq{}}\PYG{p}{)}
\PYG{+w}{                }\PYG{o}{\PYGZlt{}}\PYG{o}{\PYGZlt{}}\PYG{+w}{ }\PYG{n}{std}\PYG{o}{:}\PYG{o}{:}\PYG{n}{endl}\PYG{p}{;}
\PYG{+w}{    }\PYG{p}{\PYGZcb{}}
\PYG{+w}{  }\PYG{p}{\PYGZcb{}}
\PYG{k}{delete}\PYG{+w}{ }\PYG{p}{[}\PYG{p}{]}\PYG{+w}{ }\PYG{n}{input}\PYG{p}{;}
\end{Verbatim}
\end{codeblockbg}


%%C070073 h3

\subsection{更多内容}


%%C070074 p

本案例展示的把数据从文件读入内存的做法只是若干可选做法之一。不过与其他做法相比，它是最快的方法，尽管其他做法从面向对象的角度看可能更吸引人。比较这些做法的性能超出了本案例的范围，读者可以把它当作练习。


%%C070075 p

从文件流读取数据还有以下几种可选做法：


%%C070076 li

\begin{itemize}
\item 直接用\myttbr{std::\allowbreak{}istreambuf\_\allowbreak{}iterator}迭代器初始化一个\texttt{std::\allowbreak{}vector}（类似地，也可用于\texttt{std::\allowbreak{}string}）：


%%C070077 code

\begin{codeblockbg}
\begin{Verbatim}[commandchars=\\\{\},breaklines=true,breakanywhere=true,fontsize=\codefont,xleftmargin=17bp]
\PYG{n}{std}\PYG{o}{:}\PYG{o}{:}\PYG{n}{vector}\PYG{o}{\PYGZlt{}}\PYG{k+kt}{unsigned}\PYG{+w}{ }\PYG{k+kt}{char}\PYG{o}{\PYGZgt{}}\PYG{+w}{ }\PYG{n}{input}\PYG{p}{;}
\PYG{n}{std}\PYG{o}{:}\PYG{o}{:}\PYG{n}{ifstream}\PYG{+w}{ }\PYG{n}{ifile}\PYG{p}{(}\PYG{l+s}{\PYGZdq{}}\PYG{l+s}{sample.bin}\PYG{l+s}{\PYGZdq{}}\PYG{p}{,}\PYG{+w}{ }\PYG{n}{std}\PYG{o}{:}\PYG{o}{:}\PYG{n}{ios}\PYG{o}{:}\PYG{o}{:}\PYG{n}{binary}\PYG{p}{)}\PYG{p}{;}
\PYG{k}{if}\PYG{p}{(}\PYG{n}{ifile}\PYG{p}{.}\PYG{n}{is\PYGZus{}open}\PYG{p}{(}\PYG{p}{)}\PYG{p}{)}
\PYG{p}{\PYGZob{}}
\PYG{+w}{  }\PYG{n}{input}\PYG{+w}{ }\PYG{o}{=}\PYG{+w}{ }\PYG{n}{std}\PYG{o}{:}\PYG{o}{:}\PYG{n}{vector}\PYG{o}{\PYGZlt{}}\PYG{k+kt}{unsigned}\PYG{+w}{ }\PYG{k+kt}{char}\PYG{o}{\PYGZgt{}}\PYG{p}{(}
\PYG{+w}{     }\PYG{n}{std}\PYG{o}{:}\PYG{o}{:}\PYG{n}{istreambuf\PYGZus{}iterator}\PYG{o}{\PYGZlt{}}\PYG{k+kt}{char}\PYG{o}{\PYGZgt{}}\PYG{p}{(}\PYG{n}{ifile}\PYG{p}{)}\PYG{p}{,}
\PYG{+w}{     }\PYG{n}{std}\PYG{o}{:}\PYG{o}{:}\PYG{n}{istreambuf\PYGZus{}iterator}\PYG{o}{\PYGZlt{}}\PYG{k+kt}{char}\PYG{o}{\PYGZgt{}}\PYG{p}{(}\PYG{p}{)}\PYG{p}{)}\PYG{p}{;}
\PYG{+w}{  }\PYG{n}{ifile}\PYG{p}{.}\PYG{n}{close}\PYG{p}{(}\PYG{p}{)}\PYG{p}{;}
\PYG{p}{\PYGZcb{}}
\end{Verbatim}
\end{codeblockbg}


%%C070078 li

\item 用\myttbr{std::\allowbreak{}istreambuf\_\allowbreak{}iterator}迭代器为\texttt{std::\allowbreak{}vector}的内容赋值：


%%C070079 code

\begin{codeblockbg}
\begin{Verbatim}[commandchars=\\\{\},breaklines=true,breakanywhere=true,fontsize=\codefont,xleftmargin=17bp]
\PYG{n}{std}\PYG{o}{:}\PYG{o}{:}\PYG{n}{vector}\PYG{o}{\PYGZlt{}}\PYG{k+kt}{unsigned}\PYG{+w}{ }\PYG{k+kt}{char}\PYG{o}{\PYGZgt{}}\PYG{+w}{ }\PYG{n}{input}\PYG{p}{;}
\PYG{n}{std}\PYG{o}{:}\PYG{o}{:}\PYG{n}{ifstream}\PYG{+w}{ }\PYG{n}{ifile}\PYG{p}{(}\PYG{l+s}{\PYGZdq{}}\PYG{l+s}{sample.bin}\PYG{l+s}{\PYGZdq{}}\PYG{p}{,}\PYG{+w}{ }\PYG{n}{std}\PYG{o}{:}\PYG{o}{:}\PYG{n}{ios}\PYG{o}{:}\PYG{o}{:}\PYG{n}{binary}\PYG{p}{)}\PYG{p}{;}
\PYG{k}{if}\PYG{p}{(}\PYG{n}{ifile}\PYG{p}{.}\PYG{n}{is\PYGZus{}open}\PYG{p}{(}\PYG{p}{)}\PYG{p}{)}
\PYG{p}{\PYGZob{}}
\PYG{+w}{  }\PYG{n}{ifile}\PYG{p}{.}\PYG{n}{seekg}\PYG{p}{(}\PYG{l+m+mi}{0}\PYG{p}{,}\PYG{+w}{ }\PYG{n}{std}\PYG{o}{:}\PYG{o}{:}\PYG{n}{ios\PYGZus{}base}\PYG{o}{:}\PYG{o}{:}\PYG{n}{end}\PYG{p}{)}\PYG{p}{;}
\PYG{+w}{  }\PYG{k}{auto}\PYG{+w}{ }\PYG{n}{length}\PYG{+w}{ }\PYG{o}{=}\PYG{+w}{ }\PYG{n}{ifile}\PYG{p}{.}\PYG{n}{tellg}\PYG{p}{(}\PYG{p}{)}\PYG{p}{;}
\PYG{+w}{  }\PYG{n}{ifile}\PYG{p}{.}\PYG{n}{seekg}\PYG{p}{(}\PYG{l+m+mi}{0}\PYG{p}{,}\PYG{+w}{ }\PYG{n}{std}\PYG{o}{:}\PYG{o}{:}\PYG{n}{ios\PYGZus{}base}\PYG{o}{:}\PYG{o}{:}\PYG{n}{beg}\PYG{p}{)}\PYG{p}{;}
\PYG{+w}{  }\PYG{n}{input}\PYG{p}{.}\PYG{n}{reserve}\PYG{p}{(}\PYG{k}{static\PYGZus{}cast}\PYG{o}{\PYGZlt{}}\PYG{k+kt}{size\PYGZus{}t}\PYG{o}{\PYGZgt{}}\PYG{p}{(}\PYG{n}{length}\PYG{p}{)}\PYG{p}{)}\PYG{p}{;}
\PYG{+w}{  }\PYG{n}{input}\PYG{p}{.}\PYG{n}{assign}\PYG{p}{(}
\PYG{+w}{    }\PYG{n}{std}\PYG{o}{:}\PYG{o}{:}\PYG{n}{istreambuf\PYGZus{}iterator}\PYG{o}{\PYGZlt{}}\PYG{k+kt}{char}\PYG{o}{\PYGZgt{}}\PYG{p}{(}\PYG{n}{ifile}\PYG{p}{)}\PYG{p}{,}
\PYG{+w}{    }\PYG{n}{std}\PYG{o}{:}\PYG{o}{:}\PYG{n}{istreambuf\PYGZus{}iterator}\PYG{o}{\PYGZlt{}}\PYG{k+kt}{char}\PYG{o}{\PYGZgt{}}\PYG{p}{(}\PYG{p}{)}\PYG{p}{)}\PYG{p}{;}
\PYG{+w}{  }\PYG{n}{ifile}\PYG{p}{.}\PYG{n}{close}\PYG{p}{(}\PYG{p}{)}\PYG{p}{;}
\PYG{p}{\PYGZcb{}}
\end{Verbatim}
\end{codeblockbg}


%%C070080 li

\item 用\myttbr{std::\allowbreak{}istreambuf\_\allowbreak{}iterator}迭代器配合\myttbr{std::\allowbreak{}back\_\allowbreak{}inserter}适配器，把文件流的内容复制到vector，由该适配器向vector末尾写入：


%%C070081 code

\begin{codeblockbg}
\begin{Verbatim}[commandchars=\\\{\},breaklines=true,breakanywhere=true,fontsize=\codefont,xleftmargin=17bp]
\PYG{n}{std}\PYG{o}{:}\PYG{o}{:}\PYG{n}{vector}\PYG{o}{\PYGZlt{}}\PYG{k+kt}{unsigned}\PYG{+w}{ }\PYG{k+kt}{char}\PYG{o}{\PYGZgt{}}\PYG{+w}{ }\PYG{n}{input}\PYG{p}{;}
\PYG{n}{std}\PYG{o}{:}\PYG{o}{:}\PYG{n}{ifstream}\PYG{+w}{ }\PYG{n}{ifile}\PYG{p}{(}\PYG{l+s}{\PYGZdq{}}\PYG{l+s}{sample.bin}\PYG{l+s}{\PYGZdq{}}\PYG{p}{,}\PYG{+w}{ }\PYG{n}{std}\PYG{o}{:}\PYG{o}{:}\PYG{n}{ios}\PYG{o}{:}\PYG{o}{:}\PYG{n}{binary}\PYG{p}{)}\PYG{p}{;}
\PYG{k}{if}\PYG{p}{(}\PYG{n}{ifile}\PYG{p}{.}\PYG{n}{is\PYGZus{}open}\PYG{p}{(}\PYG{p}{)}\PYG{p}{)}
\PYG{p}{\PYGZob{}}
\PYG{+w}{  }\PYG{n}{ifile}\PYG{p}{.}\PYG{n}{seekg}\PYG{p}{(}\PYG{l+m+mi}{0}\PYG{p}{,}\PYG{+w}{ }\PYG{n}{std}\PYG{o}{:}\PYG{o}{:}\PYG{n}{ios\PYGZus{}base}\PYG{o}{:}\PYG{o}{:}\PYG{n}{end}\PYG{p}{)}\PYG{p}{;}
\PYG{+w}{  }\PYG{k}{auto}\PYG{+w}{ }\PYG{n}{length}\PYG{+w}{ }\PYG{o}{=}\PYG{+w}{ }\PYG{n}{ifile}\PYG{p}{.}\PYG{n}{tellg}\PYG{p}{(}\PYG{p}{)}\PYG{p}{;}
\PYG{+w}{  }\PYG{n}{ifile}\PYG{p}{.}\PYG{n}{seekg}\PYG{p}{(}\PYG{l+m+mi}{0}\PYG{p}{,}\PYG{+w}{ }\PYG{n}{std}\PYG{o}{:}\PYG{o}{:}\PYG{n}{ios\PYGZus{}base}\PYG{o}{:}\PYG{o}{:}\PYG{n}{beg}\PYG{p}{)}\PYG{p}{;}
\PYG{+w}{  }\PYG{n}{input}\PYG{p}{.}\PYG{n}{reserve}\PYG{p}{(}\PYG{k}{static\PYGZus{}cast}\PYG{o}{\PYGZlt{}}\PYG{k+kt}{size\PYGZus{}t}\PYG{o}{\PYGZgt{}}\PYG{p}{(}\PYG{n}{length}\PYG{p}{)}\PYG{p}{)}\PYG{p}{;}
\PYG{+w}{  }\PYG{n}{std}\PYG{o}{:}\PYG{o}{:}\PYG{n}{copy}\PYG{p}{(}\PYG{n}{std}\PYG{o}{:}\PYG{o}{:}\PYG{n}{istreambuf\PYGZus{}iterator}\PYG{o}{\PYGZlt{}}\PYG{k+kt}{char}\PYG{o}{\PYGZgt{}}\PYG{p}{(}\PYG{n}{ifile}\PYG{p}{)}\PYG{p}{,}
\PYG{+w}{            }\PYG{n}{std}\PYG{o}{:}\PYG{o}{:}\PYG{n}{istreambuf\PYGZus{}iterator}\PYG{o}{\PYGZlt{}}\PYG{k+kt}{char}\PYG{o}{\PYGZgt{}}\PYG{p}{(}\PYG{p}{)}\PYG{p}{,}
\PYG{+w}{            }\PYG{n}{std}\PYG{o}{:}\PYG{o}{:}\PYG{n}{back\PYGZus{}inserter}\PYG{p}{(}\PYG{n}{input}\PYG{p}{)}\PYG{p}{)}\PYG{p}{;}
\PYG{+w}{  }\PYG{n}{ifile}\PYG{p}{.}\PYG{n}{close}\PYG{p}{(}\PYG{p}{)}\PYG{p}{;}
\PYG{p}{\PYGZcb{}}
\end{Verbatim}
\end{codeblockbg}

\end{itemize}


%%C070082 h3

\subsection{延伸阅读}


%%C070083 li

\begin{itemize}
\item \textit{从二进制文件读写对象}


%%C070084 li

\item \textit{用I/O操纵器控制流的输出}

\end{itemize}


%%C070085 h2

\section{从二进制文件读写对象}


%%C070086 p

上一个案例展示了如何向文件写入原始数据（即无结构数据）以及从文件读取原始数据。但很多时候我们必须持久化和加载对象。上一个案例展示的写法只适用于POD类型。对于其他类型，我们必须明确决定实际写入或读取的内容，因为读写指针、虚函数表以及各种元数据不仅没有意义，而且在语义上是错误的。这类操作通常称为序列化（serialization）与反序列化（deserialization）。本案例将展示如何把POD类型和非POD类型序列化到二进制文件、以及从二进制文件反序列化。


%%C070087 h3

\subsection{准备工作}


%%C070088 p

建议你先阅读上一个案例\textit{从二进制文件读写}\textit{原始数据}，然后再继续。你还应了解什么是POD类型和非POD类型，以及如何重载运算符。


%%C070089 p

本案例的示例将使用下面给出的\texttt{foo}和\texttt{foopod}类：


%%C070090 code

\begin{codeblockbg}
\begin{Verbatim}[commandchars=\\\{\},breaklines=true,breakanywhere=true,fontsize=\codefont,xleftmargin=17bp]
\PYG{k}{class}\PYG{+w}{ }\PYG{n+nc}{foo}
\PYG{p}{\PYGZob{}}
\PYG{+w}{  }\PYG{k+kt}{int}\PYG{+w}{ }\PYG{n}{i}\PYG{p}{;}
\PYG{+w}{  }\PYG{k+kt}{char}\PYG{+w}{ }\PYG{n}{c}\PYG{p}{;}
\PYG{+w}{  }\PYG{n}{std}\PYG{o}{:}\PYG{o}{:}\PYG{n}{string}\PYG{+w}{ }\PYG{n}{s}\PYG{p}{;}
\PYG{k}{public}\PYG{o}{:}
\PYG{+w}{  }\PYG{n}{foo}\PYG{p}{(}\PYG{k+kt}{int}\PYG{+w}{ }\PYG{k}{const}\PYG{+w}{ }\PYG{n}{i}\PYG{+w}{ }\PYG{o}{=}\PYG{+w}{ }\PYG{l+m+mi}{0}\PYG{p}{,}\PYG{+w}{ }\PYG{k+kt}{char}\PYG{+w}{ }\PYG{k}{const}\PYG{+w}{ }\PYG{n}{c}\PYG{+w}{ }\PYG{o}{=}\PYG{+w}{ }\PYG{l+m+mi}{0}\PYG{p}{,}\PYG{+w}{ }\PYG{n}{std}\PYG{o}{:}\PYG{o}{:}\PYG{n}{string}\PYG{+w}{ }\PYG{k}{const}\PYG{+w}{ }\PYG{o}{\PYGZam{}}\PYG{+w}{ }\PYG{n}{s}\PYG{+w}{ }\PYG{o}{=}\PYG{+w}{ }\PYG{p}{\PYGZob{}}\PYG{p}{\PYGZcb{}}\PYG{p}{)}\PYG{o}{:}
\PYG{+w}{    }\PYG{n}{i}\PYG{p}{(}\PYG{n}{i}\PYG{p}{)}\PYG{p}{,}\PYG{+w}{ }\PYG{n}{c}\PYG{p}{(}\PYG{n}{c}\PYG{p}{)}\PYG{p}{,}\PYG{+w}{ }\PYG{n}{s}\PYG{p}{(}\PYG{n}{s}\PYG{p}{)}
\PYG{+w}{  }\PYG{p}{\PYGZob{}}\PYG{p}{\PYGZcb{}}
\PYG{n}{foo}\PYG{p}{(}\PYG{n}{foo}\PYG{+w}{ }\PYG{k}{const}\PYG{+w}{ }\PYG{o}{\PYGZam{}}\PYG{p}{)}\PYG{+w}{ }\PYG{o}{=}\PYG{+w}{ }\PYG{k}{default}\PYG{p}{;}
\PYG{n}{foo}\PYG{o}{\PYGZam{}}\PYG{+w}{ }\PYG{k}{operator}\PYG{o}{=}\PYG{p}{(}\PYG{n}{foo}\PYG{+w}{ }\PYG{k}{const}\PYG{+w}{ }\PYG{o}{\PYGZam{}}\PYG{p}{)}\PYG{+w}{ }\PYG{o}{=}\PYG{+w}{ }\PYG{k}{default}\PYG{p}{;}
\PYG{k+kt}{bool}\PYG{+w}{ }\PYG{k}{operator}\PYG{o}{=}\PYG{o}{=}\PYG{p}{(}\PYG{n}{foo}\PYG{+w}{ }\PYG{k}{const}\PYG{+w}{ }\PYG{o}{\PYGZam{}}\PYG{+w}{ }\PYG{n}{rhv}\PYG{p}{)}\PYG{+w}{ }\PYG{k}{const}
\PYG{p}{\PYGZob{}}
\PYG{+w}{  }\PYG{k}{return}\PYG{+w}{ }\PYG{n}{i}\PYG{+w}{ }\PYG{o}{=}\PYG{o}{=}\PYG{+w}{ }\PYG{n}{rhv}\PYG{p}{.}\PYG{n}{i}\PYG{+w}{ }\PYG{o}{\PYGZam{}}\PYG{o}{\PYGZam{}}
\PYG{+w}{         }\PYG{n}{c}\PYG{+w}{ }\PYG{o}{=}\PYG{o}{=}\PYG{+w}{ }\PYG{n}{rhv}\PYG{p}{.}\PYG{n}{c}\PYG{+w}{ }\PYG{o}{\PYGZam{}}\PYG{o}{\PYGZam{}}
\PYG{+w}{         }\PYG{n}{s}\PYG{+w}{ }\PYG{o}{=}\PYG{o}{=}\PYG{+w}{ }\PYG{n}{rhv}\PYG{p}{.}\PYG{n}{s}\PYG{p}{;}
\PYG{p}{\PYGZcb{}}
\PYG{+w}{      }\PYG{k+kt}{bool}\PYG{+w}{ }\PYG{k}{operator}\PYG{o}{!}\PYG{o}{=}\PYG{p}{(}\PYG{n}{foo}\PYG{+w}{ }\PYG{k}{const}\PYG{+w}{ }\PYG{o}{\PYGZam{}}\PYG{+w}{ }\PYG{n}{rhv}\PYG{p}{)}\PYG{+w}{ }\PYG{k}{const}
\PYG{+w}{  }\PYG{p}{\PYGZob{}}
\PYG{+w}{    }\PYG{k}{return}\PYG{+w}{ }\PYG{o}{!}\PYG{p}{(}\PYG{o}{*}\PYG{k}{this}\PYG{+w}{ }\PYG{o}{=}\PYG{o}{=}\PYG{+w}{ }\PYG{n}{rhv}\PYG{p}{)}\PYG{p}{;}
\PYG{+w}{  }\PYG{p}{\PYGZcb{}}
\PYG{p}{\PYGZcb{}}\PYG{p}{;}
\PYG{k}{struct}\PYG{+w}{ }\PYG{n+nc}{foopod}
\PYG{p}{\PYGZob{}}
\PYG{+w}{  }\PYG{k+kt}{bool}\PYG{+w}{ }\PYG{n}{a}\PYG{p}{;}
\PYG{+w}{  }\PYG{k+kt}{char}\PYG{+w}{ }\PYG{n}{b}\PYG{p}{;}
\PYG{+w}{  }\PYG{k+kt}{int}\PYG{+w}{ }\PYG{n}{c}\PYG{p}{[}\PYG{l+m+mi}{2}\PYG{p}{]}\PYG{p}{;}
\PYG{p}{\PYGZcb{}}\PYG{p}{;}
\PYG{k+kt}{bool}\PYG{+w}{ }\PYG{k}{operator}\PYG{o}{=}\PYG{o}{=}\PYG{p}{(}\PYG{n}{foopod}\PYG{+w}{ }\PYG{k}{const}\PYG{+w}{ }\PYG{o}{\PYGZam{}}\PYG{+w}{ }\PYG{n}{f1}\PYG{p}{,}\PYG{+w}{ }\PYG{n}{foopod}\PYG{+w}{ }\PYG{k}{const}\PYG{+w}{ }\PYG{o}{\PYGZam{}}\PYG{+w}{ }\PYG{n}{f2}\PYG{p}{)}
\PYG{p}{\PYGZob{}}
\PYG{+w}{  }\PYG{k}{return}\PYG{+w}{ }\PYG{n}{f1}\PYG{p}{.}\PYG{n}{a}\PYG{+w}{ }\PYG{o}{=}\PYG{o}{=}\PYG{+w}{ }\PYG{n}{f2}\PYG{p}{.}\PYG{n}{a}\PYG{+w}{ }\PYG{o}{\PYGZam{}}\PYG{o}{\PYGZam{}}\PYG{+w}{ }\PYG{n}{f1}\PYG{p}{.}\PYG{n}{b}\PYG{+w}{ }\PYG{o}{=}\PYG{o}{=}\PYG{+w}{ }\PYG{n}{f2}\PYG{p}{.}\PYG{n}{b}\PYG{+w}{ }\PYG{o}{\PYGZam{}}\PYG{o}{\PYGZam{}}
\PYG{+w}{         }\PYG{n}{f1}\PYG{p}{.}\PYG{n}{c}\PYG{p}{[}\PYG{l+m+mi}{0}\PYG{p}{]}\PYG{+w}{ }\PYG{o}{=}\PYG{o}{=}\PYG{+w}{ }\PYG{n}{f2}\PYG{p}{.}\PYG{n}{c}\PYG{p}{[}\PYG{l+m+mi}{0}\PYG{p}{]}\PYG{+w}{ }\PYG{o}{\PYGZam{}}\PYG{o}{\PYGZam{}}\PYG{+w}{ }\PYG{n}{f1}\PYG{p}{.}\PYG{n}{c}\PYG{p}{[}\PYG{l+m+mi}{1}\PYG{p}{]}\PYG{+w}{ }\PYG{o}{=}\PYG{o}{=}\PYG{+w}{ }\PYG{n}{f2}\PYG{p}{.}\PYG{n}{c}\PYG{p}{[}\PYG{l+m+mi}{1}\PYG{p}{]}\PYG{p}{;}
\PYG{p}{\PYGZcb{}}
\end{Verbatim}
\end{codeblockbg}


%%C070091 h3

\subsection{具体做法}


%%C070092 p

要序列化/反序列化不含指针的POD类型，可使用\texttt{ofstream::\allowbreak{}write()}和\texttt{ifstream::\allowbreak{}read()}，如上一个案例所示：


%%C070093 li

\begin{itemize}
\item 用\texttt{ofstream}和\texttt{write()}方法把对象序列化到二进制文件：


%%C070094 code

\begin{codeblockbg}
\begin{Verbatim}[commandchars=\\\{\},breaklines=true,breakanywhere=true,fontsize=\codefont,xleftmargin=17bp]
\PYG{n}{std}\PYG{o}{:}\PYG{o}{:}\PYG{n}{vector}\PYG{o}{\PYGZlt{}}\PYG{n}{foopod}\PYG{o}{\PYGZgt{}}\PYG{+w}{ }\PYG{n}{output}\PYG{+w}{ }\PYG{p}{\PYGZob{}}
\PYG{+w}{  }\PYG{p}{\PYGZob{}}\PYG{n+nb}{true}\PYG{p}{,}\PYG{+w}{ }\PYG{l+s+sc}{\PYGZsq{}}\PYG{l+s+sc}{1}\PYG{l+s+sc}{\PYGZsq{}}\PYG{p}{,}\PYG{+w}{ }\PYG{p}{\PYGZob{}}\PYG{l+m+mi}{1}\PYG{p}{,}\PYG{+w}{ }\PYG{l+m+mi}{2}\PYG{p}{\PYGZcb{}}\PYG{p}{\PYGZcb{}}\PYG{p}{,}
\PYG{+w}{  }\PYG{p}{\PYGZob{}}\PYG{n+nb}{true}\PYG{p}{,}\PYG{+w}{ }\PYG{l+s+sc}{\PYGZsq{}}\PYG{l+s+sc}{2}\PYG{l+s+sc}{\PYGZsq{}}\PYG{p}{,}\PYG{+w}{ }\PYG{p}{\PYGZob{}}\PYG{l+m+mi}{3}\PYG{p}{,}\PYG{+w}{ }\PYG{l+m+mi}{4}\PYG{p}{\PYGZcb{}}\PYG{p}{\PYGZcb{}}\PYG{p}{,}
\PYG{+w}{  }\PYG{p}{\PYGZob{}}\PYG{n+nb}{false}\PYG{p}{,}\PYG{+w}{ }\PYG{l+s+sc}{\PYGZsq{}}\PYG{l+s+sc}{3}\PYG{l+s+sc}{\PYGZsq{}}\PYG{p}{,}\PYG{+w}{ }\PYG{p}{\PYGZob{}}\PYG{l+m+mi}{4}\PYG{p}{,}\PYG{+w}{ }\PYG{l+m+mi}{5}\PYG{p}{\PYGZcb{}}\PYG{p}{\PYGZcb{}}
\PYG{p}{\PYGZcb{}}\PYG{p}{;}
\PYG{n}{std}\PYG{o}{:}\PYG{o}{:}\PYG{n}{ofstream}\PYG{+w}{ }\PYG{n}{ofile}\PYG{p}{(}\PYG{l+s}{\PYGZdq{}}\PYG{l+s}{sample.bin}\PYG{l+s}{\PYGZdq{}}\PYG{p}{,}\PYG{+w}{ }\PYG{n}{std}\PYG{o}{:}\PYG{o}{:}\PYG{n}{ios}\PYG{o}{:}\PYG{o}{:}\PYG{n}{binary}\PYG{p}{)}\PYG{p}{;}
\PYG{k}{if}\PYG{p}{(}\PYG{n}{ofile}\PYG{p}{.}\PYG{n}{is\PYGZus{}open}\PYG{p}{(}\PYG{p}{)}\PYG{p}{)}
\PYG{p}{\PYGZob{}}
\PYG{+w}{  }\PYG{k}{for}\PYG{p}{(}\PYG{k}{auto}\PYG{+w}{ }\PYG{k}{const}\PYG{+w}{ }\PYG{o}{\PYGZam{}}\PYG{+w}{ }\PYG{n}{value}\PYG{+w}{ }\PYG{o}{:}\PYG{+w}{ }\PYG{n}{output}\PYG{p}{)}
\PYG{+w}{  }\PYG{p}{\PYGZob{}}
\PYG{+w}{    }\PYG{n}{ofile}\PYG{p}{.}\PYG{n}{write}\PYG{p}{(}\PYG{k}{reinterpret\PYGZus{}cast}\PYG{o}{\PYGZlt{}}\PYG{k}{const}\PYG{+w}{ }\PYG{k+kt}{char}\PYG{o}{*}\PYG{o}{\PYGZgt{}}\PYG{p}{(}\PYG{o}{\PYGZam{}}\PYG{n}{value}\PYG{p}{)}\PYG{p}{,}
\PYG{+w}{                }\PYG{k}{sizeof}\PYG{p}{(}\PYG{n}{value}\PYG{p}{)}\PYG{p}{)}\PYG{p}{;}
\PYG{+w}{  }\PYG{p}{\PYGZcb{}}
\PYG{+w}{  }\PYG{n}{ofile}\PYG{p}{.}\PYG{n}{close}\PYG{p}{(}\PYG{p}{)}\PYG{p}{;}
\PYG{p}{\PYGZcb{}}
\end{Verbatim}
\end{codeblockbg}


%%C070095 li

\item 用\texttt{ifstream}和\texttt{read()}方法从二进制文件反序列化对象：

%%C070096 code

\begin{codeblockbg}
\begin{Verbatim}[commandchars=\\\{\},breaklines=true,breakanywhere=true,fontsize=\codefont,xleftmargin=17bp]
\PYG{n}{std}\PYG{o}{:}\PYG{o}{:}\PYG{n}{vector}\PYG{o}{\PYGZlt{}}\PYG{n}{foopod}\PYG{o}{\PYGZgt{}}\PYG{+w}{ }\PYG{n}{input}\PYG{p}{;}
\PYG{n}{std}\PYG{o}{:}\PYG{o}{:}\PYG{n}{ifstream}\PYG{+w}{ }\PYG{n}{ifile}\PYG{p}{(}\PYG{l+s}{\PYGZdq{}}\PYG{l+s}{sample.bin}\PYG{l+s}{\PYGZdq{}}\PYG{p}{,}\PYG{+w}{ }\PYG{n}{std}\PYG{o}{:}\PYG{o}{:}\PYG{n}{ios}\PYG{o}{:}\PYG{o}{:}\PYG{n}{binary}\PYG{p}{)}\PYG{p}{;}
\PYG{k}{if}\PYG{p}{(}\PYG{n}{ifile}\PYG{p}{.}\PYG{n}{is\PYGZus{}open}\PYG{p}{(}\PYG{p}{)}\PYG{p}{)}
\PYG{p}{\PYGZob{}}
\PYG{+w}{  }\PYG{k}{while}\PYG{p}{(}\PYG{n+nb}{true}\PYG{p}{)}
\PYG{+w}{  }\PYG{p}{\PYGZob{}}
\PYG{+w}{    }\PYG{n}{foopod}\PYG{+w}{ }\PYG{n}{value}\PYG{p}{;}
\PYG{+w}{    }\PYG{n}{ifile}\PYG{p}{.}\PYG{n}{read}\PYG{p}{(}\PYG{k}{reinterpret\PYGZus{}cast}\PYG{o}{\PYGZlt{}}\PYG{k+kt}{char}\PYG{o}{*}\PYG{o}{\PYGZgt{}}\PYG{p}{(}\PYG{o}{\PYGZam{}}\PYG{n}{value}\PYG{p}{)}\PYG{p}{,}
\PYG{+w}{               }\PYG{k}{sizeof}\PYG{p}{(}\PYG{n}{value}\PYG{p}{)}\PYG{p}{)}\PYG{p}{;}
\PYG{k}{if}\PYG{p}{(}\PYG{n}{ifile}\PYG{p}{.}\PYG{n}{fail}\PYG{p}{(}\PYG{p}{)}\PYG{+w}{ }\PYG{o}{|}\PYG{o}{|}\PYG{+w}{ }\PYG{n}{ifile}\PYG{p}{.}\PYG{n}{eof}\PYG{p}{(}\PYG{p}{)}\PYG{p}{)}\PYG{+w}{ }\PYG{k}{break}\PYG{p}{;}
\PYG{+w}{    }\PYG{n}{input}\PYG{p}{.}\PYG{n}{push\PYGZus{}back}\PYG{p}{(}\PYG{n}{value}\PYG{p}{)}\PYG{p}{;}
\PYG{p}{\PYGZcb{}}
\PYG{+w}{  }\PYG{n}{ifile}\PYG{p}{.}\PYG{n}{close}\PYG{p}{(}\PYG{p}{)}\PYG{p}{;}
\PYG{p}{\PYGZcb{}}
\end{Verbatim}
\end{codeblockbg}

\end{itemize}


%%C070097 p

要序列化非POD类型（或包含指针的POD类型），你必须把各个数据成员的值显式写入文件；反序列化时，则必须按同样的顺序从文件显式读入这些数据成员。为说明这一点，我们来看前面定义的\texttt{foo}类：


%%C070098 li

\begin{itemize}
\item 添加一个名为\texttt{write()}的成员函数，用来序列化该类的对象。该方法接受一个\texttt{ofstream}的引用，并返回一个\texttt{bool}，表示操作是否成功：


%%C070099 code

\begin{codeblockbg}
\begin{Verbatim}[commandchars=\\\{\},breaklines=true,breakanywhere=true,fontsize=\codefont,xleftmargin=17bp]
\PYG{k+kt}{bool}\PYG{+w}{ }\PYG{n+nf}{write}\PYG{p}{(}\PYG{n}{std}\PYG{o}{:}\PYG{o}{:}\PYG{n}{ofstream}\PYG{o}{\PYGZam{}}\PYG{+w}{ }\PYG{n}{ofile}\PYG{p}{)}\PYG{+w}{ }\PYG{k}{const}
\PYG{p}{\PYGZob{}}
\PYG{+w}{  }\PYG{n}{ofile}\PYG{p}{.}\PYG{n}{write}\PYG{p}{(}\PYG{k}{reinterpret\PYGZus{}cast}\PYG{o}{\PYGZlt{}}\PYG{k}{const}\PYG{+w}{ }\PYG{k+kt}{char}\PYG{o}{*}\PYG{o}{\PYGZgt{}}\PYG{p}{(}\PYG{o}{\PYGZam{}}\PYG{n}{i}\PYG{p}{)}\PYG{p}{,}\PYG{+w}{ }\PYG{k}{sizeof}\PYG{p}{(}\PYG{n}{i}\PYG{p}{)}\PYG{p}{)}\PYG{p}{;}
\PYG{+w}{  }\PYG{n}{ofile}\PYG{p}{.}\PYG{n}{write}\PYG{p}{(}\PYG{o}{\PYGZam{}}\PYG{n}{c}\PYG{p}{,}\PYG{+w}{ }\PYG{k}{sizeof}\PYG{p}{(}\PYG{n}{c}\PYG{p}{)}\PYG{p}{)}\PYG{p}{;}
\PYG{+w}{  }\PYG{k}{auto}\PYG{+w}{ }\PYG{n}{size}\PYG{+w}{ }\PYG{o}{=}\PYG{+w}{ }\PYG{k}{static\PYGZus{}cast}\PYG{o}{\PYGZlt{}}\PYG{k+kt}{int}\PYG{o}{\PYGZgt{}}\PYG{p}{(}\PYG{n}{s}\PYG{p}{.}\PYG{n}{size}\PYG{p}{(}\PYG{p}{)}\PYG{p}{)}\PYG{p}{;}
\PYG{+w}{  }\PYG{n}{ofile}\PYG{p}{.}\PYG{n}{write}\PYG{p}{(}\PYG{k}{reinterpret\PYGZus{}cast}\PYG{o}{\PYGZlt{}}\PYG{k+kt}{char}\PYG{o}{*}\PYG{o}{\PYGZgt{}}\PYG{p}{(}\PYG{o}{\PYGZam{}}\PYG{n}{size}\PYG{p}{)}\PYG{p}{,}\PYG{+w}{ }\PYG{k}{sizeof}\PYG{p}{(}\PYG{n}{size}\PYG{p}{)}\PYG{p}{)}\PYG{p}{;}
\PYG{+w}{  }\PYG{n}{ofile}\PYG{p}{.}\PYG{n}{write}\PYG{p}{(}\PYG{n}{s}\PYG{p}{.}\PYG{n}{data}\PYG{p}{(}\PYG{p}{)}\PYG{p}{,}\PYG{+w}{ }\PYG{n}{s}\PYG{p}{.}\PYG{n}{size}\PYG{p}{(}\PYG{p}{)}\PYG{p}{)}\PYG{p}{;}
\PYG{+w}{  }\PYG{k}{return}\PYG{+w}{ }\PYG{o}{!}\PYG{n}{ofile}\PYG{p}{.}\PYG{n}{fail}\PYG{p}{(}\PYG{p}{)}\PYG{p}{;}
\PYG{p}{\PYGZcb{}}
\end{Verbatim}
\end{codeblockbg}


%%C070100 li

\item 添加一个名为\texttt{read()}的成员函数，用来反序列化该类的对象。该方法接受一个\texttt{ifstream}的引用，并返回一个\texttt{bool}，表示操作是否成功：


%%C070101 code

\begin{codeblockbg}
\begin{Verbatim}[commandchars=\\\{\},breaklines=true,breakanywhere=true,fontsize=\codefont,xleftmargin=17bp]
\PYG{k+kt}{bool}\PYG{+w}{ }\PYG{n+nf}{read}\PYG{p}{(}\PYG{n}{std}\PYG{o}{:}\PYG{o}{:}\PYG{n}{ifstream}\PYG{o}{\PYGZam{}}\PYG{+w}{ }\PYG{n}{ifile}\PYG{p}{)}
\PYG{p}{\PYGZob{}}
\PYG{+w}{  }\PYG{n}{ifile}\PYG{p}{.}\PYG{n}{read}\PYG{p}{(}\PYG{k}{reinterpret\PYGZus{}cast}\PYG{o}{\PYGZlt{}}\PYG{k+kt}{char}\PYG{o}{*}\PYG{o}{\PYGZgt{}}\PYG{p}{(}\PYG{o}{\PYGZam{}}\PYG{n}{i}\PYG{p}{)}\PYG{p}{,}\PYG{+w}{ }\PYG{k}{sizeof}\PYG{p}{(}\PYG{n}{i}\PYG{p}{)}\PYG{p}{)}\PYG{p}{;}
\PYG{+w}{  }\PYG{n}{ifile}\PYG{p}{.}\PYG{n}{read}\PYG{p}{(}\PYG{o}{\PYGZam{}}\PYG{n}{c}\PYG{p}{,}\PYG{+w}{ }\PYG{k}{sizeof}\PYG{p}{(}\PYG{n}{c}\PYG{p}{)}\PYG{p}{)}\PYG{p}{;}
\PYG{+w}{  }\PYG{k}{auto}\PYG{+w}{ }\PYG{n}{size}\PYG{+w}{ }\PYG{p}{\PYGZob{}}\PYG{l+m+mi}{0}\PYG{p}{\PYGZcb{}}\PYG{p}{;}
\PYG{+w}{  }\PYG{n}{ifile}\PYG{p}{.}\PYG{n}{read}\PYG{p}{(}\PYG{k}{reinterpret\PYGZus{}cast}\PYG{o}{\PYGZlt{}}\PYG{k+kt}{char}\PYG{o}{*}\PYG{o}{\PYGZgt{}}\PYG{p}{(}\PYG{o}{\PYGZam{}}\PYG{n}{size}\PYG{p}{)}\PYG{p}{,}\PYG{+w}{ }\PYG{k}{sizeof}\PYG{p}{(}\PYG{n}{size}\PYG{p}{)}\PYG{p}{)}\PYG{p}{;}
\PYG{+w}{  }\PYG{n}{s}\PYG{p}{.}\PYG{n}{resize}\PYG{p}{(}\PYG{n}{size}\PYG{p}{)}\PYG{p}{;}
\PYG{+w}{  }\PYG{n}{ifile}\PYG{p}{.}\PYG{n}{read}\PYG{p}{(}\PYG{k}{reinterpret\PYGZus{}cast}\PYG{o}{\PYGZlt{}}\PYG{k+kt}{char}\PYG{o}{*}\PYG{o}{\PYGZgt{}}\PYG{p}{(}\PYG{o}{\PYGZam{}}\PYG{n}{s}\PYG{p}{.}\PYG{n}{front}\PYG{p}{(}\PYG{p}{)}\PYG{p}{)}\PYG{p}{,}\PYG{+w}{ }\PYG{n}{size}\PYG{p}{)}\PYG{p}{;}
\PYG{+w}{  }\PYG{k}{return}\PYG{+w}{ }\PYG{o}{!}\PYG{n}{ifile}\PYG{p}{.}\PYG{n}{fail}\PYG{p}{(}\PYG{p}{)}\PYG{p}{;}
\PYG{p}{\PYGZcb{}}
\end{Verbatim}
\end{codeblockbg}

\end{itemize}


%%C070102 p

除上面示例的\texttt{write()}与\texttt{read()}成员函数之外，另一种做法是重载\texttt{operator<<}与\texttt{operator>>}。要做到这一点，你需要执行以下步骤：


%%C070103 p

为非成员\texttt{operator<<}与\texttt{operator>>}添加\texttt{friend}声明，声明加在要序列化/反序列化的


%%C070104 p

类中（本例中是\texttt{foo}类）：


%%C070105 code

\begin{codeblockbg}
\begin{Verbatim}[commandchars=\\\{\},breaklines=true,breakanywhere=true,fontsize=\codefont,xleftmargin=17bp]
\PYG{k}{friend}\PYG{+w}{ }\PYG{n}{std}\PYG{o}{:}\PYG{o}{:}\PYG{n}{ofstream}\PYG{o}{\PYGZam{}}\PYG{+w}{ }\PYG{k}{operator}\PYG{o}{\PYGZlt{}}\PYG{o}{\PYGZlt{}}\PYG{p}{(}\PYG{n}{std}\PYG{o}{:}\PYG{o}{:}\PYG{n}{ofstream}\PYG{o}{\PYGZam{}}\PYG{+w}{ }\PYG{n}{ofile}\PYG{p}{,}
\PYG{+w}{                                 }\PYG{n}{foo}\PYG{+w}{ }\PYG{k}{const}\PYG{o}{\PYGZam{}}\PYG{+w}{ }\PYG{n}{f}\PYG{p}{)}\PYG{p}{;}
\PYG{k}{friend}\PYG{+w}{ }\PYG{n}{std}\PYG{o}{:}\PYG{o}{:}\PYG{n}{ifstream}\PYG{o}{\PYGZam{}}\PYG{+w}{ }\PYG{k}{operator}\PYG{o}{\PYGZgt{}}\PYG{o}{\PYGZgt{}}\PYG{p}{(}\PYG{n}{std}\PYG{o}{:}\PYG{o}{:}\PYG{n}{ifstream}\PYG{o}{\PYGZam{}}\PYG{+w}{ }\PYG{n}{ifile}\PYG{p}{,}
\PYG{+w}{                                 }\PYG{n}{foo}\PYG{o}{\PYGZam{}}\PYG{+w}{ }\PYG{n}{f}\PYG{p}{)}\PYG{p}{;}
\end{Verbatim}
\end{codeblockbg}


%%C070106 p

为你的类重载\texttt{operator<<}：


%%C070107 code

\begin{codeblockbg}
\begin{Verbatim}[commandchars=\\\{\},breaklines=true,breakanywhere=true,fontsize=\codefont,xleftmargin=17bp]
\PYG{n}{std}\PYG{o}{:}\PYG{o}{:}\PYG{n}{ofstream}\PYG{o}{\PYGZam{}}\PYG{+w}{ }\PYG{k}{operator}\PYG{o}{\PYGZlt{}}\PYG{o}{\PYGZlt{}}\PYG{p}{(}\PYG{n}{std}\PYG{o}{:}\PYG{o}{:}\PYG{n}{ofstream}\PYG{o}{\PYGZam{}}\PYG{+w}{ }\PYG{n}{ofile}\PYG{p}{,}\PYG{+w}{ }\PYG{n}{foo}\PYG{+w}{ }\PYG{k}{const}\PYG{o}{\PYGZam{}}\PYG{+w}{ }\PYG{n}{f}\PYG{p}{)}
\PYG{p}{\PYGZob{}}
\PYG{+w}{  }\PYG{n}{ofile}\PYG{p}{.}\PYG{n}{write}\PYG{p}{(}\PYG{k}{reinterpret\PYGZus{}cast}\PYG{o}{\PYGZlt{}}\PYG{k}{const}\PYG{+w}{ }\PYG{k+kt}{char}\PYG{o}{*}\PYG{o}{\PYGZgt{}}\PYG{p}{(}\PYG{o}{\PYGZam{}}\PYG{n}{f}\PYG{p}{.}\PYG{n}{i}\PYG{p}{)}\PYG{p}{,}
\PYG{+w}{              }\PYG{k}{sizeof}\PYG{p}{(}\PYG{n}{f}\PYG{p}{.}\PYG{n}{i}\PYG{p}{)}\PYG{p}{)}\PYG{p}{;}
\PYG{+w}{  }\PYG{n}{ofile}\PYG{p}{.}\PYG{n}{write}\PYG{p}{(}\PYG{o}{\PYGZam{}}\PYG{n}{f}\PYG{p}{.}\PYG{n}{c}\PYG{p}{,}\PYG{+w}{ }\PYG{k}{sizeof}\PYG{p}{(}\PYG{n}{f}\PYG{p}{.}\PYG{n}{c}\PYG{p}{)}\PYG{p}{)}\PYG{p}{;}
\PYG{+w}{  }\PYG{k}{auto}\PYG{+w}{ }\PYG{n}{size}\PYG{+w}{ }\PYG{o}{=}\PYG{+w}{ }\PYG{k}{static\PYGZus{}cast}\PYG{o}{\PYGZlt{}}\PYG{k+kt}{int}\PYG{o}{\PYGZgt{}}\PYG{p}{(}\PYG{n}{f}\PYG{p}{.}\PYG{n}{s}\PYG{p}{.}\PYG{n}{size}\PYG{p}{(}\PYG{p}{)}\PYG{p}{)}\PYG{p}{;}
\PYG{+w}{  }\PYG{n}{ofile}\PYG{p}{.}\PYG{n}{write}\PYG{p}{(}\PYG{k}{reinterpret\PYGZus{}cast}\PYG{o}{\PYGZlt{}}\PYG{k+kt}{char}\PYG{o}{*}\PYG{o}{\PYGZgt{}}\PYG{p}{(}\PYG{o}{\PYGZam{}}\PYG{n}{size}\PYG{p}{)}\PYG{p}{,}\PYG{+w}{ }\PYG{k}{sizeof}\PYG{p}{(}\PYG{n}{size}\PYG{p}{)}\PYG{p}{)}\PYG{p}{;}
\PYG{+w}{  }\PYG{n}{ofile}\PYG{p}{.}\PYG{n}{write}\PYG{p}{(}\PYG{n}{f}\PYG{p}{.}\PYG{n}{s}\PYG{p}{.}\PYG{n}{data}\PYG{p}{(}\PYG{p}{)}\PYG{p}{,}\PYG{+w}{ }\PYG{n}{f}\PYG{p}{.}\PYG{n}{s}\PYG{p}{.}\PYG{n}{size}\PYG{p}{(}\PYG{p}{)}\PYG{p}{)}\PYG{p}{;}
\PYG{+w}{  }\PYG{k}{return}\PYG{+w}{ }\PYG{n}{ofile}\PYG{p}{;}
\PYG{p}{\PYGZcb{}}
\end{Verbatim}
\end{codeblockbg}


%%C070108 p

为你的类重载\texttt{operator>>}：


%%C070109 code

\begin{codeblockbg}
\begin{Verbatim}[commandchars=\\\{\},breaklines=true,breakanywhere=true,fontsize=\codefont,xleftmargin=17bp]
\PYG{n}{std}\PYG{o}{:}\PYG{o}{:}\PYG{n}{ifstream}\PYG{o}{\PYGZam{}}\PYG{+w}{ }\PYG{k}{operator}\PYG{o}{\PYGZgt{}}\PYG{o}{\PYGZgt{}}\PYG{p}{(}\PYG{n}{std}\PYG{o}{:}\PYG{o}{:}\PYG{n}{ifstream}\PYG{o}{\PYGZam{}}\PYG{+w}{ }\PYG{n}{ifile}\PYG{p}{,}\PYG{+w}{ }\PYG{n}{foo}\PYG{o}{\PYGZam{}}\PYG{+w}{ }\PYG{n}{f}\PYG{p}{)}
\PYG{p}{\PYGZob{}}
\PYG{+w}{  }\PYG{n}{ifile}\PYG{p}{.}\PYG{n}{read}\PYG{p}{(}\PYG{k}{reinterpret\PYGZus{}cast}\PYG{o}{\PYGZlt{}}\PYG{k+kt}{char}\PYG{o}{*}\PYG{o}{\PYGZgt{}}\PYG{p}{(}\PYG{o}{\PYGZam{}}\PYG{n}{f}\PYG{p}{.}\PYG{n}{i}\PYG{p}{)}\PYG{p}{,}\PYG{+w}{ }\PYG{k}{sizeof}\PYG{p}{(}\PYG{n}{f}\PYG{p}{.}\PYG{n}{i}\PYG{p}{)}\PYG{p}{)}\PYG{p}{;}
\PYG{+w}{  }\PYG{n}{ifile}\PYG{p}{.}\PYG{n}{read}\PYG{p}{(}\PYG{o}{\PYGZam{}}\PYG{n}{f}\PYG{p}{.}\PYG{n}{c}\PYG{p}{,}\PYG{+w}{ }\PYG{k}{sizeof}\PYG{p}{(}\PYG{n}{f}\PYG{p}{.}\PYG{n}{c}\PYG{p}{)}\PYG{p}{)}\PYG{p}{;}
\PYG{+w}{  }\PYG{k}{auto}\PYG{+w}{ }\PYG{n}{size}\PYG{+w}{ }\PYG{p}{\PYGZob{}}\PYG{l+m+mi}{0}\PYG{p}{\PYGZcb{}}\PYG{p}{;}
\PYG{+w}{  }\PYG{n}{ifile}\PYG{p}{.}\PYG{n}{read}\PYG{p}{(}\PYG{k}{reinterpret\PYGZus{}cast}\PYG{o}{\PYGZlt{}}\PYG{k+kt}{char}\PYG{o}{*}\PYG{o}{\PYGZgt{}}\PYG{p}{(}\PYG{o}{\PYGZam{}}\PYG{n}{size}\PYG{p}{)}\PYG{p}{,}\PYG{+w}{ }\PYG{k}{sizeof}\PYG{p}{(}\PYG{n}{size}\PYG{p}{)}\PYG{p}{)}\PYG{p}{;}
\PYG{+w}{  }\PYG{n}{f}\PYG{p}{.}\PYG{n}{s}\PYG{p}{.}\PYG{n}{resize}\PYG{p}{(}\PYG{n}{size}\PYG{p}{)}\PYG{p}{;}
\PYG{+w}{  }\PYG{n}{ifile}\PYG{p}{.}\PYG{n}{read}\PYG{p}{(}\PYG{k}{reinterpret\PYGZus{}cast}\PYG{o}{\PYGZlt{}}\PYG{k+kt}{char}\PYG{o}{*}\PYG{o}{\PYGZgt{}}\PYG{p}{(}\PYG{o}{\PYGZam{}}\PYG{n}{f}\PYG{p}{.}\PYG{n}{s}\PYG{p}{.}\PYG{n}{front}\PYG{p}{(}\PYG{p}{)}\PYG{p}{)}\PYG{p}{,}\PYG{+w}{ }\PYG{n}{size}\PYG{p}{)}\PYG{p}{;}
\PYG{+w}{  }\PYG{k}{return}\PYG{+w}{ }\PYG{n}{ifile}\PYG{p}{;}
\PYG{p}{\PYGZcb{}}
\end{Verbatim}
\end{codeblockbg}


%%C070110 h3

\subsection{工作原理}


%%C070111 p

无论序列化的是整个对象（POD类型）还是只序列化其中的一部分，我们用的都是上一个案例中讨论过的流类：\texttt{ofstream}用于输出文件流，\texttt{ifstream}用于输入文件流。关于用这些标准类读写数据的细节，那个案例已经讲过，这里不再重复。


%%C070112 p

当你把对象序列化到文件、从文件反序列化时，应避免把指针的值写入文件，也绝不能从文件读取指针值，因为指针代表内存地址，跨进程没有意义，甚至在同一进程中稍后一段时间也没有意义。相反，你应当写入指针所指向的数据，并把数据读入指针所指向的对象。这是一条通用原则；在实践中，你可能会遇到源数据中有多个指针指向同一对象的情况，这时你可能希望只写入一份副本，并以相应的方式处理读取。


%%C070113 p

如果要序列化的对象是POD类型，完全可以像讨论原始数据时那样做。本案例的示例中序列化了一串\texttt{foopod}类型的对象。反序列化时，我们在循环中从文件流读取，直到读到文件末尾或发生失败。这里的读取方式看起来可能有些反直觉，但换一种写法会让最后读到的值被重复添加：


%%C070114 li

\begin{itemize}
\item 读取在一个无限循环中进行。


%%C070115 li

\item 在循环中执行一次读取操作。


%%C070116 li

\item 检查是否失败或到达文件末尾，若出现任一情况，就退出无限循环。


%%C070117 li

\item 把读到的值加入输入序列，然后继续循环。

\end{itemize}


%%C070118 p

如果改用带退出条件的循环来读取，而该条件检查的是文件末尾标志位，也就是\myttbr{while(\allowbreak{}!ifile.\allowbreak{}eof())}，那么最后一个值会被加入输入序列两次。原因在于：读到最后一个值时，文件末尾其实还没有遇到（文件末尾标记位于文件最后一个字节之后）。只有到下一次读取尝试时才会到达该标记，因此那时才置位流的eof位。然而输入变量里仍然是最后一个值，因为它没有被任何新数据覆盖，于是这个值第二次被加入输入向量。


%%C070119 p

如果要序列化和反序列化的对象是非POD类型，就不能把这些对象当作原始数据来读写。例如，这类对象可能有虚表。把虚表写入文件不会出问题，尽管它并没有什么价值；但是从文件读取、从而覆盖对象的虚表，会对对象和程序造成灾难性后果。


%%C070120 p

序列化和反序列化非POD类型有多种可选做法，其中一些已在上一节演示：要么提供显式的写入与读取方法，要么重载标准的\texttt{<<}与\texttt{>>}运算符。第二种做法的优点是，能让你的类用于泛型代码，而这些代码正是用这些运算符把对象写入流文件、从流文件读出的。


%%C070121 note

\begin{myNote}
当你打算序列化和反序列化对象时，最好从一开始就考虑为数据做版本管理，以免日后数据结构变化时出现问题。具体该如何做版本管理，超出了本案例的范围。
\end{myNote}


%%C070122 h3

\subsection{延伸阅读}


%%C070123 li

\begin{itemize}
\item \textit{从二进制文件读写原始数据}


%%C070124 li

\item \textit{用I/O操纵器控制流的输出}

\end{itemize}


%%C070125 h2

\section{为流使用本地化设置}


%%C070126 p

向流写入或从流读取的方式可能取决于语言与地区设置，例如数字、时间值或货币值的写入与解析，以及字符串的比较（collating）。C++输入输出库通过\textit{locales}与\textit{facets}提供了一套处理国际化特性的通用机制。本案例将学习如何用区域设置（locale）控制输入输出流的行为。


%%C070127 h3

\subsection{准备工作}


%%C070128 p

本案例所有示例都使用预定义的控制台流对象\texttt{std::\allowbreak{}cout}。不过，同样的做法适用于所有输入输出流对象。此外，本案例的示例中还要用到下列对象和lambda函数：


%%C070129 code

\begin{codeblockbg}
\begin{Verbatim}[commandchars=\\\{\},breaklines=true,breakanywhere=true,fontsize=\codefont,xleftmargin=17bp]
\PYG{k}{auto}\PYG{+w}{ }\PYG{n}{now}\PYG{+w}{ }\PYG{o}{=}\PYG{+w}{ }\PYG{n}{std}\PYG{o}{:}\PYG{o}{:}\PYG{n}{chrono}\PYG{o}{:}\PYG{o}{:}\PYG{n}{system\PYGZus{}clock}\PYG{o}{:}\PYG{o}{:}\PYG{n}{now}\PYG{p}{(}\PYG{p}{)}\PYG{p}{;}
\PYG{k}{auto}\PYG{+w}{ }\PYG{n}{stime}\PYG{+w}{ }\PYG{o}{=}\PYG{+w}{ }\PYG{n}{std}\PYG{o}{:}\PYG{o}{:}\PYG{n}{chrono}\PYG{o}{:}\PYG{o}{:}\PYG{n}{system\PYGZus{}clock}\PYG{o}{:}\PYG{o}{:}\PYG{n}{to\PYGZus{}time\PYGZus{}t}\PYG{p}{(}\PYG{n}{now}\PYG{p}{)}\PYG{p}{;}
\PYG{k}{auto}\PYG{+w}{ }\PYG{n}{ltime}\PYG{+w}{ }\PYG{o}{=}\PYG{+w}{ }\PYG{n}{std}\PYG{o}{:}\PYG{o}{:}\PYG{n}{localtime}\PYG{p}{(}\PYG{o}{\PYGZam{}}\PYG{n}{stime}\PYG{p}{)}\PYG{p}{;}
\PYG{n}{std}\PYG{o}{:}\PYG{o}{:}\PYG{n}{vector}\PYG{o}{\PYGZlt{}}\PYG{n}{std}\PYG{o}{:}\PYG{o}{:}\PYG{n}{string}\PYG{o}{\PYGZgt{}}\PYG{+w}{ }\PYG{n}{names}
\PYG{+w}{  }\PYG{p}{\PYGZob{}}\PYG{l+s}{\PYGZdq{}}\PYG{l+s}{John}\PYG{l+s}{\PYGZdq{}}\PYG{p}{,}\PYG{+w}{ }\PYG{l+s}{\PYGZdq{}}\PYG{l+s}{adele}\PYG{l+s}{\PYGZdq{}}\PYG{p}{,}\PYG{+w}{ }\PYG{l+s}{\PYGZdq{}}\PYG{l+s}{Øivind}\PYG{l+s}{\PYGZdq{}}\PYG{p}{,}\PYG{+w}{ }\PYG{l+s}{\PYGZdq{}}\PYG{l+s}{François}\PYG{l+s}{\PYGZdq{}}\PYG{p}{,}\PYG{+w}{ }\PYG{l+s}{\PYGZdq{}}\PYG{l+s}{Robert}\PYG{l+s}{\PYGZdq{}}\PYG{p}{,}\PYG{+w}{ }\PYG{l+s}{\PYGZdq{}}\PYG{l+s}{Åke}\PYG{l+s}{\PYGZdq{}}\PYG{p}{\PYGZcb{}}\PYG{p}{;}
\PYG{k}{auto}\PYG{+w}{ }\PYG{n}{sort\PYGZus{}and\PYGZus{}print}\PYG{+w}{ }\PYG{o}{=}\PYG{+w}{ }\PYG{p}{[}\PYG{p}{]}\PYG{p}{(}\PYG{n}{std}\PYG{o}{:}\PYG{o}{:}\PYG{n}{vector}\PYG{o}{\PYGZlt{}}\PYG{n}{std}\PYG{o}{:}\PYG{o}{:}\PYG{n}{string}\PYG{o}{\PYGZgt{}}\PYG{+w}{ }\PYG{n}{v}\PYG{p}{,}
\PYG{+w}{                         }\PYG{n}{std}\PYG{o}{:}\PYG{o}{:}\PYG{n}{locale}\PYG{+w}{ }\PYG{k}{const}\PYG{+w}{ }\PYG{o}{\PYGZam{}}\PYG{+w}{ }\PYG{n}{loc}\PYG{p}{)}
\PYG{p}{\PYGZob{}}
\PYG{+w}{  }\PYG{n}{std}\PYG{o}{:}\PYG{o}{:}\PYG{n}{sort}\PYG{p}{(}\PYG{n}{v}\PYG{p}{.}\PYG{n}{begin}\PYG{p}{(}\PYG{p}{)}\PYG{p}{,}\PYG{+w}{ }\PYG{n}{v}\PYG{p}{.}\PYG{n}{end}\PYG{p}{(}\PYG{p}{)}\PYG{p}{,}\PYG{+w}{ }\PYG{n}{loc}\PYG{p}{)}\PYG{p}{;}
\PYG{+w}{  }\PYG{k}{for}\PYG{+w}{ }\PYG{p}{(}\PYG{k}{auto}\PYG{+w}{ }\PYG{k}{const}\PYG{+w}{ }\PYG{o}{\PYGZam{}}\PYG{+w}{ }\PYG{n}{s}\PYG{+w}{ }\PYG{o}{:}\PYG{+w}{ }\PYG{n}{v}\PYG{p}{)}\PYG{+w}{ }\PYG{n}{std}\PYG{o}{:}\PYG{o}{:}\PYG{n}{cout}\PYG{+w}{ }\PYG{o}{\PYGZlt{}}\PYG{o}{\PYGZlt{}}\PYG{+w}{ }\PYG{n}{s}\PYG{+w}{ }\PYG{o}{\PYGZlt{}}\PYG{o}{\PYGZlt{}}\PYG{+w}{ }\PYG{l+s+sc}{\PYGZsq{}}\PYG{l+s+sc}{ }\PYG{l+s+sc}{\PYGZsq{}}\PYG{p}{;}
\PYG{+w}{  }\PYG{n}{std}\PYG{o}{:}\PYG{o}{:}\PYG{n}{cout}\PYG{+w}{ }\PYG{o}{\PYGZlt{}}\PYG{o}{\PYGZlt{}}\PYG{+w}{ }\PYG{n}{std}\PYG{o}{:}\PYG{o}{:}\PYG{n}{endl}\PYG{p}{;}
\PYG{p}{\PYGZcb{}}\PYG{p}{;}
\end{Verbatim}
\end{codeblockbg}


%%C070130 p

本案例使用的区域设置名称（\textit{en\_US.utf8}、\textit{de\_DE.utf8}等）是UNIX系统上的写法。下表列出它们在Windows系统上的对应名称：


%%C070131 table

\myTableBegin
\begin{xltabular}{\textwidth}{@{}>{\hsize=0.336\hsize\raggedright\arraybackslash}X>{\hsize=0.664\hsize\raggedright\arraybackslash}X@{}}
\rowcolor{tabhead} \textcolor{white}{\textbf{UNIX}} & \textcolor{white}{\textbf{Windows}} \\
\endfirsthead
\rowcolor{tabhead} \textcolor{white}{\textbf{UNIX}} & \textcolor{white}{\textbf{Windows}} \\
\endhead
en\_US.utf8 & English\_US.1252 \\
en\_GB.utf8 & English\_UK.1252 \\
de\_DE.utf8 & German\_Germany.1252 \\
sv\_SE.utf8 & Swedish\_Sweden.1252 \\
\end{xltabular}
\myTableEnd


%%C070132 h3

\subsection{具体做法}


%%C070133 p

要控制流的本地化设置，你需要按以下步骤操作：


%%C070134 li

\begin{itemize}
\item 用\texttt{std::\allowbreak{}locale}类表示本地化设置。构造locale对象有多种方式，包括以下几种：


%%C070135 li

\begin{itemize}
\item 默认构造，从而使用全局区域设置（程序启动时默认为\textit{C}区域设置）。


%%C070136 li

\item 由区域设置名称构造，例如\textit{C}、\textit{POSIX}、\textit{en\_US.utf8}等，前提是操作系统支持该名称。


%%C070137 li

\item 由另一个区域设置构造，但替换掉指定的facet。


%%C070138 li

\item 由另一个区域设置构造，但指定类别中的所有facet都从另一个指定的区域设置复制：


%%C070139 code

\begin{codeblockbg}
\begin{Verbatim}[commandchars=\\\{\},breaklines=true,breakanywhere=true,fontsize=\codefont,xleftmargin=17bp]
\PYG{c+c1}{// 默认构造}
\PYG{k}{auto}\PYG{+w}{ }\PYG{n}{loc\PYGZus{}def}\PYG{+w}{ }\PYG{o}{=}\PYG{+w}{ }\PYG{n}{std}\PYG{o}{:}\PYG{o}{:}\PYG{n}{locale}\PYG{+w}{ }\PYG{p}{\PYGZob{}}\PYG{p}{\PYGZcb{}}\PYG{p}{;}
\PYG{c+c1}{// 由名称构造}
\PYG{k}{auto}\PYG{+w}{ }\PYG{n}{loc\PYGZus{}us}\PYG{+w}{ }\PYG{o}{=}\PYG{+w}{ }\PYG{n}{std}\PYG{o}{:}\PYG{o}{:}\PYG{n}{locale}\PYG{+w}{ }\PYG{p}{\PYGZob{}}\PYG{l+s}{\PYGZdq{}}\PYG{l+s}{en\PYGZus{}US.utf8}\PYG{l+s}{\PYGZdq{}}\PYG{p}{\PYGZcb{}}\PYG{p}{;}
\PYG{c+c1}{// 由另一个区域设置构造，但替换一个 facet}
\PYG{k}{auto}\PYG{+w}{ }\PYG{n}{loc1}\PYG{+w}{ }\PYG{o}{=}\PYG{+w}{ }\PYG{n}{std}\PYG{o}{:}\PYG{o}{:}\PYG{n}{locale}\PYG{+w}{ }\PYG{p}{\PYGZob{}}\PYG{n}{loc\PYGZus{}def}\PYG{p}{,}\PYG{+w}{ }\PYG{k}{new}\PYG{+w}{ }\PYG{n}{std}\PYG{o}{:}\PYG{o}{:}\PYG{n}{collate}\PYG{o}{\PYGZlt{}}\PYG{k+kt}{wchar\PYGZus{}t}\PYG{o}{\PYGZgt{}}\PYG{p}{\PYGZcb{}}\PYG{p}{;}
\PYG{c+c1}{// 由另一个区域设置构造，但替换某类别中的 facet}
\PYG{k}{auto}\PYG{+w}{ }\PYG{n}{loc2}\PYG{+w}{ }\PYG{o}{=}\PYG{+w}{ }\PYG{n}{std}\PYG{o}{:}\PYG{o}{:}\PYG{n}{locale}\PYG{+w}{ }\PYG{p}{\PYGZob{}}\PYG{n}{loc\PYGZus{}def}\PYG{p}{,}\PYG{+w}{ }\PYG{n}{loc\PYGZus{}us}\PYG{p}{,}
\PYG{+w}{                         }\PYG{n}{std}\PYG{o}{:}\PYG{o}{:}\PYG{n}{locale}\PYG{o}{:}\PYG{o}{:}\PYG{n}{collate}\PYG{p}{\PYGZcb{}}\PYG{p}{;}
\end{Verbatim}
\end{codeblockbg}


%%C070140 li

\item 要获得默认\textit{C}区域设置的副本，用\myttbr{std::\allowbreak{}locale::\allowbreak{}classic()}静态方法：


%%C070141 code

\begin{codeblockbg}
\begin{Verbatim}[commandchars=\\\{\},breaklines=true,breakanywhere=true,fontsize=\codefont,xleftmargin=17bp]
\PYG{k}{auto}\PYG{+w}{ }\PYG{n}{loc}\PYG{+w}{ }\PYG{o}{=}\PYG{+w}{ }\PYG{n}{std}\PYG{o}{:}\PYG{o}{:}\PYG{n}{locale}\PYG{o}{:}\PYG{o}{:}\PYG{n}{classic}\PYG{p}{(}\PYG{p}{)}\PYG{p}{;}
\end{Verbatim}
\end{codeblockbg}


%%C070142 li

\item 要修改每次默认构造locale时所复制的默认区域设置，用\myttbr{std::\allowbreak{}locale::\allowbreak{}global()}静态方法：


%%C070143 code

\begin{codeblockbg}
\begin{Verbatim}[commandchars=\\\{\},breaklines=true,breakanywhere=true,fontsize=\codefont,xleftmargin=17bp]
\PYG{n}{std}\PYG{o}{:}\PYG{o}{:}\PYG{n}{locale}\PYG{o}{:}\PYG{o}{:}\PYG{n}{global}\PYG{p}{(}\PYG{n}{std}\PYG{o}{:}\PYG{o}{:}\PYG{n}{locale}\PYG{p}{(}\PYG{l+s}{\PYGZdq{}}\PYG{l+s}{en\PYGZus{}US.utf8}\PYG{l+s}{\PYGZdq{}}\PYG{p}{)}\PYG{p}{)}\PYG{p}{;}
\end{Verbatim}
\end{codeblockbg}


%%C070144 li

\item 用\texttt{imbue()}方法修改输入输出流的当前区域设置：


%%C070145 code

\begin{codeblockbg}
\begin{Verbatim}[commandchars=\\\{\},breaklines=true,breakanywhere=true,fontsize=\codefont,xleftmargin=17bp]
\PYG{n}{std}\PYG{o}{:}\PYG{o}{:}\PYG{n}{cout}\PYG{p}{.}\PYG{n}{imbue}\PYG{p}{(}\PYG{n}{std}\PYG{o}{:}\PYG{o}{:}\PYG{n}{locale}\PYG{p}{(}\PYG{l+s}{\PYGZdq{}}\PYG{l+s}{en\PYGZus{}US.utf8}\PYG{l+s}{\PYGZdq{}}\PYG{p}{)}\PYG{p}{)}\PYG{p}{;}
\end{Verbatim}
\end{codeblockbg}

\end{itemize}

\end{itemize}


%%C070146 p

下面的列表给出了使用各种区域设置的示例：


%%C070147 li

\begin{itemize}
\item 使用由名称指定的某个区域设置。本例中用的是德语区域设置：


%%C070148 code

\begin{codeblockbg}
\begin{Verbatim}[commandchars=\\\{\},breaklines=true,breakanywhere=true,fontsize=\codefont,xleftmargin=17bp]
\PYG{k}{auto}\PYG{+w}{ }\PYG{n}{loc}\PYG{+w}{ }\PYG{o}{=}\PYG{+w}{ }\PYG{n}{std}\PYG{o}{:}\PYG{o}{:}\PYG{n}{locale}\PYG{p}{(}\PYG{l+s}{\PYGZdq{}}\PYG{l+s}{de\PYGZus{}DE.utf8}\PYG{l+s}{\PYGZdq{}}\PYG{p}{)}\PYG{p}{;}
\PYG{n}{std}\PYG{o}{:}\PYG{o}{:}\PYG{n}{cout}\PYG{p}{.}\PYG{n}{imbue}\PYG{p}{(}\PYG{n}{loc}\PYG{p}{)}\PYG{p}{;}
\PYG{n}{std}\PYG{o}{:}\PYG{o}{:}\PYG{n}{cout}\PYG{+w}{ }\PYG{o}{\PYGZlt{}}\PYG{o}{\PYGZlt{}}\PYG{+w}{ }\PYG{l+m+mf}{1000.50}\PYG{+w}{ }\PYG{o}{\PYGZlt{}}\PYG{o}{\PYGZlt{}}\PYG{+w}{ }\PYG{n}{std}\PYG{o}{:}\PYG{o}{:}\PYG{n}{endl}\PYG{p}{;}
\PYG{c+c1}{// 1.000,5}
\PYG{n}{std}\PYG{o}{:}\PYG{o}{:}\PYG{n}{cout}\PYG{+w}{ }\PYG{o}{\PYGZlt{}}\PYG{o}{\PYGZlt{}}\PYG{+w}{ }\PYG{n}{std}\PYG{o}{:}\PYG{o}{:}\PYG{n}{showbase}\PYG{+w}{ }\PYG{o}{\PYGZlt{}}\PYG{o}{\PYGZlt{}}\PYG{+w}{ }\PYG{n}{std}\PYG{o}{:}\PYG{o}{:}\PYG{n}{put\PYGZus{}money}\PYG{p}{(}\PYG{l+m+mi}{1050}\PYG{p}{)}
\PYG{+w}{          }\PYG{o}{\PYGZlt{}}\PYG{o}{\PYGZlt{}}\PYG{+w}{ }\PYG{n}{std}\PYG{o}{:}\PYG{o}{:}\PYG{n}{endl}\PYG{p}{;}
\PYG{c+c1}{// 10,50 €}
\PYG{n}{std}\PYG{o}{:}\PYG{o}{:}\PYG{n}{cout}\PYG{+w}{ }\PYG{o}{\PYGZlt{}}\PYG{o}{\PYGZlt{}}\PYG{+w}{ }\PYG{n}{std}\PYG{o}{:}\PYG{o}{:}\PYG{n}{put\PYGZus{}time}\PYG{p}{(}\PYG{n}{ltime}\PYG{p}{,}\PYG{+w}{ }\PYG{l+s}{\PYGZdq{}}\PYG{l+s}{\PYGZpc{}c}\PYG{l+s}{\PYGZdq{}}\PYG{p}{)}\PYG{+w}{ }\PYG{o}{\PYGZlt{}}\PYG{o}{\PYGZlt{}}\PYG{+w}{ }\PYG{n}{std}\PYG{o}{:}\PYG{o}{:}\PYG{n}{endl}\PYG{p}{;}
\PYG{c+c1}{// So 04 Dez 2016 17:54:06 JST}
\PYG{n}{sort\PYGZus{}and\PYGZus{}print}\PYG{p}{(}\PYG{n}{names}\PYG{p}{,}\PYG{+w}{ }\PYG{n}{loc}\PYG{p}{)}\PYG{p}{;}
\PYG{c+c1}{// adele Åke François John Øivind Robert}
\end{Verbatim}
\end{codeblockbg}


%%C070149 li

\item 使用对应用户设置的区域设置（由系统定义）。做法是用空字符串构造一个\texttt{std::\allowbreak{}locale}对象：


%%C070150 code

\begin{codeblockbg}
\begin{Verbatim}[commandchars=\\\{\},breaklines=true,breakanywhere=true,fontsize=\codefont,xleftmargin=17bp]
\PYG{k}{auto}\PYG{+w}{ }\PYG{n}{loc}\PYG{+w}{ }\PYG{o}{=}\PYG{+w}{ }\PYG{n}{std}\PYG{o}{:}\PYG{o}{:}\PYG{n}{locale}\PYG{p}{(}\PYG{l+s}{\PYGZdq{}}\PYG{l+s}{\PYGZdq{}}\PYG{p}{)}\PYG{p}{;}
\PYG{n}{std}\PYG{o}{:}\PYG{o}{:}\PYG{n}{cout}\PYG{p}{.}\PYG{n}{imbue}\PYG{p}{(}\PYG{n}{loc}\PYG{p}{)}\PYG{p}{;}
\PYG{n}{std}\PYG{o}{:}\PYG{o}{:}\PYG{n}{cout}\PYG{+w}{ }\PYG{o}{\PYGZlt{}}\PYG{o}{\PYGZlt{}}\PYG{+w}{ }\PYG{l+m+mf}{1000.50}\PYG{+w}{ }\PYG{o}{\PYGZlt{}}\PYG{o}{\PYGZlt{}}\PYG{+w}{ }\PYG{n}{std}\PYG{o}{:}\PYG{o}{:}\PYG{n}{endl}\PYG{p}{;}
\PYG{c+c1}{// 1,000.5}
\PYG{n}{std}\PYG{o}{:}\PYG{o}{:}\PYG{n}{cout}\PYG{+w}{ }\PYG{o}{\PYGZlt{}}\PYG{o}{\PYGZlt{}}\PYG{+w}{ }\PYG{n}{std}\PYG{o}{:}\PYG{o}{:}\PYG{n}{showbase}\PYG{+w}{ }\PYG{o}{\PYGZlt{}}\PYG{o}{\PYGZlt{}}\PYG{+w}{ }\PYG{n}{std}\PYG{o}{:}\PYG{o}{:}\PYG{n}{put\PYGZus{}money}\PYG{p}{(}\PYG{l+m+mi}{1050}\PYG{p}{)}
\PYG{+w}{          }\PYG{o}{\PYGZlt{}}\PYG{o}{\PYGZlt{}}\PYG{+w}{ }\PYG{n}{std}\PYG{o}{:}\PYG{o}{:}\PYG{n}{endl}\PYG{p}{;}
\PYG{c+c1}{// \PYGZdl{}10.50}
\PYG{n}{std}\PYG{o}{:}\PYG{o}{:}\PYG{n}{cout}\PYG{+w}{ }\PYG{o}{\PYGZlt{}}\PYG{o}{\PYGZlt{}}\PYG{+w}{ }\PYG{n}{std}\PYG{o}{:}\PYG{o}{:}\PYG{n}{put\PYGZus{}time}\PYG{p}{(}\PYG{n}{ltime}\PYG{p}{,}\PYG{+w}{ }\PYG{l+s}{\PYGZdq{}}\PYG{l+s}{\PYGZpc{}c}\PYG{l+s}{\PYGZdq{}}\PYG{p}{)}\PYG{+w}{ }\PYG{o}{\PYGZlt{}}\PYG{o}{\PYGZlt{}}\PYG{+w}{ }\PYG{n}{std}\PYG{o}{:}\PYG{o}{:}\PYG{n}{endl}\PYG{p}{;}
\PYG{c+c1}{// Sun 04 Dec 2016 05:54:06 PM JST}
\PYG{n}{sort\PYGZus{}and\PYGZus{}print}\PYG{p}{(}\PYG{n}{names}\PYG{p}{,}\PYG{+w}{ }\PYG{n}{loc}\PYG{p}{)}\PYG{p}{;}
\PYG{c+c1}{// adele Åke François John Øivind Robert}
\end{Verbatim}
\end{codeblockbg}


%%C070151 li

\item 设置并使用全局区域设置：


%%C070152 code

\begin{codeblockbg}
\begin{Verbatim}[commandchars=\\\{\},breaklines=true,breakanywhere=true,fontsize=\codefont,xleftmargin=17bp]
\PYG{n}{std}\PYG{o}{:}\PYG{o}{:}\PYG{n}{locale}\PYG{o}{:}\PYG{o}{:}\PYG{n}{global}\PYG{p}{(}\PYG{n}{std}\PYG{o}{:}\PYG{o}{:}\PYG{n}{locale}\PYG{p}{(}\PYG{l+s}{\PYGZdq{}}\PYG{l+s}{sv\PYGZus{}SE.utf8}\PYG{l+s}{\PYGZdq{}}\PYG{p}{)}\PYG{p}{)}\PYG{p}{;}\PYG{+w}{ }\PYG{c+c1}{// 设置全局}
\PYG{k}{auto}\PYG{+w}{ }\PYG{n}{loc}\PYG{+w}{ }\PYG{o}{=}\PYG{+w}{ }\PYG{n}{std}\PYG{o}{:}\PYG{o}{:}\PYG{n}{locale}\PYG{p}{\PYGZob{}}\PYG{p}{\PYGZcb{}}\PYG{p}{;}\PYG{+w}{                       }\PYG{c+c1}{// 使用全局}
\PYG{n}{std}\PYG{o}{:}\PYG{o}{:}\PYG{n}{cout}\PYG{p}{.}\PYG{n}{imbue}\PYG{p}{(}\PYG{n}{loc}\PYG{p}{)}\PYG{p}{;}
\PYG{n}{std}\PYG{o}{:}\PYG{o}{:}\PYG{n}{cout}\PYG{+w}{ }\PYG{o}{\PYGZlt{}}\PYG{o}{\PYGZlt{}}\PYG{+w}{ }\PYG{l+m+mf}{1000.50}\PYG{+w}{ }\PYG{o}{\PYGZlt{}}\PYG{o}{\PYGZlt{}}\PYG{+w}{ }\PYG{n}{std}\PYG{o}{:}\PYG{o}{:}\PYG{n}{endl}\PYG{p}{;}
\PYG{c+c1}{// 1 000,5}
\PYG{n}{std}\PYG{o}{:}\PYG{o}{:}\PYG{n}{cout}\PYG{+w}{ }\PYG{o}{\PYGZlt{}}\PYG{o}{\PYGZlt{}}\PYG{+w}{ }\PYG{n}{std}\PYG{o}{:}\PYG{o}{:}\PYG{n}{showbase}\PYG{+w}{ }\PYG{o}{\PYGZlt{}}\PYG{o}{\PYGZlt{}}\PYG{+w}{ }\PYG{n}{std}\PYG{o}{:}\PYG{o}{:}\PYG{n}{put\PYGZus{}money}\PYG{p}{(}\PYG{l+m+mi}{1050}\PYG{p}{)}
\PYG{+w}{          }\PYG{o}{\PYGZlt{}}\PYG{o}{\PYGZlt{}}\PYG{+w}{ }\PYG{n}{std}\PYG{o}{:}\PYG{o}{:}\PYG{n}{endl}\PYG{p}{;}
\PYG{c+c1}{// 10,50 kr}
\PYG{n}{std}\PYG{o}{:}\PYG{o}{:}\PYG{n}{cout}\PYG{+w}{ }\PYG{o}{\PYGZlt{}}\PYG{o}{\PYGZlt{}}\PYG{+w}{ }\PYG{n}{std}\PYG{o}{:}\PYG{o}{:}\PYG{n}{put\PYGZus{}time}\PYG{p}{(}\PYG{n}{ltime}\PYG{p}{,}\PYG{+w}{ }\PYG{l+s}{\PYGZdq{}}\PYG{l+s}{\PYGZpc{}c}\PYG{l+s}{\PYGZdq{}}\PYG{p}{)}\PYG{+w}{ }\PYG{o}{\PYGZlt{}}\PYG{o}{\PYGZlt{}}\PYG{+w}{ }\PYG{n}{std}\PYG{o}{:}\PYG{o}{:}\PYG{n}{endl}\PYG{p}{;}
\PYG{c+c1}{// sön 4 dec 2016 18:02:29}
\PYG{n}{sort\PYGZus{}and\PYGZus{}print}\PYG{p}{(}\PYG{n}{names}\PYG{p}{,}\PYG{+w}{ }\PYG{n}{loc}\PYG{p}{)}\PYG{p}{;}
\PYG{c+c1}{// adele François John Robert Åke Øivind}
\end{Verbatim}
\end{codeblockbg}


%%C070153 li

\item 使用默认的\textit{C}区域设置：


%%C070154 code

\begin{codeblockbg}
\begin{Verbatim}[commandchars=\\\{\},breaklines=true,breakanywhere=true,fontsize=\codefont,xleftmargin=17bp]
\PYG{k}{auto}\PYG{+w}{ }\PYG{n}{loc}\PYG{+w}{ }\PYG{o}{=}\PYG{+w}{ }\PYG{n}{std}\PYG{o}{:}\PYG{o}{:}\PYG{n}{locale}\PYG{o}{:}\PYG{o}{:}\PYG{n}{classic}\PYG{p}{(}\PYG{p}{)}\PYG{p}{;}
\PYG{n}{std}\PYG{o}{:}\PYG{o}{:}\PYG{n}{cout}\PYG{p}{.}\PYG{n}{imbue}\PYG{p}{(}\PYG{n}{loc}\PYG{p}{)}\PYG{p}{;}
\PYG{n}{std}\PYG{o}{:}\PYG{o}{:}\PYG{n}{cout}\PYG{+w}{ }\PYG{o}{\PYGZlt{}}\PYG{o}{\PYGZlt{}}\PYG{+w}{ }\PYG{l+m+mf}{1000.50}\PYG{+w}{ }\PYG{o}{\PYGZlt{}}\PYG{o}{\PYGZlt{}}\PYG{+w}{ }\PYG{n}{std}\PYG{o}{:}\PYG{o}{:}\PYG{n}{endl}\PYG{p}{;}
\PYG{c+c1}{// 1000.5}
\PYG{n}{std}\PYG{o}{:}\PYG{o}{:}\PYG{n}{cout}\PYG{+w}{ }\PYG{o}{\PYGZlt{}}\PYG{o}{\PYGZlt{}}\PYG{+w}{ }\PYG{n}{std}\PYG{o}{:}\PYG{o}{:}\PYG{n}{showbase}\PYG{+w}{ }\PYG{o}{\PYGZlt{}}\PYG{o}{\PYGZlt{}}\PYG{+w}{ }\PYG{n}{std}\PYG{o}{:}\PYG{o}{:}\PYG{n}{put\PYGZus{}money}\PYG{p}{(}\PYG{l+m+mi}{1050}\PYG{p}{)}
\PYG{+w}{          }\PYG{o}{\PYGZlt{}}\PYG{o}{\PYGZlt{}}\PYG{+w}{ }\PYG{n}{std}\PYG{o}{:}\PYG{o}{:}\PYG{n}{endl}\PYG{p}{;}
\PYG{c+c1}{// 1050}
\PYG{n}{std}\PYG{o}{:}\PYG{o}{:}\PYG{n}{cout}\PYG{+w}{ }\PYG{o}{\PYGZlt{}}\PYG{o}{\PYGZlt{}}\PYG{+w}{ }\PYG{n}{std}\PYG{o}{:}\PYG{o}{:}\PYG{n}{put\PYGZus{}time}\PYG{p}{(}\PYG{n}{ltime}\PYG{p}{,}\PYG{+w}{ }\PYG{l+s}{\PYGZdq{}}\PYG{l+s}{\PYGZpc{}c}\PYG{l+s}{\PYGZdq{}}\PYG{p}{)}\PYG{+w}{ }\PYG{o}{\PYGZlt{}}\PYG{o}{\PYGZlt{}}\PYG{+w}{ }\PYG{n}{std}\PYG{o}{:}\PYG{o}{:}\PYG{n}{endl}\PYG{p}{;}
\PYG{c+c1}{// Sun Dec 4 17:55:14 2016}
\PYG{n}{sort\PYGZus{}and\PYGZus{}print}\PYG{p}{(}\PYG{n}{names}\PYG{p}{,}\PYG{+w}{ }\PYG{n}{loc}\PYG{p}{)}\PYG{p}{;}
\PYG{c+c1}{// François John Robert adele Åke Øivind}
\end{Verbatim}
\end{codeblockbg}

\end{itemize}


%%C070155 h3

\subsection{工作原理}


%%C070156 p

区域设置对象实际上并不存储本地化设置。区域设置是facet的异构容器。facet是定义本地化与国际化设置的对象。标准规定了每个区域设置必须包含的facet列表。除此之外，区域设置还可以包含任意用户自定义的facet。下面是标准定义的全部facet列表：


%%C070157 table

\myTableBegin
\begin{xltabular}{\textwidth}{@{}>{\hsize=0.500\hsize\raggedright\arraybackslash}X>{\hsize=0.500\hsize\raggedright\arraybackslash}X@{}}
\rowcolor{tabhead} \textcolor{white}{\myttbr{std::\allowbreak{}collate<\allowbreak{}char>}} & \textcolor{white}{\myttbr{std::\allowbreak{}collate<\allowbreak{}wchar\_\allowbreak{}t>}} \\
\endfirsthead
\rowcolor{tabhead} \textcolor{white}{\myttbr{std::\allowbreak{}collate<\allowbreak{}char>}} & \textcolor{white}{\myttbr{std::\allowbreak{}collate<\allowbreak{}wchar\_\allowbreak{}t>}} \\
\endhead
\myttbr{std::\allowbreak{}ctype<\allowbreak{}char>} & \myttbr{std::\allowbreak{}ctype<\allowbreak{}wchar\_\allowbreak{}t>} \\
\myttbr{std::\allowbreak{}codecvt<\allowbreak{}char,char,mbstate\_\allowbreak{}t>} \myttbr{std::\allowbreak{}codecvt<\allowbreak{}char16\_\allowbreak{}t,char,mbstate\_\allowbreak{}t>} & \myttbr{std::\allowbreak{}codecvt<\allowbreak{}char32\_\allowbreak{}t,char,mbstate\_\allowbreak{}t>} \myttbr{std::\allowbreak{}codecvt<\allowbreak{}wchar\_\allowbreak{}t,char,mbstate\_\allowbreak{}t>} \\
\myttbr{std::\allowbreak{}moneypunct<\allowbreak{}char>} \myttbr{std::\allowbreak{}moneypunct<\allowbreak{}char,true>} & \myttbr{std::\allowbreak{}moneypunct<\allowbreak{}wchar\_\allowbreak{}t>} \myttbr{std::\allowbreak{}moneypunct<\allowbreak{}wchar\_\allowbreak{}t,true>} \\
\myttbr{std::\allowbreak{}money\_\allowbreak{}get<\allowbreak{}char>} & \myttbr{std::\allowbreak{}money\_\allowbreak{}get<\allowbreak{}wchar\_\allowbreak{}t>} \\
\myttbr{std::\allowbreak{}money\_\allowbreak{}put<\allowbreak{}char>} & \myttbr{std::\allowbreak{}money\_\allowbreak{}put<\allowbreak{}wchar\_\allowbreak{}t>} \\
\myttbr{std::\allowbreak{}numpunct<\allowbreak{}char>} & \myttbr{std::\allowbreak{}numpunct<\allowbreak{}wchar\_\allowbreak{}t>} \\
\myttbr{std::\allowbreak{}num\_\allowbreak{}get<\allowbreak{}char>} & \myttbr{std::\allowbreak{}num\_\allowbreak{}get<\allowbreak{}wchar\_\allowbreak{}t>} \\
\myttbr{std::\allowbreak{}num\_\allowbreak{}put<\allowbreak{}char>} & \myttbr{std::\allowbreak{}num\_\allowbreak{}put<\allowbreak{}wchar\_\allowbreak{}t>} \\
\myttbr{std::\allowbreak{}time\_\allowbreak{}get<\allowbreak{}char>} & \myttbr{std::\allowbreak{}time\_\allowbreak{}get<\allowbreak{}wchar\_\allowbreak{}t>} \\
\myttbr{std::\allowbreak{}time\_\allowbreak{}put<\allowbreak{}char>} & \myttbr{std::\allowbreak{}time\_\allowbreak{}put<\allowbreak{}wchar\_\allowbreak{}t>} \\
\myttbr{std::\allowbreak{}messages<\allowbreak{}char>} & \myttbr{std::\allowbreak{}messages<\allowbreak{}wchar\_\allowbreak{}t>} \\
\end{xltabular}
\myTableEnd


%%C070158 p

区域设置是不可变对象，其中保存的facet对象也是不可变的。区域设置的实现是一个带引用计数的数组，数组元素是指向facet的带引用计数的指针。该数组以\texttt{std::\allowbreak{}locale::\allowbreak{}id}为索引；所有facet都必须派生自基类\texttt{std::\allowbreak{}locale::\allowbreak{}facet}，并且必须有一个公有静态成员，其类型为\texttt{std::\allowbreak{}locale::\allowbreak{}id}，名字为\texttt{id}。


%%C070159 p

创建区域设置对象只能通过某个重载构造函数，或者通过\texttt{combine()}方法——顾名思义，该方法把当前区域设置与一个新的、编译期可识别的facet组合起来，返回一个新的区域设置对象。另一方面，可以用\texttt{std::\allowbreak{}has\_\allowbreak{}facet()}函数模板判断某个区域设置是否包含特定的facet，也可以用\texttt{std::\allowbreak{}use\_\allowbreak{}facet()}函数模板取得对某个区域设置所实现的facet的引用。


%%C070160 p

前面的示例中，我们对一个字符串向量排序，并把区域设置对象作为第三个实参传给\texttt{std::\allowbreak{}sort()}这个通用算法。这第三个实参本应是比较函数对象。之所以能传区域设置对象，是因为\texttt{std::\allowbreak{}locale}有一个\texttt{operator()}，它用自己的collate facet按字典序比较两个字符串。这其实是\texttt{std::\allowbreak{}locale}直接提供的唯一本地化功能；不过它内部所做的是调用collate facet的\texttt{compare()}方法，由该方法按facet的规则完成字符串比较。


%%C070161 p

每个程序在启动时都会创建一个全局区域设置。这个全局区域设置的内容会被复制到每个默认构造的区域设置中。可以用静态方法\myttbr{std::\allowbreak{}locale::\allowbreak{}global()}替换全局区域设置。默认情况下，全局区域设置是\textit{C}区域设置，它等同于ANCI C中同名的区域设置。该区域设置是为处理简单的英文文本而准备的，C++把它作为默认设置，以提供与C的兼容性。可以用静态方法\myttbr{std::\allowbreak{}locale::\allowbreak{}classic()}获得对它的引用。


%%C070162 p

默认情况下，所有流都用经典区域设置来写出或解析文本。不过，可以用流的\texttt{imbue()}方法修改流所使用的区域设置。它是\texttt{std::\allowbreak{}ios\_\allowbreak{}base}类的成员，而这个类是所有输入输出流的基类。与之配套的成员是\texttt{getloc()}方法，它返回当前流区域设置的副本。


%%C070163 note

\begin{myNote}
前面的示例中，我们修改的是\texttt{std::\allowbreak{}cout}流对象的区域设置。在实践中，你可能希望为标准C流关联的所有流对象设置同一个区域设置：\texttt{cin}、\texttt{cout}、\texttt{cerr}和\texttt{clog}（或\texttt{wcin}、\texttt{wcout}、\texttt{wcerr}和\texttt{wclog}）。
\end{myNote}


%%C070164 h3

\subsection{延伸阅读}


%%C070165 li

\begin{itemize}
\item \textit{用I/O操纵器控制流的输出}


%%C070166 li

\item \textit{使用货币I/O操纵器}


%%C070167 li

\item \textit{使用时间I/O操纵器}

\end{itemize}


%%C070168 h2

\section{用I/O操纵器控制流的输出}


%%C070169 p

除了基于流的输入输出库之外，标准库还提供了一系列辅助函数，称为操纵器（manipulator），它们通过\texttt{operator<<}与\texttt{operator>>}控制输入流或输出流。本案例将介绍其中一些操纵器，并用几个格式化控制台输出的示例演示它们的用法；后续案例中还会继续介绍更多操纵器。


%%C070170 h3

\subsection{准备工作}


%%C070171 p

I/O操纵器位于\texttt{std}命名空间，分别在头文件\texttt{<\allowbreak{}ios>}、\texttt{<\allowbreak{}istream>}、\texttt{<\allowbreak{}ostream>}和\texttt{<\allowbreak{}iomanip>}中声明。本案例只讨论\texttt{<\allowbreak{}ios>}与\texttt{<\allowbreak{}iomanip>}中的部分操纵器。


%%C070172 h3

\subsection{具体做法}


%%C070173 p

下列操纵器可用于控制流的输出或输入：


%%C070174 li

\begin{itemize}
\item \texttt{boolalpha}与\texttt{noboolalpha}启用和禁用布尔值的文本表示：


%%C070175 code

\begin{codeblockbg}
\begin{Verbatim}[commandchars=\\\{\},breaklines=true,breakanywhere=true,fontsize=\codefont,xleftmargin=17bp]
\PYG{n}{std}\PYG{o}{:}\PYG{o}{:}\PYG{n}{cout}\PYG{+w}{ }\PYG{o}{\PYGZlt{}}\PYG{o}{\PYGZlt{}}\PYG{+w}{ }\PYG{n}{std}\PYG{o}{:}\PYG{o}{:}\PYG{n}{boolalpha}\PYG{+w}{ }\PYG{o}{\PYGZlt{}}\PYG{o}{\PYGZlt{}}\PYG{+w}{ }\PYG{n+nb}{true}\PYG{+w}{ }\PYG{o}{\PYGZlt{}}\PYG{o}{\PYGZlt{}}\PYG{+w}{ }\PYG{n}{std}\PYG{o}{:}\PYG{o}{:}\PYG{n}{endl}\PYG{p}{;}\PYG{+w}{    }\PYG{c+c1}{// true}
\PYG{n}{std}\PYG{o}{:}\PYG{o}{:}\PYG{n}{cout}\PYG{+w}{ }\PYG{o}{\PYGZlt{}}\PYG{o}{\PYGZlt{}}\PYG{+w}{ }\PYG{n+nb}{false}\PYG{+w}{ }\PYG{o}{\PYGZlt{}}\PYG{o}{\PYGZlt{}}\PYG{+w}{ }\PYG{n}{std}\PYG{o}{:}\PYG{o}{:}\PYG{n}{endl}\PYG{p}{;}\PYG{+w}{                     }\PYG{c+c1}{// false}
\PYG{n}{std}\PYG{o}{:}\PYG{o}{:}\PYG{n}{cout}\PYG{+w}{ }\PYG{o}{\PYGZlt{}}\PYG{o}{\PYGZlt{}}\PYG{+w}{ }\PYG{n}{std}\PYG{o}{:}\PYG{o}{:}\PYG{n}{noboolalpha}\PYG{+w}{ }\PYG{o}{\PYGZlt{}}\PYG{o}{\PYGZlt{}}\PYG{+w}{ }\PYG{n+nb}{false}\PYG{+w}{ }\PYG{o}{\PYGZlt{}}\PYG{o}{\PYGZlt{}}\PYG{+w}{ }\PYG{n}{std}\PYG{o}{:}\PYG{o}{:}\PYG{n}{endl}\PYG{p}{;}\PYG{+w}{ }\PYG{c+c1}{// 0}
\end{Verbatim}
\end{codeblockbg}

%%C070176 li

\item \texttt{left}、\texttt{right}和\texttt{internal}影响填充字符的对齐方式；\texttt{left}和\texttt{right}作用于所有文本，但\texttt{internal}只作用于整数、浮点数和货币输出：


%%C070177 code

\begin{codeblockbg}
\begin{Verbatim}[commandchars=\\\{\},breaklines=true,breakanywhere=true,fontsize=\codefont,xleftmargin=17bp]
\PYG{n}{std}\PYG{o}{:}\PYG{o}{:}\PYG{n}{cout}\PYG{+w}{ }\PYG{o}{\PYGZlt{}}\PYG{o}{\PYGZlt{}}\PYG{+w}{ }\PYG{n}{std}\PYG{o}{:}\PYG{o}{:}\PYG{n}{right}\PYG{+w}{ }\PYG{o}{\PYGZlt{}}\PYG{o}{\PYGZlt{}}\PYG{+w}{ }\PYG{n}{std}\PYG{o}{:}\PYG{o}{:}\PYG{n}{setw}\PYG{p}{(}\PYG{l+m+mi}{10}\PYG{p}{)}\PYG{+w}{ }\PYG{o}{\PYGZlt{}}\PYG{o}{\PYGZlt{}}\PYG{+w}{ }\PYG{l+s}{\PYGZdq{}}\PYG{l+s}{right}\PYG{l+s}{\PYGZdq{}}
\PYG{+w}{          }\PYG{o}{\PYGZlt{}}\PYG{o}{\PYGZlt{}}\PYG{+w}{ }\PYG{n}{std}\PYG{o}{:}\PYG{o}{:}\PYG{n}{endl}\PYG{p}{;}
\PYG{n}{std}\PYG{o}{:}\PYG{o}{:}\PYG{n}{cout}\PYG{+w}{ }\PYG{o}{\PYGZlt{}}\PYG{o}{\PYGZlt{}}\PYG{+w}{ }\PYG{n}{std}\PYG{o}{:}\PYG{o}{:}\PYG{n}{setw}\PYG{p}{(}\PYG{l+m+mi}{10}\PYG{p}{)}\PYG{+w}{ }\PYG{o}{\PYGZlt{}}\PYG{o}{\PYGZlt{}}\PYG{+w}{ }\PYG{l+s}{\PYGZdq{}}\PYG{l+s}{text}\PYG{l+s}{\PYGZdq{}}\PYG{+w}{ }\PYG{o}{\PYGZlt{}}\PYG{o}{\PYGZlt{}}\PYG{+w}{ }\PYG{n}{std}\PYG{o}{:}\PYG{o}{:}\PYG{n}{endl}\PYG{p}{;}
\PYG{n}{std}\PYG{o}{:}\PYG{o}{:}\PYG{n}{cout}\PYG{+w}{ }\PYG{o}{\PYGZlt{}}\PYG{o}{\PYGZlt{}}\PYG{+w}{ }\PYG{n}{std}\PYG{o}{:}\PYG{o}{:}\PYG{n}{left}\PYG{+w}{ }\PYG{o}{\PYGZlt{}}\PYG{o}{\PYGZlt{}}\PYG{+w}{ }\PYG{n}{std}\PYG{o}{:}\PYG{o}{:}\PYG{n}{setw}\PYG{p}{(}\PYG{l+m+mi}{10}\PYG{p}{)}\PYG{+w}{ }\PYG{o}{\PYGZlt{}}\PYG{o}{\PYGZlt{}}\PYG{+w}{ }\PYG{l+s}{\PYGZdq{}}\PYG{l+s}{left}\PYG{l+s}{\PYGZdq{}}\PYG{+w}{ }\PYG{o}{\PYGZlt{}}\PYG{o}{\PYGZlt{}}\PYG{+w}{ }\PYG{n}{std}\PYG{o}{:}\PYG{o}{:}\PYG{n}{endl}\PYG{p}{;}
\end{Verbatim}
\end{codeblockbg}


%%C070178 li

\item \texttt{fixed}、\texttt{scientific}、\texttt{hexfloat}和\texttt{defaultfloat}改变浮点类型所用的格式（输入流和输出流都适用）。后两个从C++11起才可用：


%%C070179 code

\begin{codeblockbg}
\begin{Verbatim}[commandchars=\\\{\},breaklines=true,breakanywhere=true,fontsize=\codefont,xleftmargin=17bp]
\PYG{n}{std}\PYG{o}{:}\PYG{o}{:}\PYG{n}{cout}\PYG{+w}{ }\PYG{o}{\PYGZlt{}}\PYG{o}{\PYGZlt{}}\PYG{+w}{ }\PYG{n}{std}\PYG{o}{:}\PYG{o}{:}\PYG{n}{fixed}\PYG{+w}{ }\PYG{o}{\PYGZlt{}}\PYG{o}{\PYGZlt{}}\PYG{+w}{ }\PYG{l+m+mf}{0.25}\PYG{+w}{ }\PYG{o}{\PYGZlt{}}\PYG{o}{\PYGZlt{}}\PYG{+w}{ }\PYG{n}{std}\PYG{o}{:}\PYG{o}{:}\PYG{n}{endl}\PYG{p}{;}
\PYG{c+c1}{// 0.250000}
\PYG{n}{std}\PYG{o}{:}\PYG{o}{:}\PYG{n}{cout}\PYG{+w}{ }\PYG{o}{\PYGZlt{}}\PYG{o}{\PYGZlt{}}\PYG{+w}{ }\PYG{n}{std}\PYG{o}{:}\PYG{o}{:}\PYG{n}{scientific}\PYG{+w}{ }\PYG{o}{\PYGZlt{}}\PYG{o}{\PYGZlt{}}\PYG{+w}{ }\PYG{l+m+mf}{0.25}\PYG{+w}{ }\PYG{o}{\PYGZlt{}}\PYG{o}{\PYGZlt{}}\PYG{+w}{ }\PYG{n}{std}\PYG{o}{:}\PYG{o}{:}\PYG{n}{endl}\PYG{p}{;}
\PYG{c+c1}{// 2.500000e\PYGZhy{}01}
\PYG{n}{std}\PYG{o}{:}\PYG{o}{:}\PYG{n}{cout}\PYG{+w}{ }\PYG{o}{\PYGZlt{}}\PYG{o}{\PYGZlt{}}\PYG{+w}{ }\PYG{n}{std}\PYG{o}{:}\PYG{o}{:}\PYG{n}{hexfloat}\PYG{+w}{ }\PYG{o}{\PYGZlt{}}\PYG{o}{\PYGZlt{}}\PYG{+w}{ }\PYG{l+m+mf}{0.25}\PYG{+w}{ }\PYG{o}{\PYGZlt{}}\PYG{o}{\PYGZlt{}}\PYG{+w}{ }\PYG{n}{std}\PYG{o}{:}\PYG{o}{:}\PYG{n}{endl}\PYG{p}{;}
\PYG{c+c1}{// 0x1p\PYGZhy{}2}
\PYG{n}{std}\PYG{o}{:}\PYG{o}{:}\PYG{n}{cout}\PYG{+w}{ }\PYG{o}{\PYGZlt{}}\PYG{o}{\PYGZlt{}}\PYG{+w}{ }\PYG{n}{std}\PYG{o}{:}\PYG{o}{:}\PYG{n}{defaultfloat}\PYG{+w}{ }\PYG{o}{\PYGZlt{}}\PYG{o}{\PYGZlt{}}\PYG{+w}{ }\PYG{l+m+mf}{0.25}\PYG{+w}{ }\PYG{o}{\PYGZlt{}}\PYG{o}{\PYGZlt{}}\PYG{+w}{ }\PYG{n}{std}\PYG{o}{:}\PYG{o}{:}\PYG{n}{endl}\PYG{p}{;}
\PYG{c+c1}{// 0.25}
\end{Verbatim}
\end{codeblockbg}


%%C070180 li

\item \texttt{dec}、\texttt{hex}和\texttt{oct}控制整数类型所用的进制（输入流和输出流都适用）：


%%C070181 code

\begin{codeblockbg}
\begin{Verbatim}[commandchars=\\\{\},breaklines=true,breakanywhere=true,fontsize=\codefont,xleftmargin=17bp]
\PYG{n}{std}\PYG{o}{:}\PYG{o}{:}\PYG{n}{cout}\PYG{+w}{ }\PYG{o}{\PYGZlt{}}\PYG{o}{\PYGZlt{}}\PYG{+w}{ }\PYG{n}{std}\PYG{o}{:}\PYG{o}{:}\PYG{n}{oct}\PYG{+w}{ }\PYG{o}{\PYGZlt{}}\PYG{o}{\PYGZlt{}}\PYG{+w}{ }\PYG{l+m+mi}{42}\PYG{+w}{ }\PYG{o}{\PYGZlt{}}\PYG{o}{\PYGZlt{}}\PYG{+w}{ }\PYG{n}{std}\PYG{o}{:}\PYG{o}{:}\PYG{n}{endl}\PYG{p}{;}\PYG{+w}{ }\PYG{c+c1}{// 52}
\PYG{n}{std}\PYG{o}{:}\PYG{o}{:}\PYG{n}{cout}\PYG{+w}{ }\PYG{o}{\PYGZlt{}}\PYG{o}{\PYGZlt{}}\PYG{+w}{ }\PYG{n}{std}\PYG{o}{:}\PYG{o}{:}\PYG{n}{hex}\PYG{+w}{ }\PYG{o}{\PYGZlt{}}\PYG{o}{\PYGZlt{}}\PYG{+w}{ }\PYG{l+m+mi}{42}\PYG{+w}{ }\PYG{o}{\PYGZlt{}}\PYG{o}{\PYGZlt{}}\PYG{+w}{ }\PYG{n}{std}\PYG{o}{:}\PYG{o}{:}\PYG{n}{endl}\PYG{p}{;}\PYG{+w}{ }\PYG{c+c1}{// 2a}
\PYG{n}{std}\PYG{o}{:}\PYG{o}{:}\PYG{n}{cout}\PYG{+w}{ }\PYG{o}{\PYGZlt{}}\PYG{o}{\PYGZlt{}}\PYG{+w}{ }\PYG{n}{std}\PYG{o}{:}\PYG{o}{:}\PYG{n}{dec}\PYG{+w}{ }\PYG{o}{\PYGZlt{}}\PYG{o}{\PYGZlt{}}\PYG{+w}{ }\PYG{l+m+mi}{42}\PYG{+w}{ }\PYG{o}{\PYGZlt{}}\PYG{o}{\PYGZlt{}}\PYG{+w}{ }\PYG{n}{std}\PYG{o}{:}\PYG{o}{:}\PYG{n}{endl}\PYG{p}{;}\PYG{+w}{ }\PYG{c+c1}{// 42}
\end{Verbatim}
\end{codeblockbg}


%%C070182 li

\item \texttt{setw}改变下一个输入或输出字段的宽度。默认宽度为0。


%%C070183 li

\item \texttt{setfill}改变输出流的填充字符；该字符用于填充后续字段，直到达到指定宽度。默认填充字符为空白：


%%C070184 code

\begin{codeblockbg}
\begin{Verbatim}[commandchars=\\\{\},breaklines=true,breakanywhere=true,fontsize=\codefont,xleftmargin=17bp]
\PYG{n}{std}\PYG{o}{:}\PYG{o}{:}\PYG{n}{cout}\PYG{+w}{ }\PYG{o}{\PYGZlt{}}\PYG{o}{\PYGZlt{}}\PYG{+w}{ }\PYG{n}{std}\PYG{o}{:}\PYG{o}{:}\PYG{n}{right}
\PYG{+w}{          }\PYG{o}{\PYGZlt{}}\PYG{o}{\PYGZlt{}}\PYG{+w}{ }\PYG{n}{std}\PYG{o}{:}\PYG{o}{:}\PYG{n}{setfill}\PYG{p}{(}\PYG{l+s+sc}{\PYGZsq{}}\PYG{l+s+sc}{.}\PYG{l+s+sc}{\PYGZsq{}}\PYG{p}{)}\PYG{+w}{ }\PYG{o}{\PYGZlt{}}\PYG{o}{\PYGZlt{}}\PYG{+w}{ }\PYG{n}{std}\PYG{o}{:}\PYG{o}{:}\PYG{n}{setw}\PYG{p}{(}\PYG{l+m+mi}{10}\PYG{p}{)}
\PYG{+w}{          }\PYG{o}{\PYGZlt{}}\PYG{o}{\PYGZlt{}}\PYG{+w}{ }\PYG{l+s}{\PYGZdq{}}\PYG{l+s}{right}\PYG{l+s}{\PYGZdq{}}\PYG{+w}{ }\PYG{o}{\PYGZlt{}}\PYG{o}{\PYGZlt{}}\PYG{+w}{ }\PYG{n}{std}\PYG{o}{:}\PYG{o}{:}\PYG{n}{endl}\PYG{p}{;}
\PYG{c+c1}{// .....right}
\end{Verbatim}
\end{codeblockbg}


%%C070185 li

\item \texttt{setprecision}改变输入流和输出流中浮点类型的小数精度（即生成多少位数字）。默认精度为6：


%%C070186 code

\begin{codeblockbg}
\begin{Verbatim}[commandchars=\\\{\},breaklines=true,breakanywhere=true,fontsize=\codefont,xleftmargin=17bp]
\PYG{n}{std}\PYG{o}{:}\PYG{o}{:}\PYG{n}{cout}\PYG{+w}{ }\PYG{o}{\PYGZlt{}}\PYG{o}{\PYGZlt{}}\PYG{+w}{ }\PYG{n}{std}\PYG{o}{:}\PYG{o}{:}\PYG{n}{fixed}\PYG{+w}{ }\PYG{o}{\PYGZlt{}}\PYG{o}{\PYGZlt{}}\PYG{+w}{ }\PYG{n}{std}\PYG{o}{:}\PYG{o}{:}\PYG{n}{setprecision}\PYG{p}{(}\PYG{l+m+mi}{2}\PYG{p}{)}\PYG{+w}{ }\PYG{o}{\PYGZlt{}}\PYG{o}{\PYGZlt{}}\PYG{+w}{ }\PYG{l+m+mf}{12.345}
\PYG{+w}{          }\PYG{o}{\PYGZlt{}}\PYG{o}{\PYGZlt{}}\PYG{+w}{ }\PYG{n}{std}\PYG{o}{:}\PYG{o}{:}\PYG{n}{endl}\PYG{p}{;}
\PYG{c+c1}{// 12.35}
\end{Verbatim}
\end{codeblockbg}

\end{itemize}


%%C070187 h3

\subsection{工作原理}


%%C070188 p

上面列出的所有I/O操纵器（除只作用于下一个输出字段的\texttt{setw}之外）都会影响流，此后所有连续的写入或读取操作都沿用最后指定的格式，直到再次使用另一个操纵器。


%%C070189 p

其中一些操纵器不带参数调用，例如\texttt{boolalpha}/\texttt{noboolalpha}或\texttt{dec}/\texttt{hex}/\texttt{oct}。这些操纵器是只接受一个参数的函数，该参数是对字符串的引用，函数返回对同一流的引用：


%%C070190 code

\begin{codeblockbg}
\begin{Verbatim}[commandchars=\\\{\},breaklines=true,breakanywhere=true,fontsize=\codefont,xleftmargin=17bp]
\PYG{n}{std}\PYG{o}{:}\PYG{o}{:}\PYG{n}{ios\PYGZus{}base}\PYG{o}{\PYGZam{}}\PYG{+w}{ }\PYG{n+nf}{hex}\PYG{p}{(}\PYG{n}{std}\PYG{o}{:}\PYG{o}{:}\PYG{n}{ios\PYGZus{}base}\PYG{o}{\PYGZam{}}\PYG{+w}{ }\PYG{n}{str}\PYG{p}{)}\PYG{p}{;}
\end{Verbatim}
\end{codeblockbg}


%%C070191 p

像\myttbr{std::\allowbreak{}cout << std::\allowbreak{}hex}这样的表达式之所以可行，是因为\myttbr{basic\_\allowbreak{}ostream::\allowbreak{}operator<<}和\myttbr{basic\_\allowbreak{}istream::\allowbreak{}operator>>}都有接受这类函数指针的特殊重载。


%%C070192 p

另一些操纵器（包括这里没有提到的一些）要带参数调用。它们是接受一个或多个参数、并返回未指明类型对象的函数：


%%C070193 code

\begin{codeblockbg}
\begin{Verbatim}[commandchars=\\\{\},breaklines=true,breakanywhere=true,fontsize=\codefont,xleftmargin=17bp]
\PYG{k}{template}\PYG{o}{\PYGZlt{}}\PYG{k}{class}\PYG{+w}{ }\PYG{n+nc}{CharT}\PYG{o}{\PYGZgt{}}
\PYG{c+cm}{/*未指明*/}\PYG{+w}{ }\PYG{n}{setfill}\PYG{p}{(}\PYG{n}{CharT}\PYG{+w}{ }\PYG{n}{c}\PYG{p}{)}\PYG{p}{;}
\end{Verbatim}
\end{codeblockbg}


%%C070194 p

为了更清楚地说明这些操纵器的用法，我们来看两个把内容格式化输出到控制台的例子。


%%C070195 p

第一个例子中，我们要按以下要求列出某本书的目录：


%%C070196 li

\begin{itemize}
\item 章节号右对齐，并用罗马数字表示。


%%C070197 li

\item 章节标题左对齐，标题到页码之间的剩余空间用点号填满。


%%C070198 li

\item 章节的起始页码右对齐。

\end{itemize}


%%C070199 p

这个例子中，我们使用下面这些类和辅助函数：


%%C070200 code

\begin{codeblockbg}
\begin{Verbatim}[commandchars=\\\{\},breaklines=true,breakanywhere=true,fontsize=\codefont,xleftmargin=17bp]
\PYG{k}{struct}\PYG{+w}{ }\PYG{n+nc}{Chapter}
\PYG{p}{\PYGZob{}}
\PYG{+w}{  }\PYG{k+kt}{int}\PYG{+w}{ }\PYG{n}{Number}\PYG{p}{;}
\PYG{+w}{  }\PYG{n}{std}\PYG{o}{:}\PYG{o}{:}\PYG{n}{string}\PYG{+w}{ }\PYG{n}{Title}\PYG{p}{;}
\PYG{+w}{  }\PYG{k+kt}{int}\PYG{+w}{ }\PYG{n}{Page}\PYG{p}{;}
\PYG{p}{\PYGZcb{}}\PYG{p}{;}
\PYG{k}{struct}\PYG{+w}{ }\PYG{n+nc}{BookPart}
\PYG{p}{\PYGZob{}}
\PYG{+w}{  }\PYG{n}{std}\PYG{o}{:}\PYG{o}{:}\PYG{n}{string}\PYG{+w}{ }\PYG{n}{Title}\PYG{p}{;}
\PYG{+w}{  }\PYG{n}{std}\PYG{o}{:}\PYG{o}{:}\PYG{n}{vector}\PYG{o}{\PYGZlt{}}\PYG{n}{Chapter}\PYG{o}{\PYGZgt{}}\PYG{+w}{ }\PYG{n}{Chapters}\PYG{p}{;}
\PYG{p}{\PYGZcb{}}\PYG{p}{;}
\PYG{k}{struct}\PYG{+w}{ }\PYG{n+nc}{Book}
\PYG{p}{\PYGZob{}}
\PYG{+w}{  }\PYG{n}{std}\PYG{o}{:}\PYG{o}{:}\PYG{n}{string}\PYG{+w}{ }\PYG{n}{Title}\PYG{p}{;}
\PYG{+w}{  }\PYG{n}{std}\PYG{o}{:}\PYG{o}{:}\PYG{n}{vector}\PYG{o}{\PYGZlt{}}\PYG{n}{BookPart}\PYG{o}{\PYGZgt{}}\PYG{+w}{ }\PYG{n}{Parts}\PYG{p}{;}
\PYG{p}{\PYGZcb{}}\PYG{p}{;}
\PYG{n}{std}\PYG{o}{:}\PYG{o}{:}\PYG{n}{string}\PYG{+w}{ }\PYG{n+nf}{to\PYGZus{}roman}\PYG{p}{(}\PYG{k+kt}{unsigned}\PYG{+w}{ }\PYG{k+kt}{int}\PYG{+w}{ }\PYG{n}{value}\PYG{p}{)}
\PYG{p}{\PYGZob{}}
\PYG{+w}{  }\PYG{k}{struct}\PYG{+w}{ }\PYG{n+nc}{roman\PYGZus{}t}\PYG{+w}{ }\PYG{p}{\PYGZob{}}\PYG{+w}{ }\PYG{k+kt}{unsigned}\PYG{+w}{ }\PYG{k+kt}{int}\PYG{+w}{ }\PYG{n}{value}\PYG{p}{;}\PYG{+w}{ }\PYG{k+kt}{char}\PYG{+w}{ }\PYG{k}{const}\PYG{o}{*}\PYG{+w}{ }\PYG{n}{numeral}\PYG{p}{;}\PYG{+w}{ }\PYG{p}{\PYGZcb{}}\PYG{p}{;}
\PYG{+w}{  }\PYG{k}{const}\PYG{+w}{ }\PYG{k}{static}\PYG{+w}{ }\PYG{n}{roman\PYGZus{}t}\PYG{+w}{ }\PYG{n}{rarr}\PYG{p}{[}\PYG{l+m+mi}{13}\PYG{p}{]}\PYG{+w}{ }\PYG{o}{=}
\PYG{+w}{  }\PYG{p}{\PYGZob{}}
\PYG{+w}{    }\PYG{p}{\PYGZob{}}\PYG{l+m+mi}{1000}\PYG{p}{,}\PYG{+w}{ }\PYG{l+s}{\PYGZdq{}}\PYG{l+s}{M}\PYG{l+s}{\PYGZdq{}}\PYG{p}{\PYGZcb{}}\PYG{p}{,}\PYG{+w}{ }\PYG{p}{\PYGZob{}}\PYG{l+m+mi}{900}\PYG{p}{,}\PYG{+w}{ }\PYG{l+s}{\PYGZdq{}}\PYG{l+s}{CM}\PYG{l+s}{\PYGZdq{}}\PYG{p}{\PYGZcb{}}\PYG{p}{,}\PYG{+w}{ }\PYG{p}{\PYGZob{}}\PYG{l+m+mi}{500}\PYG{p}{,}\PYG{+w}{ }\PYG{l+s}{\PYGZdq{}}\PYG{l+s}{D}\PYG{l+s}{\PYGZdq{}}\PYG{p}{\PYGZcb{}}\PYG{p}{,}\PYG{+w}{ }\PYG{p}{\PYGZob{}}\PYG{l+m+mi}{400}\PYG{p}{,}\PYG{+w}{ }\PYG{l+s}{\PYGZdq{}}\PYG{l+s}{CD}\PYG{l+s}{\PYGZdq{}}\PYG{p}{\PYGZcb{}}\PYG{p}{,}
\PYG{+w}{    }\PYG{p}{\PYGZob{}}\PYG{l+m+mi}{100}\PYG{p}{,}\PYG{+w}{ }\PYG{l+s}{\PYGZdq{}}\PYG{l+s}{C}\PYG{l+s}{\PYGZdq{}}\PYG{p}{\PYGZcb{}}\PYG{p}{,}\PYG{+w}{ }\PYG{p}{\PYGZob{}}\PYG{+w}{ }\PYG{l+m+mi}{90}\PYG{p}{,}\PYG{+w}{ }\PYG{l+s}{\PYGZdq{}}\PYG{l+s}{XC}\PYG{l+s}{\PYGZdq{}}\PYG{p}{\PYGZcb{}}\PYG{p}{,}\PYG{+w}{ }\PYG{p}{\PYGZob{}}\PYG{+w}{ }\PYG{l+m+mi}{50}\PYG{p}{,}\PYG{+w}{ }\PYG{l+s}{\PYGZdq{}}\PYG{l+s}{L}\PYG{l+s}{\PYGZdq{}}\PYG{p}{\PYGZcb{}}\PYG{p}{,}\PYG{+w}{ }\PYG{p}{\PYGZob{}}\PYG{+w}{ }\PYG{l+m+mi}{40}\PYG{p}{,}\PYG{+w}{ }\PYG{l+s}{\PYGZdq{}}\PYG{l+s}{XL}\PYG{l+s}{\PYGZdq{}}\PYG{p}{\PYGZcb{}}\PYG{p}{,}
\PYG{+w}{    }\PYG{p}{\PYGZob{}}\PYG{+w}{ }\PYG{l+m+mi}{10}\PYG{p}{,}\PYG{+w}{ }\PYG{l+s}{\PYGZdq{}}\PYG{l+s}{X}\PYG{l+s}{\PYGZdq{}}\PYG{p}{\PYGZcb{}}\PYG{p}{,}\PYG{+w}{ }\PYG{p}{\PYGZob{}}\PYG{+w}{ }\PYG{l+m+mi}{9}\PYG{p}{,}\PYG{+w}{ }\PYG{l+s}{\PYGZdq{}}\PYG{l+s}{IX}\PYG{l+s}{\PYGZdq{}}\PYG{p}{\PYGZcb{}}\PYG{p}{,}\PYG{+w}{ }\PYG{p}{\PYGZob{}}\PYG{+w}{ }\PYG{l+m+mi}{5}\PYG{p}{,}\PYG{+w}{ }\PYG{l+s}{\PYGZdq{}}\PYG{l+s}{V}\PYG{l+s}{\PYGZdq{}}\PYG{p}{\PYGZcb{}}\PYG{p}{,}\PYG{+w}{ }\PYG{p}{\PYGZob{}}\PYG{+w}{ }\PYG{l+m+mi}{4}\PYG{p}{,}\PYG{+w}{ }\PYG{l+s}{\PYGZdq{}}\PYG{l+s}{IV}\PYG{l+s}{\PYGZdq{}}\PYG{p}{\PYGZcb{}}\PYG{p}{,}
\PYG{+w}{    }\PYG{p}{\PYGZob{}}\PYG{+w}{ }\PYG{l+m+mi}{1}\PYG{p}{,}\PYG{+w}{ }\PYG{l+s}{\PYGZdq{}}\PYG{l+s}{I}\PYG{l+s}{\PYGZdq{}}\PYG{p}{\PYGZcb{}}
\PYG{+w}{  }\PYG{p}{\PYGZcb{}}\PYG{p}{;}
\PYG{n}{std}\PYG{o}{:}\PYG{o}{:}\PYG{n}{string}\PYG{+w}{ }\PYG{n}{result}\PYG{p}{;}
\PYG{k}{for}\PYG{+w}{ }\PYG{p}{(}\PYG{k}{auto}\PYG{+w}{ }\PYG{k}{const}\PYG{+w}{ }\PYG{o}{\PYGZam{}}\PYG{+w}{ }\PYG{n}{number}\PYG{+w}{ }\PYG{o}{:}\PYG{+w}{ }\PYG{n}{rarr}\PYG{p}{)}
\PYG{p}{\PYGZob{}}
\PYG{+w}{  }\PYG{k}{while}\PYG{+w}{ }\PYG{p}{(}\PYG{n}{value}\PYG{+w}{ }\PYG{o}{\PYGZgt{}}\PYG{o}{=}\PYG{+w}{ }\PYG{n}{number}\PYG{p}{.}\PYG{n}{value}\PYG{p}{)}
\PYG{+w}{  }\PYG{p}{\PYGZob{}}
\PYG{+w}{    }\PYG{n}{result}\PYG{+w}{ }\PYG{o}{+}\PYG{o}{=}\PYG{+w}{ }\PYG{n}{number}\PYG{p}{.}\PYG{n}{numeral}\PYG{p}{;}
\PYG{+w}{    }\PYG{n}{value}\PYG{+w}{ }\PYG{o}{\PYGZhy{}}\PYG{o}{=}\PYG{+w}{ }\PYG{n}{number}\PYG{p}{.}\PYG{n}{value}\PYG{p}{;}
\PYG{+w}{  }\PYG{p}{\PYGZcb{}}
\PYG{p}{\PYGZcb{}}
\PYG{+w}{  }\PYG{k}{return}\PYG{+w}{ }\PYG{n}{result}\PYG{p}{;}
\PYG{p}{\PYGZcb{}}
\end{Verbatim}
\end{codeblockbg}


%%C070201 p

\texttt{print\_\allowbreak{}toc()}函数如以下代码所示，它接受一个\texttt{Book}参数，并按上述要求把内容打印到控制台。为此我们用到：


%%C070202 li

\begin{itemize}
\item \texttt{std::\allowbreak{}left}和\texttt{std::\allowbreak{}right}指定文本对齐方式。


%%C070203 li

\item \texttt{std::\allowbreak{}setw}指定每个输出字段的宽度。


%%C070204 li

\item \texttt{std::\allowbreak{}fill}指定填充字符（章节号用空格，章节标题用点号）：


%%C070205 code

\begin{codeblockbg}
\begin{Verbatim}[commandchars=\\\{\},breaklines=true,breakanywhere=true,fontsize=\codefont,xleftmargin=17bp]
\PYG{k+kt}{void}\PYG{+w}{ }\PYG{n+nf}{print\PYGZus{}toc}\PYG{p}{(}\PYG{n}{Book}\PYG{+w}{ }\PYG{k}{const}\PYG{+w}{ }\PYG{o}{\PYGZam{}}\PYG{+w}{ }\PYG{n}{book}\PYG{p}{)}
\PYG{p}{\PYGZob{}}
\PYG{+w}{  }\PYG{n}{std}\PYG{o}{:}\PYG{o}{:}\PYG{n}{cout}\PYG{+w}{ }\PYG{o}{\PYGZlt{}}\PYG{o}{\PYGZlt{}}\PYG{+w}{ }\PYG{n}{book}\PYG{p}{.}\PYG{n}{Title}\PYG{+w}{ }\PYG{o}{\PYGZlt{}}\PYG{o}{\PYGZlt{}}\PYG{+w}{ }\PYG{n}{std}\PYG{o}{:}\PYG{o}{:}\PYG{n}{endl}\PYG{p}{;}
\PYG{+w}{  }\PYG{k}{for}\PYG{p}{(}\PYG{k}{auto}\PYG{+w}{ }\PYG{k}{const}\PYG{+w}{ }\PYG{o}{\PYGZam{}}\PYG{+w}{ }\PYG{n}{part}\PYG{+w}{ }\PYG{o}{:}\PYG{+w}{ }\PYG{n}{book}\PYG{p}{.}\PYG{n}{Parts}\PYG{p}{)}
\PYG{p}{\PYGZob{}}
\PYG{+w}{  }\PYG{n}{std}\PYG{o}{:}\PYG{o}{:}\PYG{n}{cout}\PYG{+w}{ }\PYG{o}{\PYGZlt{}}\PYG{o}{\PYGZlt{}}\PYG{+w}{ }\PYG{n}{std}\PYG{o}{:}\PYG{o}{:}\PYG{n}{left}\PYG{+w}{ }\PYG{o}{\PYGZlt{}}\PYG{o}{\PYGZlt{}}\PYG{+w}{ }\PYG{n}{std}\PYG{o}{:}\PYG{o}{:}\PYG{n}{setw}\PYG{p}{(}\PYG{l+m+mi}{15}\PYG{p}{)}\PYG{+w}{ }\PYG{o}{\PYGZlt{}}\PYG{o}{\PYGZlt{}}\PYG{+w}{ }\PYG{n}{std}\PYG{o}{:}\PYG{o}{:}\PYG{n}{setfill}\PYG{p}{(}\PYG{l+s+sc}{\PYGZsq{}}\PYG{l+s+sc}{ }\PYG{l+s+sc}{\PYGZsq{}}\PYG{p}{)}
\PYG{+w}{            }\PYG{o}{\PYGZlt{}}\PYG{o}{\PYGZlt{}}\PYG{+w}{ }\PYG{n}{part}\PYG{p}{.}\PYG{n}{Title}\PYG{+w}{ }\PYG{o}{\PYGZlt{}}\PYG{o}{\PYGZlt{}}\PYG{+w}{ }\PYG{n}{std}\PYG{o}{:}\PYG{o}{:}\PYG{n}{endl}\PYG{p}{;}
\PYG{+w}{  }\PYG{n}{std}\PYG{o}{:}\PYG{o}{:}\PYG{n}{cout}\PYG{+w}{ }\PYG{o}{\PYGZlt{}}\PYG{o}{\PYGZlt{}}\PYG{+w}{ }\PYG{n}{std}\PYG{o}{:}\PYG{o}{:}\PYG{n}{left}\PYG{+w}{ }\PYG{o}{\PYGZlt{}}\PYG{o}{\PYGZlt{}}\PYG{+w}{ }\PYG{n}{std}\PYG{o}{:}\PYG{o}{:}\PYG{n}{setw}\PYG{p}{(}\PYG{l+m+mi}{15}\PYG{p}{)}\PYG{+w}{ }\PYG{o}{\PYGZlt{}}\PYG{o}{\PYGZlt{}}\PYG{+w}{ }\PYG{n}{std}\PYG{o}{:}\PYG{o}{:}\PYG{n}{setfill}\PYG{p}{(}\PYG{l+s+sc}{\PYGZsq{}}\PYG{l+s+sc}{\PYGZhy{}}\PYG{l+s+sc}{\PYGZsq{}}\PYG{p}{)}
\PYG{+w}{            }\PYG{o}{\PYGZlt{}}\PYG{o}{\PYGZlt{}}\PYG{+w}{ }\PYG{l+s+sc}{\PYGZsq{}}\PYG{l+s+sc}{\PYGZhy{}}\PYG{l+s+sc}{\PYGZsq{}}\PYG{+w}{ }\PYG{o}{\PYGZlt{}}\PYG{o}{\PYGZlt{}}\PYG{+w}{ }\PYG{n}{std}\PYG{o}{:}\PYG{o}{:}\PYG{n}{endl}\PYG{p}{;}
\PYG{+w}{    }\PYG{k}{for}\PYG{p}{(}\PYG{k}{auto}\PYG{+w}{ }\PYG{k}{const}\PYG{+w}{ }\PYG{o}{\PYGZam{}}\PYG{+w}{ }\PYG{n}{chapter}\PYG{+w}{ }\PYG{o}{:}\PYG{+w}{ }\PYG{n}{part}\PYG{p}{.}\PYG{n}{Chapters}\PYG{p}{)}
\PYG{+w}{    }\PYG{p}{\PYGZob{}}
\PYG{+w}{      }\PYG{n}{std}\PYG{o}{:}\PYG{o}{:}\PYG{n}{cout}\PYG{+w}{ }\PYG{o}{\PYGZlt{}}\PYG{o}{\PYGZlt{}}\PYG{+w}{ }\PYG{n}{std}\PYG{o}{:}\PYG{o}{:}\PYG{n}{right}\PYG{+w}{ }\PYG{o}{\PYGZlt{}}\PYG{o}{\PYGZlt{}}\PYG{+w}{ }\PYG{n}{std}\PYG{o}{:}\PYG{o}{:}\PYG{n}{setw}\PYG{p}{(}\PYG{l+m+mi}{4}\PYG{p}{)}\PYG{+w}{ }\PYG{o}{\PYGZlt{}}\PYG{o}{\PYGZlt{}}\PYG{+w}{ }\PYG{n}{std}\PYG{o}{:}\PYG{o}{:}\PYG{n}{setfill}\PYG{p}{(}\PYG{l+s+sc}{\PYGZsq{}}\PYG{l+s+sc}{ }\PYG{l+s+sc}{\PYGZsq{}}\PYG{p}{)}
\PYG{+w}{                }\PYG{o}{\PYGZlt{}}\PYG{o}{\PYGZlt{}}\PYG{+w}{ }\PYG{n}{to\PYGZus{}roman}\PYG{p}{(}\PYG{n}{chapter}\PYG{p}{.}\PYG{n}{Number}\PYG{p}{)}\PYG{+w}{ }\PYG{o}{\PYGZlt{}}\PYG{o}{\PYGZlt{}}\PYG{+w}{ }\PYG{l+s+sc}{\PYGZsq{}}\PYG{l+s+sc}{ }\PYG{l+s+sc}{\PYGZsq{}}\PYG{p}{;}
\PYG{+w}{      }\PYG{n}{std}\PYG{o}{:}\PYG{o}{:}\PYG{n}{cout}\PYG{+w}{ }\PYG{o}{\PYGZlt{}}\PYG{o}{\PYGZlt{}}\PYG{+w}{ }\PYG{n}{std}\PYG{o}{:}\PYG{o}{:}\PYG{n}{left}\PYG{+w}{ }\PYG{o}{\PYGZlt{}}\PYG{o}{\PYGZlt{}}\PYG{+w}{ }\PYG{n}{std}\PYG{o}{:}\PYG{o}{:}\PYG{n}{setw}\PYG{p}{(}\PYG{l+m+mi}{35}\PYG{p}{)}\PYG{+w}{ }\PYG{o}{\PYGZlt{}}\PYG{o}{\PYGZlt{}}\PYG{+w}{ }\PYG{n}{std}\PYG{o}{:}\PYG{o}{:}\PYG{n}{setfill}\PYG{p}{(}\PYG{l+s+sc}{\PYGZsq{}}\PYG{l+s+sc}{.}\PYG{l+s+sc}{\PYGZsq{}}\PYG{p}{)}
\PYG{+w}{                }\PYG{o}{\PYGZlt{}}\PYG{o}{\PYGZlt{}}\PYG{+w}{ }\PYG{n}{chapter}\PYG{p}{.}\PYG{n}{Title}\PYG{p}{;}
\PYG{+w}{      }\PYG{n}{std}\PYG{o}{:}\PYG{o}{:}\PYG{n}{cout}\PYG{+w}{ }\PYG{o}{\PYGZlt{}}\PYG{o}{\PYGZlt{}}\PYG{+w}{ }\PYG{n}{std}\PYG{o}{:}\PYG{o}{:}\PYG{n}{right}\PYG{+w}{ }\PYG{o}{\PYGZlt{}}\PYG{o}{\PYGZlt{}}\PYG{+w}{ }\PYG{n}{std}\PYG{o}{:}\PYG{o}{:}\PYG{n}{setw}\PYG{p}{(}\PYG{l+m+mi}{3}\PYG{p}{)}\PYG{+w}{ }\PYG{o}{\PYGZlt{}}\PYG{o}{\PYGZlt{}}\PYG{+w}{ }\PYG{n}{std}\PYG{o}{:}\PYG{o}{:}\PYG{n}{setfill}\PYG{p}{(}\PYG{l+s+sc}{\PYGZsq{}}\PYG{l+s+sc}{.}\PYG{l+s+sc}{\PYGZsq{}}\PYG{p}{)}
\PYG{+w}{                }\PYG{o}{\PYGZlt{}}\PYG{o}{\PYGZlt{}}\PYG{+w}{ }\PYG{n}{chapter}\PYG{p}{.}\PYG{n}{Page}\PYG{+w}{ }\PYG{o}{\PYGZlt{}}\PYG{o}{\PYGZlt{}}\PYG{+w}{ }\PYG{n}{std}\PYG{o}{:}\PYG{o}{:}\PYG{n}{endl}\PYG{p}{;}
\PYG{+w}{    }\PYG{p}{\PYGZcb{}}
\PYG{+w}{  }\PYG{p}{\PYGZcb{}}
\PYG{p}{\PYGZcb{}}
\end{Verbatim}
\end{codeblockbg}

\end{itemize}


%%C070206 p

下面是使用该方法的示例，其中\texttt{Book}对象描述的是《The Fellowship of the Ring》一书的目录：


%%C070207 code

\begin{codeblockbg}
\begin{Verbatim}[commandchars=\\\{\},breaklines=true,breakanywhere=true,fontsize=\codefont,xleftmargin=17bp]
\PYG{k}{auto}\PYG{+w}{ }\PYG{n}{book}\PYG{+w}{ }\PYG{o}{=}\PYG{+w}{ }\PYG{n}{Book}
\PYG{p}{\PYGZob{}}
\PYG{+w}{  }\PYG{l+s}{\PYGZdq{}}\PYG{l+s}{THE FELLOWSHIP OF THE RING}\PYG{l+s}{\PYGZdq{}}\PYG{n}{s}\PYG{p}{,}
\PYG{+w}{  }\PYG{p}{\PYGZob{}}
\PYG{+w}{    }\PYG{p}{\PYGZob{}}
\PYG{+w}{      }\PYG{l+s}{\PYGZdq{}}\PYG{l+s}{BOOK ONE}\PYG{l+s}{\PYGZdq{}}\PYG{n}{s}\PYG{p}{,}
\PYG{+w}{      }\PYG{p}{\PYGZob{}}
\PYG{+w}{        }\PYG{p}{\PYGZob{}}\PYG{l+m+mi}{1}\PYG{p}{,}\PYG{+w}{ }\PYG{l+s}{\PYGZdq{}}\PYG{l+s}{A Long\PYGZhy{}expected Party}\PYG{l+s}{\PYGZdq{}}\PYG{n}{s}\PYG{p}{,}\PYG{+w}{ }\PYG{l+m+mi}{21}\PYG{p}{\PYGZcb{}}\PYG{p}{,}
\PYG{+w}{        }\PYG{p}{\PYGZob{}}\PYG{l+m+mi}{2}\PYG{p}{,}\PYG{+w}{ }\PYG{l+s}{\PYGZdq{}}\PYG{l+s}{The Shadow of the Past}\PYG{l+s}{\PYGZdq{}}\PYG{n}{s}\PYG{p}{,}\PYG{+w}{ }\PYG{l+m+mi}{42}\PYG{p}{\PYGZcb{}}\PYG{p}{,}
\PYG{+w}{        }\PYG{p}{\PYGZob{}}\PYG{l+m+mi}{3}\PYG{p}{,}\PYG{+w}{ }\PYG{l+s}{\PYGZdq{}}\PYG{l+s}{Three Is Company}\PYG{l+s}{\PYGZdq{}}\PYG{n}{s}\PYG{p}{,}\PYG{+w}{ }\PYG{l+m+mi}{65}\PYG{p}{\PYGZcb{}}\PYG{p}{,}
\PYG{+w}{        }\PYG{p}{\PYGZob{}}\PYG{l+m+mi}{4}\PYG{p}{,}\PYG{+w}{ }\PYG{l+s}{\PYGZdq{}}\PYG{l+s}{A Short Cut to Mushrooms}\PYG{l+s}{\PYGZdq{}}\PYG{n}{s}\PYG{p}{,}\PYG{+w}{ }\PYG{l+m+mi}{86}\PYG{p}{\PYGZcb{}}\PYG{p}{,}
\PYG{+w}{        }\PYG{p}{\PYGZob{}}\PYG{l+m+mi}{5}\PYG{p}{,}\PYG{+w}{ }\PYG{l+s}{\PYGZdq{}}\PYG{l+s}{A Conspiracy Unmasked}\PYG{l+s}{\PYGZdq{}}\PYG{n}{s}\PYG{p}{,}\PYG{+w}{ }\PYG{l+m+mi}{98}\PYG{p}{\PYGZcb{}}\PYG{p}{,}
\PYG{+w}{        }\PYG{p}{\PYGZob{}}\PYG{l+m+mi}{6}\PYG{p}{,}\PYG{+w}{ }\PYG{l+s}{\PYGZdq{}}\PYG{l+s}{The Old Forest}\PYG{l+s}{\PYGZdq{}}\PYG{n}{s}\PYG{p}{,}\PYG{+w}{ }\PYG{l+m+mi}{109}\PYG{p}{\PYGZcb{}}\PYG{p}{,}
\PYG{+w}{        }\PYG{p}{\PYGZob{}}\PYG{l+m+mi}{7}\PYG{p}{,}\PYG{+w}{ }\PYG{l+s}{\PYGZdq{}}\PYG{l+s}{In the House of Tom Bombadil}\PYG{l+s}{\PYGZdq{}}\PYG{n}{s}\PYG{p}{,}\PYG{+w}{ }\PYG{l+m+mi}{123}\PYG{p}{\PYGZcb{}}\PYG{p}{,}
\PYG{+w}{        }\PYG{p}{\PYGZob{}}\PYG{l+m+mi}{8}\PYG{p}{,}\PYG{+w}{ }\PYG{l+s}{\PYGZdq{}}\PYG{l+s}{Fog on the Barrow\PYGZhy{}downs}\PYG{l+s}{\PYGZdq{}}\PYG{n}{s}\PYG{p}{,}\PYG{+w}{ }\PYG{l+m+mi}{135}\PYG{p}{\PYGZcb{}}\PYG{p}{,}
\PYG{+w}{        }\PYG{p}{\PYGZob{}}\PYG{l+m+mi}{9}\PYG{p}{,}\PYG{+w}{ }\PYG{l+s}{\PYGZdq{}}\PYG{l+s}{At the Sign of The Prancing Pony}\PYG{l+s}{\PYGZdq{}}\PYG{n}{s}\PYG{p}{,}\PYG{+w}{ }\PYG{l+m+mi}{149}\PYG{p}{\PYGZcb{}}\PYG{p}{,}
\PYG{+w}{        }\PYG{p}{\PYGZob{}}\PYG{l+m+mi}{10}\PYG{p}{,}\PYG{+w}{ }\PYG{l+s}{\PYGZdq{}}\PYG{l+s}{Strider}\PYG{l+s}{\PYGZdq{}}\PYG{n}{s}\PYG{p}{,}\PYG{+w}{ }\PYG{l+m+mi}{163}\PYG{p}{\PYGZcb{}}\PYG{p}{,}
\PYG{+w}{        }\PYG{p}{\PYGZob{}}\PYG{l+m+mi}{11}\PYG{p}{,}\PYG{+w}{ }\PYG{l+s}{\PYGZdq{}}\PYG{l+s}{A Knife in the Dark}\PYG{l+s}{\PYGZdq{}}\PYG{n}{s}\PYG{p}{,}\PYG{+w}{ }\PYG{l+m+mi}{176}\PYG{p}{\PYGZcb{}}\PYG{p}{,}
\PYG{+w}{        }\PYG{p}{\PYGZob{}}\PYG{l+m+mi}{12}\PYG{p}{,}\PYG{+w}{ }\PYG{l+s}{\PYGZdq{}}\PYG{l+s}{Flight to the Ford}\PYG{l+s}{\PYGZdq{}}\PYG{n}{s}\PYG{p}{,}\PYG{+w}{ }\PYG{l+m+mi}{197}\PYG{p}{\PYGZcb{}}\PYG{p}{,}
\PYG{+w}{      }\PYG{p}{\PYGZcb{}}\PYG{p}{,}
\PYG{+w}{    }\PYG{p}{\PYGZcb{}}\PYG{p}{,}
\PYG{+w}{    }\PYG{p}{\PYGZob{}}
\PYG{+w}{      }\PYG{l+s}{\PYGZdq{}}\PYG{l+s}{BOOK TWO}\PYG{l+s}{\PYGZdq{}}\PYG{n}{s}\PYG{p}{,}
\PYG{+w}{      }\PYG{p}{\PYGZob{}}
\PYG{+w}{        }\PYG{p}{\PYGZob{}}\PYG{l+m+mi}{1}\PYG{p}{,}\PYG{+w}{ }\PYG{l+s}{\PYGZdq{}}\PYG{l+s}{Many Meetings}\PYG{l+s}{\PYGZdq{}}\PYG{n}{s}\PYG{p}{,}\PYG{+w}{ }\PYG{l+m+mi}{219}\PYG{p}{\PYGZcb{}}\PYG{p}{,}
\PYG{+w}{        }\PYG{p}{\PYGZob{}}\PYG{l+m+mi}{2}\PYG{p}{,}\PYG{+w}{ }\PYG{l+s}{\PYGZdq{}}\PYG{l+s}{The Council of Elrond}\PYG{l+s}{\PYGZdq{}}\PYG{n}{s}\PYG{p}{,}\PYG{+w}{ }\PYG{l+m+mi}{239}\PYG{p}{\PYGZcb{}}\PYG{p}{,}
\PYG{+w}{        }\PYG{p}{\PYGZob{}}\PYG{l+m+mi}{3}\PYG{p}{,}\PYG{+w}{ }\PYG{l+s}{\PYGZdq{}}\PYG{l+s}{The Ring Goes South}\PYG{l+s}{\PYGZdq{}}\PYG{n}{s}\PYG{p}{,}\PYG{+w}{ }\PYG{l+m+mi}{272}\PYG{p}{\PYGZcb{}}\PYG{p}{,}
\PYG{+w}{        }\PYG{p}{\PYGZob{}}\PYG{l+m+mi}{4}\PYG{p}{,}\PYG{+w}{ }\PYG{l+s}{\PYGZdq{}}\PYG{l+s}{A Journey in the Dark}\PYG{l+s}{\PYGZdq{}}\PYG{n}{s}\PYG{p}{,}\PYG{+w}{ }\PYG{l+m+mi}{295}\PYG{p}{\PYGZcb{}}\PYG{p}{,}
\PYG{+w}{        }\PYG{p}{\PYGZob{}}\PYG{l+m+mi}{5}\PYG{p}{,}\PYG{+w}{ }\PYG{l+s}{\PYGZdq{}}\PYG{l+s}{The Bridge of Khazad\PYGZhy{}dum}\PYG{l+s}{\PYGZdq{}}\PYG{n}{s}\PYG{p}{,}\PYG{+w}{ }\PYG{l+m+mi}{321}\PYG{p}{\PYGZcb{}}\PYG{p}{,}
\PYG{+w}{        }\PYG{p}{\PYGZob{}}\PYG{l+m+mi}{6}\PYG{p}{,}\PYG{+w}{ }\PYG{l+s}{\PYGZdq{}}\PYG{l+s}{Lothlorien}\PYG{l+s}{\PYGZdq{}}\PYG{n}{s}\PYG{p}{,}\PYG{+w}{ }\PYG{l+m+mi}{333}\PYG{p}{\PYGZcb{}}\PYG{p}{,}
\PYG{+w}{        }\PYG{p}{\PYGZob{}}\PYG{l+m+mi}{7}\PYG{p}{,}\PYG{+w}{ }\PYG{l+s}{\PYGZdq{}}\PYG{l+s}{The Mirror of Galadriel}\PYG{l+s}{\PYGZdq{}}\PYG{n}{s}\PYG{p}{,}\PYG{+w}{ }\PYG{l+m+mi}{353}\PYG{p}{\PYGZcb{}}\PYG{p}{,}
\PYG{+w}{        }\PYG{p}{\PYGZob{}}\PYG{l+m+mi}{8}\PYG{p}{,}\PYG{+w}{ }\PYG{l+s}{\PYGZdq{}}\PYG{l+s}{Farewell to Lorien}\PYG{l+s}{\PYGZdq{}}\PYG{n}{s}\PYG{p}{,}\PYG{+w}{ }\PYG{l+m+mi}{367}\PYG{p}{\PYGZcb{}}\PYG{p}{,}
\PYG{+w}{        }\PYG{p}{\PYGZob{}}\PYG{l+m+mi}{9}\PYG{p}{,}\PYG{+w}{ }\PYG{l+s}{\PYGZdq{}}\PYG{l+s}{The Great River}\PYG{l+s}{\PYGZdq{}}\PYG{n}{s}\PYG{p}{,}\PYG{+w}{ }\PYG{l+m+mi}{380}\PYG{p}{\PYGZcb{}}\PYG{p}{,}
\PYG{+w}{        }\PYG{p}{\PYGZob{}}\PYG{l+m+mi}{10}\PYG{p}{,}\PYG{+w}{ }\PYG{l+s}{\PYGZdq{}}\PYG{l+s}{The Breaking of the Fellowship}\PYG{l+s}{\PYGZdq{}}\PYG{n}{s}\PYG{p}{,}\PYG{+w}{ }\PYG{l+m+mi}{390}\PYG{p}{\PYGZcb{}}\PYG{p}{,}
\PYG{+w}{      }\PYG{p}{\PYGZcb{}}\PYG{p}{,}
\PYG{+w}{    }\PYG{p}{\PYGZcb{}}\PYG{p}{,}
\PYG{+w}{  }\PYG{p}{\PYGZcb{}}
\PYG{p}{\PYGZcb{}}\PYG{p}{;}
\PYG{n}{print\PYGZus{}toc}\PYG{p}{(}\PYG{n}{book}\PYG{p}{)}\PYG{p}{;}
\end{Verbatim}
\end{codeblockbg}


%%C070208 p

这时输出如下：


%%C070209 code

\begin{codeblockbg}
\begin{Verbatim}[commandchars=\\\{\},breaklines=true,breakanywhere=true,fontsize=\codefont,xleftmargin=17bp]
\PYG{n}{THE}\PYG{+w}{ }\PYG{n}{FELLOWSHIP}\PYG{+w}{ }\PYG{n}{OF}\PYG{+w}{ }\PYG{n}{THE}\PYG{+w}{ }\PYG{n}{RING}
\PYG{n}{BOOK}\PYG{+w}{ }\PYG{n}{ONE}
\PYG{o}{\PYGZhy{}}\PYG{o}{\PYGZhy{}}\PYG{o}{\PYGZhy{}}\PYG{o}{\PYGZhy{}}\PYG{o}{\PYGZhy{}}\PYG{o}{\PYGZhy{}}\PYG{o}{\PYGZhy{}}\PYG{o}{\PYGZhy{}}\PYG{o}{\PYGZhy{}}\PYG{o}{\PYGZhy{}}\PYG{o}{\PYGZhy{}}\PYG{o}{\PYGZhy{}}\PYG{o}{\PYGZhy{}}\PYG{o}{\PYGZhy{}}\PYG{o}{\PYGZhy{}}
\PYG{+w}{   }\PYG{n}{I}\PYG{+w}{ }\PYG{n}{A}\PYG{+w}{ }\PYG{n}{Long}\PYG{o}{\PYGZhy{}}\PYG{n}{expected}\PYG{+w}{ }\PYG{n}{Party}\PYG{p}{.}\PYG{p}{.}\PYG{p}{.}\PYG{p}{.}\PYG{p}{.}\PYG{p}{.}\PYG{p}{.}\PYG{p}{.}\PYG{p}{.}\PYG{p}{.}\PYG{p}{.}\PYG{p}{.}\PYG{p}{.}\PYG{p}{.}\PYG{l+m+mf}{.21}
\PYG{+w}{  }\PYG{n}{II}\PYG{+w}{ }\PYG{n}{The}\PYG{+w}{ }\PYG{n}{Shadow}\PYG{+w}{ }\PYG{n}{of}\PYG{+w}{ }\PYG{n}{the}\PYG{+w}{ }\PYG{n}{Past}\PYG{p}{.}\PYG{p}{.}\PYG{p}{.}\PYG{p}{.}\PYG{p}{.}\PYG{p}{.}\PYG{p}{.}\PYG{p}{.}\PYG{p}{.}\PYG{p}{.}\PYG{p}{.}\PYG{p}{.}\PYG{p}{.}\PYG{l+m+mf}{.42}
\PYG{+w}{ }\PYG{n}{III}\PYG{+w}{ }\PYG{n}{Three}\PYG{+w}{ }\PYG{n}{Is}\PYG{+w}{ }\PYG{n}{Company}\PYG{p}{.}\PYG{p}{.}\PYG{p}{.}\PYG{p}{.}\PYG{p}{.}\PYG{p}{.}\PYG{p}{.}\PYG{p}{.}\PYG{p}{.}\PYG{p}{.}\PYG{p}{.}\PYG{p}{.}\PYG{p}{.}\PYG{p}{.}\PYG{p}{.}\PYG{p}{.}\PYG{p}{.}\PYG{p}{.}\PYG{p}{.}\PYG{l+m+mf}{.65}
\PYG{+w}{  }\PYG{n}{IV}\PYG{+w}{ }\PYG{n}{A}\PYG{+w}{ }\PYG{n}{Short}\PYG{+w}{ }\PYG{n}{Cut}\PYG{+w}{ }\PYG{n}{to}\PYG{+w}{ }\PYG{n}{Mushrooms}\PYG{p}{.}\PYG{p}{.}\PYG{p}{.}\PYG{p}{.}\PYG{p}{.}\PYG{p}{.}\PYG{p}{.}\PYG{p}{.}\PYG{p}{.}\PYG{p}{.}\PYG{p}{.}\PYG{l+m+mf}{.86}
\PYG{+w}{   }\PYG{n}{V}\PYG{+w}{ }\PYG{n}{A}\PYG{+w}{ }\PYG{n}{Conspiracy}\PYG{+w}{ }\PYG{n}{Unmasked}\PYG{p}{.}\PYG{p}{.}\PYG{p}{.}\PYG{p}{.}\PYG{p}{.}\PYG{p}{.}\PYG{p}{.}\PYG{p}{.}\PYG{p}{.}\PYG{p}{.}\PYG{p}{.}\PYG{p}{.}\PYG{p}{.}\PYG{p}{.}\PYG{l+m+mf}{.98}
\PYG{+w}{  }\PYG{n}{VI}\PYG{+w}{ }\PYG{n}{The}\PYG{+w}{ }\PYG{n}{Old}\PYG{+w}{ }\PYG{n}{Forest}\PYG{p}{.}\PYG{p}{.}\PYG{p}{.}\PYG{p}{.}\PYG{p}{.}\PYG{p}{.}\PYG{p}{.}\PYG{p}{.}\PYG{p}{.}\PYG{p}{.}\PYG{p}{.}\PYG{p}{.}\PYG{p}{.}\PYG{p}{.}\PYG{p}{.}\PYG{p}{.}\PYG{p}{.}\PYG{p}{.}\PYG{p}{.}\PYG{p}{.}\PYG{l+m+mf}{.109}
\PYG{+w}{ }\PYG{n}{VII}\PYG{+w}{ }\PYG{n}{In}\PYG{+w}{ }\PYG{n}{the}\PYG{+w}{ }\PYG{n}{House}\PYG{+w}{ }\PYG{n}{of}\PYG{+w}{ }\PYG{n}{Tom}\PYG{+w}{ }\PYG{n}{Bombadil}\PYG{p}{.}\PYG{p}{.}\PYG{p}{.}\PYG{p}{.}\PYG{p}{.}\PYG{p}{.}\PYG{l+m+mf}{.123}
\PYG{n}{VIII}\PYG{+w}{ }\PYG{n}{Fog}\PYG{+w}{ }\PYG{n}{on}\PYG{+w}{ }\PYG{n}{the}\PYG{+w}{ }\PYG{n}{Barrow}\PYG{o}{\PYGZhy{}}\PYG{n}{downs}\PYG{p}{.}\PYG{p}{.}\PYG{p}{.}\PYG{p}{.}\PYG{p}{.}\PYG{p}{.}\PYG{p}{.}\PYG{p}{.}\PYG{p}{.}\PYG{p}{.}\PYG{p}{.}\PYG{l+m+mf}{.135}
\PYG{+w}{  }\PYG{n}{IX}\PYG{+w}{ }\PYG{n}{At}\PYG{+w}{ }\PYG{n}{the}\PYG{+w}{ }\PYG{n}{Sign}\PYG{+w}{ }\PYG{n}{of}\PYG{+w}{ }\PYG{n}{The}\PYG{+w}{ }\PYG{n}{Prancing}\PYG{+w}{ }\PYG{n}{Pony}\PYG{p}{.}\PYG{p}{.}\PYG{l+m+mf}{.149}
\PYG{+w}{   }\PYG{n}{X}\PYG{+w}{ }\PYG{n}{Strider}\PYG{p}{.}\PYG{p}{.}\PYG{p}{.}\PYG{p}{.}\PYG{p}{.}\PYG{p}{.}\PYG{p}{.}\PYG{p}{.}\PYG{p}{.}\PYG{p}{.}\PYG{p}{.}\PYG{p}{.}\PYG{p}{.}\PYG{p}{.}\PYG{p}{.}\PYG{p}{.}\PYG{p}{.}\PYG{p}{.}\PYG{p}{.}\PYG{p}{.}\PYG{p}{.}\PYG{p}{.}\PYG{p}{.}\PYG{p}{.}\PYG{p}{.}\PYG{p}{.}\PYG{p}{.}\PYG{l+m+mf}{.163}
\PYG{+w}{  }\PYG{n}{XI}\PYG{+w}{ }\PYG{n}{A}\PYG{+w}{ }\PYG{n}{Knife}\PYG{+w}{ }\PYG{n}{in}\PYG{+w}{ }\PYG{n}{the}\PYG{+w}{ }\PYG{n}{Dark}\PYG{p}{.}\PYG{p}{.}\PYG{p}{.}\PYG{p}{.}\PYG{p}{.}\PYG{p}{.}\PYG{p}{.}\PYG{p}{.}\PYG{p}{.}\PYG{p}{.}\PYG{p}{.}\PYG{p}{.}\PYG{p}{.}\PYG{p}{.}\PYG{p}{.}\PYG{l+m+mf}{.176}
\PYG{+w}{ }\PYG{n}{XII}\PYG{+w}{ }\PYG{n}{Flight}\PYG{+w}{ }\PYG{n}{to}\PYG{+w}{ }\PYG{n}{the}\PYG{+w}{ }\PYG{n}{Ford}\PYG{p}{.}\PYG{p}{.}\PYG{p}{.}\PYG{p}{.}\PYG{p}{.}\PYG{p}{.}\PYG{p}{.}\PYG{p}{.}\PYG{p}{.}\PYG{p}{.}\PYG{p}{.}\PYG{p}{.}\PYG{p}{.}\PYG{p}{.}\PYG{p}{.}\PYG{p}{.}\PYG{l+m+mf}{.197}
\PYG{n}{BOOK}\PYG{+w}{ }\PYG{n}{TWO}
\PYG{o}{\PYGZhy{}}\PYG{o}{\PYGZhy{}}\PYG{o}{\PYGZhy{}}\PYG{o}{\PYGZhy{}}\PYG{o}{\PYGZhy{}}\PYG{o}{\PYGZhy{}}\PYG{o}{\PYGZhy{}}\PYG{o}{\PYGZhy{}}\PYG{o}{\PYGZhy{}}\PYG{o}{\PYGZhy{}}\PYG{o}{\PYGZhy{}}\PYG{o}{\PYGZhy{}}\PYG{o}{\PYGZhy{}}\PYG{o}{\PYGZhy{}}\PYG{o}{\PYGZhy{}}
\PYG{+w}{   }\PYG{n}{I}\PYG{+w}{ }\PYG{n}{Many}\PYG{+w}{ }\PYG{n}{Meetings}\PYG{p}{.}\PYG{p}{.}\PYG{p}{.}\PYG{p}{.}\PYG{p}{.}\PYG{p}{.}\PYG{p}{.}\PYG{p}{.}\PYG{p}{.}\PYG{p}{.}\PYG{p}{.}\PYG{p}{.}\PYG{p}{.}\PYG{p}{.}\PYG{p}{.}\PYG{p}{.}\PYG{p}{.}\PYG{p}{.}\PYG{p}{.}\PYG{p}{.}\PYG{p}{.}\PYG{l+m+mf}{.219}
\PYG{+w}{  }\PYG{n}{II}\PYG{+w}{ }\PYG{n}{The}\PYG{+w}{ }\PYG{n}{Council}\PYG{+w}{ }\PYG{n}{of}\PYG{+w}{ }\PYG{n}{Elrond}\PYG{p}{.}\PYG{p}{.}\PYG{p}{.}\PYG{p}{.}\PYG{p}{.}\PYG{p}{.}\PYG{p}{.}\PYG{p}{.}\PYG{p}{.}\PYG{p}{.}\PYG{p}{.}\PYG{p}{.}\PYG{p}{.}\PYG{l+m+mf}{.239}
\PYG{+w}{ }\PYG{n}{III}\PYG{+w}{ }\PYG{n}{The}\PYG{+w}{ }\PYG{n}{Ring}\PYG{+w}{ }\PYG{n}{Goes}\PYG{+w}{ }\PYG{n}{South}\PYG{p}{.}\PYG{p}{.}\PYG{p}{.}\PYG{p}{.}\PYG{p}{.}\PYG{p}{.}\PYG{p}{.}\PYG{p}{.}\PYG{p}{.}\PYG{p}{.}\PYG{p}{.}\PYG{p}{.}\PYG{p}{.}\PYG{p}{.}\PYG{p}{.}\PYG{l+m+mf}{.272}
\PYG{+w}{  }\PYG{n}{IV}\PYG{+w}{ }\PYG{n}{A}\PYG{+w}{ }\PYG{n}{Journey}\PYG{+w}{ }\PYG{n}{in}\PYG{+w}{ }\PYG{n}{the}\PYG{+w}{ }\PYG{n}{Dark}\PYG{p}{.}\PYG{p}{.}\PYG{p}{.}\PYG{p}{.}\PYG{p}{.}\PYG{p}{.}\PYG{p}{.}\PYG{p}{.}\PYG{p}{.}\PYG{p}{.}\PYG{p}{.}\PYG{p}{.}\PYG{p}{.}\PYG{l+m+mf}{.295}
\PYG{+w}{   }\PYG{n}{V}\PYG{+w}{ }\PYG{n}{The}\PYG{+w}{ }\PYG{n}{Bridge}\PYG{+w}{ }\PYG{n}{of}\PYG{+w}{ }\PYG{n}{Khazad}\PYG{o}{\PYGZhy{}}\PYG{n}{dum}\PYG{p}{.}\PYG{p}{.}\PYG{p}{.}\PYG{p}{.}\PYG{p}{.}\PYG{p}{.}\PYG{p}{.}\PYG{p}{.}\PYG{p}{.}\PYG{p}{.}\PYG{l+m+mf}{.321}
\PYG{+w}{  }\PYG{n}{VI}\PYG{+w}{ }\PYG{n}{Lothlorien}\PYG{p}{.}\PYG{p}{.}\PYG{p}{.}\PYG{p}{.}\PYG{p}{.}\PYG{p}{.}\PYG{p}{.}\PYG{p}{.}\PYG{p}{.}\PYG{p}{.}\PYG{p}{.}\PYG{p}{.}\PYG{p}{.}\PYG{p}{.}\PYG{p}{.}\PYG{p}{.}\PYG{p}{.}\PYG{p}{.}\PYG{p}{.}\PYG{p}{.}\PYG{p}{.}\PYG{p}{.}\PYG{p}{.}\PYG{p}{.}\PYG{l+m+mf}{.333}
\PYG{+w}{ }\PYG{n}{VII}\PYG{+w}{ }\PYG{n}{The}\PYG{+w}{ }\PYG{n}{Mirror}\PYG{+w}{ }\PYG{n}{of}\PYG{+w}{ }\PYG{n}{Galadriel}\PYG{p}{.}\PYG{p}{.}\PYG{p}{.}\PYG{p}{.}\PYG{p}{.}\PYG{p}{.}\PYG{p}{.}\PYG{p}{.}\PYG{p}{.}\PYG{p}{.}\PYG{p}{.}\PYG{l+m+mf}{.353}
\PYG{n}{VIII}\PYG{+w}{ }\PYG{n}{Farewell}\PYG{+w}{ }\PYG{n}{to}\PYG{+w}{ }\PYG{n}{Lorien}\PYG{p}{.}\PYG{p}{.}\PYG{p}{.}\PYG{p}{.}\PYG{p}{.}\PYG{p}{.}\PYG{p}{.}\PYG{p}{.}\PYG{p}{.}\PYG{p}{.}\PYG{p}{.}\PYG{p}{.}\PYG{p}{.}\PYG{p}{.}\PYG{p}{.}\PYG{p}{.}\PYG{l+m+mf}{.367}
\PYG{+w}{  }\PYG{n}{IX}\PYG{+w}{ }\PYG{n}{The}\PYG{+w}{ }\PYG{n}{Great}\PYG{+w}{ }\PYG{n}{River}\PYG{p}{.}\PYG{p}{.}\PYG{p}{.}\PYG{p}{.}\PYG{p}{.}\PYG{p}{.}\PYG{p}{.}\PYG{p}{.}\PYG{p}{.}\PYG{p}{.}\PYG{p}{.}\PYG{p}{.}\PYG{p}{.}\PYG{p}{.}\PYG{p}{.}\PYG{p}{.}\PYG{p}{.}\PYG{p}{.}\PYG{p}{.}\PYG{l+m+mf}{.380}
\PYG{+w}{   }\PYG{n}{X}\PYG{+w}{ }\PYG{n}{The}\PYG{+w}{ }\PYG{n}{Breaking}\PYG{+w}{ }\PYG{n}{of}\PYG{+w}{ }\PYG{n}{the}\PYG{+w}{ }\PYG{n}{Fellowship}\PYG{p}{.}\PYG{p}{.}\PYG{p}{.}\PYG{p}{.}\PYG{l+m+mf}{.390}
\end{Verbatim}
\end{codeblockbg}


%%C070210 p

第二个例子要输出一张按营收排名的全球最大公司表格。表中包含公司名称、行业、营收（十亿美元）、营收增减、营收增长率、员工人数和所属国家等列。这个例子中，我们使用下面的类：


%%C070211 code

\begin{codeblockbg}
\begin{Verbatim}[commandchars=\\\{\},breaklines=true,breakanywhere=true,fontsize=\codefont,xleftmargin=17bp]
\PYG{k}{struct}\PYG{+w}{ }\PYG{n+nc}{Company}
\PYG{p}{\PYGZob{}}
\PYG{+w}{  }\PYG{n}{std}\PYG{o}{:}\PYG{o}{:}\PYG{n}{string}\PYG{+w}{ }\PYG{n}{Name}\PYG{p}{;}
\PYG{+w}{  }\PYG{n}{std}\PYG{o}{:}\PYG{o}{:}\PYG{n}{string}\PYG{+w}{ }\PYG{n}{Industry}\PYG{p}{;}
\PYG{+w}{  }\PYG{k+kt}{double}\PYG{+w}{      }\PYG{n}{Revenue}\PYG{p}{;}
\PYG{+w}{  }\PYG{k+kt}{bool}\PYG{+w}{        }\PYG{n}{RevenueIncrease}\PYG{p}{;}
\PYG{+w}{  }\PYG{k+kt}{double}\PYG{+w}{      }\PYG{n}{Growth}\PYG{p}{;}
\PYG{+w}{  }\PYG{k+kt}{int}\PYG{+w}{         }\PYG{n}{Employees}\PYG{p}{;}
\PYG{+w}{  }\PYG{n}{std}\PYG{o}{:}\PYG{o}{:}\PYG{n}{string}\PYG{+w}{ }\PYG{n}{Country}\PYG{p}{;}
\PYG{p}{\PYGZcb{}}\PYG{p}{;}
\end{Verbatim}
\end{codeblockbg}


%%C070212 p

\texttt{print\_\allowbreak{}companies()}函数如以下代码所示，它在上一例的基础上又用了几个额外的操纵器：


%%C070213 li

\begin{itemize}
\item \texttt{std::\allowbreak{}boolalpha}把布尔值显示为\texttt{true}和\texttt{false}，而不是\texttt{1}和\texttt{0}。


%%C070214 li

\item \texttt{std::\allowbreak{}fixed}表示定点浮点表示法，随后由\texttt{std::\allowbreak{}defaultfloat}恢复默认的浮点表示法。


%%C070215 li

\item \texttt{std::\allowbreak{}setprecision}指定输出中显示的小数位数。与\texttt{std::\allowbreak{}fixed}配合，可把增长率字段表示为带一位小数的定点数。


%%C070216 code

\begin{codeblockbg}
\begin{Verbatim}[commandchars=\\\{\},breaklines=true,breakanywhere=true,fontsize=\codefont,xleftmargin=17bp]
\PYG{k+kt}{void}\PYG{+w}{ }\PYG{n+nf}{print\PYGZus{}companies}\PYG{p}{(}\PYG{n}{std}\PYG{o}{:}\PYG{o}{:}\PYG{n}{vector}\PYG{o}{\PYGZlt{}}\PYG{n}{Company}\PYG{o}{\PYGZgt{}}\PYG{+w}{ }\PYG{k}{const}\PYG{+w}{ }\PYG{o}{\PYGZam{}}\PYG{+w}{ }\PYG{n}{companies}\PYG{p}{)}
\PYG{p}{\PYGZob{}}
\PYG{+w}{  }\PYG{k}{for}\PYG{p}{(}\PYG{k}{auto}\PYG{+w}{ }\PYG{k}{const}\PYG{+w}{ }\PYG{o}{\PYGZam{}}\PYG{+w}{ }\PYG{n}{company}\PYG{+w}{ }\PYG{o}{:}\PYG{+w}{ }\PYG{n}{companies}\PYG{p}{)}
\PYG{+w}{  }\PYG{p}{\PYGZob{}}
\PYG{+w}{    }\PYG{n}{std}\PYG{o}{:}\PYG{o}{:}\PYG{n}{cout}\PYG{+w}{ }\PYG{o}{\PYGZlt{}}\PYG{o}{\PYGZlt{}}\PYG{+w}{ }\PYG{n}{std}\PYG{o}{:}\PYG{o}{:}\PYG{n}{left}\PYG{+w}{ }\PYG{o}{\PYGZlt{}}\PYG{o}{\PYGZlt{}}\PYG{+w}{ }\PYG{n}{std}\PYG{o}{:}\PYG{o}{:}\PYG{n}{setw}\PYG{p}{(}\PYG{l+m+mi}{26}\PYG{p}{)}\PYG{+w}{ }\PYG{o}{\PYGZlt{}}\PYG{o}{\PYGZlt{}}\PYG{+w}{ }\PYG{n}{std}\PYG{o}{:}\PYG{o}{:}\PYG{n}{setfill}\PYG{p}{(}\PYG{l+s+sc}{\PYGZsq{}}\PYG{l+s+sc}{ }\PYG{l+s+sc}{\PYGZsq{}}\PYG{p}{)}
\PYG{+w}{              }\PYG{o}{\PYGZlt{}}\PYG{o}{\PYGZlt{}}\PYG{+w}{ }\PYG{n}{company}\PYG{p}{.}\PYG{n}{Name}\PYG{p}{;}
\PYG{+w}{    }\PYG{n}{std}\PYG{o}{:}\PYG{o}{:}\PYG{n}{cout}\PYG{+w}{ }\PYG{o}{\PYGZlt{}}\PYG{o}{\PYGZlt{}}\PYG{+w}{ }\PYG{n}{std}\PYG{o}{:}\PYG{o}{:}\PYG{n}{left}\PYG{+w}{ }\PYG{o}{\PYGZlt{}}\PYG{o}{\PYGZlt{}}\PYG{+w}{ }\PYG{n}{std}\PYG{o}{:}\PYG{o}{:}\PYG{n}{setw}\PYG{p}{(}\PYG{l+m+mi}{18}\PYG{p}{)}\PYG{+w}{ }\PYG{o}{\PYGZlt{}}\PYG{o}{\PYGZlt{}}\PYG{+w}{ }\PYG{n}{std}\PYG{o}{:}\PYG{o}{:}\PYG{n}{setfill}\PYG{p}{(}\PYG{l+s+sc}{\PYGZsq{}}\PYG{l+s+sc}{ }\PYG{l+s+sc}{\PYGZsq{}}\PYG{p}{)}
\PYG{+w}{              }\PYG{o}{\PYGZlt{}}\PYG{o}{\PYGZlt{}}\PYG{+w}{ }\PYG{n}{company}\PYG{p}{.}\PYG{n}{Industry}\PYG{p}{;}
\PYG{+w}{    }\PYG{n}{std}\PYG{o}{:}\PYG{o}{:}\PYG{n}{cout}\PYG{+w}{ }\PYG{o}{\PYGZlt{}}\PYG{o}{\PYGZlt{}}\PYG{+w}{ }\PYG{n}{std}\PYG{o}{:}\PYG{o}{:}\PYG{n}{left}\PYG{+w}{ }\PYG{o}{\PYGZlt{}}\PYG{o}{\PYGZlt{}}\PYG{+w}{ }\PYG{n}{std}\PYG{o}{:}\PYG{o}{:}\PYG{n}{setw}\PYG{p}{(}\PYG{l+m+mi}{5}\PYG{p}{)}\PYG{+w}{ }\PYG{o}{\PYGZlt{}}\PYG{o}{\PYGZlt{}}\PYG{+w}{ }\PYG{n}{std}\PYG{o}{:}\PYG{o}{:}\PYG{n}{setfill}\PYG{p}{(}\PYG{l+s+sc}{\PYGZsq{}}\PYG{l+s+sc}{ }\PYG{l+s+sc}{\PYGZsq{}}\PYG{p}{)}
\PYG{+w}{              }\PYG{o}{\PYGZlt{}}\PYG{o}{\PYGZlt{}}\PYG{+w}{ }\PYG{n}{company}\PYG{p}{.}\PYG{n}{Revenue}\PYG{p}{;}
\PYG{+w}{    }\PYG{n}{std}\PYG{o}{:}\PYG{o}{:}\PYG{n}{cout}\PYG{+w}{ }\PYG{o}{\PYGZlt{}}\PYG{o}{\PYGZlt{}}\PYG{+w}{ }\PYG{n}{std}\PYG{o}{:}\PYG{o}{:}\PYG{n}{left}\PYG{+w}{ }\PYG{o}{\PYGZlt{}}\PYG{o}{\PYGZlt{}}\PYG{+w}{ }\PYG{n}{std}\PYG{o}{:}\PYG{o}{:}\PYG{n}{setw}\PYG{p}{(}\PYG{l+m+mi}{5}\PYG{p}{)}\PYG{+w}{ }\PYG{o}{\PYGZlt{}}\PYG{o}{\PYGZlt{}}\PYG{+w}{ }\PYG{n}{std}\PYG{o}{:}\PYG{o}{:}\PYG{n}{setfill}\PYG{p}{(}\PYG{l+s+sc}{\PYGZsq{}}\PYG{l+s+sc}{ }\PYG{l+s+sc}{\PYGZsq{}}\PYG{p}{)}
\PYG{+w}{              }\PYG{o}{\PYGZlt{}}\PYG{o}{\PYGZlt{}}\PYG{+w}{ }\PYG{n}{std}\PYG{o}{:}\PYG{o}{:}\PYG{n}{boolalpha}\PYG{+w}{ }\PYG{o}{\PYGZlt{}}\PYG{o}{\PYGZlt{}}\PYG{+w}{ }\PYG{n}{company}\PYG{p}{.}\PYG{n}{RevenueIncrease}
\PYG{+w}{              }\PYG{o}{\PYGZlt{}}\PYG{o}{\PYGZlt{}}\PYG{+w}{ }\PYG{n}{std}\PYG{o}{:}\PYG{o}{:}\PYG{n}{noboolalpha}\PYG{p}{;}
\PYG{+w}{    }\PYG{n}{std}\PYG{o}{:}\PYG{o}{:}\PYG{n}{cout}\PYG{+w}{ }\PYG{o}{\PYGZlt{}}\PYG{o}{\PYGZlt{}}\PYG{+w}{ }\PYG{n}{std}\PYG{o}{:}\PYG{o}{:}\PYG{n}{right}\PYG{+w}{ }\PYG{o}{\PYGZlt{}}\PYG{o}{\PYGZlt{}}\PYG{+w}{ }\PYG{n}{std}\PYG{o}{:}\PYG{o}{:}\PYG{n}{setw}\PYG{p}{(}\PYG{l+m+mi}{5}\PYG{p}{)}\PYG{+w}{ }\PYG{o}{\PYGZlt{}}\PYG{o}{\PYGZlt{}}\PYG{+w}{ }\PYG{n}{std}\PYG{o}{:}\PYG{o}{:}\PYG{n}{setfill}\PYG{p}{(}\PYG{l+s+sc}{\PYGZsq{}}\PYG{l+s+sc}{ }\PYG{l+s+sc}{\PYGZsq{}}\PYG{p}{)}
\PYG{+w}{              }\PYG{o}{\PYGZlt{}}\PYG{o}{\PYGZlt{}}\PYG{+w}{ }\PYG{n}{std}\PYG{o}{:}\PYG{o}{:}\PYG{n}{fixed}\PYG{+w}{ }\PYG{o}{\PYGZlt{}}\PYG{o}{\PYGZlt{}}\PYG{+w}{ }\PYG{n}{std}\PYG{o}{:}\PYG{o}{:}\PYG{n}{setprecision}\PYG{p}{(}\PYG{l+m+mi}{1}\PYG{p}{)}\PYG{+w}{ }\PYG{o}{\PYGZlt{}}\PYG{o}{\PYGZlt{}}\PYG{+w}{ }\PYG{n}{company}\PYG{p}{.}\PYG{n}{Growth}
\PYG{+w}{              }\PYG{o}{\PYGZlt{}}\PYG{o}{\PYGZlt{}}\PYG{+w}{ }\PYG{n}{std}\PYG{o}{:}\PYG{o}{:}\PYG{n}{defaultfloat}\PYG{+w}{ }\PYG{o}{\PYGZlt{}}\PYG{o}{\PYGZlt{}}\PYG{+w}{ }\PYG{n}{std}\PYG{o}{:}\PYG{o}{:}\PYG{n}{setprecision}\PYG{p}{(}\PYG{l+m+mi}{6}\PYG{p}{)}\PYG{+w}{ }\PYG{o}{\PYGZlt{}}\PYG{o}{\PYGZlt{}}\PYG{+w}{ }\PYG{l+s+sc}{\PYGZsq{}}\PYG{l+s+sc}{ }\PYG{l+s+sc}{\PYGZsq{}}\PYG{p}{;}
\PYG{+w}{    }\PYG{n}{std}\PYG{o}{:}\PYG{o}{:}\PYG{n}{cout}\PYG{+w}{ }\PYG{o}{\PYGZlt{}}\PYG{o}{\PYGZlt{}}\PYG{+w}{ }\PYG{n}{std}\PYG{o}{:}\PYG{o}{:}\PYG{n}{right}\PYG{+w}{ }\PYG{o}{\PYGZlt{}}\PYG{o}{\PYGZlt{}}\PYG{+w}{ }\PYG{n}{std}\PYG{o}{:}\PYG{o}{:}\PYG{n}{setw}\PYG{p}{(}\PYG{l+m+mi}{8}\PYG{p}{)}\PYG{+w}{ }\PYG{o}{\PYGZlt{}}\PYG{o}{\PYGZlt{}}\PYG{+w}{ }\PYG{n}{std}\PYG{o}{:}\PYG{o}{:}\PYG{n}{setfill}\PYG{p}{(}\PYG{l+s+sc}{\PYGZsq{}}\PYG{l+s+sc}{ }\PYG{l+s+sc}{\PYGZsq{}}\PYG{p}{)}
\PYG{+w}{              }\PYG{o}{\PYGZlt{}}\PYG{o}{\PYGZlt{}}\PYG{+w}{ }\PYG{n}{company}\PYG{p}{.}\PYG{n}{Employees}\PYG{+w}{ }\PYG{o}{\PYGZlt{}}\PYG{o}{\PYGZlt{}}\PYG{+w}{ }\PYG{l+s+sc}{\PYGZsq{}}\PYG{l+s+sc}{ }\PYG{l+s+sc}{\PYGZsq{}}\PYG{p}{;}
\PYG{+w}{    }\PYG{n}{std}\PYG{o}{:}\PYG{o}{:}\PYG{n}{cout}\PYG{+w}{ }\PYG{o}{\PYGZlt{}}\PYG{o}{\PYGZlt{}}\PYG{+w}{ }\PYG{n}{std}\PYG{o}{:}\PYG{o}{:}\PYG{n}{left}\PYG{+w}{ }\PYG{o}{\PYGZlt{}}\PYG{o}{\PYGZlt{}}\PYG{+w}{ }\PYG{n}{std}\PYG{o}{:}\PYG{o}{:}\PYG{n}{setw}\PYG{p}{(}\PYG{l+m+mi}{2}\PYG{p}{)}\PYG{+w}{ }\PYG{o}{\PYGZlt{}}\PYG{o}{\PYGZlt{}}\PYG{+w}{ }\PYG{n}{std}\PYG{o}{:}\PYG{o}{:}\PYG{n}{setfill}\PYG{p}{(}\PYG{l+s+sc}{\PYGZsq{}}\PYG{l+s+sc}{ }\PYG{l+s+sc}{\PYGZsq{}}\PYG{p}{)}
\PYG{+w}{              }\PYG{o}{\PYGZlt{}}\PYG{o}{\PYGZlt{}}\PYG{+w}{ }\PYG{n}{company}\PYG{p}{.}\PYG{n}{Country}
\PYG{+w}{              }\PYG{o}{\PYGZlt{}}\PYG{o}{\PYGZlt{}}\PYG{+w}{ }\PYG{n}{std}\PYG{o}{:}\PYG{o}{:}\PYG{n}{endl}\PYG{p}{;}
\PYG{+w}{  }\PYG{p}{\PYGZcb{}}
\PYG{p}{\PYGZcb{}}
\end{Verbatim}
\end{codeblockbg}

\end{itemize}


%%C070217 p

下面是调用该方法的示例。这里的数据来自维基百科（https://\allowbreak{}en.\allowbreak{}wikipedia.\allowbreak{}org/\allowbreak{}wiki/\allowbreak{}List\_\allowbreak{}of\_\allowbreak{}largest\_\allowbreak{}companies\_\allowbreak{}by\_\allowbreak{}revenue,asof2016）：


%%C070218 code

\begin{codeblockbg}
\begin{Verbatim}[commandchars=\\\{\},breaklines=true,breakanywhere=true,fontsize=\codefont,xleftmargin=17bp]
\PYG{n}{std}\PYG{o}{:}\PYG{o}{:}\PYG{n}{vector}\PYG{o}{\PYGZlt{}}\PYG{n}{Company}\PYG{o}{\PYGZgt{}}\PYG{+w}{ }\PYG{n}{companies}
\PYG{p}{\PYGZob{}}
\PYG{+w}{  }\PYG{p}{\PYGZob{}}\PYG{l+s}{\PYGZdq{}}\PYG{l+s}{Walmart}\PYG{l+s}{\PYGZdq{}}\PYG{n}{s}\PYG{p}{,}\PYG{+w}{ }\PYG{l+s}{\PYGZdq{}}\PYG{l+s}{Retail}\PYG{l+s}{\PYGZdq{}}\PYG{n}{s}\PYG{p}{,}\PYG{+w}{ }\PYG{l+m+mi}{482}\PYG{p}{,}\PYG{+w}{ }\PYG{n+nb}{false}\PYG{p}{,}\PYG{+w}{ }\PYG{l+m+mf}{0.71}\PYG{p}{,}
\PYG{+w}{   }\PYG{l+m+mi}{2300000}\PYG{p}{,}\PYG{+w}{ }\PYG{l+s}{\PYGZdq{}}\PYG{l+s}{US}\PYG{l+s}{\PYGZdq{}}\PYG{n}{s}\PYG{p}{\PYGZcb{}}\PYG{p}{,}
\PYG{+w}{  }\PYG{p}{\PYGZob{}}\PYG{l+s}{\PYGZdq{}}\PYG{l+s}{State Grid}\PYG{l+s}{\PYGZdq{}}\PYG{n}{s}\PYG{p}{,}\PYG{+w}{ }\PYG{l+s}{\PYGZdq{}}\PYG{l+s}{Electric utility}\PYG{l+s}{\PYGZdq{}}\PYG{n}{s}\PYG{p}{,}\PYG{+w}{ }\PYG{l+m+mi}{330}\PYG{p}{,}\PYG{+w}{ }\PYG{n+nb}{false}\PYG{p}{,}\PYG{+w}{ }\PYG{l+m+mf}{2.91}\PYG{p}{,}
\PYG{+w}{   }\PYG{l+m+mi}{927839}\PYG{p}{,}\PYG{+w}{ }\PYG{l+s}{\PYGZdq{}}\PYG{l+s}{China}\PYG{l+s}{\PYGZdq{}}\PYG{n}{s}\PYG{p}{\PYGZcb{}}\PYG{p}{,}
\PYG{+w}{  }\PYG{p}{\PYGZob{}}\PYG{l+s}{\PYGZdq{}}\PYG{l+s}{Saudi Aramco}\PYG{l+s}{\PYGZdq{}}\PYG{n}{s}\PYG{p}{,}\PYG{+w}{ }\PYG{l+s}{\PYGZdq{}}\PYG{l+s}{Oil and gas}\PYG{l+s}{\PYGZdq{}}\PYG{n}{s}\PYG{p}{,}\PYG{+w}{ }\PYG{l+m+mi}{311}\PYG{p}{,}\PYG{+w}{ }\PYG{n+nb}{true}\PYG{p}{,}\PYG{+w}{ }\PYG{l+m+mf}{40.11}\PYG{p}{,}
\PYG{+w}{   }\PYG{l+m+mi}{65266}\PYG{p}{,}\PYG{+w}{ }\PYG{l+s}{\PYGZdq{}}\PYG{l+s}{SA}\PYG{l+s}{\PYGZdq{}}\PYG{n}{s}\PYG{p}{\PYGZcb{}}\PYG{p}{,}
\PYG{+w}{  }\PYG{p}{\PYGZob{}}\PYG{l+s}{\PYGZdq{}}\PYG{l+s}{China National Petroleum}\PYG{l+s}{\PYGZdq{}}\PYG{n}{s}\PYG{p}{,}\PYG{+w}{ }\PYG{l+s}{\PYGZdq{}}\PYG{l+s}{Oil and gas}\PYG{l+s}{\PYGZdq{}}\PYG{n}{s}\PYG{p}{,}\PYG{+w}{ }\PYG{l+m+mi}{299}\PYG{p}{,}
\PYG{+w}{   }\PYG{n+nb}{false}\PYG{p}{,}\PYG{+w}{ }\PYG{l+m+mf}{30.21}\PYG{p}{,}\PYG{+w}{ }\PYG{l+m+mi}{1589508}\PYG{p}{,}\PYG{+w}{ }\PYG{l+s}{\PYGZdq{}}\PYG{l+s}{China}\PYG{l+s}{\PYGZdq{}}\PYG{n}{s}\PYG{p}{\PYGZcb{}}\PYG{p}{,}
\PYG{+w}{  }\PYG{p}{\PYGZob{}}\PYG{l+s}{\PYGZdq{}}\PYG{l+s}{Sinopec Group}\PYG{l+s}{\PYGZdq{}}\PYG{n}{s}\PYG{p}{,}\PYG{+w}{ }\PYG{l+s}{\PYGZdq{}}\PYG{l+s}{Oil and gas}\PYG{l+s}{\PYGZdq{}}\PYG{n}{s}\PYG{p}{,}\PYG{+w}{ }\PYG{l+m+mi}{294}\PYG{p}{,}\PYG{+w}{ }\PYG{n+nb}{false}\PYG{p}{,}\PYG{+w}{ }\PYG{l+m+mf}{34.11}\PYG{p}{,}
\PYG{+w}{   }\PYG{l+m+mi}{810538}\PYG{p}{,}\PYG{+w}{ }\PYG{l+s}{\PYGZdq{}}\PYG{l+s}{China}\PYG{l+s}{\PYGZdq{}}\PYG{n}{s}\PYG{p}{\PYGZcb{}}\PYG{p}{,}
\PYG{p}{\PYGZcb{}}\PYG{p}{;}
\PYG{n}{print\PYGZus{}companies}\PYG{p}{(}\PYG{n}{companies}\PYG{p}{)}\PYG{p}{;}
\end{Verbatim}
\end{codeblockbg}


%%C070219 p

这里的输出是表格形式，如下所示；你可以把它当作练习，试着加上表头和网格线：


%%C070220 code

\begin{codeblockbg}
\begin{Verbatim}[commandchars=\\\{\},breaklines=true,breakanywhere=true,fontsize=\codefont,xleftmargin=17bp]
\PYG{n}{Walmart}\PYG{+w}{                   }\PYG{n}{Retail}\PYG{+w}{            }\PYG{l+m+mi}{482}\PYG{+w}{  }\PYG{n+nb}{false}\PYG{+w}{  }\PYG{l+m+mf}{0.7}\PYG{+w}{  }\PYG{l+m+mi}{2300000}\PYG{+w}{ }\PYG{n}{US}
\PYG{n}{State}\PYG{+w}{ }\PYG{n}{Grid}\PYG{+w}{                }\PYG{n}{Electric}\PYG{+w}{ }\PYG{n}{utility}\PYG{+w}{  }\PYG{l+m+mi}{330}\PYG{+w}{  }\PYG{n+nb}{false}\PYG{+w}{  }\PYG{l+m+mf}{2.9}\PYG{+w}{   }\PYG{l+m+mi}{927839}\PYG{+w}{ }\PYG{n}{China}
\PYG{n}{Saudi}\PYG{+w}{ }\PYG{n}{Aramco}\PYG{+w}{              }\PYG{n}{Oil}\PYG{+w}{ }\PYG{k}{and}\PYG{+w}{ }\PYG{n}{gas}\PYG{+w}{       }\PYG{l+m+mi}{311}\PYG{+w}{  }\PYG{n+nb}{true}\PYG{+w}{  }\PYG{l+m+mf}{40.1}\PYG{+w}{    }\PYG{l+m+mi}{65266}\PYG{+w}{ }\PYG{n}{SA}
\PYG{n}{China}\PYG{+w}{ }\PYG{n}{National}\PYG{+w}{ }\PYG{n}{Petroleum}\PYG{+w}{  }\PYG{n}{Oil}\PYG{+w}{ }\PYG{k}{and}\PYG{+w}{ }\PYG{n}{gas}\PYG{+w}{       }\PYG{l+m+mi}{299}\PYG{+w}{  }\PYG{n+nb}{false}\PYG{+w}{ }\PYG{l+m+mf}{30.2}\PYG{+w}{  }\PYG{l+m+mi}{1589508}\PYG{+w}{ }\PYG{n}{China}
\PYG{n}{Sinopec}\PYG{+w}{ }\PYG{n}{Group}\PYG{+w}{             }\PYG{n}{Oil}\PYG{+w}{ }\PYG{k}{and}\PYG{+w}{ }\PYG{n}{gas}\PYG{+w}{       }\PYG{l+m+mi}{294}\PYG{+w}{  }\PYG{n+nb}{false}\PYG{+w}{ }\PYG{l+m+mf}{34.1}\PYG{+w}{   }\PYG{l+m+mi}{810538}\PYG{+w}{ }\PYG{n}{China}
\end{Verbatim}
\end{codeblockbg}


%%C070221 h3

\subsection{延伸阅读}


%%C070222 li

\begin{itemize}
\item \textit{从二进制文件读写原始数据}


%%C070223 li

\item \textit{使用货币I/O操纵器}


%%C070224 li

\item \textit{使用时间I/O操纵器}

\end{itemize}


%%C070225 h2

\section{使用货币I/O操纵器}


%%C070226 p

上一个案例中，我们看过一些可用于控制输入输出流的操纵器，讨论的都是数值和文本。本案例将介绍如何用标准操纵器读写货币值。


%%C070227 h3

\subsection{准备工作}


%%C070228 p

你应当熟悉区域设置（locale）以及如何为流设置区域设置。这个话题在\textit{为流使用本地化设置}案例中讨论过。建议先读那个案例再继续。


%%C070229 p

本案例讨论的操纵器位于\texttt{<\allowbreak{}iomanip>}头文件的\texttt{std}命名空间中。


%%C070230 h3

\subsection{具体做法}


%%C070231 p

要把货币值写入输出流，应当按以下步骤操作：


%%C070232 li

\begin{itemize}
\item 设置用于控制货币格式的区域设置：


%%C070233 code

\begin{codeblockbg}
\begin{Verbatim}[commandchars=\\\{\},breaklines=true,breakanywhere=true,fontsize=\codefont,xleftmargin=17bp]
\PYG{n}{std}\PYG{o}{:}\PYG{o}{:}\PYG{n}{cout}\PYG{p}{.}\PYG{n}{imbue}\PYG{p}{(}\PYG{n}{std}\PYG{o}{:}\PYG{o}{:}\PYG{n}{locale}\PYG{p}{(}\PYG{l+s}{\PYGZdq{}}\PYG{l+s}{en\PYGZus{}GB.utf8}\PYG{l+s}{\PYGZdq{}}\PYG{p}{)}\PYG{p}{)}\PYG{p}{;}
\end{Verbatim}
\end{codeblockbg}


%%C070234 li

\item 金额使用\texttt{long double}或\texttt{std::\allowbreak{}basic\_\allowbreak{}string}值：


%%C070235 code

\begin{codeblockbg}
\begin{Verbatim}[commandchars=\\\{\},breaklines=true,breakanywhere=true,fontsize=\codefont,xleftmargin=17bp]
\PYG{k+kt}{long}\PYG{+w}{ }\PYG{k+kt}{double}\PYG{+w}{ }\PYG{n}{mon}\PYG{+w}{ }\PYG{o}{=}\PYG{+w}{ }\PYG{l+m+mf}{12345.67}\PYG{p}{;}
\PYG{n}{std}\PYG{o}{:}\PYG{o}{:}\PYG{n}{string}\PYG{+w}{ }\PYG{n}{smon}\PYG{+w}{ }\PYG{o}{=}\PYG{+w}{ }\PYG{l+s}{\PYGZdq{}}\PYG{l+s}{12345.67}\PYG{l+s}{\PYGZdq{}}\PYG{p}{;}
\end{Verbatim}
\end{codeblockbg}


%%C070236 li

\item 使用只带一个参数（货币值）的\texttt{std::\allowbreak{}put\_\allowbreak{}money}操纵器，以货币符号（如果有）显示该值：


%%C070237 code

\begin{codeblockbg}
\begin{Verbatim}[commandchars=\\\{\},breaklines=true,breakanywhere=true,fontsize=\codefont,xleftmargin=17bp]
\PYG{n}{std}\PYG{o}{:}\PYG{o}{:}\PYG{n}{cout}\PYG{+w}{ }\PYG{o}{\PYGZlt{}}\PYG{o}{\PYGZlt{}}\PYG{+w}{ }\PYG{n}{std}\PYG{o}{:}\PYG{o}{:}\PYG{n}{showbase}\PYG{+w}{ }\PYG{o}{\PYGZlt{}}\PYG{o}{\PYGZlt{}}\PYG{+w}{ }\PYG{n}{std}\PYG{o}{:}\PYG{o}{:}\PYG{n}{put\PYGZus{}money}\PYG{p}{(}\PYG{n}{mon}\PYG{p}{)}
\PYG{+w}{          }\PYG{o}{\PYGZlt{}}\PYG{o}{\PYGZlt{}}\PYG{+w}{ }\PYG{n}{std}\PYG{o}{:}\PYG{o}{:}\PYG{n}{endl}\PYG{p}{;}\PYG{+w}{ }\PYG{c+c1}{// £123.46}
\PYG{n}{std}\PYG{o}{:}\PYG{o}{:}\PYG{n}{cout}\PYG{+w}{ }\PYG{o}{\PYGZlt{}}\PYG{o}{\PYGZlt{}}\PYG{+w}{ }\PYG{n}{std}\PYG{o}{:}\PYG{o}{:}\PYG{n}{showbase}\PYG{+w}{ }\PYG{o}{\PYGZlt{}}\PYG{o}{\PYGZlt{}}\PYG{+w}{ }\PYG{n}{std}\PYG{o}{:}\PYG{o}{:}\PYG{n}{put\PYGZus{}money}\PYG{p}{(}\PYG{n}{smon}\PYG{p}{)}
\PYG{+w}{          }\PYG{o}{\PYGZlt{}}\PYG{o}{\PYGZlt{}}\PYG{+w}{ }\PYG{n}{std}\PYG{o}{:}\PYG{o}{:}\PYG{n}{endl}\PYG{p}{;}\PYG{+w}{ }\PYG{c+c1}{// £123.46}
\end{Verbatim}
\end{codeblockbg}


%%C070238 li

\item 使用带两个参数的\texttt{std::\allowbreak{}put\_\allowbreak{}money}：货币值和设为\texttt{true}的布尔标志，表示使用国际货币字符串：


%%C070239 code

\begin{codeblockbg}
\begin{Verbatim}[commandchars=\\\{\},breaklines=true,breakanywhere=true,fontsize=\codefont,xleftmargin=17bp]
\PYG{n}{std}\PYG{o}{:}\PYG{o}{:}\PYG{n}{cout}\PYG{+w}{ }\PYG{o}{\PYGZlt{}}\PYG{o}{\PYGZlt{}}\PYG{+w}{ }\PYG{n}{std}\PYG{o}{:}\PYG{o}{:}\PYG{n}{showbase}\PYG{+w}{ }\PYG{o}{\PYGZlt{}}\PYG{o}{\PYGZlt{}}\PYG{+w}{ }\PYG{n}{std}\PYG{o}{:}\PYG{o}{:}\PYG{n}{put\PYGZus{}money}\PYG{p}{(}\PYG{n}{mon}\PYG{p}{,}\PYG{+w}{ }\PYG{n+nb}{true}\PYG{p}{)}
\PYG{+w}{          }\PYG{o}{\PYGZlt{}}\PYG{o}{\PYGZlt{}}\PYG{+w}{ }\PYG{n}{std}\PYG{o}{:}\PYG{o}{:}\PYG{n}{endl}\PYG{p}{;}\PYG{+w}{ }\PYG{c+c1}{// GBP 123.46}
\PYG{n}{std}\PYG{o}{:}\PYG{o}{:}\PYG{n}{cout}\PYG{+w}{ }\PYG{o}{\PYGZlt{}}\PYG{o}{\PYGZlt{}}\PYG{+w}{ }\PYG{n}{std}\PYG{o}{:}\PYG{o}{:}\PYG{n}{showbase}\PYG{+w}{ }\PYG{o}{\PYGZlt{}}\PYG{o}{\PYGZlt{}}\PYG{+w}{ }\PYG{n}{std}\PYG{o}{:}\PYG{o}{:}\PYG{n}{put\PYGZus{}money}\PYG{p}{(}\PYG{n}{smon}\PYG{p}{,}\PYG{+w}{ }\PYG{n+nb}{true}\PYG{p}{)}
\PYG{+w}{          }\PYG{o}{\PYGZlt{}}\PYG{o}{\PYGZlt{}}\PYG{+w}{ }\PYG{n}{std}\PYG{o}{:}\PYG{o}{:}\PYG{n}{endl}\PYG{p}{;}\PYG{+w}{ }\PYG{c+c1}{// GBP 123.46}
\end{Verbatim}
\end{codeblockbg}

\end{itemize}


%%C070240 p

要从输入流读取货币值，应当按以下步骤操作：


%%C070241 li

\begin{itemize}
\item 设置用于控制货币格式的区域设置：


%%C070242 code

\begin{codeblockbg}
\begin{Verbatim}[commandchars=\\\{\},breaklines=true,breakanywhere=true,fontsize=\codefont,xleftmargin=17bp]
\PYG{n}{std}\PYG{o}{:}\PYG{o}{:}\PYG{n}{istringstream}\PYG{+w}{ }\PYG{n}{stext}\PYG{p}{(}\PYG{l+s}{\PYGZdq{}}\PYG{l+s}{\PYGZdl{}123.45 123.45 USD}\PYG{l+s}{\PYGZdq{}}\PYG{p}{)}\PYG{p}{;}
\PYG{n}{stext}\PYG{p}{.}\PYG{n}{imbue}\PYG{p}{(}\PYG{n}{std}\PYG{o}{:}\PYG{o}{:}\PYG{n}{locale}\PYG{p}{(}\PYG{l+s}{\PYGZdq{}}\PYG{l+s}{en\PYGZus{}US.utf8}\PYG{l+s}{\PYGZdq{}}\PYG{p}{)}\PYG{p}{)}\PYG{p}{;}
\end{Verbatim}
\end{codeblockbg}


%%C070243 li

\item 用\texttt{long double}或\texttt{std::\allowbreak{}basic\_\allowbreak{}string}值从输入流读取金额：


%%C070244 code

\begin{codeblockbg}
\begin{Verbatim}[commandchars=\\\{\},breaklines=true,breakanywhere=true,fontsize=\codefont,xleftmargin=17bp]
\PYG{k+kt}{long}\PYG{+w}{ }\PYG{k+kt}{double}\PYG{+w}{ }\PYG{n}{v1}\PYG{p}{;}
\PYG{n}{std}\PYG{o}{:}\PYG{o}{:}\PYG{n}{string}\PYG{+w}{ }\PYG{n}{v2}\PYG{p}{;}
\end{Verbatim}
\end{codeblockbg}


%%C070245 li

\item 如果输入流中可能出现货币符号，使用只带一个参数（用于存放货币值的变量）的\texttt{std::\allowbreak{}get\_\allowbreak{}money()}：


%%C070246 code

\begin{codeblockbg}
\begin{Verbatim}[commandchars=\\\{\},breaklines=true,breakanywhere=true,fontsize=\codefont,xleftmargin=17bp]
\PYG{n}{stext}\PYG{+w}{ }\PYG{o}{\PYGZgt{}}\PYG{o}{\PYGZgt{}}\PYG{+w}{ }\PYG{n}{std}\PYG{o}{:}\PYG{o}{:}\PYG{n}{get\PYGZus{}money}\PYG{p}{(}\PYG{n}{v1}\PYG{p}{)}\PYG{+w}{ }\PYG{o}{\PYGZgt{}}\PYG{o}{\PYGZgt{}}\PYG{+w}{ }\PYG{n}{std}\PYG{o}{:}\PYG{o}{:}\PYG{n}{get\PYGZus{}money}\PYG{p}{(}\PYG{n}{v2}\PYG{p}{)}\PYG{p}{;}
\PYG{c+c1}{// v1 = 12345, v2 = \PYGZdq{}12345\PYGZdq{}}
\end{Verbatim}
\end{codeblockbg}


%%C070247 li

\item 使用带两个参数的\texttt{std::\allowbreak{}get\_\allowbreak{}money()}：用于存放货币值的变量和设为\texttt{true}的布尔标志，表示存在国际货币字符串：


%%C070248 code

\begin{codeblockbg}
\begin{Verbatim}[commandchars=\\\{\},breaklines=true,breakanywhere=true,fontsize=\codefont,xleftmargin=17bp]
\PYG{n}{stext}\PYG{+w}{ }\PYG{o}{\PYGZgt{}}\PYG{o}{\PYGZgt{}}\PYG{+w}{ }\PYG{n}{std}\PYG{o}{:}\PYG{o}{:}\PYG{n}{get\PYGZus{}money}\PYG{p}{(}\PYG{n}{v1}\PYG{p}{,}\PYG{+w}{ }\PYG{n+nb}{true}\PYG{p}{)}\PYG{+w}{ }\PYG{o}{\PYGZgt{}}\PYG{o}{\PYGZgt{}}\PYG{+w}{ }\PYG{n}{std}\PYG{o}{:}\PYG{o}{:}\PYG{n}{get\PYGZus{}money}\PYG{p}{(}\PYG{n}{v2}\PYG{p}{,}\PYG{+w}{ }\PYG{n+nb}{true}\PYG{p}{)}\PYG{p}{;}
\PYG{c+c1}{// v1 = 0, v2 = \PYGZdq{}12345\PYGZdq{}}
\end{Verbatim}
\end{codeblockbg}

\end{itemize}


%%C070249 h3

\subsection{工作原理}


%%C070250 p

\texttt{put\_\allowbreak{}money()}和\texttt{get\_\allowbreak{}money()}操纵器非常相似：它们都是函数模板，接受一个参数，表示要写入输出流的货币值，或用于存放从输入流读入的货币值的变量，还接受第二个可选参数，用于指示是否使用国际货币字符串。默认选择货币符号（如果可用）。\texttt{put\_\allowbreak{}money()}利用\texttt{std::\allowbreak{}money\_\allowbreak{}put()}facet的设置输出货币值，\texttt{get\_\allowbreak{}money()}利用\texttt{std::\allowbreak{}money\_\allowbreak{}get()}facet解析货币值。两个操纵器函数模板都返回未指明类型的对象。这些函数不抛出异常：


%%C070251 code

\begin{codeblockbg}
\begin{Verbatim}[commandchars=\\\{\},breaklines=true,breakanywhere=true,fontsize=\codefont,xleftmargin=17bp]
\PYG{k}{template}\PYG{+w}{ }\PYG{o}{\PYGZlt{}}\PYG{k}{class}\PYG{+w}{ }\PYG{n+nc}{MoneyT}\PYG{o}{\PYGZgt{}}
\PYG{c+cm}{/*未指明*/}\PYG{+w}{ }\PYG{n}{put\PYGZus{}money}\PYG{p}{(}\PYG{k}{const}\PYG{+w}{ }\PYG{n}{MoneyT}\PYG{o}{\PYGZam{}}\PYG{+w}{ }\PYG{n}{mon}\PYG{p}{,}\PYG{+w}{ }\PYG{k+kt}{bool}\PYG{+w}{ }\PYG{n}{intl}\PYG{+w}{ }\PYG{o}{=}\PYG{+w}{ }\PYG{n+nb}{false}\PYG{p}{)}\PYG{p}{;}
\PYG{k}{template}\PYG{+w}{ }\PYG{o}{\PYGZlt{}}\PYG{k}{class}\PYG{+w}{ }\PYG{n+nc}{MoneyT}\PYG{o}{\PYGZgt{}}
\PYG{c+cm}{/*未指明*/}\PYG{+w}{ }\PYG{n}{get\PYGZus{}money}\PYG{p}{(}\PYG{n}{MoneyT}\PYG{o}{\PYGZam{}}\PYG{+w}{ }\PYG{n}{mon}\PYG{p}{,}\PYG{+w}{ }\PYG{k+kt}{bool}\PYG{+w}{ }\PYG{n}{intl}\PYG{+w}{ }\PYG{o}{=}\PYG{+w}{ }\PYG{n+nb}{false}\PYG{p}{)}\PYG{p}{;}
\end{Verbatim}
\end{codeblockbg}


%%C070252 p

两个操纵器函数都要求货币值为\texttt{long double}或\texttt{std::\allowbreak{}basic\_\allowbreak{}string}。


%%C070253 note

\begin{myNote}[tip]
不过要注意，货币值是按所用区域设置定义的货币最小单位的整数形式存储的。以美元为例，\$100.00存储为10000.0，而1美分即\$0.01存储为1.0。
\end{myNote}


%%C070254 p

向输出流写入货币值时，如果要显示货币符号或国际货币字符串，就要使用\texttt{std::\allowbreak{}showbase}操纵器。它通常用于指示数值基数的前缀（如十六进制的0x），但对货币值而言，它指示是否显示货币符号或货币字符串：


%%C070255 code

\begin{codeblockbg}
\begin{Verbatim}[commandchars=\\\{\},breaklines=true,breakanywhere=true,fontsize=\codefont,xleftmargin=17bp]
\PYG{n}{std}\PYG{o}{:}\PYG{o}{:}\PYG{n}{cout}\PYG{+w}{ }\PYG{o}{\PYGZlt{}}\PYG{o}{\PYGZlt{}}\PYG{+w}{ }\PYG{n}{std}\PYG{o}{:}\PYG{o}{:}\PYG{n}{put\PYGZus{}money}\PYG{p}{(}\PYG{l+m+mf}{12345.67}\PYG{p}{)}
\PYG{+w}{          }\PYG{o}{\PYGZlt{}}\PYG{o}{\PYGZlt{}}\PYG{+w}{ }\PYG{n}{std}\PYG{o}{:}\PYG{o}{:}\PYG{n}{endl}\PYG{p}{;}\PYG{+w}{ }\PYG{c+c1}{// 输出123.46}
\PYG{n}{std}\PYG{o}{:}\PYG{o}{:}\PYG{n}{cout}\PYG{+w}{ }\PYG{o}{\PYGZlt{}}\PYG{o}{\PYGZlt{}}\PYG{+w}{ }\PYG{n}{std}\PYG{o}{:}\PYG{o}{:}\PYG{n}{showbase}\PYG{+w}{ }\PYG{o}{\PYGZlt{}}\PYG{o}{\PYGZlt{}}\PYG{+w}{ }\PYG{n}{std}\PYG{o}{:}\PYG{o}{:}\PYG{n}{put\PYGZus{}money}\PYG{p}{(}\PYG{l+m+mf}{12345.67}\PYG{p}{)}
\PYG{+w}{          }\PYG{o}{\PYGZlt{}}\PYG{o}{\PYGZlt{}}\PYG{+w}{ }\PYG{n}{std}\PYG{o}{:}\PYG{o}{:}\PYG{n}{endl}\PYG{p}{;}\PYG{+w}{ }\PYG{c+c1}{// 输出£123.46}
\end{Verbatim}
\end{codeblockbg}


%%C070256 h3

\subsection{延伸阅读}


%%C070257 li

\begin{itemize}
\item \textit{用I/O操纵器控制流的输出}


%%C070258 li

\item \textit{使用时间I/O操纵器}

\end{itemize}


%%C070259 h2

\section{使用时间I/O操纵器}


%%C070260 p

与上一个案例讨论的货币输入输出操纵器类似，C++11标准提供了用于控制时间值写入流和从流中读取的操纵器，这类时间值以保存日历日期和时间的\texttt{std::\allowbreak{}tm}对象表示。本案例将学习如何使用这些时间操纵器。


%%C070261 h3

\subsection{准备工作}


%%C070262 p

时间I/O操纵器使用的时间值以\texttt{std::\allowbreak{}tm}值表示。你应当熟悉\texttt{<\allowbreak{}ctime>}头文件中的这个结构。


%%C070263 p

你还应当熟悉区域设置（locale）以及如何为流设置区域设置。这个话题在\textit{为流使用本地化设置}案例中讨论过。建议先读那个案例再继续。


%%C070264 p

本案例讨论的操纵器位于\texttt{<\allowbreak{}iomanip>}头文件的\texttt{std}命名空间中。


%%C070265 h3

\subsection{具体做法}


%%C070266 p

要把时间值写入输出流，应当执行以下步骤：


%%C070267 p

获取与给定时间相对应的日历日期和时间值。有多种


%%C070268 p

方法可以做到这一点。以下是把当前时间转换为以日历日期和时间表示的本地时间的示例：

%%C070269 code

\begin{codeblockbg}
\begin{Verbatim}[commandchars=\\\{\},breaklines=true,breakanywhere=true,fontsize=\codefont,xleftmargin=17bp]
\PYG{k}{auto}\PYG{+w}{ }\PYG{n}{now}\PYG{+w}{ }\PYG{o}{=}\PYG{+w}{ }\PYG{n}{std}\PYG{o}{:}\PYG{o}{:}\PYG{n}{chrono}\PYG{o}{:}\PYG{o}{:}\PYG{n}{system\PYGZus{}clock}\PYG{o}{:}\PYG{o}{:}\PYG{n}{now}\PYG{p}{(}\PYG{p}{)}\PYG{p}{;}
\PYG{k}{auto}\PYG{+w}{ }\PYG{n}{stime}\PYG{+w}{ }\PYG{o}{=}\PYG{+w}{ }\PYG{n}{std}\PYG{o}{:}\PYG{o}{:}\PYG{n}{chrono}\PYG{o}{:}\PYG{o}{:}\PYG{n}{system\PYGZus{}clock}\PYG{o}{:}\PYG{o}{:}\PYG{n}{to\PYGZus{}time\PYGZus{}t}\PYG{p}{(}\PYG{n}{now}\PYG{p}{)}\PYG{p}{;}
\PYG{k}{auto}\PYG{+w}{ }\PYG{n}{ltime}\PYG{+w}{ }\PYG{o}{=}\PYG{+w}{ }\PYG{n}{std}\PYG{o}{:}\PYG{o}{:}\PYG{n}{localtime}\PYG{p}{(}\PYG{o}{\PYGZam{}}\PYG{n}{stime}\PYG{p}{)}\PYG{p}{;}
\PYG{k}{auto}\PYG{+w}{ }\PYG{n}{ttime}\PYG{+w}{ }\PYG{o}{=}\PYG{+w}{ }\PYG{n}{std}\PYG{o}{:}\PYG{o}{:}\PYG{n}{time}\PYG{p}{(}\PYG{n+nb}{nullptr}\PYG{p}{)}\PYG{p}{;}
\PYG{k}{auto}\PYG{+w}{ }\PYG{n}{ltime}\PYG{+w}{ }\PYG{o}{=}\PYG{+w}{ }\PYG{n}{std}\PYG{o}{:}\PYG{o}{:}\PYG{n}{localtime}\PYG{p}{(}\PYG{o}{\PYGZam{}}\PYG{n}{ttime}\PYG{p}{)}\PYG{p}{;}
\end{Verbatim}
\end{codeblockbg}


%%C070270 p

使用\texttt{std::\allowbreak{}put\_\allowbreak{}time()}，传入指向表示


%%C070271 p

日历日期与时间的\texttt{std::\allowbreak{}tm}对象的指针，再传入指向以空字符结尾的字符串的指针来表示格式。C++11标准提供了很长的可用格式列表，可在http://\allowbreak{}en.\allowbreak{}cppreference.\allowbreak{}com/\allowbreak{}w/\allowbreak{}cpp/\allowbreak{}io/\allowbreak{}manip/\allowbreak{}put\_\allowbreak{}time查阅。


%%C070272 p

若要按某个特定


%%C070273 p

区域设置的规则写出标准日期与时间字符串，先调用\texttt{imbue()}为流设置区域设置，再使用\texttt{std::\allowbreak{}put\_\allowbreak{}time()}操纵器：


%%C070274 code

\begin{codeblockbg}
\begin{Verbatim}[commandchars=\\\{\},breaklines=true,breakanywhere=true,fontsize=\codefont,xleftmargin=17bp]
\PYG{n}{std}\PYG{o}{:}\PYG{o}{:}\PYG{n}{cout}\PYG{p}{.}\PYG{n}{imbue}\PYG{p}{(}\PYG{n}{std}\PYG{o}{:}\PYG{o}{:}\PYG{n}{locale}\PYG{p}{(}\PYG{l+s}{\PYGZdq{}}\PYG{l+s}{en\PYGZus{}GB.utf8}\PYG{l+s}{\PYGZdq{}}\PYG{p}{)}\PYG{p}{)}\PYG{p}{;}
\PYG{n}{std}\PYG{o}{:}\PYG{o}{:}\PYG{n}{cout}\PYG{+w}{ }\PYG{o}{\PYGZlt{}}\PYG{o}{\PYGZlt{}}\PYG{+w}{ }\PYG{n}{std}\PYG{o}{:}\PYG{o}{:}\PYG{n}{put\PYGZus{}time}\PYG{p}{(}\PYG{n}{ltime}\PYG{p}{,}\PYG{+w}{ }\PYG{l+s}{\PYGZdq{}}\PYG{l+s}{\PYGZpc{}c}\PYG{l+s}{\PYGZdq{}}\PYG{p}{)}\PYG{+w}{ }\PYG{o}{\PYGZlt{}}\PYG{o}{\PYGZlt{}}\PYG{+w}{ }\PYG{n}{std}\PYG{o}{:}\PYG{o}{:}\PYG{n}{endl}\PYG{p}{;}
\PYG{c+c1}{// Sun 04 Dec 2016 05:26:47 JST}
\end{Verbatim}
\end{codeblockbg}


%%C070275 p

以下是受支持的时间格式示例：


%%C070276 li

\begin{itemize}
\item ISO 8601日期格式\texttt{“\%F”}或\texttt{“\%Y-\%m-\%d”}：


%%C070277 code

\begin{codeblockbg}
\begin{Verbatim}[commandchars=\\\{\},breaklines=true,breakanywhere=true,fontsize=\codefont,xleftmargin=17bp]
\PYG{+w}{        }\PYG{n}{std}\PYG{o}{:}\PYG{o}{:}\PYG{n}{cout}\PYG{+w}{ }\PYG{o}{\PYGZlt{}}\PYG{o}{\PYGZlt{}}\PYG{+w}{ }\PYG{n}{std}\PYG{o}{:}\PYG{o}{:}\PYG{n}{put\PYGZus{}time}\PYG{p}{(}\PYG{n}{ltime}\PYG{p}{,}\PYG{+w}{ }\PYG{l+s}{\PYGZdq{}}\PYG{l+s}{\PYGZpc{}F}\PYG{l+s}{\PYGZdq{}}\PYG{p}{)}\PYG{+w}{ }\PYG{o}{\PYGZlt{}}\PYG{o}{\PYGZlt{}}\PYG{+w}{ }\PYG{n}{std}\PYG{o}{:}\PYG{o}{:}\PYG{n}{endl}\PYG{p}{;}
\PYG{c+c1}{// 2016\PYGZhy{}12\PYGZhy{}04}
\end{Verbatim}
\end{codeblockbg}


%%C070278 li

\item ISO 8601时间格式\texttt{“\%T”}：


%%C070279 code

\begin{codeblockbg}
\begin{Verbatim}[commandchars=\\\{\},breaklines=true,breakanywhere=true,fontsize=\codefont,xleftmargin=17bp]
\PYG{n}{std}\PYG{o}{:}\PYG{o}{:}\PYG{n}{cout}\PYG{+w}{ }\PYG{o}{\PYGZlt{}}\PYG{o}{\PYGZlt{}}\PYG{+w}{ }\PYG{n}{std}\PYG{o}{:}\PYG{o}{:}\PYG{n}{put\PYGZus{}time}\PYG{p}{(}\PYG{n}{ltime}\PYG{p}{,}\PYG{+w}{ }\PYG{l+s}{\PYGZdq{}}\PYG{l+s}{\PYGZpc{}T}\PYG{l+s}{\PYGZdq{}}\PYG{p}{)}\PYG{+w}{ }\PYG{o}{\PYGZlt{}}\PYG{o}{\PYGZlt{}}\PYG{+w}{ }\PYG{n}{std}\PYG{o}{:}\PYG{o}{:}\PYG{n}{endl}\PYG{p}{;}
\PYG{c+c1}{// 05:26:47}
\end{Verbatim}
\end{codeblockbg}


%%C070280 li

\item ISO 8601的日期时间组合UTC格式\texttt{“\%FT\%T\%z”}：


%%C070281 code

\begin{codeblockbg}
\begin{Verbatim}[commandchars=\\\{\},breaklines=true,breakanywhere=true,fontsize=\codefont,xleftmargin=17bp]
\PYG{n}{std}\PYG{o}{:}\PYG{o}{:}\PYG{n}{cout}\PYG{+w}{ }\PYG{o}{\PYGZlt{}}\PYG{o}{\PYGZlt{}}\PYG{+w}{ }\PYG{n}{std}\PYG{o}{:}\PYG{o}{:}\PYG{n}{put\PYGZus{}time}\PYG{p}{(}\PYG{n}{ltime}\PYG{p}{,}\PYG{+w}{ }\PYG{l+s}{\PYGZdq{}}\PYG{l+s}{\PYGZpc{}FT\PYGZpc{}T\PYGZpc{}z}\PYG{l+s}{\PYGZdq{}}\PYG{p}{)}\PYG{+w}{ }\PYG{o}{\PYGZlt{}}\PYG{o}{\PYGZlt{}}\PYG{+w}{ }\PYG{n}{std}\PYG{o}{:}\PYG{o}{:}\PYG{n}{endl}\PYG{p}{;}
\PYG{c+c1}{// 2016\PYGZhy{}12\PYGZhy{}04T05:26:47+0900}
\end{Verbatim}
\end{codeblockbg}


%%C070282 li

\item ISO 8601周格式\texttt{“\%Y-\allowbreak{}W\%V”}：


%%C070283 code

\begin{codeblockbg}
\begin{Verbatim}[commandchars=\\\{\},breaklines=true,breakanywhere=true,fontsize=\codefont,xleftmargin=17bp]
\PYG{n}{std}\PYG{o}{:}\PYG{o}{:}\PYG{n}{cout}\PYG{+w}{ }\PYG{o}{\PYGZlt{}}\PYG{o}{\PYGZlt{}}\PYG{+w}{ }\PYG{n}{std}\PYG{o}{:}\PYG{o}{:}\PYG{n}{put\PYGZus{}time}\PYG{p}{(}\PYG{n}{ltime}\PYG{p}{,}\PYG{+w}{ }\PYG{l+s}{\PYGZdq{}}\PYG{l+s}{\PYGZpc{}Y\PYGZhy{}W\PYGZpc{}V}\PYG{l+s}{\PYGZdq{}}\PYG{p}{)}\PYG{+w}{ }\PYG{o}{\PYGZlt{}}\PYG{o}{\PYGZlt{}}\PYG{+w}{ }\PYG{n}{std}\PYG{o}{:}\PYG{o}{:}\PYG{n}{endl}\PYG{p}{;}
\PYG{c+c1}{// 2016\PYGZhy{}W48}
\end{Verbatim}
\end{codeblockbg}


%%C070284 li

\item ISO 8601带周数的日期格式\texttt{“\%Y-\allowbreak{}W\%V-\%u”}：


%%C070285 code

\begin{codeblockbg}
\begin{Verbatim}[commandchars=\\\{\},breaklines=true,breakanywhere=true,fontsize=\codefont,xleftmargin=17bp]
\PYG{n}{std}\PYG{o}{:}\PYG{o}{:}\PYG{n}{cout}\PYG{+w}{ }\PYG{o}{\PYGZlt{}}\PYG{o}{\PYGZlt{}}\PYG{+w}{ }\PYG{n}{std}\PYG{o}{:}\PYG{o}{:}\PYG{n}{put\PYGZus{}time}\PYG{p}{(}\PYG{n}{ltime}\PYG{p}{,}\PYG{+w}{ }\PYG{l+s}{\PYGZdq{}}\PYG{l+s}{\PYGZpc{}Y\PYGZhy{}W\PYGZpc{}V\PYGZhy{}\PYGZpc{}u}\PYG{l+s}{\PYGZdq{}}\PYG{p}{)}\PYG{+w}{ }\PYG{o}{\PYGZlt{}}\PYG{o}{\PYGZlt{}}\PYG{+w}{ }\PYG{n}{std}\PYG{o}{:}\PYG{o}{:}\PYG{n}{endl}\PYG{p}{;}
\PYG{c+c1}{// 2016\PYGZhy{}W48\PYGZhy{}7}
\end{Verbatim}
\end{codeblockbg}


%%C070286 li

\item ISO 8601序数日期格式\texttt{“\%Y-\%j”}：


%%C070287 code

\begin{codeblockbg}
\begin{Verbatim}[commandchars=\\\{\},breaklines=true,breakanywhere=true,fontsize=\codefont,xleftmargin=17bp]
\PYG{n}{std}\PYG{o}{:}\PYG{o}{:}\PYG{n}{cout}\PYG{+w}{ }\PYG{o}{\PYGZlt{}}\PYG{o}{\PYGZlt{}}\PYG{+w}{ }\PYG{n}{std}\PYG{o}{:}\PYG{o}{:}\PYG{n}{put\PYGZus{}time}\PYG{p}{(}\PYG{n}{ltime}\PYG{p}{,}\PYG{+w}{ }\PYG{l+s}{\PYGZdq{}}\PYG{l+s}{\PYGZpc{}Y\PYGZhy{}\PYGZpc{}j}\PYG{l+s}{\PYGZdq{}}\PYG{p}{)}\PYG{+w}{ }\PYG{o}{\PYGZlt{}}\PYG{o}{\PYGZlt{}}\PYG{+w}{ }\PYG{n}{std}\PYG{o}{:}\PYG{o}{:}\PYG{n}{endl}\PYG{p}{;}
\PYG{c+c1}{// 2016\PYGZhy{}339}
\end{Verbatim}
\end{codeblockbg}

\end{itemize}


%%C070288 p

要从输入流读取时间值，应按以下步骤操作：


%%C070289 p

声明一个\texttt{std::\allowbreak{}tm}类型的对象，用来保存从输入流


%%C070290 p

中读取的时间值：


%%C070291 code

\begin{codeblockbg}
\begin{Verbatim}[commandchars=\\\{\},breaklines=true,breakanywhere=true,fontsize=\codefont,xleftmargin=17bp]
\PYG{k}{auto}\PYG{+w}{ }\PYG{n}{time}\PYG{+w}{ }\PYG{o}{=}\PYG{+w}{ }\PYG{n}{std}\PYG{o}{:}\PYG{o}{:}\PYG{n}{tm}\PYG{+w}{ }\PYG{p}{\PYGZob{}}\PYG{p}{\PYGZcb{}}\PYG{p}{;}
\end{Verbatim}
\end{codeblockbg}


%%C070292 p

使用\texttt{std::\allowbreak{}get\_\allowbreak{}time()}，传入指向将要保存时间值的\texttt{std::\allowbreak{}tm}对象的指针


%%C070293 p

以及指向以空字符结尾的字符串的指针来表示格式。可能的格式列表可在http://\allowbreak{}en.\allowbreak{}cppreference.\allowbreak{}com/\allowbreak{}w/\allowbreak{}cpp/\allowbreak{}io/\allowbreak{}manip/\allowbreak{}get\_\allowbreak{}time查阅。下面的示例解析一个ISO 8601日期时间组合值：


%%C070294 code

\begin{codeblockbg}
\begin{Verbatim}[commandchars=\\\{\},breaklines=true,breakanywhere=true,fontsize=\codefont,xleftmargin=17bp]
\PYG{n}{std}\PYG{o}{:}\PYG{o}{:}\PYG{n}{istringstream}\PYG{+w}{ }\PYG{n}{stext}\PYG{p}{(}\PYG{l+s}{\PYGZdq{}}\PYG{l+s}{2016\PYGZhy{}12\PYGZhy{}04T05:26:47+0900}\PYG{l+s}{\PYGZdq{}}\PYG{p}{)}\PYG{p}{;}
\PYG{n}{stext}\PYG{+w}{ }\PYG{o}{\PYGZgt{}}\PYG{o}{\PYGZgt{}}\PYG{+w}{ }\PYG{n}{std}\PYG{o}{:}\PYG{o}{:}\PYG{n}{get\PYGZus{}time}\PYG{p}{(}\PYG{o}{\PYGZam{}}\PYG{n}{time}\PYG{p}{,}\PYG{+w}{ }\PYG{l+s}{\PYGZdq{}}\PYG{l+s}{\PYGZpc{}Y\PYGZhy{}\PYGZpc{}m\PYGZhy{}\PYGZpc{}dT\PYGZpc{}H:\PYGZpc{}M:\PYGZpc{}S}\PYG{l+s}{\PYGZdq{}}\PYG{p}{)}\PYG{p}{;}
\PYG{k}{if}\PYG{+w}{ }\PYG{p}{(}\PYG{o}{!}\PYG{n}{stext}\PYG{p}{.}\PYG{n}{fail}\PYG{p}{(}\PYG{p}{)}\PYG{p}{)}\PYG{+w}{ }\PYG{p}{\PYGZob{}}\PYG{+w}{ }\PYG{c+cm}{/* 做点什么 */}\PYG{+w}{ }\PYG{p}{\PYGZcb{}}
\end{Verbatim}
\end{codeblockbg}


%%C070295 p

若要按某个特定


%%C070296 p

区域设置的规则读取标准日期与时间字符串，先调用\texttt{imblue()}为流设置区域设置，再使用\texttt{std::\allowbreak{}get\_\allowbreak{}time()}操纵器：


%%C070297 code

\begin{codeblockbg}
\begin{Verbatim}[commandchars=\\\{\},breaklines=true,breakanywhere=true,fontsize=\codefont,xleftmargin=17bp]
\PYG{n}{std}\PYG{o}{:}\PYG{o}{:}\PYG{n}{istringstream}\PYG{+w}{ }\PYG{n}{stext}\PYG{p}{(}\PYG{l+s}{\PYGZdq{}}\PYG{l+s}{Sun 04 Dec 2016 05:35:30 JST}\PYG{l+s}{\PYGZdq{}}\PYG{p}{)}\PYG{p}{;}
\PYG{n}{stext}\PYG{p}{.}\PYG{n}{imbue}\PYG{p}{(}\PYG{n}{std}\PYG{o}{:}\PYG{o}{:}\PYG{n}{locale}\PYG{p}{(}\PYG{l+s}{\PYGZdq{}}\PYG{l+s}{en\PYGZus{}GB.utf8}\PYG{l+s}{\PYGZdq{}}\PYG{p}{)}\PYG{p}{)}\PYG{p}{;}
\PYG{n}{stext}\PYG{+w}{ }\PYG{o}{\PYGZgt{}}\PYG{o}{\PYGZgt{}}\PYG{+w}{ }\PYG{n}{std}\PYG{o}{:}\PYG{o}{:}\PYG{n}{get\PYGZus{}time}\PYG{p}{(}\PYG{o}{\PYGZam{}}\PYG{n}{time}\PYG{p}{,}\PYG{+w}{ }\PYG{l+s}{\PYGZdq{}}\PYG{l+s}{\PYGZpc{}c}\PYG{l+s}{\PYGZdq{}}\PYG{p}{)}\PYG{p}{;}
\PYG{k}{if}\PYG{+w}{ }\PYG{p}{(}\PYG{n}{stext}\PYG{p}{.}\PYG{n}{fail}\PYG{p}{(}\PYG{p}{)}\PYG{p}{)}\PYG{+w}{ }\PYG{p}{\PYGZob{}}\PYG{+w}{ }\PYG{c+cm}{/* 做别的事 */}\PYG{+w}{ }\PYG{p}{\PYGZcb{}}
\end{Verbatim}
\end{codeblockbg}


%%C070298 h3

\subsection{How it works...}


%%C070299 p

用于时间值的两个操纵器\texttt{put\_\allowbreak{}time()}和\texttt{get\_\allowbreak{}time()}非常相似：它们都是带两个参数的函数模板。第一个参数是指向\texttt{std::\allowbreak{}tm}对象的指针，该对象表示日历日期与时间，其值为要写入流或从流中读取的值。第二个参数是指向以空字符结尾的字符串的指针，表示时间文本的格式。\texttt{put\_\allowbreak{}time()}使用\texttt{std::\allowbreak{}time\_\allowbreak{}put()}facet输出日期时间值，\texttt{get\_\allowbreak{}time()}使用\texttt{std::\allowbreak{}time\_\allowbreak{}get()}facet解析日期时间值。两个操纵器函数模板都返回未指明类型的对象。这些函数不抛异常：


%%C070300 code

\begin{codeblockbg}
\begin{Verbatim}[commandchars=\\\{\},breaklines=true,breakanywhere=true,fontsize=\codefont,xleftmargin=17bp]
\PYG{k}{template}\PYG{o}{\PYGZlt{}}\PYG{k}{class}\PYG{+w}{ }\PYG{n+nc}{CharT}\PYG{o}{\PYGZgt{}}
\PYG{c+cm}{/*unspecified*/}\PYG{+w}{ }\PYG{n}{put\PYGZus{}time}\PYG{p}{(}\PYG{k}{const}\PYG{+w}{ }\PYG{n}{std}\PYG{o}{:}\PYG{o}{:}\PYG{n}{tm}\PYG{o}{*}\PYG{+w}{ }\PYG{n}{tmb}\PYG{p}{,}\PYG{+w}{ }\PYG{k}{const}\PYG{+w}{ }\PYG{n}{CharT}\PYG{o}{*}\PYG{+w}{ }\PYG{n}{fmt}\PYG{p}{)}\PYG{p}{;}
\PYG{k}{template}\PYG{o}{\PYGZlt{}}\PYG{k}{class}\PYG{+w}{ }\PYG{n+nc}{CharT}\PYG{o}{\PYGZgt{}}
\PYG{c+cm}{/*unspecified*/}\PYG{+w}{ }\PYG{n}{get\PYGZus{}time}\PYG{p}{(}\PYG{n}{std}\PYG{o}{:}\PYG{o}{:}\PYG{n}{tm}\PYG{o}{*}\PYG{+w}{ }\PYG{n}{tmb}\PYG{p}{,}\PYG{+w}{ }\PYG{k}{const}\PYG{+w}{ }\PYG{n}{CharT}\PYG{o}{*}\PYG{+w}{ }\PYG{n}{fmt}\PYG{p}{)}\PYG{p}{;}
\end{Verbatim}
\end{codeblockbg}


%%C070301 note

\begin{myNote}[tip]
用\texttt{put\_\allowbreak{}time()}把日期时间值写入输出流所得到的字符串，与调用\texttt{std::\allowbreak{}strftime()}或\texttt{std::\allowbreak{}wcsftime()}得到的结果相同。
\end{myNote}


%%C070302 p

标准定义了很长的转换说明符列表，用于组成格式字符串。这些说明符以\texttt{\%}开头，某些情况下后面还跟一个\texttt{E}或\texttt{0}。其中有些说明符是等价的，例如\texttt{\%F}等价于\texttt{\%Y-\%m-\%d}（即ISO 8601日期格式），\texttt{\%T}等价于\texttt{\%H:\%M:\%S}（即ISO 8601时间格式）。本案例的示例只涉及少量转换说明符，且都与ISO 8601日期时间格式有关。完整的转换说明符列表请查阅C++标准，或访问前面提到的链接。


%%C070303 note

\begin{myNote}
要注意，\texttt{put\_\allowbreak{}time()}支持的转换说明符并非都受\texttt{get\_\allowbreak{}time()}支持。例如\texttt{z}（ISO 8601格式中相对UTC的偏移）和\texttt{Z}（时区名称或缩写）说明符就只能用于\texttt{put\_\allowbreak{}time()}。
\end{myNote}


%%C070304 code

\begin{codeblockbg}
\begin{Verbatim}[commandchars=\\\{\},breaklines=true,breakanywhere=true,fontsize=\codefont,xleftmargin=17bp]
\PYG{n}{std}\PYG{o}{:}\PYG{o}{:}\PYG{n}{istringstream}\PYG{+w}{ }\PYG{n}{stext}\PYG{p}{(}\PYG{l+s}{\PYGZdq{}}\PYG{l+s}{2016\PYGZhy{}12\PYGZhy{}04T05:26:47+0900}\PYG{l+s}{\PYGZdq{}}\PYG{p}{)}\PYG{p}{;}
\PYG{k}{auto}\PYG{+w}{ }\PYG{n}{time}\PYG{+w}{ }\PYG{o}{=}\PYG{+w}{ }\PYG{n}{std}\PYG{o}{:}\PYG{o}{:}\PYG{n}{tm}\PYG{+w}{ }\PYG{p}{\PYGZob{}}\PYG{p}{\PYGZcb{}}\PYG{p}{;}
\PYG{n}{stext}\PYG{+w}{ }\PYG{o}{\PYGZgt{}}\PYG{o}{\PYGZgt{}}\PYG{+w}{ }\PYG{n}{std}\PYG{o}{:}\PYG{o}{:}\PYG{n}{get\PYGZus{}time}\PYG{p}{(}\PYG{o}{\PYGZam{}}\PYG{n}{time}\PYG{p}{,}\PYG{+w}{ }\PYG{l+s}{\PYGZdq{}}\PYG{l+s}{\PYGZpc{}Y\PYGZhy{}\PYGZpc{}m\PYGZhy{}\PYGZpc{}dT\PYGZpc{}H:\PYGZpc{}M:\PYGZpc{}S\PYGZpc{}z}\PYG{l+s}{\PYGZdq{}}\PYG{p}{)}\PYG{p}{;}\PYG{+w}{ }\PYG{c+c1}{// 失败}
\PYG{n}{stext}\PYG{+w}{ }\PYG{o}{\PYGZgt{}}\PYG{o}{\PYGZgt{}}\PYG{+w}{ }\PYG{n}{std}\PYG{o}{:}\PYG{o}{:}\PYG{n}{get\PYGZus{}time}\PYG{p}{(}\PYG{o}{\PYGZam{}}\PYG{n}{time}\PYG{p}{,}\PYG{+w}{ }\PYG{l+s}{\PYGZdq{}}\PYG{l+s}{\PYGZpc{}Y\PYGZhy{}\PYGZpc{}m\PYGZhy{}\PYGZpc{}dT\PYGZpc{}H:\PYGZpc{}M:\PYGZpc{}S}\PYG{l+s}{\PYGZdq{}}\PYG{p}{)}\PYG{p}{;}\PYG{+w}{   }\PYG{c+c1}{// 成功}
\end{Verbatim}
\end{codeblockbg}


%%C070305 p

某些转换说明符表示的文本与区域设置有关。所有以\texttt{E}或\texttt{0}开头的说明符都与区域设置有关。要为流设置特定的区域设置，使用\texttt{imbue()}方法，如上一节的示例所示。


%%C070306 h3

\subsection{See also}


%%C070307 li

\begin{itemize}
\item \textit{用I/O操纵器控制流的输出}


%%C070308 li

\item \textit{使用货币I/O操纵器}

\end{itemize}


%%C070309 h2

\section{处理文件系统路径}


%%C070310 p

C++17标准的一个重要新增内容是\texttt{filesystem}库，它使开发者能在层次化文件系统（如Windows或POSIX文件系统）中处理路径、文件和目录。该标准库是在\texttt{boost.\allowbreak{}filesystem}库的基础上发展而来的。接下来的几个案例将探讨库中那些使我们能对文件和目录执行操作（例如创建、移动或删除），以及查询属性和搜索的特性。不过，先看一看库是如何处理路径的，这很重要。


%%C070311 h3

\subsection{Getting ready}


%%C070312 p

本案例中，大部分示例使用Windows路径。在随书代码中，所有示例都提供了Windows和POSIX两种版本。


%%C070313 p

\texttt{filesystem}库位于\texttt{<\allowbreak{}filesystem>}头文件中的\texttt{std::\allowbreak{}filesystem}命名空间。为简化代码，所有示例都使用下面的命名空间别名：


%%C070314 code

\begin{codeblockbg}
\begin{Verbatim}[commandchars=\\\{\},breaklines=true,breakanywhere=true,fontsize=\codefont,xleftmargin=17bp]
\PYG{k}{namespace}\PYG{+w}{ }\PYG{n+nn}{fs}\PYG{+w}{ }\PYG{o}{=}\PYG{+w}{ }\PYG{n+nn}{std}\PYG{o}{:}\PYG{o}{:}\PYG{n+nn}{filesystem}\PYG{p}{;}
\end{Verbatim}
\end{codeblockbg}


%%C070315 note

\begin{myNote}
在写作本书时，所有主流编译器都提供了该库的实现，不过它仍被视为实验性的，因此被放在以此命名的命名空间中。正因如此，该库实际使用的命名空间是\myttbr{std::\allowbreak{}experimental::\allowbreak{}filesystem}，GCC和Clang中实际使用的头文件是\myttbr{<\allowbreak{}experimental/\allowbreak{}filesystem>}。
\end{myNote}


%%C070316 p

指向文件系统组件（文件、目录、硬链接或软链接）的路径由\texttt{path}类表示。


%%C070317 h3

\subsection{How to do it...}


%%C070318 p

以下是对路径最常见的操作：


%%C070319 li

\begin{itemize}
\item 用构造函数、赋值运算符或\texttt{assign()}方法创建路径：


%%C070320 code

\begin{codeblockbg}
\begin{Verbatim}[commandchars=\\\{\},breaklines=true,breakanywhere=true,fontsize=\codefont,xleftmargin=17bp]
\PYG{c+c1}{// Windows}
\PYG{k}{auto}\PYG{+w}{ }\PYG{n}{path}\PYG{+w}{ }\PYG{o}{=}\PYG{+w}{ }\PYG{n}{fs}\PYG{o}{:}\PYG{o}{:}\PYG{n}{path}\PYG{p}{\PYGZob{}}\PYG{l+s}{\PYGZdq{}}\PYG{l+s}{C:}\PYG{l+s+se}{\PYGZbs{}\PYGZbs{}}\PYG{l+s}{Users}\PYG{l+s+se}{\PYGZbs{}\PYGZbs{}}\PYG{l+s}{Marius}\PYG{l+s+se}{\PYGZbs{}\PYGZbs{}}\PYG{l+s}{Documents}\PYG{l+s}{\PYGZdq{}}\PYG{p}{\PYGZcb{}}\PYG{p}{;}
\PYG{c+c1}{// POSIX}
\PYG{k}{auto}\PYG{+w}{ }\PYG{n}{path}\PYG{+w}{ }\PYG{o}{=}\PYG{+w}{ }\PYG{n}{fs}\PYG{o}{:}\PYG{o}{:}\PYG{n}{path}\PYG{p}{\PYGZob{}}\PYG{+w}{ }\PYG{l+s}{\PYGZdq{}}\PYG{l+s}{/home/marius/docs}\PYG{l+s}{\PYGZdq{}}\PYG{+w}{ }\PYG{p}{\PYGZcb{}}\PYG{p}{;}
\end{Verbatim}
\end{codeblockbg}


%%C070321 li

\item 用成员\texttt{operator /\allowbreak{}=}、非成员\texttt{operator /}或\texttt{append()}方法向路径追加元素，追加时包含目录分隔符：


%%C070322 code

\begin{codeblockbg}
\begin{Verbatim}[commandchars=\\\{\},breaklines=true,breakanywhere=true,fontsize=\codefont,xleftmargin=17bp]
\PYG{n}{path}\PYG{+w}{ }\PYG{o}{/}\PYG{o}{=}\PYG{+w}{ }\PYG{l+s}{\PYGZdq{}}\PYG{l+s}{Book}\PYG{l+s}{\PYGZdq{}}\PYG{p}{;}
\PYG{n}{path}\PYG{+w}{ }\PYG{o}{=}\PYG{+w}{ }\PYG{n}{path}\PYG{+w}{ }\PYG{o}{/}\PYG{+w}{ }\PYG{l+s}{\PYGZdq{}}\PYG{l+s}{Modern}\PYG{l+s}{\PYGZdq{}}\PYG{+w}{ }\PYG{o}{/}\PYG{+w}{ }\PYG{l+s}{\PYGZdq{}}\PYG{l+s}{Cpp}\PYG{l+s}{\PYGZdq{}}\PYG{p}{;}
\PYG{n}{path}\PYG{p}{.}\PYG{n}{append}\PYG{p}{(}\PYG{l+s}{\PYGZdq{}}\PYG{l+s}{Programming}\PYG{l+s}{\PYGZdq{}}\PYG{p}{)}\PYG{p}{;}
\PYG{c+c1}{// Windows: C:\PYGZbs{}Users\PYGZbs{}Marius\PYGZbs{}Documents\PYGZbs{}Book\PYGZbs{}Modern\PYGZbs{}Cpp\PYGZbs{}Programming}
\PYG{c+c1}{// POSIX:   /home/marius/docs/Book/Modern/Cpp/Programming}
\end{Verbatim}
\end{codeblockbg}


%%C070323 li

\item 用成员\texttt{operator +=}、非成员\texttt{operator +}或\texttt{concat()}方法向路径连接元素，连接时不包含目录分隔符：


%%C070324 code

\begin{codeblockbg}
\begin{Verbatim}[commandchars=\\\{\},breaklines=true,breakanywhere=true,fontsize=\codefont,xleftmargin=17bp]
\PYG{k}{auto}\PYG{+w}{ }\PYG{n}{path}\PYG{+w}{ }\PYG{o}{=}\PYG{+w}{ }\PYG{n}{fs}\PYG{o}{:}\PYG{o}{:}\PYG{n}{path}\PYG{p}{\PYGZob{}}\PYG{+w}{ }\PYG{l+s}{\PYGZdq{}}\PYG{l+s}{C:}\PYG{l+s+se}{\PYGZbs{}\PYGZbs{}}\PYG{l+s}{Users}\PYG{l+s+se}{\PYGZbs{}\PYGZbs{}}\PYG{l+s}{Marius}\PYG{l+s+se}{\PYGZbs{}\PYGZbs{}}\PYG{l+s}{Documents}\PYG{l+s}{\PYGZdq{}}\PYG{+w}{ }\PYG{p}{\PYGZcb{}}\PYG{p}{;}
\PYG{n}{path}\PYG{+w}{ }\PYG{o}{+}\PYG{o}{=}\PYG{+w}{ }\PYG{l+s}{\PYGZdq{}}\PYG{l+s}{Book}\PYG{l+s}{\PYGZdq{}}\PYG{p}{;}
\PYG{n}{path}\PYG{p}{.}\PYG{n}{concat}\PYG{p}{(}\PYG{l+s}{\PYGZdq{}}\PYG{l+s}{Modern}\PYG{l+s}{\PYGZdq{}}\PYG{p}{)}\PYG{p}{;}
\PYG{c+c1}{// path = C:\PYGZbs{}Users\PYGZbs{}Marius\PYGZbs{}Documents\PYGZbs{}Book\PYGZbs{}Modern}
\end{Verbatim}
\end{codeblockbg}


%%C070325 li

\item 用\texttt{root\_\allowbreak{}name()}、\texttt{root\_\allowbreak{}dir()}、\texttt{filename()}、\texttt{stem()}、\texttt{extension()}等成员函数把路径的元素分解成各个部分，如根名称、根目录、父路径、文件名、扩展名等（下面的示例展示了全部这些函数）：


%%C070326 code

\begin{codeblockbg}
\begin{Verbatim}[commandchars=\\\{\},breaklines=true,breakanywhere=true,fontsize=\codefont,xleftmargin=17bp]
\PYG{k}{auto}\PYG{+w}{ }\PYG{n}{path}\PYG{+w}{ }\PYG{o}{=}
\PYG{+w}{  }\PYG{n}{fs}\PYG{o}{:}\PYG{o}{:}\PYG{n}{path}\PYG{p}{\PYGZob{}}\PYG{l+s}{\PYGZdq{}}\PYG{l+s}{C:}\PYG{l+s+se}{\PYGZbs{}\PYGZbs{}}\PYG{l+s}{Users}\PYG{l+s+se}{\PYGZbs{}\PYGZbs{}}\PYG{l+s}{Marius}\PYG{l+s+se}{\PYGZbs{}\PYGZbs{}}\PYG{l+s}{Documents}\PYG{l+s+se}{\PYGZbs{}\PYGZbs{}}\PYG{l+s}{sample.file.txt}\PYG{l+s}{\PYGZdq{}}\PYG{p}{\PYGZcb{}}\PYG{p}{;}
\PYG{n}{std}\PYG{o}{:}\PYG{o}{:}\PYG{n}{cout}
\PYG{+w}{  }\PYG{o}{\PYGZlt{}}\PYG{o}{\PYGZlt{}}\PYG{+w}{ }\PYG{l+s}{\PYGZdq{}}\PYG{l+s}{root: }\PYG{l+s}{\PYGZdq{}}\PYG{+w}{ }\PYG{o}{\PYGZlt{}}\PYG{o}{\PYGZlt{}}\PYG{+w}{ }\PYG{n}{path}\PYG{p}{.}\PYG{n}{root\PYGZus{}name}\PYG{p}{(}\PYG{p}{)}\PYG{+w}{ }\PYG{o}{\PYGZlt{}}\PYG{o}{\PYGZlt{}}\PYG{+w}{ }\PYG{n}{std}\PYG{o}{:}\PYG{o}{:}\PYG{n}{endl}
\PYG{+w}{  }\PYG{o}{\PYGZlt{}}\PYG{o}{\PYGZlt{}}\PYG{+w}{ }\PYG{l+s}{\PYGZdq{}}\PYG{l+s}{root dir: }\PYG{l+s}{\PYGZdq{}}\PYG{+w}{ }\PYG{o}{\PYGZlt{}}\PYG{o}{\PYGZlt{}}\PYG{+w}{ }\PYG{n}{path}\PYG{p}{.}\PYG{n}{root\PYGZus{}directory}\PYG{p}{(}\PYG{p}{)}\PYG{+w}{ }\PYG{o}{\PYGZlt{}}\PYG{o}{\PYGZlt{}}\PYG{+w}{ }\PYG{n}{std}\PYG{o}{:}\PYG{o}{:}\PYG{n}{endl}
\PYG{+w}{  }\PYG{o}{\PYGZlt{}}\PYG{o}{\PYGZlt{}}\PYG{+w}{ }\PYG{l+s}{\PYGZdq{}}\PYG{l+s}{root path: }\PYG{l+s}{\PYGZdq{}}\PYG{+w}{ }\PYG{o}{\PYGZlt{}}\PYG{o}{\PYGZlt{}}\PYG{+w}{ }\PYG{n}{path}\PYG{p}{.}\PYG{n}{root\PYGZus{}path}\PYG{p}{(}\PYG{p}{)}\PYG{+w}{ }\PYG{o}{\PYGZlt{}}\PYG{o}{\PYGZlt{}}\PYG{+w}{ }\PYG{n}{std}\PYG{o}{:}\PYG{o}{:}\PYG{n}{endl}
\PYG{+w}{  }\PYG{o}{\PYGZlt{}}\PYG{o}{\PYGZlt{}}\PYG{+w}{ }\PYG{l+s}{\PYGZdq{}}\PYG{l+s}{rel path: }\PYG{l+s}{\PYGZdq{}}\PYG{+w}{ }\PYG{o}{\PYGZlt{}}\PYG{o}{\PYGZlt{}}\PYG{+w}{ }\PYG{n}{path}\PYG{p}{.}\PYG{n}{relative\PYGZus{}path}\PYG{p}{(}\PYG{p}{)}\PYG{+w}{ }\PYG{o}{\PYGZlt{}}\PYG{o}{\PYGZlt{}}\PYG{+w}{ }\PYG{n}{std}\PYG{o}{:}\PYG{o}{:}\PYG{n}{endl}
\PYG{+w}{  }\PYG{o}{\PYGZlt{}}\PYG{o}{\PYGZlt{}}\PYG{+w}{ }\PYG{l+s}{\PYGZdq{}}\PYG{l+s}{parent path: }\PYG{l+s}{\PYGZdq{}}\PYG{+w}{ }\PYG{o}{\PYGZlt{}}\PYG{o}{\PYGZlt{}}\PYG{+w}{ }\PYG{n}{path}\PYG{p}{.}\PYG{n}{parent\PYGZus{}path}\PYG{p}{(}\PYG{p}{)}\PYG{+w}{ }\PYG{o}{\PYGZlt{}}\PYG{o}{\PYGZlt{}}\PYG{+w}{ }\PYG{n}{std}\PYG{o}{:}\PYG{o}{:}\PYG{n}{endl}
\PYG{o}{\PYGZlt{}}\PYG{o}{\PYGZlt{}}\PYG{+w}{ }\PYG{l+s}{\PYGZdq{}}\PYG{l+s}{filename: }\PYG{l+s}{\PYGZdq{}}\PYG{+w}{ }\PYG{o}{\PYGZlt{}}\PYG{o}{\PYGZlt{}}\PYG{+w}{ }\PYG{n}{path}\PYG{p}{.}\PYG{n}{filename}\PYG{p}{(}\PYG{p}{)}\PYG{+w}{ }\PYG{o}{\PYGZlt{}}\PYG{o}{\PYGZlt{}}\PYG{+w}{ }\PYG{n}{std}\PYG{o}{:}\PYG{o}{:}\PYG{n}{endl}
\PYG{o}{\PYGZlt{}}\PYG{o}{\PYGZlt{}}\PYG{+w}{ }\PYG{l+s}{\PYGZdq{}}\PYG{l+s}{stem: }\PYG{l+s}{\PYGZdq{}}\PYG{+w}{ }\PYG{o}{\PYGZlt{}}\PYG{o}{\PYGZlt{}}\PYG{+w}{ }\PYG{n}{path}\PYG{p}{.}\PYG{n}{stem}\PYG{p}{(}\PYG{p}{)}\PYG{+w}{ }\PYG{o}{\PYGZlt{}}\PYG{o}{\PYGZlt{}}\PYG{+w}{ }\PYG{n}{std}\PYG{o}{:}\PYG{o}{:}\PYG{n}{endl}
\PYG{o}{\PYGZlt{}}\PYG{o}{\PYGZlt{}}\PYG{+w}{ }\PYG{l+s}{\PYGZdq{}}\PYG{l+s}{extension: }\PYG{l+s}{\PYGZdq{}}\PYG{+w}{ }\PYG{o}{\PYGZlt{}}\PYG{o}{\PYGZlt{}}\PYG{+w}{ }\PYG{n}{path}\PYG{p}{.}\PYG{n}{extension}\PYG{p}{(}\PYG{p}{)}\PYG{+w}{ }\PYG{o}{\PYGZlt{}}\PYG{o}{\PYGZlt{}}\PYG{+w}{ }\PYG{n}{std}\PYG{o}{:}\PYG{o}{:}\PYG{n}{endl}\PYG{p}{;}
\end{Verbatim}
\end{codeblockbg}


%%C070327 li

\item 用\texttt{has\_\allowbreak{}root\_\allowbreak{}name()}、\myttbr{has\_\allowbreak{}root\_\allowbreak{}directory()}、\texttt{has\_\allowbreak{}filename()}、\texttt{has\_\allowbreak{}stem()}和\texttt{has\_\allowbreak{}extension()}等成员函数查询路径是否含有某些部分（下面的示例展示了全部这些函数）：


%%C070328 code

\begin{codeblockbg}
\begin{Verbatim}[commandchars=\\\{\},breaklines=true,breakanywhere=true,fontsize=\codefont,xleftmargin=17bp]
\PYG{k}{auto}\PYG{+w}{ }\PYG{n}{path}\PYG{+w}{ }\PYG{o}{=}
\PYG{+w}{  }\PYG{n}{fs}\PYG{o}{:}\PYG{o}{:}\PYG{n}{path}\PYG{p}{\PYGZob{}}\PYG{l+s}{\PYGZdq{}}\PYG{l+s}{C:}\PYG{l+s+se}{\PYGZbs{}\PYGZbs{}}\PYG{l+s}{Users}\PYG{l+s+se}{\PYGZbs{}\PYGZbs{}}\PYG{l+s}{Marius}\PYG{l+s+se}{\PYGZbs{}\PYGZbs{}}\PYG{l+s}{Documents}\PYG{l+s+se}{\PYGZbs{}\PYGZbs{}}\PYG{l+s}{sample.file.txt}\PYG{l+s}{\PYGZdq{}}\PYG{p}{\PYGZcb{}}\PYG{p}{;}
\PYG{n}{std}\PYG{o}{:}\PYG{o}{:}\PYG{n}{cout}
\PYG{+w}{  }\PYG{o}{\PYGZlt{}}\PYG{o}{\PYGZlt{}}\PYG{+w}{ }\PYG{l+s}{\PYGZdq{}}\PYG{l+s}{has root: }\PYG{l+s}{\PYGZdq{}}\PYG{+w}{ }\PYG{o}{\PYGZlt{}}\PYG{o}{\PYGZlt{}}\PYG{+w}{ }\PYG{n}{path}\PYG{p}{.}\PYG{n}{has\PYGZus{}root\PYGZus{}name}\PYG{p}{(}\PYG{p}{)}\PYG{+w}{ }\PYG{o}{\PYGZlt{}}\PYG{o}{\PYGZlt{}}\PYG{+w}{ }\PYG{n}{std}\PYG{o}{:}\PYG{o}{:}\PYG{n}{endl}
\PYG{+w}{  }\PYG{o}{\PYGZlt{}}\PYG{o}{\PYGZlt{}}\PYG{+w}{ }\PYG{l+s}{\PYGZdq{}}\PYG{l+s}{has root dir: }\PYG{l+s}{\PYGZdq{}}\PYG{+w}{ }\PYG{o}{\PYGZlt{}}\PYG{o}{\PYGZlt{}}\PYG{+w}{ }\PYG{n}{path}\PYG{p}{.}\PYG{n}{has\PYGZus{}root\PYGZus{}directory}\PYG{p}{(}\PYG{p}{)}\PYG{+w}{ }\PYG{o}{\PYGZlt{}}\PYG{o}{\PYGZlt{}}\PYG{+w}{ }\PYG{n}{std}\PYG{o}{:}\PYG{o}{:}\PYG{n}{endl}
\PYG{+w}{  }\PYG{o}{\PYGZlt{}}\PYG{o}{\PYGZlt{}}\PYG{+w}{ }\PYG{l+s}{\PYGZdq{}}\PYG{l+s}{has root path: }\PYG{l+s}{\PYGZdq{}}\PYG{+w}{ }\PYG{o}{\PYGZlt{}}\PYG{o}{\PYGZlt{}}\PYG{+w}{ }\PYG{n}{path}\PYG{p}{.}\PYG{n}{has\PYGZus{}root\PYGZus{}path}\PYG{p}{(}\PYG{p}{)}\PYG{+w}{ }\PYG{o}{\PYGZlt{}}\PYG{o}{\PYGZlt{}}\PYG{+w}{ }\PYG{n}{std}\PYG{o}{:}\PYG{o}{:}\PYG{n}{endl}
\PYG{+w}{  }\PYG{o}{\PYGZlt{}}\PYG{o}{\PYGZlt{}}\PYG{+w}{ }\PYG{l+s}{\PYGZdq{}}\PYG{l+s}{has rel path: }\PYG{l+s}{\PYGZdq{}}\PYG{+w}{ }\PYG{o}{\PYGZlt{}}\PYG{o}{\PYGZlt{}}\PYG{+w}{ }\PYG{n}{path}\PYG{p}{.}\PYG{n}{has\PYGZus{}relative\PYGZus{}path}\PYG{p}{(}\PYG{p}{)}\PYG{+w}{ }\PYG{o}{\PYGZlt{}}\PYG{o}{\PYGZlt{}}\PYG{+w}{ }\PYG{n}{std}\PYG{o}{:}\PYG{o}{:}\PYG{n}{endl}
\PYG{+w}{  }\PYG{o}{\PYGZlt{}}\PYG{o}{\PYGZlt{}}\PYG{+w}{ }\PYG{l+s}{\PYGZdq{}}\PYG{l+s}{has parent path: }\PYG{l+s}{\PYGZdq{}}\PYG{+w}{ }\PYG{o}{\PYGZlt{}}\PYG{o}{\PYGZlt{}}\PYG{+w}{ }\PYG{n}{path}\PYG{p}{.}\PYG{n}{has\PYGZus{}parent\PYGZus{}path}\PYG{p}{(}\PYG{p}{)}\PYG{+w}{ }\PYG{o}{\PYGZlt{}}\PYG{o}{\PYGZlt{}}\PYG{+w}{ }\PYG{n}{std}\PYG{o}{:}\PYG{o}{:}\PYG{n}{endl}
\PYG{+w}{  }\PYG{o}{\PYGZlt{}}\PYG{o}{\PYGZlt{}}\PYG{+w}{ }\PYG{l+s}{\PYGZdq{}}\PYG{l+s}{has filename: }\PYG{l+s}{\PYGZdq{}}\PYG{+w}{ }\PYG{o}{\PYGZlt{}}\PYG{o}{\PYGZlt{}}\PYG{+w}{ }\PYG{n}{path}\PYG{p}{.}\PYG{n}{has\PYGZus{}filename}\PYG{p}{(}\PYG{p}{)}\PYG{+w}{ }\PYG{o}{\PYGZlt{}}\PYG{o}{\PYGZlt{}}\PYG{+w}{ }\PYG{n}{std}\PYG{o}{:}\PYG{o}{:}\PYG{n}{endl}
\PYG{+w}{  }\PYG{o}{\PYGZlt{}}\PYG{o}{\PYGZlt{}}\PYG{+w}{ }\PYG{l+s}{\PYGZdq{}}\PYG{l+s}{has stem: }\PYG{l+s}{\PYGZdq{}}\PYG{+w}{ }\PYG{o}{\PYGZlt{}}\PYG{o}{\PYGZlt{}}\PYG{+w}{ }\PYG{n}{path}\PYG{p}{.}\PYG{n}{has\PYGZus{}stem}\PYG{p}{(}\PYG{p}{)}\PYG{+w}{ }\PYG{o}{\PYGZlt{}}\PYG{o}{\PYGZlt{}}\PYG{+w}{ }\PYG{n}{std}\PYG{o}{:}\PYG{o}{:}\PYG{n}{endl}
\PYG{+w}{  }\PYG{o}{\PYGZlt{}}\PYG{o}{\PYGZlt{}}\PYG{+w}{ }\PYG{l+s}{\PYGZdq{}}\PYG{l+s}{has extension: }\PYG{l+s}{\PYGZdq{}}\PYG{+w}{ }\PYG{o}{\PYGZlt{}}\PYG{o}{\PYGZlt{}}\PYG{+w}{ }\PYG{n}{path}\PYG{p}{.}\PYG{n}{has\PYGZus{}extension}\PYG{p}{(}\PYG{p}{)}\PYG{+w}{ }\PYG{o}{\PYGZlt{}}\PYG{o}{\PYGZlt{}}\PYG{+w}{ }\PYG{n}{std}\PYG{o}{:}\PYG{o}{:}\PYG{n}{endl}\PYG{p}{;}
\end{Verbatim}
\end{codeblockbg}


%%C070329 li

\item 检查路径是相对路径还是绝对路径：


%%C070330 code

\begin{codeblockbg}
\begin{Verbatim}[commandchars=\\\{\},breaklines=true,breakanywhere=true,fontsize=\codefont,xleftmargin=17bp]
\PYG{k}{auto}\PYG{+w}{ }\PYG{n}{path2}\PYG{+w}{ }\PYG{o}{=}\PYG{+w}{ }\PYG{n}{fs}\PYG{o}{:}\PYG{o}{:}\PYG{n}{path}\PYG{p}{\PYGZob{}}\PYG{+w}{ }\PYG{l+s}{\PYGZdq{}}\PYG{l+s}{marius}\PYG{l+s+se}{\PYGZbs{}\PYGZbs{}}\PYG{l+s}{temp}\PYG{l+s}{\PYGZdq{}}\PYG{+w}{ }\PYG{p}{\PYGZcb{}}\PYG{p}{;}
\PYG{n}{std}\PYG{o}{:}\PYG{o}{:}\PYG{n}{cout}
\PYG{+w}{  }\PYG{o}{\PYGZlt{}}\PYG{o}{\PYGZlt{}}\PYG{+w}{ }\PYG{l+s}{\PYGZdq{}}\PYG{l+s}{absolute: }\PYG{l+s}{\PYGZdq{}}\PYG{+w}{ }\PYG{o}{\PYGZlt{}}\PYG{o}{\PYGZlt{}}\PYG{+w}{ }\PYG{n}{path1}\PYG{p}{.}\PYG{n}{is\PYGZus{}absolute}\PYG{p}{(}\PYG{p}{)}\PYG{+w}{ }\PYG{o}{\PYGZlt{}}\PYG{o}{\PYGZlt{}}\PYG{+w}{ }\PYG{n}{std}\PYG{o}{:}\PYG{o}{:}\PYG{n}{endl}
\PYG{+w}{  }\PYG{o}{\PYGZlt{}}\PYG{o}{\PYGZlt{}}\PYG{+w}{ }\PYG{l+s}{\PYGZdq{}}\PYG{l+s}{absolute: }\PYG{l+s}{\PYGZdq{}}\PYG{+w}{ }\PYG{o}{\PYGZlt{}}\PYG{o}{\PYGZlt{}}\PYG{+w}{ }\PYG{n}{path2}\PYG{p}{.}\PYG{n}{is\PYGZus{}absolute}\PYG{p}{(}\PYG{p}{)}\PYG{+w}{ }\PYG{o}{\PYGZlt{}}\PYG{o}{\PYGZlt{}}\PYG{+w}{ }\PYG{n}{std}\PYG{o}{:}\PYG{o}{:}\PYG{n}{endl}\PYG{p}{;}
\end{Verbatim}
\end{codeblockbg}


%%C070331 li

\item 修改路径的各个部分，例如用\myttbr{replace\_\allowbreak{}filename()}和\texttt{remove\_\allowbreak{}filename()}修改文件名，用\myttbr{replace\_\allowbreak{}extension()}修改扩展名：


%%C070332 code

\begin{codeblockbg}
\begin{Verbatim}[commandchars=\\\{\},breaklines=true,breakanywhere=true,fontsize=\codefont,xleftmargin=17bp]
\PYG{k}{auto}\PYG{+w}{ }\PYG{n}{path}\PYG{+w}{ }\PYG{o}{=}
\PYG{+w}{  }\PYG{n}{fs}\PYG{o}{:}\PYG{o}{:}\PYG{n}{path}\PYG{p}{\PYGZob{}}\PYG{l+s}{\PYGZdq{}}\PYG{l+s}{C:}\PYG{l+s+se}{\PYGZbs{}\PYGZbs{}}\PYG{l+s}{Users}\PYG{l+s+se}{\PYGZbs{}\PYGZbs{}}\PYG{l+s}{Marius}\PYG{l+s+se}{\PYGZbs{}\PYGZbs{}}\PYG{l+s}{Documents}\PYG{l+s+se}{\PYGZbs{}\PYGZbs{}}\PYG{l+s}{sample.file.txt}\PYG{l+s}{\PYGZdq{}}\PYG{p}{\PYGZcb{}}\PYG{p}{;}
\PYG{n}{path}\PYG{p}{.}\PYG{n}{replace\PYGZus{}filename}\PYG{p}{(}\PYG{l+s}{\PYGZdq{}}\PYG{l+s}{output}\PYG{l+s}{\PYGZdq{}}\PYG{p}{)}\PYG{p}{;}
\PYG{n}{path}\PYG{p}{.}\PYG{n}{replace\PYGZus{}extension}\PYG{p}{(}\PYG{l+s}{\PYGZdq{}}\PYG{l+s}{.log}\PYG{l+s}{\PYGZdq{}}\PYG{p}{)}\PYG{p}{;}
\PYG{c+c1}{// path = C:\PYGZbs{}Users\PYGZbs{}Marius\PYGZbs{}Documents\PYGZbs{}output.log}
\PYG{n}{path}\PYG{p}{.}\PYG{n}{remove\PYGZus{}filename}\PYG{p}{(}\PYG{p}{)}\PYG{p}{;}
\PYG{c+c1}{// path = C:\PYGZbs{}Users\PYGZbs{}Marius\PYGZbs{}Documents}
\end{Verbatim}
\end{codeblockbg}


%%C070333 li

\item 把目录分隔符转换为系统首选的分隔符：


%%C070334 code

\begin{codeblockbg}
\begin{Verbatim}[commandchars=\\\{\},breaklines=true,breakanywhere=true,fontsize=\codefont,xleftmargin=17bp]
\PYG{c+c1}{// Windows}
\PYG{k}{auto}\PYG{+w}{ }\PYG{n}{path}\PYG{+w}{ }\PYG{o}{=}\PYG{+w}{ }\PYG{n}{fs}\PYG{o}{:}\PYG{o}{:}\PYG{n}{path}\PYG{p}{\PYGZob{}}\PYG{l+s}{\PYGZdq{}}\PYG{l+s}{Users/Marius/Documents}\PYG{l+s}{\PYGZdq{}}\PYG{p}{\PYGZcb{}}\PYG{p}{;}
\PYG{n}{path}\PYG{p}{.}\PYG{n}{make\PYGZus{}preferred}\PYG{p}{(}\PYG{p}{)}\PYG{p}{;}
\PYG{c+c1}{// path = Users\PYGZbs{}Marius\PYGZbs{}Documents}
\PYG{c+c1}{// POSIX}
\PYG{k}{auto}\PYG{+w}{ }\PYG{n}{path}\PYG{+w}{ }\PYG{o}{=}\PYG{+w}{ }\PYG{n}{fs}\PYG{o}{:}\PYG{o}{:}\PYG{n}{path}\PYG{p}{\PYGZob{}}\PYG{+w}{ }\PYG{l+s}{\PYGZdq{}}\PYG{l+s}{homedocs}\PYG{l+s}{\PYGZdq{}}\PYG{+w}{ }\PYG{p}{\PYGZcb{}}\PYG{p}{;}
\PYG{n}{path}\PYG{p}{.}\PYG{n}{make\PYGZus{}preferred}\PYG{p}{(}\PYG{p}{)}\PYG{p}{;}
\PYG{c+c1}{// path = /home/marius/docs}
\end{Verbatim}
\end{codeblockbg}

\end{itemize}


%%C070335 h3

\subsection{How it works...}


%%C070336 p

\myttbr{std::\allowbreak{}filesystem::\allowbreak{}path}类用于为指向文件系统组件的路径建模。不过它只处理语法，并不验证路径所表示的组件（如文件或目录）是否存在。


%%C070337 p

该库为路径定义了一套可移植的通用语法，能适应POSIX或Windows等各种文件系统，包括Microsoft Windows的\textit{UNC}（\textit{Universal} \textit{Naming Convention}，通用命名约定）格式。这两类文件系统在几个关键方面存在差异：


%%C070338 li

\begin{itemize}
\item POSIX系统只有一棵目录树，没有根名称，只有一个名为\texttt{/}的根目录，也只有一个当前目录，并用\texttt{/}作为目录分隔符。路径表示为以空字符结尾的\texttt{char}字符串，采用UTF-8编码。


%%C070339 li

\item Windows系统有多棵目录树，每棵树都有根名称（如\texttt{C:}）、根目录（如\texttt{C:}）和一个当前目录。路径表示为以空字符结尾的宽字符字符串，采用UTF-16编码。

\end{itemize}


%%C070340 p

\texttt{filesystem}库中定义的路径名具有如下语法：


%%C070341 li

\begin{itemize}
\item 可选的根名称（\texttt{C:}或\texttt{//\allowbreak{}localhost}）


%%C070342 li

\item 可选的根目录


%%C070343 li

\item 零个或多个文件名（可指向文件、目录、硬链接或符号链接），或目录分隔符

\end{itemize}


%%C070344 p

有两个特殊文件名会被识别：单点（\texttt{.}）表示当前目录，双点（\texttt{..}）表示父目录。目录分隔符可以重复出现，此时会当作单个分隔符处理（也就是说，\texttt{/\allowbreak{}home////\allowbreak{}docs}与\texttt{/\allowbreak{}home/\allowbreak{}marius/\allowbreak{}docs}相同）。不含多余的当前目录名（\texttt{.}）、不含多余的父目录名（\texttt{..}），也不含多余目录分隔符的路径，称为处于规范形式（normal form）。


%%C070345 p

上一节介绍的路径操作只是最常见的路径操作。不过，实现还定义了额外的查询方法与修改方法、迭代器、非成员比较运算符等。下面的示例遍历路径的各个部分并打印到控制台：


%%C070346 code

\begin{codeblockbg}
\begin{Verbatim}[commandchars=\\\{\},breaklines=true,breakanywhere=true,fontsize=\codefont,xleftmargin=17bp]
\PYG{k}{auto}\PYG{+w}{ }\PYG{n}{path}\PYG{+w}{ }\PYG{o}{=}
\PYG{+w}{  }\PYG{n}{fs}\PYG{o}{:}\PYG{o}{:}\PYG{n}{path}\PYG{p}{\PYGZob{}}\PYG{+w}{ }\PYG{l+s}{\PYGZdq{}}\PYG{l+s}{C:}\PYG{l+s+se}{\PYGZbs{}\PYGZbs{}}\PYG{l+s}{Users}\PYG{l+s+se}{\PYGZbs{}\PYGZbs{}}\PYG{l+s}{Marius}\PYG{l+s+se}{\PYGZbs{}\PYGZbs{}}\PYG{l+s}{Documents}\PYG{l+s+se}{\PYGZbs{}\PYGZbs{}}\PYG{l+s}{sample.file.txt}\PYG{l+s}{\PYGZdq{}}\PYG{+w}{ }\PYG{p}{\PYGZcb{}}\PYG{p}{;}
\PYG{k}{for}\PYG{+w}{ }\PYG{p}{(}\PYG{k}{auto}\PYG{+w}{ }\PYG{k}{const}\PYG{+w}{ }\PYG{o}{\PYGZam{}}\PYG{+w}{ }\PYG{n}{part}\PYG{+w}{ }\PYG{o}{:}\PYG{+w}{ }\PYG{n}{path}\PYG{p}{)}
\PYG{p}{\PYGZob{}}
\PYG{+w}{  }\PYG{n}{std}\PYG{o}{:}\PYG{o}{:}\PYG{n}{cout}\PYG{+w}{ }\PYG{o}{\PYGZlt{}}\PYG{o}{\PYGZlt{}}\PYG{+w}{ }\PYG{n}{part}\PYG{+w}{ }\PYG{o}{\PYGZlt{}}\PYG{o}{\PYGZlt{}}\PYG{+w}{ }\PYG{n}{std}\PYG{o}{:}\PYG{o}{:}\PYG{n}{endl}\PYG{p}{;}
\PYG{p}{\PYGZcb{}}
\end{Verbatim}
\end{codeblockbg}


%%C070347 p

结果如下：


%%C070348 code

\begin{codeblockbg}
\begin{Verbatim}[commandchars=\\\{\},breaklines=true,breakanywhere=true,fontsize=\codefont,xleftmargin=17bp]
\PYG{n+nl}{C}\PYG{p}{:}
\PYG{n}{Users}
\PYG{n}{Marius}
\PYG{n}{Documents}
\PYG{n}{sample}\PYG{p}{.}\PYG{n}{file}\PYG{p}{.}\PYG{n}{txt}
\end{Verbatim}
\end{codeblockbg}


%%C070349 p

在本例中，\texttt{sample.\allowbreak{}file.\allowbreak{}txt}是文件名。它基本上是从最后一个目录分隔符到路径末尾的部分，也就是成员函数\texttt{filename()}对该路径返回的内容。该文件的扩展名是\texttt{.\allowbreak{}txt}，即\texttt{extension()}成员函数返回的字符串。要获取不含扩展名的文件名，可以用另一个名为\texttt{stem()}的成员函数。在本例中，该方法返回的字符串是\texttt{sample.\allowbreak{}file}。所有这些方法以及其他分解方法，都有一个同名的对应查询方法，只是加上\texttt{has\_}前缀，例如\texttt{has\_\allowbreak{}filename()}、\texttt{has\_\allowbreak{}stem()}和\texttt{has\_\allowbreak{}extension()}。这些方法都返回\texttt{bool}值，表示路径是否含有相应的部分。


%%C070350 h3

\subsection{See also}


%%C070351 li

\begin{itemize}
\item \textit{创建、复制与删除文件和目录}


%%C070352 li

\item \textit{检查现有文件或目录的属性}

\end{itemize}


%%C070353 h2

\section{创建、复制与删除文件和目录}


%%C070354 p

针对文件的操作（如复制、移动和删除）以及针对目录的操作（如创建、重命名和删除）都由\texttt{filesystem}库支持。文件和目录用路径来标识（路径可以是绝对路径、规范路径或相对路径），这一主题前面的案例已经讨论过。本案例将介绍上述操作所对应的标准函数，以及它们的工作方式。


%%C070355 h3

\subsection{Getting ready}


%%C070356 p

继续之前，你应先阅读\textit{处理文件系统路径}案例。该案例中的入门说明在这里同样适用。不过，本案例的所有示例都与平台无关。


%%C070357 p

下面的所有示例都使用以下变量，并假定当前路径在Windows上是\texttt{C:\textbackslash{}Users\textbackslash{}Marius\textbackslash{}Documents}，在POSIX系统上是\texttt{/\allowbreak{}home/\allowbreak{}marius/\allowbreak{}docs}。同时假定当前路径的\texttt{temp}子目录中存在名为\texttt{sample.\allowbreak{}txt}的文件（如\texttt{C:\textbackslash{}Users\textbackslash{}Marius\textbackslash{}Documents\textbackslash{}temp\textbackslash{}sample.txt}或\myttbr{/\allowbreak{}home/\allowbreak{}marius/\allowbreak{}docs/\allowbreak{}temp/\allowbreak{}sample.\allowbreak{}txt}）：


%%C070358 code

\begin{codeblockbg}
\begin{Verbatim}[commandchars=\\\{\},breaklines=true,breakanywhere=true,fontsize=\codefont,xleftmargin=17bp]
\PYG{k}{auto}\PYG{+w}{ }\PYG{n}{err}\PYG{+w}{ }\PYG{o}{=}\PYG{+w}{ }\PYG{n}{std}\PYG{o}{:}\PYG{o}{:}\PYG{n}{error\PYGZus{}code}\PYG{p}{\PYGZob{}}\PYG{p}{\PYGZcb{}}\PYG{p}{;}
\PYG{k}{auto}\PYG{+w}{ }\PYG{n}{basepath}\PYG{+w}{ }\PYG{o}{=}\PYG{+w}{ }\PYG{n}{fs}\PYG{o}{:}\PYG{o}{:}\PYG{n}{current\PYGZus{}path}\PYG{p}{(}\PYG{p}{)}\PYG{p}{;}
\PYG{k}{auto}\PYG{+w}{ }\PYG{n}{path}\PYG{+w}{ }\PYG{o}{=}\PYG{+w}{ }\PYG{n}{basepath}\PYG{+w}{ }\PYG{o}{/}\PYG{+w}{ }\PYG{l+s}{\PYGZdq{}}\PYG{l+s}{temp}\PYG{l+s}{\PYGZdq{}}\PYG{p}{;}
\PYG{k}{auto}\PYG{+w}{ }\PYG{n}{filepath}\PYG{+w}{ }\PYG{o}{=}\PYG{+w}{ }\PYG{n}{path}\PYG{+w}{ }\PYG{o}{/}\PYG{+w}{ }\PYG{l+s}{\PYGZdq{}}\PYG{l+s}{sample.txt}\PYG{l+s}{\PYGZdq{}}\PYG{p}{;}
\end{Verbatim}
\end{codeblockbg}


%%C070359 h3

\subsection{How to do it...}


%%C070360 p

用以下库函数对目录执行操作：


%%C070361 li

\begin{itemize}
\item 要创建新目录，使用\myttbr{create\_\allowbreak{}directory()}。若目录已存在，该方法什么也不做，但它不会递归创建目录：


%%C070362 code

\begin{codeblockbg}
\begin{Verbatim}[commandchars=\\\{\},breaklines=true,breakanywhere=true,fontsize=\codefont,xleftmargin=17bp]
\PYG{k}{auto}\PYG{+w}{ }\PYG{n}{success}\PYG{+w}{ }\PYG{o}{=}\PYG{+w}{ }\PYG{n}{fs}\PYG{o}{:}\PYG{o}{:}\PYG{n}{create\PYGZus{}directory}\PYG{p}{(}\PYG{n}{path}\PYG{p}{,}\PYG{+w}{ }\PYG{n}{err}\PYG{p}{)}\PYG{p}{;}
\end{Verbatim}
\end{codeblockbg}


%%C070363 li

\item 要递归创建新目录，使用\myttbr{create\_\allowbreak{}directories()}：


%%C070364 code

\begin{codeblockbg}
\begin{Verbatim}[commandchars=\\\{\},breaklines=true,breakanywhere=true,fontsize=\codefont,xleftmargin=17bp]
\PYG{k}{auto}\PYG{+w}{ }\PYG{n}{temp}\PYG{+w}{ }\PYG{o}{=}\PYG{+w}{ }\PYG{n}{path}\PYG{+w}{ }\PYG{o}{/}\PYG{+w}{ }\PYG{l+s}{\PYGZdq{}}\PYG{l+s}{tmp1}\PYG{l+s}{\PYGZdq{}}\PYG{+w}{ }\PYG{o}{/}\PYG{+w}{ }\PYG{l+s}{\PYGZdq{}}\PYG{l+s}{tmp2}\PYG{l+s}{\PYGZdq{}}\PYG{+w}{ }\PYG{o}{/}\PYG{+w}{ }\PYG{l+s}{\PYGZdq{}}\PYG{l+s}{tmp3}\PYG{l+s}{\PYGZdq{}}\PYG{p}{;}
\PYG{k}{auto}\PYG{+w}{ }\PYG{n}{success}\PYG{+w}{ }\PYG{o}{=}\PYG{+w}{ }\PYG{n}{fs}\PYG{o}{:}\PYG{o}{:}\PYG{n}{create\PYGZus{}directories}\PYG{p}{(}\PYG{n}{temp}\PYG{p}{,}\PYG{+w}{ }\PYG{n}{err}\PYG{p}{)}\PYG{p}{;}
\end{Verbatim}
\end{codeblockbg}


%%C070365 li

\item 要移动已有目录，使用\texttt{rename()}：


%%C070366 code

\begin{codeblockbg}
\begin{Verbatim}[commandchars=\\\{\},breaklines=true,breakanywhere=true,fontsize=\codefont,xleftmargin=17bp]
\PYG{k}{auto}\PYG{+w}{ }\PYG{n}{temp}\PYG{+w}{ }\PYG{o}{=}\PYG{+w}{ }\PYG{n}{path}\PYG{+w}{ }\PYG{o}{/}\PYG{+w}{ }\PYG{l+s}{\PYGZdq{}}\PYG{l+s}{tmp1}\PYG{l+s}{\PYGZdq{}}\PYG{+w}{ }\PYG{o}{/}\PYG{+w}{ }\PYG{l+s}{\PYGZdq{}}\PYG{l+s}{tmp2}\PYG{l+s}{\PYGZdq{}}\PYG{+w}{ }\PYG{o}{/}\PYG{+w}{ }\PYG{l+s}{\PYGZdq{}}\PYG{l+s}{tmp3}\PYG{l+s}{\PYGZdq{}}\PYG{p}{;}
\PYG{k}{auto}\PYG{+w}{ }\PYG{n}{newtemp}\PYG{+w}{ }\PYG{o}{=}\PYG{+w}{ }\PYG{n}{path}\PYG{+w}{ }\PYG{o}{/}\PYG{+w}{ }\PYG{l+s}{\PYGZdq{}}\PYG{l+s}{tmp1}\PYG{l+s}{\PYGZdq{}}\PYG{+w}{ }\PYG{o}{/}\PYG{+w}{ }\PYG{l+s}{\PYGZdq{}}\PYG{l+s}{tmp3}\PYG{l+s}{\PYGZdq{}}\PYG{p}{;}
\PYG{n}{fs}\PYG{o}{:}\PYG{o}{:}\PYG{n}{rename}\PYG{p}{(}\PYG{n}{temp}\PYG{p}{,}\PYG{+w}{ }\PYG{n}{newtemp}\PYG{p}{,}\PYG{+w}{ }\PYG{n}{err}\PYG{p}{)}\PYG{p}{;}
\PYG{k}{if}\PYG{+w}{ }\PYG{p}{(}\PYG{n}{err}\PYG{p}{)}\PYG{+w}{ }\PYG{n}{std}\PYG{o}{:}\PYG{o}{:}\PYG{n}{cout}\PYG{+w}{ }\PYG{o}{\PYGZlt{}}\PYG{o}{\PYGZlt{}}\PYG{+w}{ }\PYG{n}{err}\PYG{p}{.}\PYG{n}{message}\PYG{p}{(}\PYG{p}{)}\PYG{+w}{ }\PYG{o}{\PYGZlt{}}\PYG{o}{\PYGZlt{}}\PYG{+w}{ }\PYG{n}{std}\PYG{o}{:}\PYG{o}{:}\PYG{n}{endl}\PYG{p}{;}
\end{Verbatim}
\end{codeblockbg}


%%C070367 li

\item 要重命名已有目录，同样使用\texttt{rename()}：


%%C070368 code

\begin{codeblockbg}
\begin{Verbatim}[commandchars=\\\{\},breaklines=true,breakanywhere=true,fontsize=\codefont,xleftmargin=17bp]
\PYG{k}{auto}\PYG{+w}{ }\PYG{n}{temp}\PYG{+w}{ }\PYG{o}{=}\PYG{+w}{ }\PYG{n}{path}\PYG{+w}{ }\PYG{o}{/}\PYG{+w}{ }\PYG{l+s}{\PYGZdq{}}\PYG{l+s}{tmp1}\PYG{l+s}{\PYGZdq{}}\PYG{+w}{ }\PYG{o}{/}\PYG{+w}{ }\PYG{l+s}{\PYGZdq{}}\PYG{l+s}{tmp3}\PYG{l+s}{\PYGZdq{}}\PYG{p}{;}
\PYG{k}{auto}\PYG{+w}{ }\PYG{n}{newtemp}\PYG{+w}{ }\PYG{o}{=}\PYG{+w}{ }\PYG{n}{path}\PYG{+w}{ }\PYG{o}{/}\PYG{+w}{ }\PYG{l+s}{\PYGZdq{}}\PYG{l+s}{tmp1}\PYG{l+s}{\PYGZdq{}}\PYG{+w}{ }\PYG{o}{/}\PYG{+w}{ }\PYG{l+s}{\PYGZdq{}}\PYG{l+s}{tmp4}\PYG{l+s}{\PYGZdq{}}\PYG{p}{;}
\PYG{n}{fs}\PYG{o}{:}\PYG{o}{:}\PYG{n}{rename}\PYG{p}{(}\PYG{n}{temp}\PYG{p}{,}\PYG{+w}{ }\PYG{n}{newtemp}\PYG{p}{,}\PYG{+w}{ }\PYG{n}{err}\PYG{p}{)}\PYG{p}{;}
\PYG{k}{if}\PYG{+w}{ }\PYG{p}{(}\PYG{n}{err}\PYG{p}{)}\PYG{+w}{ }\PYG{n}{std}\PYG{o}{:}\PYG{o}{:}\PYG{n}{cout}\PYG{+w}{ }\PYG{o}{\PYGZlt{}}\PYG{o}{\PYGZlt{}}\PYG{+w}{ }\PYG{n}{err}\PYG{p}{.}\PYG{n}{message}\PYG{p}{(}\PYG{p}{)}\PYG{+w}{ }\PYG{o}{\PYGZlt{}}\PYG{o}{\PYGZlt{}}\PYG{+w}{ }\PYG{n}{std}\PYG{o}{:}\PYG{o}{:}\PYG{n}{endl}\PYG{p}{;}
\end{Verbatim}
\end{codeblockbg}


%%C070369 li

\item 要复制已有目录，使用\texttt{copy()}。要递归复制目录的全部内容，使用\myttbr{copy\_\allowbreak{}options::\allowbreak{}recursive}标志：


%%C070370 code

\begin{codeblockbg}
\begin{Verbatim}[commandchars=\\\{\},breaklines=true,breakanywhere=true,fontsize=\codefont,xleftmargin=17bp]
\PYG{n}{fs}\PYG{o}{:}\PYG{o}{:}\PYG{n}{copy}\PYG{p}{(}\PYG{n}{basepath}\PYG{+w}{ }\PYG{o}{/}\PYG{+w}{ }\PYG{l+s}{\PYGZdq{}}\PYG{l+s}{temp}\PYG{l+s}{\PYGZdq{}}\PYG{p}{,}\PYG{+w}{ }\PYG{n}{basepath}\PYG{+w}{ }\PYG{o}{/}\PYG{+w}{ }\PYG{l+s}{\PYGZdq{}}\PYG{l+s}{temp2}\PYG{l+s}{\PYGZdq{}}\PYG{p}{,}
\PYG{+w}{         }\PYG{n}{fs}\PYG{o}{:}\PYG{o}{:}\PYG{n}{copy\PYGZus{}options}\PYG{o}{:}\PYG{o}{:}\PYG{n}{recursive}\PYG{p}{,}\PYG{+w}{ }\PYG{n}{err}\PYG{p}{)}\PYG{p}{;}
\PYG{k}{if}\PYG{+w}{ }\PYG{p}{(}\PYG{n}{err}\PYG{p}{)}\PYG{+w}{ }\PYG{n}{std}\PYG{o}{:}\PYG{o}{:}\PYG{n}{cout}\PYG{+w}{ }\PYG{o}{\PYGZlt{}}\PYG{o}{\PYGZlt{}}\PYG{+w}{ }\PYG{n}{err}\PYG{p}{.}\PYG{n}{message}\PYG{p}{(}\PYG{p}{)}\PYG{+w}{ }\PYG{o}{\PYGZlt{}}\PYG{o}{\PYGZlt{}}\PYG{+w}{ }\PYG{n}{std}\PYG{o}{:}\PYG{o}{:}\PYG{n}{endl}\PYG{p}{;}
\end{Verbatim}
\end{codeblockbg}


%%C070371 li

\item 要为目录创建符号链接，使用\myttbr{create\_\allowbreak{}directory\_\allowbreak{}symlink()}：


%%C070372 code

\begin{codeblockbg}
\begin{Verbatim}[commandchars=\\\{\},breaklines=true,breakanywhere=true,fontsize=\codefont,xleftmargin=17bp]
\PYG{k}{auto}\PYG{+w}{ }\PYG{n}{linkdir}\PYG{+w}{ }\PYG{o}{=}\PYG{+w}{ }\PYG{n}{basepath}\PYG{+w}{ }\PYG{o}{/}\PYG{+w}{ }\PYG{l+s}{\PYGZdq{}}\PYG{l+s}{templink}\PYG{l+s}{\PYGZdq{}}\PYG{p}{;}
\PYG{n}{fs}\PYG{o}{:}\PYG{o}{:}\PYG{n}{create\PYGZus{}directory\PYGZus{}symlink}\PYG{p}{(}\PYG{n}{path}\PYG{p}{,}\PYG{+w}{ }\PYG{n}{linkdir}\PYG{p}{,}\PYG{+w}{ }\PYG{n}{err}\PYG{p}{)}\PYG{p}{;}
\PYG{k}{if}\PYG{+w}{ }\PYG{p}{(}\PYG{n}{err}\PYG{p}{)}\PYG{+w}{ }\PYG{n}{std}\PYG{o}{:}\PYG{o}{:}\PYG{n}{cout}\PYG{+w}{ }\PYG{o}{\PYGZlt{}}\PYG{o}{\PYGZlt{}}\PYG{+w}{ }\PYG{n}{err}\PYG{p}{.}\PYG{n}{message}\PYG{p}{(}\PYG{p}{)}\PYG{+w}{ }\PYG{o}{\PYGZlt{}}\PYG{o}{\PYGZlt{}}\PYG{+w}{ }\PYG{n}{std}\PYG{o}{:}\PYG{o}{:}\PYG{n}{endl}\PYG{p}{;}
\end{Verbatim}
\end{codeblockbg}


%%C070373 li

\item 要删除空目录，使用\texttt{remove()}：

%%C070374 code

\begin{codeblockbg}
\begin{Verbatim}[commandchars=\\\{\},breaklines=true,breakanywhere=true,fontsize=\codefont,xleftmargin=17bp]
\PYG{k}{auto}\PYG{+w}{ }\PYG{n}{temp}\PYG{+w}{ }\PYG{o}{=}\PYG{+w}{ }\PYG{n}{path}\PYG{+w}{ }\PYG{o}{/}\PYG{+w}{ }\PYG{l+s}{\PYGZdq{}}\PYG{l+s}{tmp1}\PYG{l+s}{\PYGZdq{}}\PYG{+w}{ }\PYG{o}{/}\PYG{+w}{ }\PYG{l+s}{\PYGZdq{}}\PYG{l+s}{tmp4}\PYG{l+s}{\PYGZdq{}}\PYG{p}{;}
\PYG{k}{auto}\PYG{+w}{ }\PYG{n}{success}\PYG{+w}{ }\PYG{o}{=}\PYG{+w}{ }\PYG{n}{fs}\PYG{o}{:}\PYG{o}{:}\PYG{n}{remove}\PYG{p}{(}\PYG{n}{temp}\PYG{p}{,}\PYG{+w}{ }\PYG{n}{err}\PYG{p}{)}\PYG{p}{;}
\end{Verbatim}
\end{codeblockbg}


%%C070375 li

\item 若要递归删除目录的全部内容以及该目录本身，请使用\texttt{remove\_\allowbreak{}all()}：


%%C070376 code

\begin{codeblockbg}
\begin{Verbatim}[commandchars=\\\{\},breaklines=true,breakanywhere=true,fontsize=\codefont,xleftmargin=17bp]
\PYG{k}{auto}\PYG{+w}{ }\PYG{n}{success}\PYG{+w}{ }\PYG{o}{=}\PYG{+w}{ }\PYG{n}{fs}\PYG{o}{:}\PYG{o}{:}\PYG{n}{remove\PYGZus{}all}\PYG{p}{(}\PYG{n}{path}\PYG{p}{,}\PYG{+w}{ }\PYG{n}{err}\PYG{p}{)}\PYG{+w}{ }\PYG{o}{!}\PYG{o}{=}
\PYG{+w}{               }\PYG{k}{static\PYGZus{}cast}\PYG{o}{\PYGZlt{}}\PYG{n}{std}\PYG{o}{:}\PYG{o}{:}\PYG{k+kt}{uintmax\PYGZus{}t}\PYG{o}{\PYGZgt{}}\PYG{p}{(}\PYG{l+m+mi}{\PYGZhy{}1}\PYG{p}{)}\PYG{p}{;}
\end{Verbatim}
\end{codeblockbg}

\end{itemize}


%%C070377 p

使用以下库函数对文件执行操作：


%%C070378 li

\begin{itemize}
\item 要复制文件，请使用\texttt{copy()}或\texttt{copy\_\allowbreak{}file()}。下一节将说明这两者的区别：


%%C070379 code

\begin{codeblockbg}
\begin{Verbatim}[commandchars=\\\{\},breaklines=true,breakanywhere=true,fontsize=\codefont,xleftmargin=17bp]
\PYG{k}{auto}\PYG{+w}{ }\PYG{n}{success}\PYG{+w}{ }\PYG{o}{=}\PYG{+w}{ }\PYG{n}{fs}\PYG{o}{:}\PYG{o}{:}\PYG{n}{copy\PYGZus{}file}\PYG{p}{(}\PYG{n}{filepath}\PYG{p}{,}\PYG{+w}{ }\PYG{n}{path}\PYG{+w}{ }\PYG{o}{/}\PYG{+w}{ }\PYG{l+s}{\PYGZdq{}}\PYG{l+s}{sample.bak}\PYG{l+s}{\PYGZdq{}}\PYG{p}{,}\PYG{+w}{ }\PYG{n}{err}\PYG{p}{)}\PYG{p}{;}
\PYG{k}{if}\PYG{+w}{ }\PYG{p}{(}\PYG{o}{!}\PYG{n}{success}\PYG{p}{)}\PYG{+w}{ }\PYG{n}{std}\PYG{o}{:}\PYG{o}{:}\PYG{n}{cout}\PYG{+w}{ }\PYG{o}{\PYGZlt{}}\PYG{o}{\PYGZlt{}}\PYG{+w}{ }\PYG{n}{err}\PYG{p}{.}\PYG{n}{message}\PYG{p}{(}\PYG{p}{)}\PYG{+w}{ }\PYG{o}{\PYGZlt{}}\PYG{o}{\PYGZlt{}}\PYG{+w}{ }\PYG{n}{std}\PYG{o}{:}\PYG{o}{:}\PYG{n}{endl}\PYG{p}{;}
\PYG{n}{fs}\PYG{o}{:}\PYG{o}{:}\PYG{n}{copy}\PYG{p}{(}\PYG{n}{filepath}\PYG{p}{,}\PYG{+w}{ }\PYG{n}{path}\PYG{+w}{ }\PYG{o}{/}\PYG{+w}{ }\PYG{l+s}{\PYGZdq{}}\PYG{l+s}{sample.cpy}\PYG{l+s}{\PYGZdq{}}\PYG{p}{,}\PYG{+w}{ }\PYG{n}{err}\PYG{p}{)}\PYG{p}{;}
\PYG{k}{if}\PYG{+w}{ }\PYG{p}{(}\PYG{n}{err}\PYG{p}{)}\PYG{+w}{ }\PYG{n}{std}\PYG{o}{:}\PYG{o}{:}\PYG{n}{cout}\PYG{+w}{ }\PYG{o}{\PYGZlt{}}\PYG{o}{\PYGZlt{}}\PYG{+w}{ }\PYG{n}{err}\PYG{p}{.}\PYG{n}{message}\PYG{p}{(}\PYG{p}{)}\PYG{+w}{ }\PYG{o}{\PYGZlt{}}\PYG{o}{\PYGZlt{}}\PYG{+w}{ }\PYG{n}{std}\PYG{o}{:}\PYG{o}{:}\PYG{n}{endl}\PYG{p}{;}
\end{Verbatim}
\end{codeblockbg}


%%C070380 li

\item 要重命名文件，请使用\texttt{rename()}：


%%C070381 code

\begin{codeblockbg}
\begin{Verbatim}[commandchars=\\\{\},breaklines=true,breakanywhere=true,fontsize=\codefont,xleftmargin=17bp]
\PYG{k}{auto}\PYG{+w}{ }\PYG{n}{newpath}\PYG{+w}{ }\PYG{o}{=}\PYG{+w}{ }\PYG{n}{path}\PYG{+w}{ }\PYG{o}{/}\PYG{+w}{ }\PYG{l+s}{\PYGZdq{}}\PYG{l+s}{sample.log}\PYG{l+s}{\PYGZdq{}}\PYG{p}{;}
\PYG{n}{fs}\PYG{o}{:}\PYG{o}{:}\PYG{n}{rename}\PYG{p}{(}\PYG{n}{filepath}\PYG{p}{,}\PYG{+w}{ }\PYG{n}{newpath}\PYG{p}{,}\PYG{+w}{ }\PYG{n}{err}\PYG{p}{)}\PYG{p}{;}
\PYG{k}{if}\PYG{+w}{ }\PYG{p}{(}\PYG{n}{err}\PYG{p}{)}\PYG{+w}{ }\PYG{n}{std}\PYG{o}{:}\PYG{o}{:}\PYG{n}{cout}\PYG{+w}{ }\PYG{o}{\PYGZlt{}}\PYG{o}{\PYGZlt{}}\PYG{+w}{ }\PYG{n}{err}\PYG{p}{.}\PYG{n}{message}\PYG{p}{(}\PYG{p}{)}\PYG{+w}{ }\PYG{o}{\PYGZlt{}}\PYG{o}{\PYGZlt{}}\PYG{+w}{ }\PYG{n}{std}\PYG{o}{:}\PYG{o}{:}\PYG{n}{endl}\PYG{p}{;}
\end{Verbatim}
\end{codeblockbg}


%%C070382 li

\item 要移动文件，请使用\texttt{rename()}：


%%C070383 code

\begin{codeblockbg}
\begin{Verbatim}[commandchars=\\\{\},breaklines=true,breakanywhere=true,fontsize=\codefont,xleftmargin=17bp]
\PYG{k}{auto}\PYG{+w}{ }\PYG{n}{newpath}\PYG{+w}{ }\PYG{o}{=}\PYG{+w}{ }\PYG{n}{path}\PYG{+w}{ }\PYG{o}{/}\PYG{+w}{ }\PYG{l+s}{\PYGZdq{}}\PYG{l+s}{sample.log}\PYG{l+s}{\PYGZdq{}}\PYG{p}{;}
\PYG{n}{fs}\PYG{o}{:}\PYG{o}{:}\PYG{n}{rename}\PYG{p}{(}\PYG{n}{newpath}\PYG{p}{,}\PYG{+w}{ }\PYG{n}{path}\PYG{+w}{ }\PYG{o}{/}\PYG{+w}{ }\PYG{l+s}{\PYGZdq{}}\PYG{l+s}{tmp1}\PYG{l+s}{\PYGZdq{}}\PYG{+w}{ }\PYG{o}{/}\PYG{+w}{ }\PYG{l+s}{\PYGZdq{}}\PYG{l+s}{sample.log}\PYG{l+s}{\PYGZdq{}}\PYG{p}{,}\PYG{+w}{ }\PYG{n}{err}\PYG{p}{)}\PYG{p}{;}
\PYG{k}{if}\PYG{+w}{ }\PYG{p}{(}\PYG{n}{err}\PYG{p}{)}\PYG{+w}{ }\PYG{n}{std}\PYG{o}{:}\PYG{o}{:}\PYG{n}{cout}\PYG{+w}{ }\PYG{o}{\PYGZlt{}}\PYG{o}{\PYGZlt{}}\PYG{+w}{ }\PYG{n}{err}\PYG{p}{.}\PYG{n}{message}\PYG{p}{(}\PYG{p}{)}\PYG{+w}{ }\PYG{o}{\PYGZlt{}}\PYG{o}{\PYGZlt{}}\PYG{+w}{ }\PYG{n}{std}\PYG{o}{:}\PYG{o}{:}\PYG{n}{endl}\PYG{p}{;}
\end{Verbatim}
\end{codeblockbg}


%%C070384 li

\item 要为文件创建符号链接，请使用\texttt{create\_\allowbreak{}symlink()}：


%%C070385 code

\begin{codeblockbg}
\begin{Verbatim}[commandchars=\\\{\},breaklines=true,breakanywhere=true,fontsize=\codefont,xleftmargin=17bp]
\PYG{k}{auto}\PYG{+w}{ }\PYG{n}{linkpath}\PYG{+w}{ }\PYG{o}{=}\PYG{+w}{ }\PYG{n}{path}\PYG{+w}{ }\PYG{o}{/}\PYG{+w}{ }\PYG{l+s}{\PYGZdq{}}\PYG{l+s}{sample.txt.link}\PYG{l+s}{\PYGZdq{}}\PYG{p}{;}
\PYG{n}{fs}\PYG{o}{:}\PYG{o}{:}\PYG{n}{create\PYGZus{}symlink}\PYG{p}{(}\PYG{n}{filepath}\PYG{p}{,}\PYG{+w}{ }\PYG{n}{linkpath}\PYG{p}{,}\PYG{+w}{ }\PYG{n}{err}\PYG{p}{)}\PYG{p}{;}
\PYG{k}{if}\PYG{+w}{ }\PYG{p}{(}\PYG{n}{err}\PYG{p}{)}\PYG{+w}{ }\PYG{n}{std}\PYG{o}{:}\PYG{o}{:}\PYG{n}{cout}\PYG{+w}{ }\PYG{o}{\PYGZlt{}}\PYG{o}{\PYGZlt{}}\PYG{+w}{ }\PYG{n}{err}\PYG{p}{.}\PYG{n}{message}\PYG{p}{(}\PYG{p}{)}\PYG{+w}{ }\PYG{o}{\PYGZlt{}}\PYG{o}{\PYGZlt{}}\PYG{+w}{ }\PYG{n}{std}\PYG{o}{:}\PYG{o}{:}\PYG{n}{endl}\PYG{p}{;}
\end{Verbatim}
\end{codeblockbg}


%%C070386 li

\item 要删除文件，请使用\texttt{remove()}：


%%C070387 code

\begin{codeblockbg}
\begin{Verbatim}[commandchars=\\\{\},breaklines=true,breakanywhere=true,fontsize=\codefont,xleftmargin=17bp]
\PYG{k}{auto}\PYG{+w}{ }\PYG{n}{success}\PYG{+w}{ }\PYG{o}{=}\PYG{+w}{ }\PYG{n}{fs}\PYG{o}{:}\PYG{o}{:}\PYG{n}{remove}\PYG{p}{(}\PYG{n}{path}\PYG{+w}{ }\PYG{o}{/}\PYG{+w}{ }\PYG{l+s}{\PYGZdq{}}\PYG{l+s}{sample.cpy}\PYG{l+s}{\PYGZdq{}}\PYG{p}{,}\PYG{+w}{ }\PYG{n}{err}\PYG{p}{)}\PYG{p}{;}
\PYG{k}{if}\PYG{+w}{ }\PYG{p}{(}\PYG{o}{!}\PYG{n}{success}\PYG{p}{)}\PYG{+w}{ }\PYG{n}{std}\PYG{o}{:}\PYG{o}{:}\PYG{n}{cout}\PYG{+w}{ }\PYG{o}{\PYGZlt{}}\PYG{o}{\PYGZlt{}}\PYG{+w}{ }\PYG{n}{err}\PYG{p}{.}\PYG{n}{message}\PYG{p}{(}\PYG{p}{)}\PYG{+w}{ }\PYG{o}{\PYGZlt{}}\PYG{o}{\PYGZlt{}}\PYG{+w}{ }\PYG{n}{std}\PYG{o}{:}\PYG{o}{:}\PYG{n}{endl}\PYG{p}{;}
\end{Verbatim}
\end{codeblockbg}

\end{itemize}


%%C070388 h3

\subsection{工作原理}


%%C070389 p

本案例中提到的所有函数，以及这里未讨论的其他类似函数，都有多个重载，可分为两类：


%%C070390 li

\begin{itemize}
\item 最后一个参数为\texttt{std::\allowbreak{}error\_\allowbreak{}code}引用的重载：这些重载不抛出异常（它们以\texttt{noexcept}说明定义），而是在发生操作系统错误时把\texttt{error\_\allowbreak{}code}对象的值设为该操作系统错误码。如果没有发生这类错误，则调用\texttt{error\_\allowbreak{}code}对象的\texttt{clear()}方法，以重置此前可能设置过的错误码。


%%C070391 li

\item 最后一个参数不是\texttt{std::\allowbreak{}error\_\allowbreak{}code}类型的重载：这些重载在出错时抛出异常。如果发生操作系统错误，它们抛出\myttbr{std::\allowbreak{}filesystem::\allowbreak{}filesystem\_\allowbreak{}error}异常。另一方面，如果内存分配失败，这些函数抛出\texttt{std::\allowbreak{}bad\_\allowbreak{}alloc}异常。

\end{itemize}


%%C070392 p

上一节中的所有示例都使用了不抛异常、而是在出错时设置错误码的重载。有些函数返回\texttt{bool}以表示成功或失败。你可以通过两种方式检查\texttt{error\_\allowbreak{}code}对象是否保存了错误码：一是检查\texttt{value()}方法返回的错误码值是否不等于零，二是使用转换运算符\texttt{operator bool}，它在同样的情况下返回\texttt{true}，否则返回\texttt{false}。要获取错误码的说明性字符串，请使用\texttt{message()}方法。


%%C070393 p

\texttt{filesystem}库中的某些函数对文件和目录都通用，\texttt{rename()}、\texttt{remove()}和\texttt{copy()}就是如此。这些函数各自的工作细节可能相当复杂，尤其是\texttt{copy()}，这超出了本案例的范围。如果你需要执行的并非这里介绍的简单操作，应查阅参考文档。


%%C070394 p

说到复制文件，有两个函数可用：\texttt{copy()}和\texttt{copy\_\allowbreak{}file()}。它们具有签名完全相同的等价重载，表面上行为也一样。然而，两者有一个重要区别（除了\texttt{copy()}对目录也适用这一点之外）：\texttt{copy\_\allowbreak{}file()}会跟随符号链接。若要避免这样做，改为复制符号链接本身，你必须使用\texttt{copy\_\allowbreak{}symlink()}，或者使用带\myttbr{copy\_\allowbreak{}options::\allowbreak{}copy\_\allowbreak{}symlinks}标志的\texttt{copy()}。\texttt{copy()}和\texttt{copy\_\allowbreak{}file()}都有一个重载接受\myttbr{std::\allowbreak{}filesystem::\allowbreak{}copy\_\allowbreak{}options}类型的参数，用于定义操作应如何执行。\texttt{copy\_\allowbreak{}options}是一个限定作用域\texttt{enum}，定义如下：


%%C070395 code

\begin{codeblockbg}
\begin{Verbatim}[commandchars=\\\{\},breaklines=true,breakanywhere=true,fontsize=\codefont,xleftmargin=17bp]
\PYG{k}{enum}\PYG{+w}{ }\PYG{k}{class}\PYG{+w}{ }\PYG{n+nc}{copy\PYGZus{}options}
\PYG{p}{\PYGZob{}}
\PYG{+w}{  }\PYG{n}{none}\PYG{+w}{ }\PYG{o}{=}\PYG{+w}{ }\PYG{l+m+mi}{0}\PYG{p}{,}
\PYG{+w}{  }\PYG{n}{skip\PYGZus{}existing}\PYG{+w}{ }\PYG{o}{=}\PYG{+w}{ }\PYG{l+m+mi}{1}\PYG{p}{,}
\PYG{+w}{  }\PYG{n}{overwrite\PYGZus{}existing}\PYG{+w}{ }\PYG{o}{=}\PYG{+w}{ }\PYG{l+m+mi}{2}\PYG{p}{,}
\PYG{+w}{  }\PYG{n}{update\PYGZus{}existing}\PYG{+w}{ }\PYG{o}{=}\PYG{+w}{ }\PYG{l+m+mi}{4}\PYG{p}{,}
\PYG{+w}{  }\PYG{n}{recursive}\PYG{+w}{ }\PYG{o}{=}\PYG{+w}{ }\PYG{l+m+mi}{8}\PYG{p}{,}
\PYG{+w}{  }\PYG{n}{copy\PYGZus{}symlinks}\PYG{+w}{ }\PYG{o}{=}\PYG{+w}{ }\PYG{l+m+mi}{16}\PYG{p}{,}
\PYG{+w}{  }\PYG{n}{skip\PYGZus{}symlinks}\PYG{+w}{ }\PYG{o}{=}\PYG{+w}{ }\PYG{l+m+mi}{32}\PYG{p}{,}
\PYG{+w}{  }\PYG{n}{directories\PYGZus{}only}\PYG{+w}{ }\PYG{o}{=}\PYG{+w}{ }\PYG{l+m+mi}{64}\PYG{p}{,}
\PYG{+w}{  }\PYG{n}{create\PYGZus{}symlinks}\PYG{+w}{ }\PYG{o}{=}\PYG{+w}{ }\PYG{l+m+mi}{128}\PYG{p}{,}
\PYG{+w}{  }\PYG{n}{create\PYGZus{}hard\PYGZus{}links}\PYG{+w}{ }\PYG{o}{=}\PYG{+w}{ }\PYG{l+m+mi}{256}
\PYG{p}{\PYGZcb{}}\PYG{p}{;}
\end{Verbatim}
\end{codeblockbg}


%%C070396 p

下表定义了这些标志如何影响复制操作（无论使用\texttt{copy()}还是\texttt{copy\_\allowbreak{}file()}）。该表取自C++17标准的27.10.10.4节：


%%C070397 table

\myTableBegin
\begin{xltabular}{\textwidth}{@{}>{\hsize=0.512\hsize\raggedright\arraybackslash}X>{\hsize=0.488\hsize\raggedright\arraybackslash}X@{}}
\rowcolor{tabhead} \textcolor{white}{\textbf{控制copy\_file函数}} \textbf{对已存在目标文件的影响的选项组} & \textcolor{white}{} \\
\endfirsthead
\rowcolor{tabhead} \textcolor{white}{\textbf{控制copy\_file函数}} \textbf{对已存在目标文件的影响的选项组} & \textcolor{white}{} \\
\endhead
\myttbr{none} & （默认）报错；文件已存在\\
\myttbr{skip\_\allowbreak{}existing} & 不覆盖已存在的文件；不报告错误\\
\myttbr{overwrite\_\allowbreak{}existing} & 覆盖已存在的文件\\
\myttbr{update\_\allowbreak{}existing} & 若已存在的文件比用于替换的文件更旧，则覆盖它\\
\textbf{控制copy函数} \textbf{对子目录的影响的选项组} &  \\
\myttbr{none} & （默认）不复制子目录\\
\end{xltabular}
\myTableEnd


%%C070398 table

\myTableBegin
\begin{xltabular}{\textwidth}{@{}>{\hsize=0.512\hsize\raggedright\arraybackslash}X>{\hsize=0.488\hsize\raggedright\arraybackslash}X@{}}
\rowcolor{tabhead} \textcolor{white}{\myttbr{recursive}} & \textcolor{white}{递归复制子目录及其内容} \\
\endfirsthead
\rowcolor{tabhead} \textcolor{white}{\myttbr{recursive}} & \textcolor{white}{递归复制子目录及其内容} \\
\endhead
\textbf{控制copy函数} \textbf{对符号链接的影响的选项组} &  \\
\myttbr{none} & （默认）跟随符号链接\\
\myttbr{copy\_\allowbreak{}symlinks} & 把符号链接复制为符号链接，而不是复制它们所指向的文件\\
\myttbr{skip\_\allowbreak{}symlinks} & 忽略符号链接\\
\textbf{控制copy函数} \textbf{选择复制形式的选项组} &  \\
\myttbr{none} & （默认）复制内容\\
\myttbr{directories\_\allowbreak{}only} & 只复制目录结构，不复制非目录文件\\
\myttbr{create\_\allowbreak{}symlinks} & 创建符号链接而不是复制文件；除非目标路径位于当前目录，否则源路径应为绝对路径\\
\myttbr{create\_\allowbreak{}hard\_\allowbreak{}links} & 创建硬链接而不是复制文件\\
\end{xltabular}
\myTableEnd


%%C070399 p

另一个需要提及的方面与符号链接有关：\myttbr{create\_\allowbreak{}directory\_\allowbreak{}symlink()}创建指向目录的符号链接，而\texttt{create\_\allowbreak{}symlink()}既可创建指向文件的符号链接，也可创建指向目录的符号链接。在POSIX系统上，对于目录而言这两者是相同的。在其他系统（如Windows）上，指向目录的符号链接的创建方式与指向文件的符号链接不同。因此，建议对目录使用\myttbr{create\_\allowbreak{}directory\_\allowbreak{}symlink()}，以便写出在所有系统上都能正确工作的代码。


%%C070400 note

\begin{myNote}[tip]
当你对文件和目录执行本案例所述的操作，并使用可能抛出异常的重载时，务必对这些调用进行try-catch。无论使用哪种重载，都应检查操作是否成功，并在失败时采取相应措施。
\end{myNote}


%%C070401 h3

\subsection{延伸阅读}


%%C070402 li

\begin{itemize}
\item \textit{处理文件系统路径}


%%C070403 li

\item \textit{清除文件内容}


%%C070404 li

\item \textit{检查现有文件或目录的属性}

\end{itemize}


%%C070405 h2

\section{清除文件内容}


%%C070406 p

复制、重命名、移动或删除文件等操作由\texttt{filesystem}库直接提供。不过，要从文件中删除内容，你必须显式执行操作。无论处理文本文件还是二进制文件，都必须实现以下模式：


%%C070407 p

创建一个临时文件。


%%C070408 p

只把原文件中你想要的内容复制到临时文件。


%%C070409 p

删除原文件。


%%C070410 p

把临时文件重命名/移动到原文件的名称/位置。


%%C070411 h3

\subsection{准备工作}


%%C070412 p

本案例将演示如何对文本文件实现前面提到的模式。为此，我们考虑删除空行或以分号（\texttt{;}）开头的行。在这个例子中，初始文件名为\texttt{sample.\allowbreak{}dat}，其中包含莎士比亚剧作的名称，但也包含空行和以分号开头的行。以下是该文件的部分内容（从开头开始）：


%%C070413 code

\begin{codeblockbg}
\begin{Verbatim}[commandchars=\\\{\},breaklines=true,breakanywhere=true,fontsize=\codefont,xleftmargin=17bp]
\PYG{p}{;}\PYG{n}{Shakespeare}\PYG{err}{\PYGZsq{}}\PYG{n}{s}\PYG{+w}{ }\PYG{n}{plays}\PYG{p}{,}\PYG{+w}{ }\PYG{n}{listed}\PYG{+w}{ }\PYG{n}{by}\PYG{+w}{ }\PYG{n}{genre}
\PYG{p}{;}\PYG{n}{TRAGEDIES}
\PYG{n}{Troilus}\PYG{+w}{ }\PYG{k}{and}\PYG{+w}{ }\PYG{n}{Cressida}
\PYG{n}{Coriolanus}
\PYG{n}{Titus}\PYG{+w}{ }\PYG{n}{Andronicus}
\PYG{n}{Romeo}\PYG{+w}{ }\PYG{k}{and}\PYG{+w}{ }\PYG{n}{Juliet}
\PYG{n}{Timon}\PYG{+w}{ }\PYG{n}{of}\PYG{+w}{ }\PYG{n}{Athens}
\PYG{n}{Julius}\PYG{+w}{ }\PYG{n}{Caesar}
\end{Verbatim}
\end{codeblockbg}


%%C070414 p

下一节列出的代码示例使用以下变量：


%%C070415 code

\begin{codeblockbg}
\begin{Verbatim}[commandchars=\\\{\},breaklines=true,breakanywhere=true,fontsize=\codefont,xleftmargin=17bp]
\PYG{k}{auto}\PYG{+w}{ }\PYG{n}{path}\PYG{+w}{ }\PYG{o}{=}\PYG{+w}{ }\PYG{n}{fs}\PYG{o}{:}\PYG{o}{:}\PYG{n}{current\PYGZus{}path}\PYG{p}{(}\PYG{p}{)}\PYG{p}{;}
\PYG{k}{auto}\PYG{+w}{ }\PYG{n}{filepath}\PYG{+w}{ }\PYG{o}{=}\PYG{+w}{ }\PYG{n}{path}\PYG{+w}{ }\PYG{o}{/}\PYG{+w}{ }\PYG{l+s}{\PYGZdq{}}\PYG{l+s}{sample.dat}\PYG{l+s}{\PYGZdq{}}\PYG{p}{;}
\PYG{k}{auto}\PYG{+w}{ }\PYG{n}{temppath}\PYG{+w}{ }\PYG{o}{=}\PYG{+w}{ }\PYG{n}{path}\PYG{+w}{ }\PYG{o}{/}\PYG{+w}{ }\PYG{l+s}{\PYGZdq{}}\PYG{l+s}{sample.tmp}\PYG{l+s}{\PYGZdq{}}\PYG{p}{;}
\PYG{k}{auto}\PYG{+w}{ }\PYG{n}{err}\PYG{+w}{ }\PYG{o}{=}\PYG{+w}{ }\PYG{n}{std}\PYG{o}{:}\PYG{o}{:}\PYG{n}{error\PYGZus{}code}\PYG{p}{\PYGZob{}}\PYG{p}{\PYGZcb{}}\PYG{p}{;}
\end{Verbatim}
\end{codeblockbg}


%%C070416 h3

\subsection{具体做法}


%%C070417 p

执行以下操作从文件中删除内容：


%%C070418 p

以读取方式打开文件：


%%C070419 code

\begin{codeblockbg}
\begin{Verbatim}[commandchars=\\\{\},breaklines=true,breakanywhere=true,fontsize=\codefont,xleftmargin=17bp]
\PYG{n}{std}\PYG{o}{:}\PYG{o}{:}\PYG{n}{ifstream}\PYG{+w}{ }\PYG{n+nf}{in}\PYG{p}{(}\PYG{n}{filepath}\PYG{p}{)}\PYG{p}{;}
\PYG{k}{if}\PYG{+w}{ }\PYG{p}{(}\PYG{o}{!}\PYG{n}{in}\PYG{p}{.}\PYG{n}{is\PYGZus{}open}\PYG{p}{(}\PYG{p}{)}\PYG{p}{)}
\PYG{p}{\PYGZob{}}
\PYG{+w}{  }\PYG{n}{std}\PYG{o}{:}\PYG{o}{:}\PYG{n}{cout}\PYG{+w}{ }\PYG{o}{\PYGZlt{}}\PYG{o}{\PYGZlt{}}\PYG{+w}{ }\PYG{l+s}{\PYGZdq{}}\PYG{l+s}{File could not be opened!}\PYG{l+s}{\PYGZdq{}}\PYG{+w}{ }\PYG{o}{\PYGZlt{}}\PYG{o}{\PYGZlt{}}\PYG{+w}{ }\PYG{n}{std}\PYG{o}{:}\PYG{o}{:}\PYG{n}{endl}\PYG{p}{;}
\PYG{+w}{  }\PYG{k}{return}\PYG{p}{;}
\PYG{p}{\PYGZcb{}}
\end{Verbatim}
\end{codeblockbg}


%%C070420 p

打开另一个临时文件用于写入；如果该文件已存在，则截断其


%%C070421 p

内容：


%%C070422 code

\begin{codeblockbg}
\begin{Verbatim}[commandchars=\\\{\},breaklines=true,breakanywhere=true,fontsize=\codefont,xleftmargin=17bp]
\PYG{n}{std}\PYG{o}{:}\PYG{o}{:}\PYG{n}{ofstream}\PYG{+w}{ }\PYG{n+nf}{out}\PYG{p}{(}\PYG{n}{temppath}\PYG{p}{,}\PYG{+w}{ }\PYG{n}{std}\PYG{o}{:}\PYG{o}{:}\PYG{n}{ios}\PYG{o}{:}\PYG{o}{:}\PYG{n}{trunc}\PYG{p}{)}\PYG{p}{;}
\PYG{k}{if}\PYG{+w}{ }\PYG{p}{(}\PYG{o}{!}\PYG{n}{out}\PYG{p}{.}\PYG{n}{is\PYGZus{}open}\PYG{p}{(}\PYG{p}{)}\PYG{p}{)}
\PYG{p}{\PYGZob{}}
\PYG{+w}{  }\PYG{n}{std}\PYG{o}{:}\PYG{o}{:}\PYG{n}{cout}\PYG{+w}{ }\PYG{o}{\PYGZlt{}}\PYG{o}{\PYGZlt{}}\PYG{+w}{ }\PYG{l+s}{\PYGZdq{}}\PYG{l+s}{Temporary file could not be created!}\PYG{l+s}{\PYGZdq{}}
\PYG{+w}{            }\PYG{o}{\PYGZlt{}}\PYG{o}{\PYGZlt{}}\PYG{+w}{ }\PYG{n}{std}\PYG{o}{:}\PYG{o}{:}\PYG{n}{endl}\PYG{p}{;}
\PYG{+w}{  }\PYG{k}{return}\PYG{p}{;}
\PYG{p}{\PYGZcb{}}
\end{Verbatim}
\end{codeblockbg}


%%C070423 p

从输入文件逐行读取，并把选定的内容复制到输出


%%C070424 p

文件：


%%C070425 code

\begin{codeblockbg}
\begin{Verbatim}[commandchars=\\\{\},breaklines=true,breakanywhere=true,fontsize=\codefont,xleftmargin=17bp]
\PYG{k}{auto}\PYG{+w}{ }\PYG{n}{line}\PYG{+w}{ }\PYG{o}{=}\PYG{+w}{ }\PYG{n}{std}\PYG{o}{:}\PYG{o}{:}\PYG{n}{string}\PYG{p}{\PYGZob{}}\PYG{p}{\PYGZcb{}}\PYG{p}{;}
\PYG{k}{while}\PYG{+w}{ }\PYG{p}{(}\PYG{n}{std}\PYG{o}{:}\PYG{o}{:}\PYG{n}{getline}\PYG{p}{(}\PYG{n}{in}\PYG{p}{,}\PYG{+w}{ }\PYG{n}{line}\PYG{p}{)}\PYG{p}{)}
\PYG{p}{\PYGZob{}}
\PYG{+w}{  }\PYG{k}{if}\PYG{+w}{ }\PYG{p}{(}\PYG{o}{!}\PYG{n}{line}\PYG{p}{.}\PYG{n}{empty}\PYG{p}{(}\PYG{p}{)}\PYG{+w}{ }\PYG{o}{\PYGZam{}}\PYG{o}{\PYGZam{}}\PYG{+w}{ }\PYG{n}{line}\PYG{p}{.}\PYG{n}{at}\PYG{p}{(}\PYG{l+m+mi}{0}\PYG{p}{)}\PYG{+w}{ }\PYG{o}{!}\PYG{o}{=}\PYG{+w}{ }\PYG{l+s+sc}{\PYGZsq{}}\PYG{l+s+sc}{;}\PYG{l+s+sc}{\PYGZsq{}}\PYG{p}{)}
\PYG{+w}{  }\PYG{p}{\PYGZob{}}
\PYG{+w}{    }\PYG{n}{out}\PYG{+w}{ }\PYG{o}{\PYGZlt{}}\PYG{o}{\PYGZlt{}}\PYG{+w}{ }\PYG{n}{line}\PYG{+w}{ }\PYG{o}{\PYGZlt{}}\PYG{o}{\PYGZlt{}}\PYG{+w}{ }\PYG{l+s+sc}{\PYGZsq{}}\PYG{l+s+sc}{n}\PYG{l+s+sc}{\PYGZsq{}}\PYG{p}{;}
\PYG{+w}{  }\PYG{p}{\PYGZcb{}}
\PYG{p}{\PYGZcb{}}
\end{Verbatim}
\end{codeblockbg}


%%C070426 p

关闭输入和输出文件：


%%C070427 code

\begin{codeblockbg}
\begin{Verbatim}[commandchars=\\\{\},breaklines=true,breakanywhere=true,fontsize=\codefont,xleftmargin=17bp]
\PYG{n}{in}\PYG{p}{.}\PYG{n}{close}\PYG{p}{(}\PYG{p}{)}\PYG{p}{;}
\PYG{n}{out}\PYG{p}{.}\PYG{n}{close}\PYG{p}{(}\PYG{p}{)}\PYG{p}{;}
\end{Verbatim}
\end{codeblockbg}


%%C070428 p

删除原文件：


%%C070429 code

\begin{codeblockbg}
\begin{Verbatim}[commandchars=\\\{\},breaklines=true,breakanywhere=true,fontsize=\codefont,xleftmargin=17bp]
\PYG{k}{auto}\PYG{+w}{ }\PYG{n}{success}\PYG{+w}{ }\PYG{o}{=}\PYG{+w}{ }\PYG{n}{fs}\PYG{o}{:}\PYG{o}{:}\PYG{n}{remove}\PYG{p}{(}\PYG{n}{filepath}\PYG{p}{,}\PYG{+w}{ }\PYG{n}{err}\PYG{p}{)}\PYG{p}{;}
\PYG{k}{if}\PYG{p}{(}\PYG{o}{!}\PYG{n}{success}\PYG{+w}{ }\PYG{o}{|}\PYG{o}{|}\PYG{+w}{ }\PYG{n}{err}\PYG{p}{)}
\PYG{p}{\PYGZob{}}
\PYG{+w}{  }\PYG{n}{std}\PYG{o}{:}\PYG{o}{:}\PYG{n}{cout}\PYG{+w}{ }\PYG{o}{\PYGZlt{}}\PYG{o}{\PYGZlt{}}\PYG{+w}{ }\PYG{n}{err}\PYG{p}{.}\PYG{n}{message}\PYG{p}{(}\PYG{p}{)}\PYG{+w}{ }\PYG{o}{\PYGZlt{}}\PYG{o}{\PYGZlt{}}\PYG{+w}{ }\PYG{n}{std}\PYG{o}{:}\PYG{o}{:}\PYG{n}{endl}\PYG{p}{;}
\PYG{+w}{  }\PYG{k}{return}\PYG{p}{;}
\PYG{p}{\PYGZcb{}}
\end{Verbatim}
\end{codeblockbg}


%%C070430 p

把临时文件重命名/移动到原文件的名称/位置：


%%C070431 code

\begin{codeblockbg}
\begin{Verbatim}[commandchars=\\\{\},breaklines=true,breakanywhere=true,fontsize=\codefont,xleftmargin=17bp]
\PYG{n}{fs}\PYG{o}{:}\PYG{o}{:}\PYG{n}{rename}\PYG{p}{(}\PYG{n}{temppath}\PYG{p}{,}\PYG{+w}{ }\PYG{n}{filepath}\PYG{p}{,}\PYG{+w}{ }\PYG{n}{err}\PYG{p}{)}\PYG{p}{;}
\PYG{k}{if}\PYG{+w}{ }\PYG{p}{(}\PYG{n}{err}\PYG{p}{)}
\PYG{p}{\PYGZob{}}
\PYG{+w}{  }\PYG{n}{std}\PYG{o}{:}\PYG{o}{:}\PYG{n}{cout}\PYG{+w}{ }\PYG{o}{\PYGZlt{}}\PYG{o}{\PYGZlt{}}\PYG{+w}{ }\PYG{n}{err}\PYG{p}{.}\PYG{n}{message}\PYG{p}{(}\PYG{p}{)}\PYG{+w}{ }\PYG{o}{\PYGZlt{}}\PYG{o}{\PYGZlt{}}\PYG{+w}{ }\PYG{n}{std}\PYG{o}{:}\PYG{o}{:}\PYG{n}{endl}\PYG{p}{;}
\PYG{p}{\PYGZcb{}}
\end{Verbatim}
\end{codeblockbg}


%%C070432 h3

\subsection{工作原理}


%%C070433 p

上述模式对二进制文件同样适用，但为方便起见，我们只讨论文本文件的例子。本例中的临时文件与原文件位于同一目录。也可以把它放在单独的目录中，例如用户临时目录。要获取临时目录的路径，可以使用\myttbr{std::\allowbreak{}filesystem::\allowbreak{}temp\_\allowbreak{}directory\_\allowbreak{}path()}。在Windows系统上，该函数返回与\texttt{GetTempPath()}相同的目录。在POSIX系统上，它返回环境变量\texttt{TMPDIR}、\texttt{TMP}、\texttt{TEMP}或\texttt{TEMPDIR}中之一指定的路径；如果这些变量都不可用，则返回路径\texttt{/\allowbreak{}tmp}。


%%C070434 p

原文件内容如何复制到临时文件，因具体情况和需要复制的内容而异。在前面的例子中，我们复制整行，除非该行为空或以分号开头。为此，我们用\texttt{std::\allowbreak{}getline()}逐行读取原文件的内容，直到没有更多行可读。复制完所有必需的内容后，应关闭文件，以便随后移动或删除它们。


%%C070435 p

要完成该操作，有以下三种选择：


%%C070436 li

\begin{itemize}
\item 删除原文件，并把临时文件重命名为与原文件相同的名称（若二者位于同一目录），或者在二者位于不同目录时把临时文件移动到原文件所在位置。这是本案例采用的做法。为此，我们用\texttt{remove()}函数删除原文件，用\texttt{rename()}把临时文件重命名为原文件名。


%%C070437 li

\item 把临时文件的内容复制到原文件（可以使用\texttt{copy()}或\texttt{copy\_\allowbreak{}file()}函数），然后删除临时文件（用\texttt{remove()}）。


%%C070438 li

\item 重命名原文件（例如更改扩展名或名称），然后用原文件名重命名/移动临时文件。

\end{itemize}


%%C070439 note

\begin{myNote}[tip]
如果采用上面第一种做法，你必须确保随后替换原文件的临时文件具有与原文件相同的文件权限，否则根据你的解决方案所处的上下文，这可能引发问题。
\end{myNote}


%%C070440 h3

\subsection{延伸阅读}


%%C070441 li

\begin{itemize}
\item \textit{创建、复制与删除文件和目录}

\end{itemize}


%%C070442 h2

\section{检查现有文件或目录的属性}


%%C070443 p

\texttt{filesystem}库提供的函数和类型使开发者能够检查文件系统对象（如文件或目录）是否存在，以及它的属性，例如类型（文件、目录、符号链接等）、最后写入时间、权限等。本案例将介绍这些类型和函数是什么，以及如何使用它们。


%%C070444 h3

\subsection{准备工作}


%%C070445 p

在继续本案例之前，你应阅读\textit{处理文件系统路径}案例。


%%C070446 p

对于以下代码示例，我们将为\texttt{std::\allowbreak{}filesystem}命名空间使用命名空间别名\texttt{fs}。\texttt{filesystem}库位于同名头文件\texttt{<\allowbreak{}filesystem>}中。此外，我们将使用这里展示的变量：\texttt{path}表示文件的路径，\texttt{err}用于接收文件系统API可能返回的操作系统错误码：


%%C070447 code

\begin{codeblockbg}
\begin{Verbatim}[commandchars=\\\{\},breaklines=true,breakanywhere=true,fontsize=\codefont,xleftmargin=17bp]
\PYG{k}{auto}\PYG{+w}{ }\PYG{n}{path}\PYG{+w}{ }\PYG{o}{=}\PYG{+w}{ }\PYG{n}{fs}\PYG{o}{:}\PYG{o}{:}\PYG{n}{current\PYGZus{}path}\PYG{p}{(}\PYG{p}{)}\PYG{+w}{ }\PYG{o}{/}\PYG{+w}{ }\PYG{l+s}{\PYGZdq{}}\PYG{l+s}{main.cpp}\PYG{l+s}{\PYGZdq{}}\PYG{p}{;}
\PYG{k}{auto}\PYG{+w}{ }\PYG{n}{err}\PYG{+w}{ }\PYG{o}{=}\PYG{+w}{ }\PYG{n}{std}\PYG{o}{:}\PYG{o}{:}\PYG{n}{error\PYGZus{}code}\PYG{p}{\PYGZob{}}\PYG{p}{\PYGZcb{}}\PYG{p}{;}
\end{Verbatim}
\end{codeblockbg}


%%C070448 h3

\subsection{具体做法}


%%C070449 p

使用以下库函数获取文件系统对象的信息：


%%C070450 li

\begin{itemize}
\item 要检查某个路径是否指向已存在的文件系统对象，请使用\texttt{exists()}：


%%C070451 code

\begin{codeblockbg}
\begin{Verbatim}[commandchars=\\\{\},breaklines=true,breakanywhere=true,fontsize=\codefont,xleftmargin=17bp]
\PYG{k}{auto}\PYG{+w}{ }\PYG{n}{exists}\PYG{+w}{ }\PYG{o}{=}\PYG{+w}{ }\PYG{n}{fs}\PYG{o}{:}\PYG{o}{:}\PYG{n}{exists}\PYG{p}{(}\PYG{n}{path}\PYG{p}{,}\PYG{+w}{ }\PYG{n}{err}\PYG{p}{)}\PYG{p}{;}
\PYG{n}{std}\PYG{o}{:}\PYG{o}{:}\PYG{n}{cout}\PYG{+w}{ }\PYG{o}{\PYGZlt{}}\PYG{o}{\PYGZlt{}}\PYG{+w}{ }\PYG{l+s}{\PYGZdq{}}\PYG{l+s}{file exists: }\PYG{l+s}{\PYGZdq{}}\PYG{+w}{ }\PYG{o}{\PYGZlt{}}\PYG{o}{\PYGZlt{}}\PYG{+w}{ }\PYG{n}{std}\PYG{o}{:}\PYG{o}{:}\PYG{n}{boolalpha}
\PYG{+w}{          }\PYG{o}{\PYGZlt{}}\PYG{o}{\PYGZlt{}}\PYG{+w}{ }\PYG{n}{exists}\PYG{+w}{ }\PYG{o}{\PYGZlt{}}\PYG{o}{\PYGZlt{}}\PYG{+w}{ }\PYG{n}{std}\PYG{o}{:}\PYG{o}{:}\PYG{n}{endl}\PYG{p}{;}
\end{Verbatim}
\end{codeblockbg}


%%C070452 li

\item 要检查两个不同路径是否指向同一文件系统对象，请使用\texttt{equivalent()}：


%%C070453 code

\begin{codeblockbg}
\begin{Verbatim}[commandchars=\\\{\},breaklines=true,breakanywhere=true,fontsize=\codefont,xleftmargin=17bp]
\PYG{k}{auto}\PYG{+w}{ }\PYG{n}{same}\PYG{+w}{ }\PYG{o}{=}\PYG{+w}{ }\PYG{n}{fs}\PYG{o}{:}\PYG{o}{:}\PYG{n}{equivalent}\PYG{p}{(}\PYG{n}{path}\PYG{p}{,}
\PYG{+w}{               }\PYG{n}{fs}\PYG{o}{:}\PYG{o}{:}\PYG{n}{current\PYGZus{}path}\PYG{p}{(}\PYG{p}{)}\PYG{+w}{ }\PYG{o}{/}\PYG{+w}{ }\PYG{l+s}{\PYGZdq{}}\PYG{l+s}{.}\PYG{l+s}{\PYGZdq{}}\PYG{+w}{ }\PYG{o}{/}\PYG{+w}{ }\PYG{l+s}{\PYGZdq{}}\PYG{l+s}{main.cpp}\PYG{l+s}{\PYGZdq{}}\PYG{p}{)}\PYG{p}{;}
\PYG{n}{std}\PYG{o}{:}\PYG{o}{:}\PYG{n}{cout}\PYG{+w}{ }\PYG{o}{\PYGZlt{}}\PYG{o}{\PYGZlt{}}\PYG{+w}{ }\PYG{l+s}{\PYGZdq{}}\PYG{l+s}{equivalent: }\PYG{l+s}{\PYGZdq{}}\PYG{+w}{ }\PYG{o}{\PYGZlt{}}\PYG{o}{\PYGZlt{}}\PYG{+w}{ }\PYG{n}{same}\PYG{+w}{ }\PYG{o}{\PYGZlt{}}\PYG{o}{\PYGZlt{}}\PYG{+w}{ }\PYG{n}{std}\PYG{o}{:}\PYG{o}{:}\PYG{n}{endl}\PYG{p}{;}
\end{Verbatim}
\end{codeblockbg}


%%C070454 li

\item 要以字节为单位获取文件大小，请使用\texttt{file\_\allowbreak{}size()}：


%%C070455 code

\begin{codeblockbg}
\begin{Verbatim}[commandchars=\\\{\},breaklines=true,breakanywhere=true,fontsize=\codefont,xleftmargin=17bp]
\PYG{k}{auto}\PYG{+w}{ }\PYG{n}{size}\PYG{+w}{ }\PYG{o}{=}\PYG{+w}{ }\PYG{n}{fs}\PYG{o}{:}\PYG{o}{:}\PYG{n}{file\PYGZus{}size}\PYG{p}{(}\PYG{n}{path}\PYG{p}{,}\PYG{+w}{ }\PYG{n}{err}\PYG{p}{)}\PYG{p}{;}
\PYG{n}{std}\PYG{o}{:}\PYG{o}{:}\PYG{n}{cout}\PYG{+w}{ }\PYG{o}{\PYGZlt{}}\PYG{o}{\PYGZlt{}}\PYG{+w}{ }\PYG{l+s}{\PYGZdq{}}\PYG{l+s}{file size: }\PYG{l+s}{\PYGZdq{}}\PYG{+w}{ }\PYG{o}{\PYGZlt{}}\PYG{o}{\PYGZlt{}}\PYG{+w}{ }\PYG{n}{size}\PYG{+w}{ }\PYG{o}{\PYGZlt{}}\PYG{o}{\PYGZlt{}}\PYG{+w}{ }\PYG{n}{std}\PYG{o}{:}\PYG{o}{:}\PYG{n}{endl}\PYG{p}{;}
\end{Verbatim}
\end{codeblockbg}


%%C070456 li

\item 要获取指向某个文件系统对象的硬链接数量，请使用\myttbr{hard\_\allowbreak{}link\_\allowbreak{}count()}：


%%C070457 code

\begin{codeblockbg}
\begin{Verbatim}[commandchars=\\\{\},breaklines=true,breakanywhere=true,fontsize=\codefont,xleftmargin=17bp]
\PYG{k}{auto}\PYG{+w}{ }\PYG{n}{links}\PYG{+w}{ }\PYG{o}{=}\PYG{+w}{ }\PYG{n}{fs}\PYG{o}{:}\PYG{o}{:}\PYG{n}{hard\PYGZus{}link\PYGZus{}count}\PYG{p}{(}\PYG{n}{path}\PYG{p}{,}\PYG{+w}{ }\PYG{n}{err}\PYG{p}{)}\PYG{p}{;}
\PYG{k}{if}\PYG{p}{(}\PYG{n}{links}\PYG{+w}{ }\PYG{o}{!}\PYG{o}{=}\PYG{+w}{ }\PYG{k}{static\PYGZus{}cast}\PYG{o}{\PYGZlt{}}\PYG{k+kt}{uintmax\PYGZus{}t}\PYG{o}{\PYGZgt{}}\PYG{p}{(}\PYG{l+m+mi}{\PYGZhy{}1}\PYG{p}{)}\PYG{p}{)}
\PYG{+w}{  }\PYG{n}{std}\PYG{o}{:}\PYG{o}{:}\PYG{n}{cout}\PYG{+w}{ }\PYG{o}{\PYGZlt{}}\PYG{o}{\PYGZlt{}}\PYG{+w}{ }\PYG{l+s}{\PYGZdq{}}\PYG{l+s}{hard links: }\PYG{l+s}{\PYGZdq{}}\PYG{+w}{ }\PYG{o}{\PYGZlt{}}\PYG{o}{\PYGZlt{}}\PYG{+w}{ }\PYG{n}{links}\PYG{+w}{ }\PYG{o}{\PYGZlt{}}\PYG{o}{\PYGZlt{}}\PYG{+w}{ }\PYG{n}{std}\PYG{o}{:}\PYG{o}{:}\PYG{n}{endl}\PYG{p}{;}
\PYG{k}{else}
\PYG{+w}{  }\PYG{n}{std}\PYG{o}{:}\PYG{o}{:}\PYG{n}{cout}\PYG{+w}{ }\PYG{o}{\PYGZlt{}}\PYG{o}{\PYGZlt{}}\PYG{+w}{ }\PYG{l+s}{\PYGZdq{}}\PYG{l+s}{hard links: error}\PYG{l+s}{\PYGZdq{}}\PYG{+w}{ }\PYG{o}{\PYGZlt{}}\PYG{o}{\PYGZlt{}}\PYG{+w}{ }\PYG{n}{std}\PYG{o}{:}\PYG{o}{:}\PYG{n}{endl}\PYG{p}{;}
\end{Verbatim}
\end{codeblockbg}


%%C070458 li

\item 要获取或设置文件系统对象的最后修改时间，请使用\myttbr{last\_\allowbreak{}write\_\allowbreak{}time()}：


%%C070459 code

\begin{codeblockbg}
\begin{Verbatim}[commandchars=\\\{\},breaklines=true,breakanywhere=true,fontsize=\codefont,xleftmargin=17bp]
\PYG{k}{auto}\PYG{+w}{ }\PYG{n}{lwt}\PYG{+w}{ }\PYG{o}{=}\PYG{+w}{ }\PYG{n}{fs}\PYG{o}{:}\PYG{o}{:}\PYG{n}{last\PYGZus{}write\PYGZus{}time}\PYG{p}{(}\PYG{n}{path}\PYG{p}{,}\PYG{+w}{ }\PYG{n}{err}\PYG{p}{)}\PYG{p}{;}
\PYG{k}{auto}\PYG{+w}{ }\PYG{n}{time}\PYG{+w}{ }\PYG{o}{=}\PYG{+w}{ }\PYG{k}{decltype}\PYG{p}{(}\PYG{n}{lwt}\PYG{p}{)}\PYG{o}{:}\PYG{o}{:}\PYG{n}{clock}\PYG{o}{:}\PYG{o}{:}\PYG{n}{to\PYGZus{}time\PYGZus{}t}\PYG{p}{(}\PYG{n}{lwt}\PYG{p}{)}\PYG{p}{;}
\PYG{k}{auto}\PYG{+w}{ }\PYG{n}{localtime}\PYG{+w}{ }\PYG{o}{=}\PYG{+w}{ }\PYG{n}{std}\PYG{o}{:}\PYG{o}{:}\PYG{n}{localtime}\PYG{p}{(}\PYG{o}{\PYGZam{}}\PYG{n}{time}\PYG{p}{)}\PYG{p}{;}
\PYG{n}{std}\PYG{o}{:}\PYG{o}{:}\PYG{n}{cout}\PYG{+w}{ }\PYG{o}{\PYGZlt{}}\PYG{o}{\PYGZlt{}}\PYG{+w}{ }\PYG{l+s}{\PYGZdq{}}\PYG{l+s}{last write time: }\PYG{l+s}{\PYGZdq{}}
\PYG{+w}{          }\PYG{o}{\PYGZlt{}}\PYG{o}{\PYGZlt{}}\PYG{+w}{ }\PYG{n}{std}\PYG{o}{:}\PYG{o}{:}\PYG{n}{put\PYGZus{}time}\PYG{p}{(}\PYG{n}{localtime}\PYG{p}{,}\PYG{+w}{ }\PYG{l+s}{\PYGZdq{}}\PYG{l+s}{\PYGZpc{}c}\PYG{l+s}{\PYGZdq{}}\PYG{p}{)}\PYG{+w}{ }\PYG{o}{\PYGZlt{}}\PYG{o}{\PYGZlt{}}\PYG{+w}{ }\PYG{n}{std}\PYG{o}{:}\PYG{o}{:}\PYG{n}{endl}\PYG{p}{;}
\end{Verbatim}
\end{codeblockbg}


%%C070460 li

\item 要获取POSIX文件属性（如类型和权限），请使用\texttt{status()}函数。该函数会跟随符号链接。要获取符号链接自身的文件属性而不跟随它，请使用\texttt{symlink\_\allowbreak{}status()}：


%%C070461 code

\begin{codeblockbg}
\begin{Verbatim}[commandchars=\\\{\},breaklines=true,breakanywhere=true,fontsize=\codefont,xleftmargin=17bp]
\PYG{k}{auto}\PYG{+w}{ }\PYG{n}{print\PYGZus{}perm}\PYG{+w}{ }\PYG{o}{=}\PYG{+w}{ }\PYG{p}{[}\PYG{p}{]}\PYG{p}{(}\PYG{n}{fs}\PYG{o}{:}\PYG{o}{:}\PYG{n}{perms}\PYG{+w}{ }\PYG{n}{p}\PYG{p}{)}
\PYG{p}{\PYGZob{}}
\PYG{+w}{  }\PYG{n}{std}\PYG{o}{:}\PYG{o}{:}\PYG{n}{cout}
\PYG{+w}{    }\PYG{o}{\PYGZlt{}}\PYG{o}{\PYGZlt{}}\PYG{+w}{ }\PYG{p}{(}\PYG{p}{(}\PYG{n}{p}\PYG{+w}{ }\PYG{o}{\PYGZam{}}\PYG{+w}{ }\PYG{n}{fs}\PYG{o}{:}\PYG{o}{:}\PYG{n}{perms}\PYG{o}{:}\PYG{o}{:}\PYG{n}{owner\PYGZus{}read}\PYG{p}{)}\PYG{+w}{ }\PYG{o}{!}\PYG{o}{=}\PYG{+w}{ }\PYG{n}{fs}\PYG{o}{:}\PYG{o}{:}\PYG{n}{perms}\PYG{o}{:}\PYG{o}{:}\PYG{n}{none}\PYG{+w}{ }\PYG{o}{?}
\PYG{+w}{       }\PYG{l+s}{\PYGZdq{}}\PYG{l+s}{r}\PYG{l+s}{\PYGZdq{}}\PYG{+w}{ }\PYG{o}{:}\PYG{+w}{ }\PYG{l+s}{\PYGZdq{}}\PYG{l+s}{\PYGZhy{}}\PYG{l+s}{\PYGZdq{}}\PYG{p}{)}
\PYG{+w}{    }\PYG{o}{\PYGZlt{}}\PYG{o}{\PYGZlt{}}\PYG{+w}{ }\PYG{p}{(}\PYG{p}{(}\PYG{n}{p}\PYG{+w}{ }\PYG{o}{\PYGZam{}}\PYG{+w}{ }\PYG{n}{fs}\PYG{o}{:}\PYG{o}{:}\PYG{n}{perms}\PYG{o}{:}\PYG{o}{:}\PYG{n}{owner\PYGZus{}write}\PYG{p}{)}\PYG{+w}{ }\PYG{o}{!}\PYG{o}{=}\PYG{+w}{ }\PYG{n}{fs}\PYG{o}{:}\PYG{o}{:}\PYG{n}{perms}\PYG{o}{:}\PYG{o}{:}\PYG{n}{none}\PYG{+w}{ }\PYG{o}{?}
\PYG{+w}{       }\PYG{l+s}{\PYGZdq{}}\PYG{l+s}{w}\PYG{l+s}{\PYGZdq{}}\PYG{+w}{ }\PYG{o}{:}\PYG{+w}{ }\PYG{l+s}{\PYGZdq{}}\PYG{l+s}{\PYGZhy{}}\PYG{l+s}{\PYGZdq{}}\PYG{p}{)}
\PYG{+w}{    }\PYG{o}{\PYGZlt{}}\PYG{o}{\PYGZlt{}}\PYG{+w}{ }\PYG{p}{(}\PYG{p}{(}\PYG{n}{p}\PYG{+w}{ }\PYG{o}{\PYGZam{}}\PYG{+w}{ }\PYG{n}{fs}\PYG{o}{:}\PYG{o}{:}\PYG{n}{perms}\PYG{o}{:}\PYG{o}{:}\PYG{n}{owner\PYGZus{}exec}\PYG{p}{)}\PYG{+w}{ }\PYG{o}{!}\PYG{o}{=}\PYG{+w}{ }\PYG{n}{fs}\PYG{o}{:}\PYG{o}{:}\PYG{n}{perms}\PYG{o}{:}\PYG{o}{:}\PYG{n}{none}\PYG{+w}{ }\PYG{o}{?}
\PYG{+w}{       }\PYG{l+s}{\PYGZdq{}}\PYG{l+s}{x}\PYG{l+s}{\PYGZdq{}}\PYG{+w}{ }\PYG{o}{:}\PYG{+w}{ }\PYG{l+s}{\PYGZdq{}}\PYG{l+s}{\PYGZhy{}}\PYG{l+s}{\PYGZdq{}}\PYG{p}{)}
\PYG{+w}{    }\PYG{o}{\PYGZlt{}}\PYG{o}{\PYGZlt{}}\PYG{+w}{ }\PYG{p}{(}\PYG{p}{(}\PYG{n}{p}\PYG{+w}{ }\PYG{o}{\PYGZam{}}\PYG{+w}{ }\PYG{n}{fs}\PYG{o}{:}\PYG{o}{:}\PYG{n}{perms}\PYG{o}{:}\PYG{o}{:}\PYG{n}{group\PYGZus{}read}\PYG{p}{)}\PYG{+w}{ }\PYG{o}{!}\PYG{o}{=}\PYG{+w}{ }\PYG{n}{fs}\PYG{o}{:}\PYG{o}{:}\PYG{n}{perms}\PYG{o}{:}\PYG{o}{:}\PYG{n}{none}\PYG{+w}{ }\PYG{o}{?}
\PYG{+w}{       }\PYG{l+s}{\PYGZdq{}}\PYG{l+s}{r}\PYG{l+s}{\PYGZdq{}}\PYG{+w}{ }\PYG{o}{:}\PYG{+w}{ }\PYG{l+s}{\PYGZdq{}}\PYG{l+s}{\PYGZhy{}}\PYG{l+s}{\PYGZdq{}}\PYG{p}{)}
\PYG{+w}{    }\PYG{o}{\PYGZlt{}}\PYG{o}{\PYGZlt{}}\PYG{+w}{ }\PYG{p}{(}\PYG{p}{(}\PYG{n}{p}\PYG{+w}{ }\PYG{o}{\PYGZam{}}\PYG{+w}{ }\PYG{n}{fs}\PYG{o}{:}\PYG{o}{:}\PYG{n}{perms}\PYG{o}{:}\PYG{o}{:}\PYG{n}{group\PYGZus{}write}\PYG{p}{)}\PYG{+w}{ }\PYG{o}{!}\PYG{o}{=}\PYG{+w}{ }\PYG{n}{fs}\PYG{o}{:}\PYG{o}{:}\PYG{n}{perms}\PYG{o}{:}\PYG{o}{:}\PYG{n}{none}\PYG{+w}{ }\PYG{o}{?}
\PYG{+w}{       }\PYG{l+s}{\PYGZdq{}}\PYG{l+s}{w}\PYG{l+s}{\PYGZdq{}}\PYG{+w}{ }\PYG{o}{:}\PYG{+w}{ }\PYG{l+s}{\PYGZdq{}}\PYG{l+s}{\PYGZhy{}}\PYG{l+s}{\PYGZdq{}}\PYG{p}{)}
\PYG{+w}{    }\PYG{o}{\PYGZlt{}}\PYG{o}{\PYGZlt{}}\PYG{+w}{ }\PYG{p}{(}\PYG{p}{(}\PYG{n}{p}\PYG{+w}{ }\PYG{o}{\PYGZam{}}\PYG{+w}{ }\PYG{n}{fs}\PYG{o}{:}\PYG{o}{:}\PYG{n}{perms}\PYG{o}{:}\PYG{o}{:}\PYG{n}{group\PYGZus{}exec}\PYG{p}{)}\PYG{+w}{ }\PYG{o}{!}\PYG{o}{=}\PYG{+w}{ }\PYG{n}{fs}\PYG{o}{:}\PYG{o}{:}\PYG{n}{perms}\PYG{o}{:}\PYG{o}{:}\PYG{n}{none}\PYG{+w}{ }\PYG{o}{?}
\PYG{+w}{       }\PYG{l+s}{\PYGZdq{}}\PYG{l+s}{x}\PYG{l+s}{\PYGZdq{}}\PYG{+w}{ }\PYG{o}{:}\PYG{+w}{ }\PYG{l+s}{\PYGZdq{}}\PYG{l+s}{\PYGZhy{}}\PYG{l+s}{\PYGZdq{}}\PYG{p}{)}
\PYG{+w}{    }\PYG{o}{\PYGZlt{}}\PYG{o}{\PYGZlt{}}\PYG{+w}{ }\PYG{p}{(}\PYG{p}{(}\PYG{n}{p}\PYG{+w}{ }\PYG{o}{\PYGZam{}}\PYG{+w}{ }\PYG{n}{fs}\PYG{o}{:}\PYG{o}{:}\PYG{n}{perms}\PYG{o}{:}\PYG{o}{:}\PYG{n}{others\PYGZus{}read}\PYG{p}{)}\PYG{+w}{ }\PYG{o}{!}\PYG{o}{=}\PYG{+w}{ }\PYG{n}{fs}\PYG{o}{:}\PYG{o}{:}\PYG{n}{perms}\PYG{o}{:}\PYG{o}{:}\PYG{n}{none}\PYG{+w}{ }\PYG{o}{?}
\PYG{+w}{       }\PYG{l+s}{\PYGZdq{}}\PYG{l+s}{r}\PYG{l+s}{\PYGZdq{}}\PYG{+w}{ }\PYG{o}{:}\PYG{+w}{ }\PYG{l+s}{\PYGZdq{}}\PYG{l+s}{\PYGZhy{}}\PYG{l+s}{\PYGZdq{}}\PYG{p}{)}
\PYG{+w}{    }\PYG{o}{\PYGZlt{}}\PYG{o}{\PYGZlt{}}\PYG{+w}{ }\PYG{p}{(}\PYG{p}{(}\PYG{n}{p}\PYG{+w}{ }\PYG{o}{\PYGZam{}}\PYG{+w}{ }\PYG{n}{fs}\PYG{o}{:}\PYG{o}{:}\PYG{n}{perms}\PYG{o}{:}\PYG{o}{:}\PYG{n}{others\PYGZus{}write}\PYG{p}{)}\PYG{+w}{ }\PYG{o}{!}\PYG{o}{=}\PYG{+w}{ }\PYG{n}{fs}\PYG{o}{:}\PYG{o}{:}\PYG{n}{perms}\PYG{o}{:}\PYG{o}{:}\PYG{n}{none}\PYG{+w}{ }\PYG{o}{?}
\PYG{+w}{       }\PYG{l+s}{\PYGZdq{}}\PYG{l+s}{w}\PYG{l+s}{\PYGZdq{}}\PYG{+w}{ }\PYG{o}{:}\PYG{+w}{ }\PYG{l+s}{\PYGZdq{}}\PYG{l+s}{\PYGZhy{}}\PYG{l+s}{\PYGZdq{}}\PYG{p}{)}
\PYG{+w}{    }\PYG{o}{\PYGZlt{}}\PYG{o}{\PYGZlt{}}\PYG{+w}{ }\PYG{p}{(}\PYG{p}{(}\PYG{n}{p}\PYG{+w}{ }\PYG{o}{\PYGZam{}}\PYG{+w}{ }\PYG{n}{fs}\PYG{o}{:}\PYG{o}{:}\PYG{n}{perms}\PYG{o}{:}\PYG{o}{:}\PYG{n}{others\PYGZus{}exec}\PYG{p}{)}\PYG{+w}{ }\PYG{o}{!}\PYG{o}{=}\PYG{+w}{ }\PYG{n}{fs}\PYG{o}{:}\PYG{o}{:}\PYG{n}{perms}\PYG{o}{:}\PYG{o}{:}\PYG{n}{none}\PYG{+w}{ }\PYG{o}{?}
\PYG{+w}{       }\PYG{l+s}{\PYGZdq{}}\PYG{l+s}{x}\PYG{l+s}{\PYGZdq{}}\PYG{+w}{ }\PYG{o}{:}\PYG{+w}{ }\PYG{l+s}{\PYGZdq{}}\PYG{l+s}{\PYGZhy{}}\PYG{l+s}{\PYGZdq{}}\PYG{p}{)}
\PYG{+w}{    }\PYG{o}{\PYGZlt{}}\PYG{o}{\PYGZlt{}}\PYG{+w}{ }\PYG{n}{std}\PYG{o}{:}\PYG{o}{:}\PYG{n}{endl}\PYG{p}{;}
\PYG{p}{\PYGZcb{}}\PYG{p}{;}
\PYG{k}{auto}\PYG{+w}{ }\PYG{n}{status}\PYG{+w}{ }\PYG{o}{=}\PYG{+w}{ }\PYG{n}{fs}\PYG{o}{:}\PYG{o}{:}\PYG{n}{status}\PYG{p}{(}\PYG{n}{path}\PYG{p}{,}\PYG{+w}{ }\PYG{n}{err}\PYG{p}{)}\PYG{p}{;}
\PYG{n}{std}\PYG{o}{:}\PYG{o}{:}\PYG{n}{cout}\PYG{+w}{ }\PYG{o}{\PYGZlt{}}\PYG{o}{\PYGZlt{}}\PYG{+w}{ }\PYG{l+s}{\PYGZdq{}}\PYG{l+s}{permissions: }\PYG{l+s}{\PYGZdq{}}\PYG{p}{;}
\PYG{n}{print\PYGZus{}perm}\PYG{p}{(}\PYG{n}{status}\PYG{p}{.}\PYG{n}{permissions}\PYG{p}{(}\PYG{p}{)}\PYG{p}{)}\PYG{p}{;}
\end{Verbatim}
\end{codeblockbg}


%%C070462 li

\item 要检查某个路径是否指向特定类型的文件系统对象（如文件、目录、符号链接等），请使用\myttbr{is\_\allowbreak{}regular\_\allowbreak{}file()}、\texttt{is\_\allowbreak{}directory()}、\texttt{is\_\allowbreak{}symlink()}等函数：


%%C070463 code

\begin{codeblockbg}
\begin{Verbatim}[commandchars=\\\{\},breaklines=true,breakanywhere=true,fontsize=\codefont,xleftmargin=17bp]
\PYG{n}{std}\PYG{o}{:}\PYG{o}{:}\PYG{n}{cout}\PYG{+w}{ }\PYG{o}{\PYGZlt{}}\PYG{o}{\PYGZlt{}}\PYG{+w}{ }\PYG{l+s}{\PYGZdq{}}\PYG{l+s}{regular file? }\PYG{l+s}{\PYGZdq{}}\PYG{+w}{ }\PYG{o}{\PYGZlt{}}\PYG{o}{\PYGZlt{}}
\PYG{+w}{          }\PYG{n}{fs}\PYG{o}{:}\PYG{o}{:}\PYG{n}{is\PYGZus{}regular\PYGZus{}file}\PYG{p}{(}\PYG{n}{path}\PYG{p}{,}\PYG{+w}{ }\PYG{n}{err}\PYG{p}{)}\PYG{+w}{ }\PYG{o}{\PYGZlt{}}\PYG{o}{\PYGZlt{}}\PYG{+w}{ }\PYG{n}{std}\PYG{o}{:}\PYG{o}{:}\PYG{n}{endl}\PYG{p}{;}
\PYG{n}{std}\PYG{o}{:}\PYG{o}{:}\PYG{n}{cout}\PYG{+w}{ }\PYG{o}{\PYGZlt{}}\PYG{o}{\PYGZlt{}}\PYG{+w}{ }\PYG{l+s}{\PYGZdq{}}\PYG{l+s}{directory? }\PYG{l+s}{\PYGZdq{}}\PYG{+w}{ }\PYG{o}{\PYGZlt{}}\PYG{o}{\PYGZlt{}}
\PYG{+w}{          }\PYG{n}{fs}\PYG{o}{:}\PYG{o}{:}\PYG{n}{is\PYGZus{}directory}\PYG{p}{(}\PYG{n}{path}\PYG{p}{,}\PYG{+w}{ }\PYG{n}{err}\PYG{p}{)}\PYG{+w}{ }\PYG{o}{\PYGZlt{}}\PYG{o}{\PYGZlt{}}\PYG{+w}{ }\PYG{n}{std}\PYG{o}{:}\PYG{o}{:}\PYG{n}{endl}\PYG{p}{;}
\PYG{n}{std}\PYG{o}{:}\PYG{o}{:}\PYG{n}{cout}\PYG{+w}{ }\PYG{o}{\PYGZlt{}}\PYG{o}{\PYGZlt{}}\PYG{+w}{ }\PYG{l+s}{\PYGZdq{}}\PYG{l+s}{char file? }\PYG{l+s}{\PYGZdq{}}\PYG{+w}{ }\PYG{o}{\PYGZlt{}}\PYG{o}{\PYGZlt{}}
\PYG{+w}{          }\PYG{n}{fs}\PYG{o}{:}\PYG{o}{:}\PYG{n}{is\PYGZus{}character\PYGZus{}file}\PYG{p}{(}\PYG{n}{path}\PYG{p}{,}\PYG{+w}{ }\PYG{n}{err}\PYG{p}{)}\PYG{+w}{ }\PYG{o}{\PYGZlt{}}\PYG{o}{\PYGZlt{}}\PYG{+w}{ }\PYG{n}{std}\PYG{o}{:}\PYG{o}{:}\PYG{n}{endl}\PYG{p}{;}
\PYG{n}{std}\PYG{o}{:}\PYG{o}{:}\PYG{n}{cout}\PYG{+w}{ }\PYG{o}{\PYGZlt{}}\PYG{o}{\PYGZlt{}}\PYG{+w}{ }\PYG{l+s}{\PYGZdq{}}\PYG{l+s}{symlink? }\PYG{l+s}{\PYGZdq{}}\PYG{+w}{ }\PYG{o}{\PYGZlt{}}\PYG{o}{\PYGZlt{}}
\PYG{+w}{          }\PYG{n}{fs}\PYG{o}{:}\PYG{o}{:}\PYG{n}{is\PYGZus{}symlink}\PYG{p}{(}\PYG{n}{path}\PYG{p}{,}\PYG{+w}{ }\PYG{n}{err}\PYG{p}{)}\PYG{+w}{ }\PYG{o}{\PYGZlt{}}\PYG{o}{\PYGZlt{}}\PYG{+w}{ }\PYG{n}{std}\PYG{o}{:}\PYG{o}{:}\PYG{n}{endl}\PYG{p}{;}
\end{Verbatim}
\end{codeblockbg}

\end{itemize}


%%C070464 h3

\subsection{工作原理}


%%C070465 p

本案例讨论的所有函数都有一个在出错时抛出异常的重载，以及一个不抛出异常、而是通过函数参数返回错误码的重载。本案例中的所有示例都采用了后一种做法。关于这些重载集的更多信息，可参见\textit{创建、复制与删除文件和目录}案例。

%%C070466 p

用于获取文件系统文件与目录信息的这些函数总体上简单直接。不过，有几点需要考虑：


%%C070467 li

\begin{itemize}
\item 检查某个文件系统对象是否存在，可以用\texttt{exists()}完成：既可以传入路径，也可以传入此前用\texttt{status()}函数取得的\myttbr{std::\allowbreak{}filesystem::\allowbreak{}file\_\allowbreak{}status}对象。


%%C070468 li

\item \texttt{equivalent()}函数判断两个文件系统对象是否具有相同的状态，这里的状态由\texttt{status()}函数取得。如果两个路径都不存在，或者两者都存在但不是文件、目录或符号链接，该函数会返回错误。指向同一文件对象的硬链接是等价的。符号链接与其目标也是等价的。


%%C070469 li

\item \texttt{file\_\allowbreak{}size()}函数只能用于获取普通文件以及指向普通文件的符号链接的大小。对于其他任何类型的文件系统对象（例如目录），该函数都会失败。该函数返回文件的大小（以字节计），出错时返回-1。如果你想判断文件是否为空，可以使用\texttt{is\_\allowbreak{}empty()}函数。它对所有类型的文件系统对象都适用，包括目录。


%%C070470 li

\item \myttbr{last\_\allowbreak{}write\_\allowbreak{}time()}函数有两组重载：一组用于获取文件系统对象的最后修改时间，另一组用于设置最后修改时间。时间用\myttbr{std::\allowbreak{}filesystem::\allowbreak{}file\_\allowbreak{}time\_\allowbreak{}type}对象表示，它本质上是\myttbr{std::\allowbreak{}chrono::\allowbreak{}time\_\allowbreak{}point}的类型别名。下面的例子把文件的最后写入时间改为比原来早30分钟：


%%C070471 code

\begin{codeblockbg}
\begin{Verbatim}[commandchars=\\\{\},breaklines=true,breakanywhere=true,fontsize=\codefont,xleftmargin=17bp]
\PYG{k}{using}\PYG{+w}{ }\PYG{k}{namespace}\PYG{+w}{ }\PYG{n+nn}{std}\PYG{o}{:}\PYG{o}{:}\PYG{n+nn}{chrono\PYGZus{}literals}\PYG{p}{;}
\PYG{k}{auto}\PYG{+w}{ }\PYG{n}{lwt}\PYG{+w}{ }\PYG{o}{=}\PYG{+w}{ }\PYG{n}{fs}\PYG{o}{:}\PYG{o}{:}\PYG{n}{last\PYGZus{}write\PYGZus{}time}\PYG{p}{(}\PYG{n}{path}\PYG{p}{,}\PYG{+w}{ }\PYG{n}{err}\PYG{p}{)}\PYG{p}{;}
\PYG{n}{fs}\PYG{o}{:}\PYG{o}{:}\PYG{n}{last\PYGZus{}write\PYGZus{}time}\PYG{p}{(}\PYG{n}{path}\PYG{p}{,}\PYG{+w}{ }\PYG{n}{lwt}\PYG{+w}{ }\PYG{o}{\PYGZhy{}}\PYG{+w}{ }\PYG{l+m+mi}{30}\PYG{n}{min}\PYG{p}{)}\PYG{p}{;}
\end{Verbatim}
\end{codeblockbg}


%%C070472 li

\item \texttt{status()}函数确定文件系统对象的类型与权限。不过，如果该文件是符号链接，返回的信息描述的是符号链接的目标。要获取符号链接自身的信息，必须使用\texttt{symlink\_\allowbreak{}status()}函数。权限定义为一个枚举，即\myttbr{std::\allowbreak{}filesystem::\allowbreak{}perms}。这个限定作用域\texttt{enum}的枚举量并非都表示权限，有些表示控制位，例如\texttt{add\_\allowbreak{}perms}表示应添加权限，\texttt{remove\_\allowbreak{}perms}表示应移除权限。\texttt{permissions()}函数可用于修改文件或目录的权限。下面的例子为文件的所有者和用户组添加全部权限：


%%C070473 code

\begin{codeblockbg}
\begin{Verbatim}[commandchars=\\\{\},breaklines=true,breakanywhere=true,fontsize=\codefont,xleftmargin=17bp]
\PYG{n}{fs}\PYG{o}{:}\PYG{o}{:}\PYG{n}{permissions}\PYG{p}{(}
\PYG{+w}{  }\PYG{n}{path}\PYG{p}{,}
\PYG{+w}{  }\PYG{n}{fs}\PYG{o}{:}\PYG{o}{:}\PYG{n}{perms}\PYG{o}{:}\PYG{o}{:}\PYG{n}{add\PYGZus{}perms}\PYG{+w}{ }\PYG{o}{|}
\PYG{+w}{  }\PYG{n}{fs}\PYG{o}{:}\PYG{o}{:}\PYG{n}{perms}\PYG{o}{:}\PYG{o}{:}\PYG{n}{owner\PYGZus{}all}\PYG{+w}{ }\PYG{o}{|}\PYG{+w}{ }\PYG{n}{fs}\PYG{o}{:}\PYG{o}{:}\PYG{n}{perms}\PYG{o}{:}\PYG{o}{:}\PYG{n}{group\PYGZus{}all}\PYG{p}{,}
\PYG{+w}{  }\PYG{n}{err}\PYG{p}{)}\PYG{p}{;}
\end{Verbatim}
\end{codeblockbg}


%%C070474 li

\item 要确定文件系统对象的类型（例如文件、目录或符号链接），有两种可选做法：先取得文件状态，再检查类型属性；或者使用现成的文件系统函数，例如\myttbr{is\_\allowbreak{}regular\_\allowbreak{}file()}、\texttt{is\_\allowbreak{}symlink()}或\texttt{is\_\allowbreak{}directory()}。下面两个例子都用于检查某个路径是否指向普通文件，它们是等价的：


%%C070475 code

\begin{codeblockbg}
\begin{Verbatim}[commandchars=\\\{\},breaklines=true,breakanywhere=true,fontsize=\codefont,xleftmargin=17bp]
\PYG{k}{auto}\PYG{+w}{ }\PYG{n}{s}\PYG{+w}{ }\PYG{o}{=}\PYG{+w}{ }\PYG{n}{fs}\PYG{o}{:}\PYG{o}{:}\PYG{n}{status}\PYG{p}{(}\PYG{n}{path}\PYG{p}{,}\PYG{+w}{ }\PYG{n}{err}\PYG{p}{)}\PYG{p}{;}
\PYG{k}{auto}\PYG{+w}{ }\PYG{n}{isfile}\PYG{+w}{ }\PYG{o}{=}\PYG{+w}{ }\PYG{n}{s}\PYG{p}{.}\PYG{n}{type}\PYG{p}{(}\PYG{p}{)}\PYG{+w}{ }\PYG{o}{=}\PYG{o}{=}\PYG{+w}{ }\PYG{n}{std}\PYG{o}{:}\PYG{o}{:}\PYG{n}{filesystem}\PYG{o}{:}\PYG{o}{:}\PYG{n}{file\PYGZus{}type}\PYG{o}{:}\PYG{o}{:}\PYG{n}{regular}\PYG{p}{;}
\PYG{k}{auto}\PYG{+w}{ }\PYG{n}{isfile}\PYG{+w}{ }\PYG{o}{=}\PYG{+w}{ }\PYG{n}{fs}\PYG{o}{:}\PYG{o}{:}\PYG{n}{is\PYGZus{}regular\PYGZus{}file}\PYG{p}{(}\PYG{n}{path}\PYG{p}{,}\PYG{+w}{ }\PYG{n}{err}\PYG{p}{)}\PYG{p}{;}
\end{Verbatim}
\end{codeblockbg}

\end{itemize}


%%C070476 h3

\subsection{延伸阅读}


%%C070477 li

\begin{itemize}
\item \textit{处理文件系统路径}


%%C070478 li

\item \textit{创建、复制与删除文件和目录}


%%C070479 li

\item \textit{枚举目录内容}

\end{itemize}


%%C070480 h2

\section{枚举目录内容}


%%C070481 p

本章到目前为止，我们考察了\texttt{filesystem}库提供的许多功能，例如处理路径、对文件和目录执行操作（创建、移动、重命名、删除等），以及查询或修改属性。处理文件系统时另一个有用的功能是遍历目录内容。文件库提供了两种目录迭代器：一种叫\myttbr{directory\_\allowbreak{}iterator}，用于迭代目录的内容；另一种叫\myttbr{recursive\_\allowbreak{}directory\_\allowbreak{}iterator}，用于递归地迭代目录及其子目录的内容。本案例将介绍如何使用它们。


%%C070482 h3

\subsection{准备工作}


%%C070483 p

本案例将处理文件系统路径，并检查文件系统对象的属性。因此，建议你先阅读\textit{处理文件系统路径}和\textit{检查现有文件或目录的属性}两个案例。


%%C070484 p

本案例考虑具有以下结构的目录：


%%C070485 p

\begin{codeblockbg}
\begin{Verbatim}[commandchars=\\\{\},breaklines=true,breakanywhere=true,fontsize=\codefont,xleftmargin=17bp,formatcom=\bfseries]
test/
├──data/
│ ├──input.dat
│ └──output.dat
├──file_1.txt
├──file_2.txt
└──file_3.log
\end{Verbatim}
\end{codeblockbg}


%%C070486 h3

\subsection{具体做法}


%%C070487 p

使用以下模式枚举目录内容：


%%C070488 li

\begin{itemize}
\item 若只迭代目录的内容而不递归访问其子目录，请使用\myttbr{directory\_\allowbreak{}iterator}：


%%C070489 code

\begin{codeblockbg}
\begin{Verbatim}[commandchars=\\\{\},breaklines=true,breakanywhere=true,fontsize=\codefont,xleftmargin=17bp]
\PYG{k+kt}{void}\PYG{+w}{ }\PYG{n+nf}{visit\PYGZus{}directory}\PYG{p}{(}\PYG{n}{fs}\PYG{o}{:}\PYG{o}{:}\PYG{n}{path}\PYG{+w}{ }\PYG{k}{const}\PYG{+w}{ }\PYG{o}{\PYGZam{}}\PYG{+w}{ }\PYG{n}{dir}\PYG{p}{)}
\PYG{p}{\PYGZob{}}
\PYG{+w}{  }\PYG{k}{if}\PYG{+w}{ }\PYG{p}{(}\PYG{n}{fs}\PYG{o}{:}\PYG{o}{:}\PYG{n}{exists}\PYG{p}{(}\PYG{n}{dir}\PYG{p}{)}\PYG{+w}{ }\PYG{o}{\PYGZam{}}\PYG{o}{\PYGZam{}}\PYG{+w}{ }\PYG{n}{fs}\PYG{o}{:}\PYG{o}{:}\PYG{n}{is\PYGZus{}directory}\PYG{p}{(}\PYG{n}{dir}\PYG{p}{)}\PYG{p}{)}
\PYG{+w}{  }\PYG{p}{\PYGZob{}}
\PYG{+w}{    }\PYG{k}{for}\PYG{+w}{ }\PYG{p}{(}\PYG{k}{auto}\PYG{+w}{ }\PYG{k}{const}\PYG{+w}{ }\PYG{o}{\PYGZam{}}\PYG{+w}{ }\PYG{n}{entry}\PYG{+w}{ }\PYG{o}{:}\PYG{+w}{ }\PYG{n}{fs}\PYG{o}{:}\PYG{o}{:}\PYG{n}{directory\PYGZus{}iterator}\PYG{p}{(}\PYG{n}{dir}\PYG{p}{)}\PYG{p}{)}
\PYG{+w}{    }\PYG{p}{\PYGZob{}}
\PYG{+w}{      }\PYG{k}{auto}\PYG{+w}{ }\PYG{n}{filename}\PYG{+w}{ }\PYG{o}{=}\PYG{+w}{ }\PYG{n}{entry}\PYG{p}{.}\PYG{n}{path}\PYG{p}{(}\PYG{p}{)}\PYG{p}{.}\PYG{n}{filename}\PYG{p}{(}\PYG{p}{)}\PYG{p}{;}
\PYG{+w}{      }\PYG{k}{if}\PYG{+w}{ }\PYG{p}{(}\PYG{n}{fs}\PYG{o}{:}\PYG{o}{:}\PYG{n}{is\PYGZus{}directory}\PYG{p}{(}\PYG{n}{entry}\PYG{p}{.}\PYG{n}{status}\PYG{p}{(}\PYG{p}{)}\PYG{p}{)}\PYG{p}{)}
\PYG{+w}{        }\PYG{n}{std}\PYG{o}{:}\PYG{o}{:}\PYG{n}{cout}\PYG{+w}{ }\PYG{o}{\PYGZlt{}}\PYG{o}{\PYGZlt{}}\PYG{+w}{ }\PYG{l+s}{\PYGZdq{}}\PYG{l+s}{[+]}\PYG{l+s}{\PYGZdq{}}\PYG{+w}{ }\PYG{o}{\PYGZlt{}}\PYG{o}{\PYGZlt{}}\PYG{+w}{ }\PYG{n}{filename}\PYG{+w}{ }\PYG{o}{\PYGZlt{}}\PYG{o}{\PYGZlt{}}\PYG{+w}{ }\PYG{n}{std}\PYG{o}{:}\PYG{o}{:}\PYG{n}{endl}\PYG{p}{;}
\PYG{+w}{      }\PYG{k}{else}\PYG{+w}{ }\PYG{k}{if}\PYG{+w}{ }\PYG{p}{(}\PYG{n}{fs}\PYG{o}{:}\PYG{o}{:}\PYG{n}{is\PYGZus{}symlink}\PYG{p}{(}\PYG{n}{entry}\PYG{p}{.}\PYG{n}{status}\PYG{p}{(}\PYG{p}{)}\PYG{p}{)}\PYG{p}{)}
\PYG{+w}{        }\PYG{n}{std}\PYG{o}{:}\PYG{o}{:}\PYG{n}{cout}\PYG{+w}{ }\PYG{o}{\PYGZlt{}}\PYG{o}{\PYGZlt{}}\PYG{+w}{ }\PYG{l+s}{\PYGZdq{}}\PYG{l+s}{[\PYGZgt{}]}\PYG{l+s}{\PYGZdq{}}\PYG{+w}{ }\PYG{o}{\PYGZlt{}}\PYG{o}{\PYGZlt{}}\PYG{+w}{ }\PYG{n}{filename}\PYG{+w}{ }\PYG{o}{\PYGZlt{}}\PYG{o}{\PYGZlt{}}\PYG{+w}{ }\PYG{n}{std}\PYG{o}{:}\PYG{o}{:}\PYG{n}{endl}\PYG{p}{;}
\PYG{+w}{      }\PYG{k}{else}\PYG{+w}{ }\PYG{k}{if}\PYG{+w}{ }\PYG{p}{(}\PYG{n}{fs}\PYG{o}{:}\PYG{o}{:}\PYG{n}{is\PYGZus{}regular\PYGZus{}file}\PYG{p}{(}\PYG{n}{entry}\PYG{p}{.}\PYG{n}{status}\PYG{p}{(}\PYG{p}{)}\PYG{p}{)}\PYG{p}{)}
\PYG{+w}{        }\PYG{n}{std}\PYG{o}{:}\PYG{o}{:}\PYG{n}{cout}\PYG{+w}{ }\PYG{o}{\PYGZlt{}}\PYG{o}{\PYGZlt{}}\PYG{+w}{ }\PYG{l+s}{\PYGZdq{}}\PYG{l+s}{ }\PYG{l+s}{\PYGZdq{}}\PYG{+w}{ }\PYG{o}{\PYGZlt{}}\PYG{o}{\PYGZlt{}}\PYG{+w}{ }\PYG{n}{filename}\PYG{+w}{ }\PYG{o}{\PYGZlt{}}\PYG{o}{\PYGZlt{}}\PYG{+w}{ }\PYG{n}{std}\PYG{o}{:}\PYG{o}{:}\PYG{n}{endl}\PYG{p}{;}
\PYG{+w}{      }\PYG{k}{else}
\PYG{+w}{        }\PYG{n}{std}\PYG{o}{:}\PYG{o}{:}\PYG{n}{cout}\PYG{+w}{ }\PYG{o}{\PYGZlt{}}\PYG{o}{\PYGZlt{}}\PYG{+w}{ }\PYG{l+s}{\PYGZdq{}}\PYG{l+s}{[?]}\PYG{l+s}{\PYGZdq{}}\PYG{+w}{ }\PYG{o}{\PYGZlt{}}\PYG{o}{\PYGZlt{}}\PYG{+w}{ }\PYG{n}{filename}\PYG{+w}{ }\PYG{o}{\PYGZlt{}}\PYG{o}{\PYGZlt{}}\PYG{+w}{ }\PYG{n}{std}\PYG{o}{:}\PYG{o}{:}\PYG{n}{endl}\PYG{p}{;}
\PYG{+w}{     }\PYG{p}{\PYGZcb{}}
\PYG{+w}{   }\PYG{p}{\PYGZcb{}}
\PYG{+w}{ }\PYG{p}{\PYGZcb{}}
\end{Verbatim}
\end{codeblockbg}


%%C070490 li

\item 若要迭代目录的全部内容（包括其子目录），且在条目的处理顺序无关紧要时，请使用\myttbr{recursive\_\allowbreak{}directory\_\allowbreak{}iterator}：


%%C070491 code

\begin{codeblockbg}
\begin{Verbatim}[commandchars=\\\{\},breaklines=true,breakanywhere=true,fontsize=\codefont,xleftmargin=17bp]
\PYG{k+kt}{void}\PYG{+w}{ }\PYG{n+nf}{visit\PYGZus{}directory\PYGZus{}rec}\PYG{p}{(}\PYG{n}{fs}\PYG{o}{:}\PYG{o}{:}\PYG{n}{path}\PYG{+w}{ }\PYG{k}{const}\PYG{+w}{ }\PYG{o}{\PYGZam{}}\PYG{+w}{ }\PYG{n}{dir}\PYG{p}{)}
\PYG{p}{\PYGZob{}}
\PYG{+w}{  }\PYG{k}{if}\PYG{+w}{ }\PYG{p}{(}\PYG{n}{fs}\PYG{o}{:}\PYG{o}{:}\PYG{n}{exists}\PYG{p}{(}\PYG{n}{dir}\PYG{p}{)}\PYG{+w}{ }\PYG{o}{\PYGZam{}}\PYG{o}{\PYGZam{}}\PYG{+w}{ }\PYG{n}{fs}\PYG{o}{:}\PYG{o}{:}\PYG{n}{is\PYGZus{}directory}\PYG{p}{(}\PYG{n}{dir}\PYG{p}{)}\PYG{p}{)}
\PYG{+w}{  }\PYG{p}{\PYGZob{}}
\PYG{+w}{    }\PYG{k}{for}\PYG{+w}{ }\PYG{p}{(}\PYG{k}{auto}\PYG{+w}{ }\PYG{k}{const}\PYG{+w}{ }\PYG{o}{\PYGZam{}}\PYG{+w}{ }\PYG{n}{entry}\PYG{+w}{ }\PYG{o}{:}
\PYG{+w}{         }\PYG{n}{fs}\PYG{o}{:}\PYG{o}{:}\PYG{n}{recursive\PYGZus{}directory\PYGZus{}iterator}\PYG{p}{(}\PYG{n}{dir}\PYG{p}{)}\PYG{p}{)}
\PYG{+w}{    }\PYG{p}{\PYGZob{}}
\PYG{+w}{      }\PYG{k}{auto}\PYG{+w}{ }\PYG{n}{filename}\PYG{+w}{ }\PYG{o}{=}\PYG{+w}{ }\PYG{n}{entry}\PYG{p}{.}\PYG{n}{path}\PYG{p}{(}\PYG{p}{)}\PYG{p}{.}\PYG{n}{filename}\PYG{p}{(}\PYG{p}{)}\PYG{p}{;}
\PYG{+w}{      }\PYG{k}{if}\PYG{+w}{ }\PYG{p}{(}\PYG{n}{fs}\PYG{o}{:}\PYG{o}{:}\PYG{n}{is\PYGZus{}directory}\PYG{p}{(}\PYG{n}{entry}\PYG{p}{.}\PYG{n}{status}\PYG{p}{(}\PYG{p}{)}\PYG{p}{)}\PYG{p}{)}
\PYG{+w}{        }\PYG{n}{std}\PYG{o}{:}\PYG{o}{:}\PYG{n}{cout}\PYG{+w}{ }\PYG{o}{\PYGZlt{}}\PYG{o}{\PYGZlt{}}\PYG{+w}{ }\PYG{l+s}{\PYGZdq{}}\PYG{l+s}{[+]}\PYG{l+s}{\PYGZdq{}}\PYG{+w}{ }\PYG{o}{\PYGZlt{}}\PYG{o}{\PYGZlt{}}\PYG{+w}{ }\PYG{n}{filename}\PYG{+w}{ }\PYG{o}{\PYGZlt{}}\PYG{o}{\PYGZlt{}}\PYG{+w}{ }\PYG{n}{std}\PYG{o}{:}\PYG{o}{:}\PYG{n}{endl}\PYG{p}{;}
\PYG{+w}{      }\PYG{k}{else}\PYG{+w}{ }\PYG{k}{if}\PYG{+w}{ }\PYG{p}{(}\PYG{n}{fs}\PYG{o}{:}\PYG{o}{:}\PYG{n}{is\PYGZus{}symlink}\PYG{p}{(}\PYG{n}{entry}\PYG{p}{.}\PYG{n}{status}\PYG{p}{(}\PYG{p}{)}\PYG{p}{)}\PYG{p}{)}
\PYG{+w}{        }\PYG{n}{std}\PYG{o}{:}\PYG{o}{:}\PYG{n}{cout}\PYG{+w}{ }\PYG{o}{\PYGZlt{}}\PYG{o}{\PYGZlt{}}\PYG{+w}{ }\PYG{l+s}{\PYGZdq{}}\PYG{l+s}{[\PYGZgt{}]}\PYG{l+s}{\PYGZdq{}}\PYG{+w}{ }\PYG{o}{\PYGZlt{}}\PYG{o}{\PYGZlt{}}\PYG{+w}{ }\PYG{n}{filename}\PYG{+w}{ }\PYG{o}{\PYGZlt{}}\PYG{o}{\PYGZlt{}}\PYG{+w}{ }\PYG{n}{std}\PYG{o}{:}\PYG{o}{:}\PYG{n}{endl}\PYG{p}{;}
\PYG{+w}{      }\PYG{k}{else}\PYG{+w}{ }\PYG{k}{if}\PYG{+w}{ }\PYG{p}{(}\PYG{n}{fs}\PYG{o}{:}\PYG{o}{:}\PYG{n}{is\PYGZus{}regular\PYGZus{}file}\PYG{p}{(}\PYG{n}{entry}\PYG{p}{.}\PYG{n}{status}\PYG{p}{(}\PYG{p}{)}\PYG{p}{)}\PYG{p}{)}
\PYG{+w}{        }\PYG{n}{std}\PYG{o}{:}\PYG{o}{:}\PYG{n}{cout}\PYG{+w}{ }\PYG{o}{\PYGZlt{}}\PYG{o}{\PYGZlt{}}\PYG{+w}{ }\PYG{l+s}{\PYGZdq{}}\PYG{l+s}{ }\PYG{l+s}{\PYGZdq{}}\PYG{+w}{ }\PYG{o}{\PYGZlt{}}\PYG{o}{\PYGZlt{}}\PYG{+w}{ }\PYG{n}{filename}\PYG{+w}{ }\PYG{o}{\PYGZlt{}}\PYG{o}{\PYGZlt{}}\PYG{+w}{ }\PYG{n}{std}\PYG{o}{:}\PYG{o}{:}\PYG{n}{endl}\PYG{p}{;}
\PYG{+w}{      }\PYG{k}{else}
\PYG{+w}{        }\PYG{n}{std}\PYG{o}{:}\PYG{o}{:}\PYG{n}{cout}\PYG{+w}{ }\PYG{o}{\PYGZlt{}}\PYG{o}{\PYGZlt{}}\PYG{+w}{ }\PYG{l+s}{\PYGZdq{}}\PYG{l+s}{[?]}\PYG{l+s}{\PYGZdq{}}\PYG{+w}{ }\PYG{o}{\PYGZlt{}}\PYG{o}{\PYGZlt{}}\PYG{+w}{ }\PYG{n}{filename}\PYG{+w}{ }\PYG{o}{\PYGZlt{}}\PYG{o}{\PYGZlt{}}\PYG{+w}{ }\PYG{n}{std}\PYG{o}{:}\PYG{o}{:}\PYG{n}{endl}\PYG{p}{;}
\PYG{+w}{    }\PYG{p}{\PYGZcb{}}
\PYG{+w}{  }\PYG{p}{\PYGZcb{}}
\PYG{p}{\PYGZcb{}}
\end{Verbatim}
\end{codeblockbg}


%%C070492 li

\item 若要以结构化方式迭代目录的全部内容（包括其子目录），例如按树形遍历，请使用与第一个例子类似的函数：它用\myttbr{directory\_\allowbreak{}iterator}迭代目录内容，但对每个子目录递归调用自身：


%%C070493 code

\begin{codeblockbg}
\begin{Verbatim}[commandchars=\\\{\},breaklines=true,breakanywhere=true,fontsize=\codefont,xleftmargin=17bp]
\PYG{k+kt}{void}\PYG{+w}{ }\PYG{n+nf}{visit\PYGZus{}directory}\PYG{p}{(}
\PYG{+w}{  }\PYG{n}{fs}\PYG{o}{:}\PYG{o}{:}\PYG{n}{path}\PYG{+w}{ }\PYG{k}{const}\PYG{+w}{ }\PYG{o}{\PYGZam{}}\PYG{+w}{ }\PYG{n}{dir}\PYG{p}{,}
\PYG{+w}{  }\PYG{k+kt}{bool}\PYG{+w}{ }\PYG{k}{const}\PYG{+w}{ }\PYG{n}{recursive}\PYG{+w}{ }\PYG{o}{=}\PYG{+w}{ }\PYG{n+nb}{false}\PYG{p}{,}
\PYG{+w}{  }\PYG{k+kt}{unsigned}\PYG{+w}{ }\PYG{k+kt}{int}\PYG{+w}{ }\PYG{k}{const}\PYG{+w}{ }\PYG{n}{level}\PYG{+w}{ }\PYG{o}{=}\PYG{+w}{ }\PYG{l+m+mi}{0}\PYG{p}{)}
\PYG{p}{\PYGZob{}}
\PYG{+w}{  }\PYG{k}{if}\PYG{+w}{ }\PYG{p}{(}\PYG{n}{fs}\PYG{o}{:}\PYG{o}{:}\PYG{n}{exists}\PYG{p}{(}\PYG{n}{dir}\PYG{p}{)}\PYG{+w}{ }\PYG{o}{\PYGZam{}}\PYG{o}{\PYGZam{}}\PYG{+w}{ }\PYG{n}{fs}\PYG{o}{:}\PYG{o}{:}\PYG{n}{is\PYGZus{}directory}\PYG{p}{(}\PYG{n}{dir}\PYG{p}{)}\PYG{p}{)}
\PYG{+w}{  }\PYG{p}{\PYGZob{}}
\PYG{+w}{    }\PYG{k}{auto}\PYG{+w}{ }\PYG{n}{lead}\PYG{+w}{ }\PYG{o}{=}\PYG{+w}{ }\PYG{n}{std}\PYG{o}{:}\PYG{o}{:}\PYG{n}{string}\PYG{p}{(}\PYG{n}{level}\PYG{o}{*}\PYG{l+m+mi}{3}\PYG{p}{,}\PYG{+w}{ }\PYG{l+s+sc}{\PYGZsq{}}\PYG{l+s+sc}{ }\PYG{l+s+sc}{\PYGZsq{}}\PYG{p}{)}\PYG{p}{;}
\PYG{+w}{    }\PYG{k}{for}\PYG{+w}{ }\PYG{p}{(}\PYG{k}{auto}\PYG{+w}{ }\PYG{k}{const}\PYG{+w}{ }\PYG{o}{\PYGZam{}}\PYG{+w}{ }\PYG{n}{entry}\PYG{+w}{ }\PYG{o}{:}\PYG{+w}{ }\PYG{n}{fs}\PYG{o}{:}\PYG{o}{:}\PYG{n}{directory\PYGZus{}iterator}\PYG{p}{(}\PYG{n}{dir}\PYG{p}{)}\PYG{p}{)}
\PYG{+w}{    }\PYG{p}{\PYGZob{}}
\PYG{+w}{      }\PYG{k}{auto}\PYG{+w}{ }\PYG{n}{filename}\PYG{+w}{ }\PYG{o}{=}\PYG{+w}{ }\PYG{n}{entry}\PYG{p}{.}\PYG{n}{path}\PYG{p}{(}\PYG{p}{)}\PYG{p}{.}\PYG{n}{filename}\PYG{p}{(}\PYG{p}{)}\PYG{p}{;}
\PYG{+w}{      }\PYG{k}{if}\PYG{+w}{ }\PYG{p}{(}\PYG{n}{fs}\PYG{o}{:}\PYG{o}{:}\PYG{n}{is\PYGZus{}directory}\PYG{p}{(}\PYG{n}{entry}\PYG{p}{.}\PYG{n}{status}\PYG{p}{(}\PYG{p}{)}\PYG{p}{)}\PYG{p}{)}
\PYG{+w}{      }\PYG{p}{\PYGZob{}}
\PYG{+w}{        }\PYG{n}{std}\PYG{o}{:}\PYG{o}{:}\PYG{n}{cout}\PYG{+w}{ }\PYG{o}{\PYGZlt{}}\PYG{o}{\PYGZlt{}}\PYG{+w}{ }\PYG{n}{lead}\PYG{+w}{ }\PYG{o}{\PYGZlt{}}\PYG{o}{\PYGZlt{}}\PYG{+w}{ }\PYG{l+s}{\PYGZdq{}}\PYG{l+s}{[+]}\PYG{l+s}{\PYGZdq{}}\PYG{+w}{ }\PYG{o}{\PYGZlt{}}\PYG{o}{\PYGZlt{}}\PYG{+w}{ }\PYG{n}{filename}\PYG{+w}{ }\PYG{o}{\PYGZlt{}}\PYG{o}{\PYGZlt{}}\PYG{+w}{ }\PYG{n}{std}\PYG{o}{:}\PYG{o}{:}\PYG{n}{endl}\PYG{p}{;}
\PYG{+w}{        }\PYG{k}{if}\PYG{p}{(}\PYG{n}{recursive}\PYG{p}{)}
\PYG{+w}{          }\PYG{n}{visit\PYGZus{}directory}\PYG{p}{(}\PYG{n}{entry}\PYG{p}{,}\PYG{+w}{ }\PYG{n}{recursive}\PYG{p}{,}\PYG{+w}{ }\PYG{n}{level}\PYG{o}{+}\PYG{l+m+mi}{1}\PYG{p}{)}\PYG{p}{;}
\PYG{+w}{      }\PYG{p}{\PYGZcb{}}
\PYG{+w}{      }\PYG{k}{else}\PYG{+w}{ }\PYG{k}{if}\PYG{+w}{ }\PYG{p}{(}\PYG{n}{fs}\PYG{o}{:}\PYG{o}{:}\PYG{n}{is\PYGZus{}symlink}\PYG{p}{(}\PYG{n}{entry}\PYG{p}{.}\PYG{n}{status}\PYG{p}{(}\PYG{p}{)}\PYG{p}{)}\PYG{p}{)}
\PYG{+w}{        }\PYG{n}{std}\PYG{o}{:}\PYG{o}{:}\PYG{n}{cout}\PYG{+w}{ }\PYG{o}{\PYGZlt{}}\PYG{o}{\PYGZlt{}}\PYG{+w}{ }\PYG{n}{lead}\PYG{+w}{ }\PYG{o}{\PYGZlt{}}\PYG{o}{\PYGZlt{}}\PYG{+w}{ }\PYG{l+s}{\PYGZdq{}}\PYG{l+s}{[\PYGZgt{}]}\PYG{l+s}{\PYGZdq{}}\PYG{+w}{ }\PYG{o}{\PYGZlt{}}\PYG{o}{\PYGZlt{}}\PYG{+w}{ }\PYG{n}{filename}\PYG{+w}{ }\PYG{o}{\PYGZlt{}}\PYG{o}{\PYGZlt{}}\PYG{+w}{ }\PYG{n}{std}\PYG{o}{:}\PYG{o}{:}\PYG{n}{endl}\PYG{p}{;}
\PYG{+w}{      }\PYG{k}{else}\PYG{+w}{ }\PYG{k}{if}\PYG{+w}{ }\PYG{p}{(}\PYG{n}{fs}\PYG{o}{:}\PYG{o}{:}\PYG{n}{is\PYGZus{}regular\PYGZus{}file}\PYG{p}{(}\PYG{n}{entry}\PYG{p}{.}\PYG{n}{status}\PYG{p}{(}\PYG{p}{)}\PYG{p}{)}\PYG{p}{)}
\PYG{+w}{        }\PYG{n}{std}\PYG{o}{:}\PYG{o}{:}\PYG{n}{cout}\PYG{+w}{ }\PYG{o}{\PYGZlt{}}\PYG{o}{\PYGZlt{}}\PYG{+w}{ }\PYG{n}{lead}\PYG{+w}{ }\PYG{o}{\PYGZlt{}}\PYG{o}{\PYGZlt{}}\PYG{+w}{ }\PYG{l+s}{\PYGZdq{}}\PYG{l+s}{ }\PYG{l+s}{\PYGZdq{}}\PYG{+w}{ }\PYG{o}{\PYGZlt{}}\PYG{o}{\PYGZlt{}}\PYG{+w}{ }\PYG{n}{filename}\PYG{+w}{ }\PYG{o}{\PYGZlt{}}\PYG{o}{\PYGZlt{}}\PYG{+w}{ }\PYG{n}{std}\PYG{o}{:}\PYG{o}{:}\PYG{n}{endl}\PYG{p}{;}
\PYG{+w}{      }\PYG{k}{else}
\PYG{+w}{        }\PYG{n}{std}\PYG{o}{:}\PYG{o}{:}\PYG{n}{cout}\PYG{+w}{ }\PYG{o}{\PYGZlt{}}\PYG{o}{\PYGZlt{}}\PYG{+w}{ }\PYG{n}{lead}\PYG{+w}{ }\PYG{o}{\PYGZlt{}}\PYG{o}{\PYGZlt{}}\PYG{+w}{ }\PYG{l+s}{\PYGZdq{}}\PYG{l+s}{[?]}\PYG{l+s}{\PYGZdq{}}\PYG{+w}{ }\PYG{o}{\PYGZlt{}}\PYG{o}{\PYGZlt{}}\PYG{+w}{ }\PYG{n}{filename}\PYG{+w}{ }\PYG{o}{\PYGZlt{}}\PYG{o}{\PYGZlt{}}\PYG{+w}{ }\PYG{n}{std}\PYG{o}{:}\PYG{o}{:}\PYG{n}{endl}\PYG{p}{;}
\PYG{+w}{    }\PYG{p}{\PYGZcb{}}
\PYG{+w}{  }\PYG{p}{\PYGZcb{}}
\PYG{p}{\PYGZcb{}}
\end{Verbatim}
\end{codeblockbg}

\end{itemize}


%%C070494 h3

\subsection{工作原理}


%%C070495 p

\myttbr{directory\_\allowbreak{}iterator}与\myttbr{recursive\_\allowbreak{}directory\_\allowbreak{}iterator}都是输入迭代器，用于迭代目录的各个条目。区别在于，前者不递归访问子目录，而后者顾名思义会这样做。两者的行为相似：


%%C070496 li

\begin{itemize}
\item 迭代顺序未作规定。


%%C070497 li

\item 每个目录条目只被访问一次。


%%C070498 li

\item 特殊路径点（\texttt{.}）和双点（\texttt{..}）会被跳过。


%%C070499 li

\item 默认构造的迭代器是尾后迭代器，两个尾后迭代器总是相等。


%%C070500 li

\item 当迭代越过最后一个目录条目后，迭代器就与尾后迭代器相等。


%%C070501 li

\item 如果在迭代器创建之后向被迭代的目录中添加或删除条目，标准并未规定会发生什么。


%%C070502 li

\item 标准为\myttbr{directory\_\allowbreak{}iterator}和\myttbr{recursive\_\allowbreak{}directory\_\allowbreak{}iterator}都定义了非成员函数\texttt{begin()}和\texttt{end()}，这使我们能在基于范围的\texttt{for}循环中使用这两种迭代器，正如前面的示例所示。

\end{itemize}


%%C070503 p

两种迭代器都有多个重载构造函数。\myttbr{recursive\_\allowbreak{}directory\_\allowbreak{}iterator}构造函数的一些重载接受\myttbr{std::\allowbreak{}filesystem::\allowbreak{}directory\_\allowbreak{}options}类型的参数，用于指定迭代的附加选项：


%%C070504 li

\begin{itemize}
\item \texttt{none}：这是默认值，不指定任何选项。


%%C070505 li

\item \myttbr{follow\_\allowbreak{}directory\_\allowbreak{}symlink}：指定迭代应跟随符号链接，而不是访问链接自身。


%%C070506 li

\item \myttbr{skip\_\allowbreak{}permission\_\allowbreak{}denied}：指定忽略并跳过那些会触发“拒绝访问”错误的目录。

\end{itemize}


%%C070507 p

两种目录迭代器所指向的元素都是\texttt{directory\_\allowbreak{}entry}类型。\texttt{path()}成员函数返回该对象所表示的文件系统对象的路径。文件系统对象的状态可以用成员函数\texttt{status()}取得，对于符号链接则用\texttt{symlink\_\allowbreak{}status()}。


%%C070508 p

前面的示例遵循一个共同的模式：


%%C070509 li

\begin{itemize}
\item 确认要迭代的路径确实存在。


%%C070510 li

\item 用基于范围的\texttt{for}循环迭代目录的所有条目。


%%C070511 li

\item 根据迭代应当采用的方式，选用\texttt{filesystem}库提供的两种目录迭代器之一。


%%C070512 li

\item 按需求处理每个条目。

\end{itemize}


%%C070513 p

在我们的示例中，只是把目录条目的名称打印到控制台。要注意，如前所述，目录内容的迭代顺序是未指定的。如果你想以结构化方式处理内容，例如把子目录及其条目按缩进显示（本案例正是如此），或者以树形显示（在其他类型的应用中），那么使用\myttbr{recursive\_\allowbreak{}directory\_\allowbreak{}iterator}并不合适。此时应在迭代过程中对每个子目录递归调用一个使用\myttbr{directory\_\allowbreak{}iterator}的函数，如上一节最后一个示例所示。


%%C070514 p

考虑本案例开头给出的目录结构（相对于当前路径），使用递归迭代器时，得到的输出如下例所示：


%%C070515 code

\begin{codeblockbg}
\begin{Verbatim}[commandchars=\\\{\},breaklines=true,breakanywhere=true,fontsize=\codefont,xleftmargin=17bp]
\PYG{n}{visit\PYGZus{}directory\PYGZus{}rec}\PYG{p}{(}\PYG{n}{fs}\PYG{o}{:}\PYG{o}{:}\PYG{n}{current\PYGZus{}path}\PYG{p}{(}\PYG{p}{)}\PYG{+w}{ }\PYG{o}{/}\PYG{+w}{ }\PYG{l+s}{\PYGZdq{}}\PYG{l+s}{test}\PYG{l+s}{\PYGZdq{}}\PYG{p}{)}\PYG{p}{;}
\PYG{p}{[}\PYG{o}{+}\PYG{p}{]}\PYG{n}{data}
\PYG{+w}{   }\PYG{n}{input}\PYG{p}{.}\PYG{n}{dat}
\PYG{+w}{   }\PYG{n}{output}\PYG{p}{.}\PYG{n}{dat}
\PYG{+w}{   }\PYG{n}{file\PYGZus{}1}\PYG{p}{.}\PYG{n}{txt}
\PYG{+w}{   }\PYG{n}{file\PYGZus{}2}\PYG{p}{.}\PYG{n}{txt}
\PYG{+w}{   }\PYG{n}{file\PYGZus{}3}\PYG{p}{.}\PYG{n}{log}
\end{Verbatim}
\end{codeblockbg}


%%C070516 p

另一方面，如下所示使用第三个示例中的递归函数时，输出按子层级有序显示，符合预期：


%%C070517 code

\begin{codeblockbg}
\begin{Verbatim}[commandchars=\\\{\},breaklines=true,breakanywhere=true,fontsize=\codefont,xleftmargin=17bp]
\PYG{n}{visit\PYGZus{}directory}\PYG{p}{(}\PYG{n}{fs}\PYG{o}{:}\PYG{o}{:}\PYG{n}{current\PYGZus{}path}\PYG{p}{(}\PYG{p}{)}\PYG{+w}{ }\PYG{o}{/}\PYG{+w}{ }\PYG{l+s}{\PYGZdq{}}\PYG{l+s}{test}\PYG{l+s}{\PYGZdq{}}\PYG{p}{,}\PYG{+w}{ }\PYG{n+nb}{true}\PYG{p}{)}\PYG{p}{;}
\PYG{p}{[}\PYG{o}{+}\PYG{p}{]}\PYG{n}{data}
\PYG{+w}{      }\PYG{n}{input}\PYG{p}{.}\PYG{n}{dat}
\PYG{+w}{      }\PYG{n}{output}\PYG{p}{.}\PYG{n}{dat}
\PYG{+w}{   }\PYG{n}{file\PYGZus{}1}\PYG{p}{.}\PYG{n}{txt}
\PYG{+w}{   }\PYG{n}{file\PYGZus{}2}\PYG{p}{.}\PYG{n}{txt}
\PYG{+w}{   }\PYG{n}{file\PYGZus{}3}\PYG{p}{.}\PYG{n}{log}
\end{Verbatim}
\end{codeblockbg}


%%C070518 h3

\subsection{更多内容}


%%C070519 p

在上一案例\textit{检查现有文件或目录的属性}中，我们讨论过\texttt{file\_\allowbreak{}size()}等函数，它返回文件的大小（以字节计）。不过，如果指定的路径是目录，该函数会失败。要确定目录的大小，我们需要递归遍历目录内容，取得普通文件或符号链接的大小并把它们相加。但必须注意检查\texttt{file\_\allowbreak{}size()}的返回值：出错时它是转换为\texttt{std::\allowbreak{}uintmax\_\allowbreak{}t}的\texttt{-\allowbreak{}1}。这个表示失败的值不应计入目录的总大小。下面的函数在出错时同样返回作为\texttt{uintmax\_\allowbreak{}t}的\texttt{-\allowbreak{}1}：


%%C070520 code

\begin{codeblockbg}
\begin{Verbatim}[commandchars=\\\{\},breaklines=true,breakanywhere=true,fontsize=\codefont,xleftmargin=17bp]
\PYG{n}{std}\PYG{o}{:}\PYG{o}{:}\PYG{k+kt}{uintmax\PYGZus{}t}\PYG{+w}{ }\PYG{n+nf}{dir\PYGZus{}size}\PYG{p}{(}\PYG{n}{fs}\PYG{o}{:}\PYG{o}{:}\PYG{n}{path}\PYG{+w}{ }\PYG{k}{const}\PYG{+w}{ }\PYG{o}{\PYGZam{}}\PYG{+w}{ }\PYG{n}{path}\PYG{p}{)}
\PYG{p}{\PYGZob{}}
\PYG{+w}{  }\PYG{k}{auto}\PYG{+w}{ }\PYG{n}{size}\PYG{+w}{ }\PYG{o}{=}\PYG{+w}{ }\PYG{k}{static\PYGZus{}cast}\PYG{o}{\PYGZlt{}}\PYG{k+kt}{uintmax\PYGZus{}t}\PYG{o}{\PYGZgt{}}\PYG{p}{(}\PYG{l+m+mi}{\PYGZhy{}1}\PYG{p}{)}\PYG{p}{;}
\PYG{+w}{  }\PYG{k}{if}\PYG{+w}{ }\PYG{p}{(}\PYG{n}{fs}\PYG{o}{:}\PYG{o}{:}\PYG{n}{exists}\PYG{p}{(}\PYG{n}{path}\PYG{p}{)}\PYG{+w}{ }\PYG{o}{\PYGZam{}}\PYG{o}{\PYGZam{}}\PYG{+w}{ }\PYG{n}{fs}\PYG{o}{:}\PYG{o}{:}\PYG{n}{is\PYGZus{}directory}\PYG{p}{(}\PYG{n}{path}\PYG{p}{)}\PYG{p}{)}
\PYG{+w}{  }\PYG{p}{\PYGZob{}}
\PYG{+w}{    }\PYG{k}{for}\PYG{+w}{ }\PYG{p}{(}\PYG{k}{auto}\PYG{+w}{ }\PYG{k}{const}\PYG{+w}{ }\PYG{o}{\PYGZam{}}\PYG{+w}{ }\PYG{n}{entry}\PYG{+w}{ }\PYG{o}{:}\PYG{+w}{ }\PYG{n}{fs}\PYG{o}{:}\PYG{o}{:}\PYG{n}{recursive\PYGZus{}directory\PYGZus{}iterator}\PYG{p}{(}\PYG{n}{path}\PYG{p}{)}\PYG{p}{)}
\PYG{+w}{    }\PYG{p}{\PYGZob{}}
\PYG{+w}{      }\PYG{k}{if}\PYG{+w}{ }\PYG{p}{(}\PYG{n}{fs}\PYG{o}{:}\PYG{o}{:}\PYG{n}{is\PYGZus{}regular\PYGZus{}file}\PYG{p}{(}\PYG{n}{entry}\PYG{p}{.}\PYG{n}{status}\PYG{p}{(}\PYG{p}{)}\PYG{p}{)}\PYG{+w}{ }\PYG{o}{|}\PYG{o}{|}
\PYG{+w}{      }\PYG{n}{fs}\PYG{o}{:}\PYG{o}{:}\PYG{n}{is\PYGZus{}symlink}\PYG{p}{(}\PYG{n}{entry}\PYG{p}{.}\PYG{n}{status}\PYG{p}{(}\PYG{p}{)}\PYG{p}{)}\PYG{p}{)}
\PYG{+w}{      }\PYG{p}{\PYGZob{}}
\PYG{+w}{        }\PYG{k}{auto}\PYG{+w}{ }\PYG{n}{err}\PYG{+w}{ }\PYG{o}{=}\PYG{+w}{ }\PYG{n}{std}\PYG{o}{:}\PYG{o}{:}\PYG{n}{error\PYGZus{}code}\PYG{p}{\PYGZob{}}\PYG{p}{\PYGZcb{}}\PYG{p}{;}
\PYG{+w}{        }\PYG{k}{auto}\PYG{+w}{ }\PYG{n}{filesize}\PYG{+w}{ }\PYG{o}{=}\PYG{+w}{ }\PYG{n}{fs}\PYG{o}{:}\PYG{o}{:}\PYG{n}{file\PYGZus{}size}\PYG{p}{(}\PYG{n}{entry}\PYG{p}{)}\PYG{p}{;}
\PYG{+w}{        }\PYG{k}{if}\PYG{+w}{ }\PYG{p}{(}\PYG{n}{filesize}\PYG{+w}{ }\PYG{o}{!}\PYG{o}{=}\PYG{+w}{ }\PYG{k}{static\PYGZus{}cast}\PYG{o}{\PYGZlt{}}\PYG{k+kt}{uintmax\PYGZus{}t}\PYG{o}{\PYGZgt{}}\PYG{p}{(}\PYG{l+m+mi}{\PYGZhy{}1}\PYG{p}{)}\PYG{p}{)}
\PYG{+w}{          }\PYG{n}{size}\PYG{+w}{ }\PYG{o}{+}\PYG{o}{=}\PYG{+w}{ }\PYG{n}{filesize}\PYG{p}{;}
\PYG{+w}{      }\PYG{p}{\PYGZcb{}}
\PYG{+w}{    }\PYG{p}{\PYGZcb{}}
\PYG{+w}{  }\PYG{p}{\PYGZcb{}}
\PYG{+w}{  }\PYG{k}{return}\PYG{+w}{ }\PYG{n}{size}\PYG{p}{;}
\PYG{p}{\PYGZcb{}}
\end{Verbatim}
\end{codeblockbg}


%%C070521 h3

\subsection{延伸阅读}


%%C070522 li

\begin{itemize}
\item \textit{检查现有文件或目录的属性}


%%C070523 li

\item \textit{查找文件}

\end{itemize}


%%C070524 h2

\section{查找文件}


%%C070525 p

在上一案例中，我们看到了如何使用\myttbr{directory\_\allowbreak{}iterator}和\myttbr{recursive\_\allowbreak{}directory\_\allowbreak{}iterator}枚举目录内容。像上一案例那样显示目录内容，只是需要这样做的场景之一。另一个主要场景是在目录中搜索特定条目，例如具有特定名称、特定扩展名等的文件。本案例将介绍如何使用目录迭代器以及前面展示的迭代模式来查找符合给定条件的文件。


%%C070526 h3

\subsection{准备工作}


%%C070527 p

关于目录迭代器的细节，你应阅读上一案例\textit{枚举目录内容}。本案例还将使用上一案例中给出的同一个测试目录结构。


%%C070528 h3

\subsection{具体做法}


%%C070529 p

要查找符合特定条件的文件，请使用以下模式，\texttt{find\_\allowbreak{}files()}函数就是它的示例：


%%C070530 p

使用\myttbr{recursive\_\allowbreak{}directory\_\allowbreak{}iterator}迭代目录的所有条目


%%C070531 p

并递归遍历其子目录。考虑普通文件（以及你可能需要处理的其他类型文件）。


%%C070532 p

使用一个函数对象（例如lambda表达式）只筛选出符合


%%C070533 p

条件的文件。把选中的条目加入一个范围（例如vector）：


%%C070534 code

\begin{codeblockbg}
\begin{Verbatim}[commandchars=\\\{\},breaklines=true,breakanywhere=true,fontsize=\codefont,xleftmargin=17bp]
\PYG{n}{std}\PYG{o}{:}\PYG{o}{:}\PYG{n}{vector}\PYG{o}{\PYGZlt{}}\PYG{n}{fs}\PYG{o}{:}\PYG{o}{:}\PYG{n}{path}\PYG{o}{\PYGZgt{}}\PYG{+w}{ }\PYG{n}{find\PYGZus{}files}\PYG{p}{(}
\PYG{+w}{   }\PYG{n}{fs}\PYG{o}{:}\PYG{o}{:}\PYG{n}{path}\PYG{+w}{ }\PYG{k}{const}\PYG{+w}{ }\PYG{o}{\PYGZam{}}\PYG{+w}{ }\PYG{n}{dir}\PYG{p}{,}
\PYG{+w}{   }\PYG{n}{std}\PYG{o}{:}\PYG{o}{:}\PYG{n}{function}\PYG{o}{\PYGZlt{}}\PYG{k+kt}{bool}\PYG{p}{(}\PYG{n}{fs}\PYG{o}{:}\PYG{o}{:}\PYG{n}{path}\PYG{+w}{ }\PYG{k}{const}\PYG{o}{\PYGZam{}}\PYG{p}{)}\PYG{o}{\PYGZgt{}}\PYG{+w}{ }\PYG{n}{filter}\PYG{p}{)}
\PYG{p}{\PYGZob{}}
\PYG{+w}{  }\PYG{k}{auto}\PYG{+w}{ }\PYG{n}{result}\PYG{+w}{ }\PYG{o}{=}\PYG{+w}{ }\PYG{n}{std}\PYG{o}{:}\PYG{o}{:}\PYG{n}{vector}\PYG{o}{\PYGZlt{}}\PYG{n}{fs}\PYG{o}{:}\PYG{o}{:}\PYG{n}{path}\PYG{o}{\PYGZgt{}}\PYG{p}{\PYGZob{}}\PYG{p}{\PYGZcb{}}\PYG{p}{;}
\PYG{k}{if}\PYG{+w}{ }\PYG{p}{(}\PYG{n}{fs}\PYG{o}{:}\PYG{o}{:}\PYG{n}{exists}\PYG{p}{(}\PYG{n}{dir}\PYG{p}{)}\PYG{p}{)}
\PYG{p}{\PYGZob{}}
\PYG{+w}{  }\PYG{k}{for}\PYG{+w}{ }\PYG{p}{(}\PYG{k}{auto}\PYG{+w}{ }\PYG{k}{const}\PYG{+w}{ }\PYG{o}{\PYGZam{}}\PYG{+w}{ }\PYG{n}{entry}\PYG{+w}{ }\PYG{o}{:}
\PYG{+w}{    }\PYG{n}{fs}\PYG{o}{:}\PYG{o}{:}\PYG{n}{recursive\PYGZus{}directory\PYGZus{}iterator}\PYG{p}{(}
\PYG{+w}{      }\PYG{n}{dir}\PYG{p}{,}
\PYG{+w}{      }\PYG{n}{fs}\PYG{o}{:}\PYG{o}{:}\PYG{n}{directory\PYGZus{}options}\PYG{o}{:}\PYG{o}{:}\PYG{n}{follow\PYGZus{}directory\PYGZus{}symlink}\PYG{p}{)}\PYG{p}{)}
\PYG{+w}{  }\PYG{p}{\PYGZob{}}
\PYG{+w}{    }\PYG{k}{if}\PYG{+w}{ }\PYG{p}{(}\PYG{n}{fs}\PYG{o}{:}\PYG{o}{:}\PYG{n}{is\PYGZus{}regular\PYGZus{}file}\PYG{p}{(}\PYG{n}{entry}\PYG{p}{)}\PYG{+w}{ }\PYG{o}{\PYGZam{}}\PYG{o}{\PYGZam{}}
\PYG{+w}{       }\PYG{n}{filter}\PYG{p}{(}\PYG{n}{entry}\PYG{p}{)}\PYG{p}{)}
\PYG{+w}{    }\PYG{p}{\PYGZob{}}
\PYG{+w}{      }\PYG{n}{result}\PYG{p}{.}\PYG{n}{push\PYGZus{}back}\PYG{p}{(}\PYG{n}{entry}\PYG{p}{)}\PYG{p}{;}
\PYG{+w}{    }\PYG{p}{\PYGZcb{}}
\PYG{+w}{  }\PYG{p}{\PYGZcb{}}
\PYG{p}{\PYGZcb{}}
\PYG{+w}{  }\PYG{k}{return}\PYG{+w}{ }\PYG{n}{result}\PYG{p}{;}
\PYG{p}{\PYGZcb{}}
\end{Verbatim}
\end{codeblockbg}


%%C070535 h3

\subsection{工作原理}


%%C070536 p

当我们要在目录中查找文件时，目录的结构以及访问其各个条目（包括子目录）的顺序大概并不重要。因此，我们可以用\myttbr{recursive\_\allowbreak{}directory\_\allowbreak{}iterator}来迭代这些条目。


%%C070537 p

函数\texttt{find\_\allowbreak{}files()}接受两个参数：一个路径，以及一个用于筛选应返回条目的函数包装器。返回类型是\texttt{filesystem::\allowbreak{}path}的vector，不过也可以是\myttbr{filesystem::\allowbreak{}directory\_\allowbreak{}entry}的vector。本例使用的递归目录迭代器不跟随符号链接，返回的是链接本身而不是其目标。这一行为可以改变：使用带\myttbr{filesystem::\allowbreak{}directory\_\allowbreak{}options}类型参数的重载构造函数，并传入\myttbr{follow\_\allowbreak{}directory\_\allowbreak{}symlink}即可。


%%C070538 p

在上面的例子中，我们只考虑普通文件，忽略其他类型的文件系统对象。谓词作用于目录条目，若返回\texttt{true}，就把该条目加入结果。


%%C070539 p

下面的例子用\texttt{find\_\allowbreak{}files()}函数查找测试目录中所有以\texttt{file\_}为前缀的文件：


%%C070540 code

\begin{codeblockbg}
\begin{Verbatim}[commandchars=\\\{\},breaklines=true,breakanywhere=true,fontsize=\codefont,xleftmargin=17bp]
\PYG{k}{auto}\PYG{+w}{ }\PYG{n}{results}\PYG{+w}{ }\PYG{o}{=}\PYG{+w}{ }\PYG{n}{find\PYGZus{}files}\PYG{p}{(}
\PYG{+w}{          }\PYG{n}{fs}\PYG{o}{:}\PYG{o}{:}\PYG{n}{current\PYGZus{}path}\PYG{p}{(}\PYG{p}{)}\PYG{+w}{ }\PYG{o}{/}\PYG{+w}{ }\PYG{l+s}{\PYGZdq{}}\PYG{l+s}{test}\PYG{l+s}{\PYGZdq{}}\PYG{p}{,}
\PYG{+w}{          }\PYG{p}{[}\PYG{p}{]}\PYG{p}{(}\PYG{n}{fs}\PYG{o}{:}\PYG{o}{:}\PYG{n}{path}\PYG{+w}{ }\PYG{k}{const}\PYG{+w}{ }\PYG{o}{\PYGZam{}}\PYG{+w}{ }\PYG{n}{p}\PYG{p}{)}\PYG{+w}{ }\PYG{p}{\PYGZob{}}
\PYG{+w}{  }\PYG{k}{auto}\PYG{+w}{ }\PYG{n}{filename}\PYG{+w}{ }\PYG{o}{=}\PYG{+w}{ }\PYG{n}{p}\PYG{p}{.}\PYG{n}{wstring}\PYG{p}{(}\PYG{p}{)}\PYG{p}{;}
\PYG{+w}{  }\PYG{k}{return}\PYG{+w}{ }\PYG{n}{filename}\PYG{p}{.}\PYG{n}{find}\PYG{p}{(}\PYG{l+s+sa}{L}\PYG{l+s}{\PYGZdq{}}\PYG{l+s}{file\PYGZus{}}\PYG{l+s}{\PYGZdq{}}\PYG{p}{)}\PYG{+w}{ }\PYG{o}{!}\PYG{o}{=}\PYG{+w}{ }\PYG{n}{std}\PYG{o}{:}\PYG{o}{:}\PYG{n}{wstring}\PYG{o}{:}\PYG{o}{:}\PYG{n}{npos}\PYG{p}{;}
\PYG{p}{\PYGZcb{}}\PYG{p}{)}\PYG{p}{;}
\PYG{k}{for}\PYG{+w}{ }\PYG{p}{(}\PYG{k}{auto}\PYG{+w}{ }\PYG{k}{const}\PYG{+w}{ }\PYG{o}{\PYGZam{}}\PYG{+w}{ }\PYG{n}{path}\PYG{+w}{ }\PYG{o}{:}\PYG{+w}{ }\PYG{n}{results}\PYG{p}{)}
\PYG{p}{\PYGZob{}}
\PYG{+w}{  }\PYG{n}{std}\PYG{o}{:}\PYG{o}{:}\PYG{n}{cout}\PYG{+w}{ }\PYG{o}{\PYGZlt{}}\PYG{o}{\PYGZlt{}}\PYG{+w}{ }\PYG{n}{path}\PYG{+w}{ }\PYG{o}{\PYGZlt{}}\PYG{o}{\PYGZlt{}}\PYG{+w}{ }\PYG{n}{std}\PYG{o}{:}\PYG{o}{:}\PYG{n}{endl}\PYG{p}{;}
\PYG{p}{\PYGZcb{}}
\end{Verbatim}
\end{codeblockbg}


%%C070541 p

执行该程序，得到的输出如下（路径相对于当前路径）：


%%C070542 code

\begin{codeblockbg}
\begin{Verbatim}[commandchars=\\\{\},breaklines=true,breakanywhere=true,fontsize=\codefont,xleftmargin=17bp]
\PYG{n}{test}\PYG{err}{\PYGZbs{}}\PYG{n}{file\PYGZus{}1}\PYG{p}{.}\PYG{n}{txt}
\PYG{n}{test}\PYG{err}{\PYGZbs{}}\PYG{n}{file\PYGZus{}2}\PYG{p}{.}\PYG{n}{txt}
\PYG{n}{test}\PYG{err}{\PYGZbs{}}\PYG{n}{file\PYGZus{}3}\PYG{p}{.}\PYG{n}{log}
\end{Verbatim}
\end{codeblockbg}


%%C070543 p

第二个示例演示如何查找具有特定扩展名的文件，这里以扩展名\texttt{.\allowbreak{}dat}为例：


%%C070544 code

\begin{codeblockbg}
\begin{Verbatim}[commandchars=\\\{\},breaklines=true,breakanywhere=true,fontsize=\codefont,xleftmargin=17bp]
\PYG{k}{auto}\PYG{+w}{ }\PYG{n}{results}\PYG{+w}{ }\PYG{o}{=}\PYG{+w}{ }\PYG{n}{find\PYGZus{}files}\PYG{p}{(}
\PYG{+w}{       }\PYG{n}{fs}\PYG{o}{:}\PYG{o}{:}\PYG{n}{current\PYGZus{}path}\PYG{p}{(}\PYG{p}{)}\PYG{+w}{ }\PYG{o}{/}\PYG{+w}{ }\PYG{l+s}{\PYGZdq{}}\PYG{l+s}{test}\PYG{l+s}{\PYGZdq{}}\PYG{p}{,}
\PYG{+w}{       }\PYG{p}{[}\PYG{p}{]}\PYG{p}{(}\PYG{n}{fs}\PYG{o}{:}\PYG{o}{:}\PYG{n}{path}\PYG{+w}{ }\PYG{k}{const}\PYG{+w}{ }\PYG{o}{\PYGZam{}}\PYG{+w}{ }\PYG{n}{p}\PYG{p}{)}\PYG{+w}{ }\PYG{p}{\PYGZob{}}
\PYG{+w}{  }\PYG{k}{return}\PYG{+w}{ }\PYG{n}{p}\PYG{p}{.}\PYG{n}{extension}\PYG{p}{(}\PYG{p}{)}\PYG{+w}{ }\PYG{o}{=}\PYG{o}{=}\PYG{+w}{ }\PYG{l+s+sa}{L}\PYG{l+s}{\PYGZdq{}}\PYG{l+s}{.dat}\PYG{l+s}{\PYGZdq{}}\PYG{p}{;}
\PYG{p}{\PYGZcb{}}\PYG{p}{)}\PYG{p}{;}
\PYG{k}{for}\PYG{+w}{ }\PYG{p}{(}\PYG{k}{auto}\PYG{+w}{ }\PYG{k}{const}\PYG{+w}{ }\PYG{o}{\PYGZam{}}\PYG{+w}{ }\PYG{n}{path}\PYG{+w}{ }\PYG{o}{:}\PYG{+w}{ }\PYG{n}{results}\PYG{p}{)}
\PYG{p}{\PYGZob{}}
\PYG{+w}{  }\PYG{n}{std}\PYG{o}{:}\PYG{o}{:}\PYG{n}{cout}\PYG{+w}{ }\PYG{o}{\PYGZlt{}}\PYG{o}{\PYGZlt{}}\PYG{+w}{ }\PYG{n}{path}\PYG{+w}{ }\PYG{o}{\PYGZlt{}}\PYG{o}{\PYGZlt{}}\PYG{+w}{ }\PYG{n}{std}\PYG{o}{:}\PYG{o}{:}\PYG{n}{endl}\PYG{p}{;}
\PYG{p}{\PYGZcb{}}
\end{Verbatim}
\end{codeblockbg}


%%C070545 p

输出同样相对于当前路径，如下所示：


%%C070546 code

\begin{codeblockbg}
\begin{Verbatim}[commandchars=\\\{\},breaklines=true,breakanywhere=true,fontsize=\codefont,xleftmargin=17bp]
\PYG{n}{test}\PYG{err}{\PYGZbs{}}\PYG{n}{datainput}\PYG{p}{.}\PYG{n}{dat}
\PYG{n}{test}\PYG{err}{\PYGZbs{}}\PYG{n}{dataoutput}\PYG{p}{.}\PYG{n}{dat}
\end{Verbatim}
\end{codeblockbg}


%%C070547 h3

\subsection{延伸阅读}


%%C070548 li

\begin{itemize}
\item \textit{检查现有文件或目录的属性}


%%C070549 li

\item \textit{枚举目录内容}

\end{itemize}

%%C070542 code

\begin{codeblockbg}
\begin{Verbatim}[commandchars=\\\{\},breaklines=true,breakanywhere=true,fontsize=\codefont,xleftmargin=17bp]
\PYG{n}{test}\PYG{err}{\PYGZbs{}}\PYG{n}{file\PYGZus{}1}\PYG{p}{.}\PYG{n}{txt}
\PYG{n}{test}\PYG{err}{\PYGZbs{}}\PYG{n}{file\PYGZus{}2}\PYG{p}{.}\PYG{n}{txt}
\PYG{n}{test}\PYG{err}{\PYGZbs{}}\PYG{n}{file\PYGZus{}3}\PYG{p}{.}\PYG{n}{log}
\end{Verbatim}
\end{codeblockbg}


%%C070543 p

第二个示例演示如何查找具有特定扩展名的文件，这里以扩展名\texttt{.\allowbreak{}dat}为例：


%%C070544 code

\begin{codeblockbg}
\begin{Verbatim}[commandchars=\\\{\},breaklines=true,breakanywhere=true,fontsize=\codefont,xleftmargin=17bp]
\PYG{k}{auto}\PYG{+w}{ }\PYG{n}{results}\PYG{+w}{ }\PYG{o}{=}\PYG{+w}{ }\PYG{n}{find\PYGZus{}files}\PYG{p}{(}
\PYG{+w}{       }\PYG{n}{fs}\PYG{o}{:}\PYG{o}{:}\PYG{n}{current\PYGZus{}path}\PYG{p}{(}\PYG{p}{)}\PYG{+w}{ }\PYG{o}{/}\PYG{+w}{ }\PYG{l+s}{\PYGZdq{}}\PYG{l+s}{test}\PYG{l+s}{\PYGZdq{}}\PYG{p}{,}
\PYG{+w}{       }\PYG{p}{[}\PYG{p}{]}\PYG{p}{(}\PYG{n}{fs}\PYG{o}{:}\PYG{o}{:}\PYG{n}{path}\PYG{+w}{ }\PYG{k}{const}\PYG{+w}{ }\PYG{o}{\PYGZam{}}\PYG{+w}{ }\PYG{n}{p}\PYG{p}{)}\PYG{+w}{ }\PYG{p}{\PYGZob{}}
\PYG{+w}{  }\PYG{k}{return}\PYG{+w}{ }\PYG{n}{p}\PYG{p}{.}\PYG{n}{extension}\PYG{p}{(}\PYG{p}{)}\PYG{+w}{ }\PYG{o}{=}\PYG{o}{=}\PYG{+w}{ }\PYG{l+s+sa}{L}\PYG{l+s}{\PYGZdq{}}\PYG{l+s}{.dat}\PYG{l+s}{\PYGZdq{}}\PYG{p}{;}
\PYG{p}{\PYGZcb{}}\PYG{p}{)}\PYG{p}{;}
\PYG{k}{for}\PYG{+w}{ }\PYG{p}{(}\PYG{k}{auto}\PYG{+w}{ }\PYG{k}{const}\PYG{+w}{ }\PYG{o}{\PYGZam{}}\PYG{+w}{ }\PYG{n}{path}\PYG{+w}{ }\PYG{o}{:}\PYG{+w}{ }\PYG{n}{results}\PYG{p}{)}
\PYG{p}{\PYGZob{}}
\PYG{+w}{  }\PYG{n}{std}\PYG{o}{:}\PYG{o}{:}\PYG{n}{cout}\PYG{+w}{ }\PYG{o}{\PYGZlt{}}\PYG{o}{\PYGZlt{}}\PYG{+w}{ }\PYG{n}{path}\PYG{+w}{ }\PYG{o}{\PYGZlt{}}\PYG{o}{\PYGZlt{}}\PYG{+w}{ }\PYG{n}{std}\PYG{o}{:}\PYG{o}{:}\PYG{n}{endl}\PYG{p}{;}
\PYG{p}{\PYGZcb{}}
\end{Verbatim}
\end{codeblockbg}


%%C070545 p

输出同样相对于当前路径，如下所示：


%%C070546 code

\begin{codeblockbg}
\begin{Verbatim}[commandchars=\\\{\},breaklines=true,breakanywhere=true,fontsize=\codefont,xleftmargin=17bp]
\PYG{n}{test}\PYG{err}{\PYGZbs{}}\PYG{n}{datainput}\PYG{p}{.}\PYG{n}{dat}
\PYG{n}{test}\PYG{err}{\PYGZbs{}}\PYG{n}{dataoutput}\PYG{p}{.}\PYG{n}{dat}
\end{Verbatim}
\end{codeblockbg}


%%C070547 h3

\subsection{延伸阅读}


%%C070548 li

\begin{itemize}
\item \textit{检查现有文件或目录的属性}


%%C070549 li

\item \textit{枚举目录内容}

\end{itemize}
